\section{The primary arithmetic object and the topology actually specified}
\label{sec:cq-source}
The starting point is Connes--Consani--Moscovici,
\emph{Zeta zeros and prolate wave operators: Semilocal adelic operators},
Section 3.6, proof-PDF pages 17--19, especially Lemma 3.3,
equations (3.16)--(3.22), and Proposition 3.6
\cite{cqCCM}. The preceding arithmetic map is
\[
 \mathcal E f(x)=x^{1/2}\sum_{n\geq1}f(nx),\qquad
 \mathcal F_\mu g(s)=\int_0^\infty g(x)x^{-is}\,\frac{dx}{x}.
\]
The input consists of even Schwartz functions with
$f(0)=\widehat f(0)=0$. The paper's two Hermite families are
\[
 \psi_\ell^+=h_{4\ell}-\frac{h_{4\ell}(0)}{h_0(0)}h_0,
 \qquad
 \psi_\ell^-=-h_{4\ell+2}
       +\frac{h_{4\ell+2}(0)}{h_2(0)}h_2\quad(\ell\geq1).
\]
These signs and the minus sign in the Mellin exponent are retained.
The completed function is
\begin{equation}\label{eq:cq-Xi}
 \Xi(s)=\xi(\tfrac12+is),\qquad
 \xi(z)=\tfrac12z(z-1)\pi^{-z/2}\Gamma(z/2)\zeta(z),
 \qquad \Xi(s)=8H_0(2s).
\end{equation}
The paper proves
\begin{equation}\label{eq:cq-image}
 \mathcal F_\mu\mathcal E\psi_\ell^\pm=R_\ell^\pm\Xi,
 \qquad P_\ell^\pm=-\tfrac12(\tfrac14+s^2)R_\ell^\pm.
\end{equation}
It also states immediately before Proposition 3.6 that these give all
polynomial multiples of $\Xi$. Here is the precise algebra behind the
last assertion. The top summands of (3.19) and (3.20) are nonzero;
their degrees are $2\ell$ and $2\ell+1$. Division by
$-\tfrac12(\tfrac14+s^2)$ gives degrees $2\ell-2$ and
$2\ell-1$, respectively. The even and odd parities remain.
Thus the ordered family $R_1^+,R_1^-,R_2^+,R_2^-,\ldots$
has one polynomial of each degree, with a nonzero leading coefficient.
Subtracting the appropriate multiple of the polynomial of largest
degree, and iterating down to degree zero, expresses every polynomial
as a finite linear combination. Linear independence follows by the
largest-degree coefficient of a finite relation. Consequently
\begin{equation}\label{eq:cq-N}
 \cqN=\overline{\operatorname{span}
 \{\mathcal F_\mu\mathcal E\psi_\ell^\pm\}}^{\,\tau}
 =\overline{\cqC[s]\Xi}^{\,\tau},\qquad
 \cqQ=(\cqE,\tau)/\cqN.
\end{equation}
Here $\tau$ denotes precisely the source's topology. No equality of
this closed space with a principal algebraic ideal is being asserted.
Multiplication $M f(s)=sf(s)$ is continuous in the topological ring;
it preserves the dense generating subspace, and hence $\cqN$.
Its induced operator $S$ is the operator of Proposition 3.6.
We use the source's spectral assertion as an explicitly imported theorem:
\begin{equation}\label{eq:cq-source-spectrum}
 \cqSpec(S)=\{\lambda\in\cqC:
             \zeta(\tfrac12+i\lambda)=0
             \text{ at a nontrivial zero}\}.
\end{equation}
Spectrum here means failure of a continuous two-sided inverse.
No proof of the entire source spectral theorem is claimed as a new
result of this manuscript.

There is a material source limitation. Page 19 calls $\cqE$ the
Hadamard topological ring of entire functions of order at most one;
it does not state seminorms or a convergence criterion. The cited
Connes--Marcolli Proposition 2.24, printed pages 379--381,
proves the polynomial-image assertion and explicitly discusses
different function-space realizations \cite{cqCM}. It does not supply
that missing specification. We therefore retain $\tau$ in every
source quotient. Below, all assertions depending on another topology
name and define that topology. The comparison with the source is an
exact quotient diagram, rather than an identification by terminology.
The bounded edition audit gives the same limitation in the primary
arXiv v1 source, lines 683--704 (Proposition 3.9), and in v2,
\texttt{mainc2m24fine.tex}, lines 668--689 (Proposition 3.6):
each states the order class, polynomial trace image, and spectral
set, without defining its topology or algebraic multiplicities.
These exact edition locators supplement the author-proof locator
in \cite{cqCCM}; no claim of byte identity with a publisher's final
subscription text is required.
Following the book's cited Sobolev construction leads to explicit
weighted Hilbert spaces in its Chapter 2, Section 9, and in
Connes's original spectral realization \cite{cqConnesSelecta}.
Section~\ref{sec:cq-sobolev-comparison} retains those spaces and
proves their exact common-test comparison with $\cqQ$, including
the original generator and the allowed zero multiplicities.
The inspected construction does not specify the topology on all
of $\cqE$ or identify it with those Hilbert topologies.

\section{A globally defined heat map on the original underlying functions}
\label{sec:cq-heatspace}
For an entire function $f$, let
$M_f(r)=\max_{|s|\leq r}|f(s)|$ and define its order by
\[
 \limsup_{r\to\infty}
       \frac{\log\log\max(e,M_f(r))}{\log r}.
\]
The zero function is included. For $1<p<2$, $a>0$ set
\[
 \|f\|_{p,a}=\sup_{s\in\cqC}|f(s)|e^{-a|s|^p}.
\]
The underlying set $\cqE$ is exactly the set where all these seminorms
are finite. Indeed order at most one gives, for $1<q<p$,
$M_f(r)\leq\exp(r^q)$ for sufficiently large $r$; the expression
$r^q-ar^p$ is bounded above. Conversely finiteness gives order at
most $p$ for every $p>1$.
Write $\gamma$ for this explicitly defined growth topology.
It is enough to take $p=1+1/j$, $j\geq2$, and $a=1/k$, $k\geq1$:
given $p>1,a>0$, select one of these $p_j<p$ and use the boundedness
of $|s|^{p_j}/k-a|s|^p$. This proves equivalence of the countable
family with all displayed seminorms. A sequence Cauchy in these
seminorms converges uniformly on compact sets to an entire function.
Taking pointwise limits in each weighted Cauchy inequality gives
convergence in that seminorm and finiteness of its norm. Thus
$(\cqE,\gamma)$ is a complete metrizable locally convex space.

\begin{cqlemma}[Derivative bounds and the heat group]\label{thm:cq-heat}
Put $D=\cqD$. Differentiation and multiplication by $s$ map the
original underlying set $\cqE$ into itself. For $m\geq1$,
\begin{equation}\label{eq:cq-cauchy}
 \|D^mf\|_{p,a}\leq
 m!\left(\frac{epa}{m}\right)^{m/p}
 \|f\|_{p,a/2^{p-1}}.
\end{equation}
For every $t\in\cqC$, the series
\begin{equation}\label{eq:cq-U}
 U_tf=\sum_{k=0}^\infty\frac{(-t/4)^k}{k!}D^{2k}f
\end{equation}
converges in $\gamma$, locally uniformly in $t$; $U_t$ is a
continuous automorphism of $(\cqE,\gamma)$ and
$U_tU_u=U_{t+u}$, $U_t^{-1}=U_{-t}$.
\end{cqlemma}
\begin{proof}
Cauchy's circle formula on a circle of radius $R$ about $s$ gives
\[
 |D^mf(s)|\leq m!R^{-m}\max_{|z-s|=R}|f(z)|.
\]
For $b=a/2^{p-1}$ the inequality
$|s+z|^p\leq2^{p-1}(|s|^p+|z|^p)$ bounds the last expression,
after multiplication by $e^{-a|s|^p}$, by
$m!R^{-m}e^{aR^p}\|f\|_{p,b}$. Choosing
$R=(m/(ap))^{1/p}$ proves \eqref{eq:cq-cauchy}.
For multiplication use
\[
 \|Mf\|_{p,a}\leq\sup_{r\geq0}re^{-ar^p/2}\|f\|_{p,a/2}.
\]
The $k$th heat summand is at most
\[
 \frac{(|t|/4)^k}{k!}(2k)!
       \left(\frac{epa}{2k}\right)^{2k/p}\|f\|_{p,b}.
\]
Since $(2k)!/k!\leq(2k)^k$, its scalar factor is at most
$[C_{p,a}|t|k^{1-2/p}]^k$, whose $k$th root tends to zero.
This proves summability, continuity, and locally uniform convergence
of every time-differentiated series by the same estimate with the
additional polynomial factors in $k$. Commutation with $D$ follows
from continuity of $D$. In the double composition series group
terms with $n=k+\ell$. The sum of the absolute scalar coefficients
at $D^{2n}f$ is $( (|t|+|u|)/4)^n/n!$; the preceding estimate
justifies regrouping. The binomial theorem gives $U_{t+u}$.
The value at zero is the identity, proving the inverse.
\end{proof}

\begin{cqlemma}[Polynomial approximation]\label{thm:cq-density}
The Taylor polynomials of any $f\in\cqE$ converge to $f$ in $\gamma$.
Multiplication of two elements is continuous in $\gamma$.
\end{cqlemma}
\begin{proof}
If $f=\sum c_ns^n$, Cauchy's coefficient bound with parameter
$b<a$ gives, for $n\geq1$,
\[
 |c_n|\leq\|f\|_{p,b}(epb/n)^{n/p}.
\]
Also $\|s^n\|_{p,a}=(n/(epa))^{n/p}$. Therefore
$\|c_ns^n\|_{p,a}\leq\|f\|_{p,b}(b/a)^{n/p}$.
The tail is bounded by a convergent geometric series.
Finally $\|fg\|_{p,a}\leq\|f\|_{p,a/2}\|g\|_{p,a/2}$.
\end{proof}

Retain the actual real heat family, rather than substitute a polynomial:
\begin{align}\label{eq:cq-realheat}
 \Phi(u)&=\sum_{n\geq1}
 (2\pi^2n^4e^{9u}-3\pi n^2e^{5u})e^{-\pi n^2e^{4u}},\notag\\
 Z_t(s)&=8H_t(2s)=8\int_0^\infty e^{tu^2}\Phi(u)\cos(2su)\,du,
 &\partial_tZ_t=-\tfrac14D^2Z_t,
 &Z_0=\Xi.
\end{align}
These are the Rodgers--Tao conventions, Section 1 \cite{cqRT}.
For clarity the analytic exchanges in this transport can be proved
directly. Since $n^2e^{4u}\geq(n^2+e^{4u})/2$, the absolutely
convergent sums $\sum n^je^{-\pi n^2/2}$ give constants $C,c>0$
such that $|\Phi(u)|\leq Ce^{9u}e^{-ce^{4u}}$ for $u\geq0$.
For $|t|\leq T$ and $|s|\leq R$, every derivative of the integrand
is bounded by a polynomial in $u$ times
$C\exp(Tu^2+(2R+9)u-ce^{4u})$, which is integrable.
The exponential series in $t$ can therefore be integrated absolutely.
Because $D^{2k}\cos(2su)=(-1)^k(2u)^{2k}\cos(2su)$, it gives
exactly $Z_t=U_t\Xi$ with the coefficient $(-t/4)^k/k!$.
In particular $Z_t$ lies in the original order-at-most-one set for
every complex $t$, by Lemma~\ref{thm:cq-heat}; it is never identically
zero because $U_t$ is invertible and $\Xi(0)\ne0$.

\section{Transport of the actual closed arithmetic image}
\label{sec:cq-moving}
Define $\tau_t$ on the same underlying vector space $\cqE$ by
declaring $U_t:(\cqE,\tau)\longrightarrow(\cqE,\tau_t)$ a
homeomorphism. This definition specifies every open set:
$V$ is $\tau_t$-open exactly when $U_{-t}V$ is $\tau$-open.
It uses the globally defined bijection just proved and does not require
unproved continuity in $\tau$. Set
\begin{equation}\label{eq:cq-Qt}
 N_t=U_t\cqN
 =\overline{\operatorname{span}\{U_t(R_\ell^\pm\Xi)\}}^{\,\tau_t},
 \qquad Q_t=(\cqE,\tau_t)/N_t.
\end{equation}
The closure equality follows because homeomorphisms carry closure to
closure. In particular $N_t$ is closed. The quotient map
\begin{equation}\label{eq:cq-Phi}
 \mathcal U_t:\cqQ\longrightarrow Q_t,\qquad [f]\longmapsto[U_tf]
\end{equation}
is well defined, bijective, and continuous, with continuous inverse
$[g]\mapsto[U_{-t}g]$. Indeed changing $f$ by $n\in\cqN$
changes $U_tf$ by $U_tn\in N_t$, and the inverse gives necessity;
the quotient universal property proves both continuity statements.

\begin{cqtheorem}[Exact arithmetic heat intertwiner]\label{thm:cq-moving}
On the entire-function set,
\begin{equation}\label{eq:cq-B}
 B_t=U_tMU_{-t}=M-\tfrac t2D.
\end{equation}
This operator preserves $N_t$, is continuous in $\tau_t$, and induces
$S_t$ on $Q_t$. It satisfies
$S_t\mathcal U_t=\mathcal U_tS$ and
\begin{equation}\label{eq:cq-specmove}
 \cqSpec(S_t)=\cqSpec(S)
 =\{\lambda:\zeta(\tfrac12+i\lambda)=0\text{ nontrivially}\}.
\end{equation}
The time-dependent trace functions are exactly
$U_t(R_\ell^\pm\Xi)=R_\ell^\pm(B_t)Z_t$.
\end{cqtheorem}
\begin{proof}
The product rule gives $D^{2k}Mf=MD^{2k}f+2kD^{2k-1}f$.
Insert this identity in the convergent heat series to obtain
$U_tM=(M-(t/2)D)U_t$, proving \eqref{eq:cq-B}.
If $n\in\cqN$, then $B_tU_tn=U_tMn\in N_t$.
Continuity is conjugation of the continuous $M$ by the defining
homeomorphism. The induced identity follows on representatives.
If $T$ is a continuous inverse of $S-\lambda$, then
$\mathcal U_tT\mathcal U_t^{-1}$ is a continuous inverse of
$S_t-\lambda$. Conjugating in the opposite direction proves the
converse, and hence the spectrum equality using
\eqref{eq:cq-source-spectrum}. Apply the polynomial identity
$U_tp(M)=p(B_t)U_t$ to $\Xi$ for the last assertion.
\end{proof}

The heat map also transports the full multiplication, with every
cross term retained. For $f,g\in\cqE$ define
\begin{equation}\label{eq:cq-star}
 f\star_t g=U_t\bigl((U_{-t}f)(U_{-t}g)\bigr)
    =\sum_{n=0}^\infty\frac{(-t/2)^n}{n!}(D^nf)(D^ng).
\end{equation}
We prove the series identity, including convergence. For
$F_t=U_tf$, $G_t=U_tg$, the series
\[
 J_t=\sum_{n=0}^\infty\frac{(-t/2)^n}{n!}
                    (D^nF_t)(D^nG_t)
\]
converges in $\gamma$, uniformly on compact time sets: applying
\eqref{eq:cq-cauchy} with output norm $(p,a/2)$ to both factors
bounds its $n$th term by
\[
 K\left[\frac T2\left(\frac{epa}{2}\right)^{2/p}
                   n^{1-2/p}\right]^n\quad(|t|\leq T),
\]
where $K$ bounds the product of the two input norms
$(p,a/2^p)$, and is finite by the heat estimate. The root test
applies. Cauchy's formula in time permits termwise time derivatives;
continuity of $D$ permits spatial derivatives. Differentiating the
coefficient contributes
$-\tfrac12\sum(-t/2)^n(D^{n+1}F_t)(D^{n+1}G_t)/n!$.
Differentiating the factors contributes the two pure second-derivative
terms with coefficient $-1/4$. The complete product rule therefore
gives $\partial_tJ_t=-D^2J_t/4$ and $J_0=fg$.
Its entire time Taylor coefficients are uniquely
$(-1/4)^kD^{2k}(fg)/k!$, so $J_t=U_t(fg)$.
Replacing the inputs by $U_{-t}f,U_{-t}g$ proves
\eqref{eq:cq-star}. Conjugation of ordinary multiplication proves
associativity, commutativity, unit $1$, and continuity for $\tau_t$,
even when continuity of the series in the source topology is unknown.
The identity $s\star_t f=sf-tf'/2$ recovers exactly $B_t$.
No ideal property of $\cqN$ beyond its stated polynomial invariance
is required for this identity or for the quotient in
Theorem~\ref{thm:cq-moving}.

The transported actual trace image has no common point-evaluation
zero at any nonzero time. Suppose $\lambda$ were such a common
zero. It would annihilate each $U_t(s^n\Xi)$. For every $a\in\cqC$,
the Taylor series of $e^{as}$ converges in $\gamma$ by
Lemma~\ref{thm:cq-density}. Multiplication, $U_t$, and evaluation
are continuous in $\gamma$, so
\[
 0=\sum_{n=0}^\infty\frac{a^n}{n!}U_t(s^n\Xi)(\lambda)
   =e^{a\lambda-ta^2/4}Z_t(\lambda-ta/2).
\]
The last equality follows from the just-proved product formula
and the entire Taylor formula for $Z_t$. At $t\ne0$ varying $a$
forces $Z_t$ to vanish everywhere, a contradiction. This argument
uses only the polynomial generators contained in $N_t$; it assumes
no source-topology convergence of the exponential series. Thus the
source arithmetic spectrum in \eqref{eq:cq-specmove} is carried by
the explicitly transported operator and relation, even though the
transported relation has no common ordinary evaluation zero.

The first two nonconstant polynomial generators make the extra terms
visible without changing any polynomial coefficient:
\begin{align}\label{eq:cq-generators}
 U_t(s\Xi)&=sZ_t-\tfrac t2Z_t',\notag\\
 U_t(s^2\Xi)&=s^2Z_t-t sZ_t'-\tfrac t2Z_t+\tfrac{t^2}4Z_t''.
\end{align}
For $L=\partial_t+\tfrac14D^2$, the complete product rule is
\begin{equation}\label{eq:cq-unit}
 L(AP)=A\bigl(P_t+\tfrac14P_{ss}\bigr)
       +\bigl(A_t+\tfrac14A_{ss}\bigr)P
       +\tfrac12A_sP_s.
\end{equation}
Thus $L(R(s)Z_t)=\tfrac14R''Z_t+\tfrac12R'Z_t'$.
The family $R(B_t)Z_t$ solves the homogeneous heat equation, whereas
the displayed expression is the full defect of replacing it by
$R(s)Z_t$. For a local factorization $Z_t=AP$ with $A\ne0$,
the same formula retains the analytic unit $A$ and its derivatives.
At a simple zero $\lambda(t)$, differentiating
$Z_t(\lambda(t))=0$ gives
\begin{equation}\label{eq:cq-zero}
 \lambda'(t)=\frac14\frac{Z_{ss}}{Z_s}(t,\lambda(t))
 =\frac14\frac{P_{ss}}{P_s}(t,\lambda(t))
    +\frac12\frac{A_s}{A}(t,\lambda(t)).
\end{equation}
This statement is on the open locus $Z_s\ne0$, on which the analytic
implicit-function theorem supplies the local curve; it claims no
chosen zero is simple. At a multiple zero the division is undefined,
and the full equation \eqref{eq:cq-unit} remains the equation to use.

\section{Derivative domains, multivalued operators, and jet extensions}
\label{sec:cq-jets}
Every derivative has the full underlying domain $\cqE$ by
Lemma~\ref{thm:cq-heat}. Continuity in the unnamed source topology is
not assumed. The following gives exact source-quotient operators and
their domains without imposing it. For $k\geq1$ put
\[
 K_k=\{n\in\cqN:D^kn\in\cqN\},\qquad
 W_k=(D^k\cqN+\cqN)/\cqN\subseteq\cqQ.
\]
The derivative correspondence is the actual linear relation
\begin{equation}\label{eq:cq-relation}
 \mathfrak D_k=\{([f],[D^kf]):f\in\cqE\}
                    \subseteq\cqQ\times\cqQ.
\end{equation}
For fixed $[f]$, its output fibre is precisely $[D^kf]+W_k$:
all representatives are $f+n$ with $n\in\cqN$, and linearity gives
all and only the stated outputs. The kernel of its presenting map
$f\mapsto([f],[D^kf])$ is exactly $K_k$. The relation has domain
all of $\cqQ$; its multivalued part at zero is exactly $W_k$.
Thus it is the graph of an operator exactly when $D^k\cqN\subseteq
\cqN$. This is a calculated equivalence, with an explicit obstruction
space, rather than an assumption of invariance.

There is no nonempty domain of quotient classes on which the same
representative formula is unambiguous if $W_k\ne0$ and every
representative of a class is admitted: addition of the same offending
$n$ gives two different outputs for every class. A chosen linear
lift is a different map. More precisely any vector subspace
$A\subseteq\cqE$ supplies a single-valued map on its quotient image
$q(A)$ exactly when $D^k(A\cap\cqN)\subseteq\cqN$; the identical
difference-of-representatives argument proves both directions.
The first heat generator is $-D^2/4$, so its ambiguity space is
$-W_2/4$, with the coefficient and sign unchanged.

\subsection{Exact algebraic spectral defects of the source relation}
\label{sec:cq-sourcealgebra}
These calculations keep the actual $\cqN$ and require only the
already proved $M\cqN\subseteq\cqN$. For every $\lambda\in\cqC$,
division by $x=s-\lambda$ gives the algebraic identities
\begin{equation}\label{eq:cq-source-kercoker}
 \begin{aligned}
 \operatorname{coker}_{\rm alg}(S-\lambda)
     &=\cqE/(\cqN+x\cqE),\\
 \ker(S-\lambda)&\simeq
        (\cqN\cap x\cqE)/(x\cqN),
        & [n]&\longmapsto[n/x].
 \end{aligned}
\end{equation}
Division belongs to the original order class: if $f(\lambda)=0$,
the filled quotient is entire, is bounded on a fixed disk, and
outside $|s|\geq2|\lambda|+1$ its modulus is at most $2|f(s)|$.
Thus $x\cqE$ is exactly the kernel of evaluation at $\lambda$.
The cokernel identity follows from the image of the induced
multiplication. For the second identity, changing $n$ by $xn_0$
changes $n/x$ by $n_0\in\cqN$; its image is in the stated
kernel. Conversely any kernel representative $f$ has
$xf\in\cqN\cap x\cqE$. Finally $[n/x]=0$ exactly when
$n\in x\cqN$. This proves every asserted kernel and image.

The cokernel has only two possibilities. When every $n\in\cqN$
vanishes at $\lambda$, it is one-dimensional, with its exact
identification given by evaluation. Otherwise retain any
$h_\lambda\in\cqN$ with $h_\lambda(\lambda)\ne0$. Then
\begin{equation}\label{eq:cq-source-division}
 \begin{aligned}
 A_\lambda f(s)&=
  \frac{f(s)-f(\lambda)h_\lambda(s)/h_\lambda(\lambda)}{s-\lambda},\\
 f&=\frac{f(\lambda)}{h_\lambda(\lambda)}h_\lambda
                                      +(s-\lambda)A_\lambda f
 \end{aligned}
\end{equation}
proves surjectivity of $S-\lambda$, and its exact remaining kernel is
\begin{equation}\label{eq:cq-source-splitting}
 \ker(S-\lambda)\simeq
 \cqN/\bigl((s-\lambda)\cqN\oplus\cqC h_\lambda\bigr),
 \qquad n\longmapsto[A_\lambda n].
\end{equation}
The map is onto because a kernel class $[g]$ is the image of
$n=(s-\lambda)g\in\cqN$. Its kernel consists exactly of
$n=(s-\lambda)n_0+c h_\lambda$, $n_0\in\cqN$, by
\eqref{eq:cq-source-division}. The sum is direct by evaluation.
In this second case algebraic invertibility is therefore equivalent
to the exact identity
$\cqN=(s-\lambda)\cqN\oplus\cqC h_\lambda$.
On that locus the representative formula $[f]\mapsto[A_\lambda f]$
is independent of representatives, since $A_\lambda\cqN\subseteq
\cqN$; multiplication in \eqref{eq:cq-source-division} proves the
right inverse, and $A_\lambda((s-\lambda)f)=f$ proves the left
inverse. These are algebraic statements; continuity of this inverse
in $\tau$ has not been inserted.

The first derivative and the heat generator have exactly the same
frozen invariance question for the source subspace:
\begin{equation}\label{eq:cq-source-D-equivalence}
 D\cqN\subseteq\cqN
 \quad\Longleftrightarrow\quad D^2\cqN\subseteq\cqN
 \quad\Longleftrightarrow\quad -\tfrac14D^2\cqN\subseteq\cqN.
\end{equation}
Indeed the product rule is
$D^2(sn)-sD^2n=2Dn$. Since $sn\in\cqN$, second-derivative
invariance puts both terms on the left in $\cqN$, proving
first-derivative invariance. Iteration proves the reverse implication;
the retained nonzero scalar $-1/4$ proves the last equivalence.
If these equivalent inclusions hold, all $D^j\Xi$ belong to
$\cqN$. At each $\lambda$ at least one is nonzero: otherwise
the Taylor series would make $\Xi$ identically zero. Thus every
$S-\lambda$ is then algebraically surjective by
\eqref{eq:cq-source-division}. Its prescribed topological spectral
points would be represented entirely by the kernel
\eqref{eq:cq-source-splitting} or a noncontinuous algebraic inverse.
This computes the precise alternatives without assigning either
one to the source's stated spectral set.

There is also a full multiplicity consequence of algebraic derivative
descent. Writing $d[f]=[Df]$ on its exact descent locus, the product
rule gives $[d,S]=I$. For $0\ne v\in\ker(S-\lambda)$,
\begin{equation}\label{eq:cq-source-weylladder}
 (S-\lambda)d^kv=-k d^{k-1}v,
 \qquad (S-\lambda)^kd^kv=(-1)^k k!v\quad(k\geq0),
\end{equation}
with zero right side in the first identity at $k=0$.
To prove induction, commute the next $d$ using
$(S-\lambda)d=d(S-\lambda)-I$; the coefficient $-k$ becomes
$-(k+1)$. Iteration proves the second identity. A finite linear
dependence among $v,dv,\ldots,d^nv$, with largest nonzero
coefficient at index $j$, becomes that coefficient times
$(-1)^j j!v$ after applying $(S-\lambda)^j$, a contradiction.
Thus every such eigenvector generates an infinite independent
generalized eigenvector sequence. No finite source multiplicity
assertion has been assumed.

Finally even a single forward nonzero heat inclusion has a calculated
consequence. Fix $t\ne0$ and $c=t/2$, and define
\begin{equation}\label{eq:cq-source-heatcommutator}
 \begin{gathered}
 A_0=U_t,\qquad A_{k+1}=MA_k-A_kM,\qquad A_k=P_k(D)U_t,\\
 P_0(X)=1,\qquad P_{k+1}(X)=cX P_k(X)-P_k'(X).
 \end{gathered}
\end{equation}
The identity follows by induction from $[M,U_t]=cDU_t$ and
$[M,D^j]=-jD^{j-1}$: commuting $M$ through $P_k(D)U_t$
contributes $-P_k'(D)U_t+cP_k(D)DU_t$.
The degree of $P_k$ is exactly $k$, with leading coefficient
$c^k\ne0$. Under $U_t\cqN\subseteq\cqN$, induction with
$M\cqN\subseteq\cqN$ gives $A_k\cqN\subseteq\cqN$;
subtracting the lower-degree terms of $P_k$ then gives exactly
\begin{equation}\label{eq:cq-source-heatconsequence}
 D^kU_t\cqN\subseteq\cqN\quad(k\geq0).
\end{equation}
In particular every derivative of the nonzero $Z_t=U_t\Xi$ lies
in $\cqN$. At every point one derivative is nonzero by the entire
Taylor theorem, so every $S-\lambda$ is algebraically surjective
in this case too. This proof uses only the one forward inclusion.
For the stronger equality $U_t\cqN=\cqN$, it also proves
$D\cqN\subseteq\cqN$, by writing each $n\in\cqN$ as
$U_tm$ in \eqref{eq:cq-source-heatconsequence}. None of these
exact implications presumes the inclusion or equality for the actual
source closure. The explicit obstruction spaces above remain the
calculated source question.


Even when $W_k$ is not closed, the smallest source-closed enlargement
of the target relation subspace is explicit:
\begin{equation}\label{eq:cq-target}
 [f]\longmapsto[D^kf]:
 \cqQ\longrightarrow
 \cqE/\overline{\cqN+D^k\cqN}^{\,\tau}.
\end{equation}
It is well defined algebraically. Equip its source with the quotient
of the graph topology $f\mapsto(f,D^kf)$ in
$(\cqE,\tau)^2$ to obtain a continuous map. This specifies a
topology change explicitly, without claiming it is the original one.
The lost part of the original target $\cqQ$ is exactly
$\overline{\cqN+D^k\cqN}^{\,\tau}/\cqN$.

Here is an extension which retains correlated derivatives rather than
quotienting them away. For $j\geq1$ define the bijection
\[
 T_j(f_0,\ldots,f_j)
   =(f_0,f_1-Df_0,\ldots,f_j-D^jf_0).
\]
Its inverse is $(f_0,g_1,\ldots,g_j)\mapsto
(f_0,g_1+Df_0,\ldots,g_j+D^jf_0)$.
Give $V_j=\cqE^{j+1}$ the topology transported by $T_j$ from
$(\cqE,\tau)\times(\cqE,\gamma)^j$. The residual factors
are explicitly given the proved Hausdorff growth topology; the
source factor keeps its original topology. For $J_jf=(f,Df,\ldots,D^jf)$,
$T_j(J_j\cqN)=\cqN\oplus0^j$, so $J_j\cqN$ is closed.
Define
\begin{equation}\label{eq:cq-jetextension}
 X_j=V_j/J_j\cqN.
\end{equation}
Then the map
\[
 [f_0,\ldots,f_j]\longmapsto
 ([f_0],f_1-Df_0,\ldots,f_j-D^jf_0)
\]
is a homeomorphism $X_j\simeq\cqQ\oplus(\cqE,\gamma)^j$; independence,
the inverse, and continuity follow directly from $T_j$ and the
quotient definition. In particular there is the split exact sequence
\begin{equation}\label{eq:cq-jetseq}
 0\longrightarrow\cqE^j\longrightarrow X_j
 \xrightarrow{[f_0,\ldots,f_j]\mapsto[f_0]}\cqQ
 \longrightarrow0,
 \qquad [f]\longmapsto[J_jf].
\end{equation}
Its relation on a source generator retains every product term:
\begin{equation}\label{eq:cq-rawjet}
 D^k(R\Xi)=\sum_{h=0}^k\binom kh R^{(h)}\Xi^{(k-h)}
                    \quad(0\leq k\leq j).
\end{equation}
The binomial formula follows by induction from the ordinary product
rule, including the two edge terms at each step.

The original spectral coordinate acts on $V_j$ by
\[
 (\mathcal M_j f)_k=s f_k+k f_{k-1}\quad(0\leq k\leq j),
 \qquad f_{-1}=0.
\]
The identity $D^k(sf)=sD^kf+kD^{k-1}f$ proves
$\mathcal M_jJ_jf=J_j(sf)$; it therefore preserves the exact
relation subspace. In the $T_j$ coordinates it acts by $S$ on
$\cqQ$ and by $g_k\mapsto sg_k+kg_{k-1}$ on the residual terms
($g_0=0$). This also proves continuity in the stated topology.
Notice that \eqref{eq:cq-jetseq} adds a free residual module:
the spectrum of the entire extension is not claimed to equal the
source spectrum. The arithmetic quotient and its split inclusion
are explicitly preserved.

\section{A fully computed analytic realization and its exact comparison}
\label{sec:cq-analytic}
We now compute the closures and the failure of descent in specified
topologies. This does not identify them with $\tau$.
Let $\kappa$ be the topology of uniform convergence on compact sets
on the same set $\cqE$, and define
\[
 N_\gamma=\overline{\cqC[s]\Xi}^{\,\gamma},\qquad
 N_\kappa=\overline{\cqC[s]\Xi}^{\,\kappa}.
\]
Evaluations of all derivatives are continuous in both topologies by
Cauchy's formula. The growth topology is finer than $\kappa$, so
$N_\gamma\subseteq N_\kappa$.
If $\lambda$ is a zero of $\Xi$ with its actual multiplicity $m_\lambda$,
every element of either closed subspace has vanishing derivatives
of orders $0,\ldots,m_\lambda-1$ at $\lambda$.

\begin{cqlemma}[Entire division with its growth and topology retained]
\label{thm:cq-division}
Let $g\in\cqE$ with $g(0)\ne0$. If $f\in\cqE$ and $h=f/g$
extends to an entire function, then $h\in\cqE$. For $1<p<2$,
$a>0$, put $b=a/(6\cdot2^p)$. There is the linear division estimate
\begin{equation}\label{eq:cq-division-bound}
 \|f/g\|_{p,a}\leq
 \left(\frac{\max(1,\|g\|_{p,b})}{|g(0)|}\right)^3
 \|f\|_{p,b}.
\end{equation}
Multiplication by $g$ is a topological linear isomorphism from
$(\cqE,\gamma)$ onto the closed subspace of functions divisible by
$g$, with its subspace topology. Compact-topology division is also
continuous.
\end{cqlemma}
\begin{proof}
Choose $R>0$ with no zero of $g$ on $|s|=R$. Factor its finitely
many zeros in the disk, retaining multiplicities. The residual is
holomorphic and zero-free there, so its log modulus is harmonic.
For $|z|<R$ the boundary mean of $\log|Re^{i\theta}-z|$ equals
$\log R$: expand $\log(1-(z/R)e^{-i\theta})$ in its uniformly
convergent series and take real parts. The mean-value formula for
the residual gives the exact Jensen identity
\[
 \frac1{2\pi}\int_0^{2\pi}\log|g(Re^{i\theta})|\,d\theta
 =\log|g(0)|+\sum_{|z|<R}m_z\log(R/|z|)\geq\log|g(0)|.
\]
There is no zero at zero. With
$\log^+x=\max(0,\log x)$ and $\log^-x=\max(0,-\log x)$,
the inequality $\log^+|h|\leq\log^+|f|+\log^-|g|$ gives
\begin{equation}\label{eq:cq-division-mean}
 \frac1{2\pi}\int_0^{2\pi}\log^+|h(Re^{i\theta})|\,d\theta
 \leq\log^+M_f(R)+\log^+M_g(R)-\log|g(0)|.
\end{equation}
The right side is nonnegative since $M_g(R)\geq|g(0)|$.
The harmonic Poisson extension of the continuous nonnegative boundary
function $\log^+|h|$ majorizes $\log|h|$ by the maximum principle
and is nonnegative; hence it majorizes $\log^+|h|$ in the disk.
For $|s|\leq r<R$ the Poisson kernel is at most $(R+r)/(R-r)$.
Take $R\downarrow2r$ through circles avoiding zeros. Continuity of
the two maximum functions gives, for every $r>0$,
\begin{equation}\label{eq:cq-division-radial}
 \log^+M_h(r)\leq
 3\bigl(\log^+M_f(2r)+\log^+M_g(2r)-\log|g(0)|\bigr).
\end{equation}
No pointwise lower bound for $g$ on a large circle has been assumed.
For $f\ne0$, put $F=\|f\|_{p,b}>0$ and apply this to $f/F,h/F$.
Then $\log^+M_{f/F}(2r)\leq b(2r)^p$, and
$\log^+M_g(2r)\leq\log\max(1,\|g\|_{p,b})+b(2r)^p$.
The coefficient of $r^p$ is $6b2^p=a$, giving
\eqref{eq:cq-division-bound}; continuity fills in $r=0$.
The zero numerator is immediate. All $p,a$ are allowed, so $h$
belongs to the original order class and division is continuous.
Multiplication is continuous by Lemma~\ref{thm:cq-density}.
The requisite vanishing jets at every zero are closed conditions;
local Taylor division identifies them with entire divisibility.
The estimate proves this closed set is exactly $g\cqE$.
For the compact topology choose $R>r$ with no boundary zero and
$\delta=\min_{|s|=R}|g(s)|>0$. The maximum principle for the filled
entire quotient gives $M_h(r)\leq\delta^{-1}M_f(R)$, proving
compact continuity too.
\end{proof}

For general nonzero $g\in\cqE$, choose and retain a point $c$ with
$g(c)\ne0$ and use disks centered at $c$. The functions and spectral
coordinate remain unchanged. The bound
$|s+c|^p\leq2^{p-1}(|s|^p+|c|^p)$ proves continuity of translations
by $c$ and $-c$ in $\gamma$; compact containment proves it in
$\kappa$. Thus all conclusions hold with the stated translated
estimates. The original center $c=0$ already applies to $\Xi$.

\begin{cqproposition}[The computed closures are the same subspace]
\label{thm:cq-equalclosures}
With their definitions unchanged,
\begin{equation}\label{eq:cq-equalclosures}
 \begin{aligned}
 N_\gamma=N_\kappa&=\Xi\cqE\\
 &=\{f\in\cqE:D^hf(\lambda)=0\text{ for every }\Xi(\lambda)=0,
                                      \ 0\leq h<m_\lambda\}.
 \end{aligned}
\end{equation}
This is equality of subsets of the original ring, not an
identification of $\gamma$, $\kappa$, or $\tau$.
\end{cqproposition}
\begin{proof}
Jet continuity puts both closures inside the displayed vanishing
set and gives $N_\gamma\subseteq N_\kappa$. Local Taylor division
makes $h=f/\Xi$ entire for any $f$ in that set; the division lemma
then puts $h$ in $\cqE$. Its Taylor polynomials $p_n$ converge in
$\gamma$, and multiplication gives $p_n\Xi\to f$ there. Thus the
vanishing set is contained in $N_\gamma$, proving all equalities.
The algebraic-ideal equality is the conclusion of this argument;
it was not substituted for an unevaluated closure.
\end{proof}

\begin{cqtheorem}[Compact closure, multiplicities, and coordinate spectrum]
\label{thm:cq-compact}
The compact closure is exactly
\begin{equation}\label{eq:cq-compactideal}
 \begin{aligned}
 N_\kappa
 &=\{f\in\cqE:f/\Xi\text{ extends to an entire function}\}\\
 &=\{f\in\cqE:D^hf(\lambda)=0\ (h<m_\lambda),\ \Xi(\lambda)=0\}.
 \end{aligned}
\end{equation}
On $Q_\kappa=(\cqE,\kappa)/N_\kappa$, multiplication by $s$ has
spectrum exactly the zeros of $\Xi$ in the original $s$ coordinate.
For a zero $\lambda$ of multiplicity $m$, there is a continuous
surjection
\[
 Q_\kappa\longrightarrow\cqC[\epsilon]/(\epsilon^m),\qquad
 [f]\longmapsto\sum_{h=0}^{m-1}\frac{D^hf(\lambda)}{h!}\epsilon^h,
\]
intertwining multiplication by $s$ with $\lambda+\epsilon$.
\end{cqtheorem}
\begin{proof}
The vanishing conditions are closed and hold on every $p\Xi$,
so the closure is contained in the last displayed set. Local Taylor
division at each isolated zero identifies that set with entire
divisibility by $\Xi$. For such an $f$, put $g=f/\Xi\in\cqO(\cqC)$.
Its Taylor polynomials $p_n$ converge uniformly on every compact;
$p_n\Xi\to f$ in $\kappa$. Every $p_n\Xi$ belongs to the original
order class even though no growth assertion about $g$ is needed.
This proves \eqref{eq:cq-compactideal} with the closure retained.
The local-jet map is independent of representatives; the polynomials
$1,(s-\lambda),\ldots,(s-\lambda)^{m-1}$ prove surjectivity.
Leibniz's rule gives the stated action, and the finite jet evaluations
are continuous.

If $\Xi(\lambda)\ne0$, define
\[
 T_\lambda f(s)=
 \frac{f(s)-f(\lambda)\Xi(s)/\Xi(\lambda)}{s-\lambda},
\]
filling in the removable value. Division by a linear factor preserves
the order bound: outside a disk about $\lambda$ the denominator is
bounded below, and inside a disk the filled entire function is bounded.
It is continuous in $\kappa$ by Cauchy's formula on a larger disk
applied to the numerator. For $n\in N_\kappa$ the quotient
$T_\lambda n/\Xi=(n/\Xi-(n/\Xi)(\lambda))/(s-\lambda)$ is entire;
thus it descends. Direct multiplication proves
$(S_\kappa-\lambda)[T_\lambda f]=[f]$.
Applying $T_\lambda$ to $(s-\lambda)f$ gives $f$ exactly, proving
the other inverse identity. If $\Xi(\lambda)=0$, evaluation at
$\lambda$ descends to a nonzero continuous functional which annihilates
the image of $S_\kappa-\lambda$. Hence that operator is not
surjective. This proves the spectrum claim.
\end{proof}

The same coordinate spectrum and the same finite-jet surjections hold
for $Q_\gamma=(\cqE,\gamma)/\Xi\cqE$. Indeed the numerator in
$T_\lambda f$ depends continuously on $f$ in $\gamma$. Division
by $s-\lambda$, on its entire-divisibility domain, is continuous by
Lemma~\ref{thm:cq-division} with a nonzero center retained when
needed. The already proved two inverse identities therefore give a
continuous resolvent in $\gamma$ as well. The evaluation obstruction
at each zero is continuous in that topology too. This proves the
claim, without transferring a resolvent through an unproved topology
identification.

The ambient topologies $\kappa$ and $\gamma$ are in fact distinct.
The original functions $f_n(s)=\exp(ns-n^2)$ tend to zero uniformly
on each compact set. On the positive real axis the maximum of
$nx-x^{3/2}$ is $4n^3/27$, attained at $x=4n^2/9$. Since
$\operatorname{Re}s\leq|s|$, that axis attains the full weighted
supremum, giving the exact value
\[
 \|f_n\|_{3/2,1}=\exp(4n^3/27-n^2)\longrightarrow\infty.
\]
Thus equality of the closed arithmetic relation subspaces and of
their spectral sets does not identify the ambient topologies. The
identity induces a continuous linear bijection $Q_\gamma\to Q_\kappa$;
the next theorem proves that its inverse is not continuous.


\subsection{The actual difference between the two analytic quotient topologies}
\label{sec:cq-quotient-topologies}
Retain the original entire-function set $\cqE=\mathcal H_{\leq1}$,
the original coordinate $s$, and
$\Xi(s)=\xi(\tfrac12+is)$ from \eqref{eq:cq-Xi}.
Let $\Lambda$ be the set of its distinct zeros and let
$m_\lambda=\operatorname{ord}_\lambda\Xi\geq1$ be each actual
multiplicity. Put
\[
 I=\Xi\cqE,\qquad
 Q_\kappa=(\cqE,\kappa)/I,\qquad
 Q_\gamma=(\cqE,\gamma)/I,\qquad
 \mathcal P_\Xi=\prod_{\lambda\in\Lambda}\cqC^{m_\lambda}.
\]
The product has its ordinary product topology. Coordinates in
$\cqC^{m_\lambda}$ are the actual derivatives of orders
$0,\ldots,m_\lambda-1$, with no factorial rescaling. Define
\begin{equation}\label{eq:cq-total-jet-map}
 J_\Xi f=(f^{(j)}(\lambda))_{
          \lambda\in\Lambda,\ 0\leq j<m_\lambda},
 \qquad \mathcal V_\Xi=J_\Xi(\cqE)\subset\mathcal P_\Xi.
\end{equation}
Every point of $\mathcal V_\Xi$ retains the entire input's arithmetic
jets. The restriction to this image is essential.

\begin{cqtheorem}[Jet topology and its completion]
\label{thm:cq-jet-topology}
The map induced by \eqref{eq:cq-total-jet-map} is a topological
linear isomorphism
\[
 \overline J_\Xi:Q_\kappa\xrightarrow{\ \cong\ }
 \mathcal V_\Xi
\]
onto the subspace topology from $\mathcal P_\Xi$.
The image is dense and proper. Consequently $Q_\kappa$ is
incomplete, and its Hausdorff completion is canonically
$\mathcal P_\Xi$ via this map. In contrast $Q_\gamma$ is a
Fr\'echet space and is complete. The identity on the same
underlying quotient vector space induces a continuous linear
bijection
\begin{equation}\label{eq:cq-topology-comparison-strict}
 \iota:Q_\gamma\longrightarrow Q_\kappa
\end{equation}
whose inverse is not continuous.
\end{cqtheorem}

\begin{proof}
We first construct each finite interpolation map with its constants.
For a finite nonempty set $F\subset\Lambda$, write
$d_F=\sum_{\lambda\in F}m_\lambda$ and set
\[
 B_{\lambda,F}(s)=
 \prod_{\substack{\nu\in F\\\nu\ne\lambda}}
 (s-\nu)^{m_\nu}.
\]
The empty product is $1$. Since
$B_{\lambda,F}(\lambda)\ne0$, the germ
$(s-\lambda)^j/(j!B_{\lambda,F}(s))$ is holomorphic at
$\lambda$. Denote its Taylor polynomial through degree
$m_\lambda-1$, in powers of the unchanged difference
$s-\lambda$, by $T_{\lambda,m_\lambda-1}$ of that germ. Put
\begin{equation}\label{eq:cq-Hermite-basis-unscaled}
 L_{\lambda,j,F}(s)=B_{\lambda,F}(s)
 T_{\lambda,m_\lambda-1}
 \left(\frac{(s-\lambda)^j}{j!B_{\lambda,F}(s)}\right),
 \qquad 0\leq j<m_\lambda.
\end{equation}
This is a polynomial of degree at most $d_F-1$.
For $\nu\in F\setminus\{\lambda\}$ it vanishes to order
at least $m_\nu$ at $\nu$. At $\lambda$ its difference from
$(s-\lambda)^j/j!$ vanishes to order at least $m_\lambda$.
Therefore, for all $\nu\in F$ and $0\leq k<m_\nu$,
\[
 L_{\lambda,j,F}^{(k)}(\nu)
   =\begin{cases}1,&(\nu,k)=(\lambda,j),\\0,&\text{otherwise}.
     \end{cases}
\]
It follows that the linear map to polynomials
\begin{equation}\label{eq:cq-finite-Hermite-lift}
 H_F(v)(s)=\sum_{\lambda\in F}\sum_{j=0}^{m_\lambda-1}
              v_{\lambda,j}L_{\lambda,j,F}(s)
\end{equation}
realizes every prescribed finite array of actual derivatives
$v=(v_{\lambda,j})$. For each $R>0$ define
\[
 M_f(R)=\max_{|s|\leq R}|f(s)|,\qquad
 C_{R,F}=\sum_{\lambda\in F}\sum_{j=0}^{m_\lambda-1}
                  M_{L_{\lambda,j,F}}(R).
\]
The triangle inequality gives the exact finite bound
\begin{equation}\label{eq:cq-Hermite-lift-bound}
 M_{H_F(v)}(R)\leq C_{R,F}
             \max_{\lambda\in F,\,j<m_\lambda}|v_{\lambda,j}|.
\end{equation}
When $F$ is empty set $H_F=0$, $C_{R,F}=0$ and the maximum
over that empty set to $0$.

The kernel of $J_\Xi$ is precisely $I$. Indeed $J_\Xi f=0$
means that $f/\Xi$ has a removable singularity at each zero of
$\Xi$ by its local Taylor expansions. Away from those zeros it
is already holomorphic. Thus it extends to an entire function,
which belongs to the original order class $\cqE$ by
Lemma~\ref{thm:cq-division}; consequently $f\in\Xi\cqE$.
The converse follows by the product rule at each zero. This proves
that $\overline J_\Xi$ is a well-defined linear bijection onto
$\mathcal V_\Xi$, with inverse sending the jet array of any
$f\in\cqE$ to its class $[f]$. The kernel calculation proves
that this inverse is independent of the choice of such $f$.

The quotient topology on $Q_\kappa$ is generated by
\begin{equation}\label{eq:cq-compact-quotient-seminorm}
 q_R([f])=\inf_{h\in\cqE}M_{f-\Xi h}(R),\qquad R>0.
\end{equation}
For completeness, these are seminorms because they are infima of
the representatives' seminorms modulo a linear subspace. The
identity
$\{q_R<\epsilon\}=\pi_\kappa(\{M_f(R)<\epsilon\})$
holds for every $\epsilon>0$, and the quotient map
$\pi_\kappa$ is open since the inverse image of the image of
an open set $U$ is the union $U+I$ of its translates.
The increasing disk seminorms generate the compact-open topology;
the displayed identity therefore proves the assertion about the
quotient topology. Integer $R\geq1$ suffice.

Fix $R>0$. Choose $R'>R$ such that the circle $|s|=R'$
contains no zero of $\Xi$, and put
$F=\{\lambda\in\Lambda:|\lambda|<R'\}$, a finite set.
For $f\in\cqE$ form
$H=H_F((f^{(j)}(\lambda))_{\lambda\in F,j<m_\lambda})$.
The function $(f-H)/\Xi$ is holomorphic throughout $|s|<R'$
by the exact finite interpolation just proved. Its Taylor
polynomials $p_n(s)$ at $s=0$ consequently converge uniformly on
$|s|\leq R$ to $(f-H)/\Xi$. In particular $p_n\in\cqE$,
and
\[
 M_{f-\Xi p_n}(R)
 \leq M_H(R)+M_\Xi(R)
        M_{(f-H)/\Xi-p_n}(R).
\]
Taking the infimum over admissible representatives and then the
limit gives
\begin{equation}\label{eq:cq-quotient-from-finite-jets}
 q_R([f])\leq M_H(R)
 \leq C_{R,F}
       \max_{\lambda\in F,\,j<m_\lambda}|f^{(j)}(\lambda)|.
\end{equation}
No entire extension or growth bound for $(f-H)/\Xi$ outside this
disk was used. In particular there is no additional division
assumption hidden in the approximation. If $F$ is empty, this
proof gives $q_R=0$, which is compatible with the other quotient
seminorms separating points.

Conversely, fix a finite set $F\subset\Lambda$, and choose
$R>\max_{\lambda\in F}|\lambda|$ and
$0<\delta<R-\max_{\lambda\in F}|\lambda|$.
Cauchy's formula on the radius-$\delta$ circle about each
$\lambda$ gives, for every $h\in\cqE$,
\[
 |f^{(j)}(\lambda)|
 =|(f-\Xi h)^{(j)}(\lambda)|
 \leq j!\delta^{-j}M_{f-\Xi h}(R),
 \qquad\lambda\in F,\quad j<m_\lambda.
\]
Here $0!=1$, so this bound includes the zeroth derivative without
altering its value. Taking infima gives
\begin{equation}\label{eq:cq-finite-jets-from-quotient}
 \max_{\lambda\in F,\,j<m_\lambda}|f^{(j)}(\lambda)|
 \leq\left(\max_{\lambda\in F,\,j<m_\lambda}
                        j!\delta^{-j}\right)q_R([f]).
\end{equation}
Equations \eqref{eq:cq-quotient-from-finite-jets} and
\eqref{eq:cq-finite-jets-from-quotient} prove continuity in both
directions, with the exact stated topologies.

We next verify that the zero set needed for incompleteness is
indeed infinite and has an enumeration escaping every compact
set. The source Hadamard product for the original even $\Xi$,
retaining its nonzero constant $\Xi(0)$, implies that finitely
many zeros would make $\Xi$ a polynomial. This is impossible:
at the original coordinate $s=-i(2n-\tfrac12)$, $n\geq1$,
the defining completion gives exactly
\begin{equation}\label{eq:cq-Xi-nonpolynomial-values}
 \Xi\bigl(-i(2n-\tfrac12)\bigr)
 =\xi(2n)=n(2n-1)(n-1)!\,\pi^{-n}\zeta(2n).
\end{equation}
For integer $n\geq2$, $\zeta(2n)>1$ by its positive Dirichlet
series, and the last $\lfloor n/2\rfloor$ factors of $(n-1)!$
are at least $n/2$. Thus the logarithm of the right hand side is
at least
\[
 \log(n(2n-1))+\lfloor n/2\rfloor\log(n/2)-n\log\pi,
\]
which exceeds $d\log(2n)$ for every fixed $d$ and all
sufficiently large $n$. Polynomial values of degree $d$ on these
points are bounded by a constant times $(2n)^d$, a contradiction.
Since $\Xi$ is nonzero and entire, its zeros are isolated, each
multiplicity is finite and each closed disk contains finitely
many distinct zeros. We may therefore enumerate all distinct
zeros as $\Lambda=\{\lambda_1,\lambda_2,\ldots\}$ with
$|\lambda_n|\to\infty$, retaining each $m_{\lambda_n}$.
This argument uses neither reality nor simplicity of any zero.

Every projection of $\mathcal V_\Xi$ onto finitely many full
zero-jet blocks is onto, by \eqref{eq:cq-finite-Hermite-lift}.
Such finite coordinate cylinders form a base of the product
topology. Consequently $\mathcal V_\Xi$ is dense in
$\mathcal P_\Xi$. It is proper: define the particular array
$v\in\mathcal P_\Xi$ by
\begin{equation}\label{eq:cq-nonrealizable-jet-array}
 v_{\lambda_n,0}=\exp(|\lambda_n|^2),\qquad
 v_{\lambda_n,j}=0\quad(1\leq j<m_{\lambda_n}).
\end{equation}
For each $f\in\cqE$ there is a finite constant $A_f>0$ with
$M_f(r)\leq A_f\exp(r^{3/2})$ for every $r\geq0$; this is
exactly the order-at-most-one bound, extended over a bounded
interval by increasing $A_f$. Were $J_\Xi f=v$, we would have
\[
 \exp(|\lambda_n|^2)
 \leq A_f\exp(|\lambda_n|^{3/2})\quad\text{for all }n,
\]
which contradicts $|\lambda_n|\to\infty$.

The incompleteness can be exhibited by an actual sequence of
polynomial classes. Set $F_n=\{\lambda_1,\ldots,\lambda_n\}$,
$H_n=H_{F_n}(v|_{F_n})$, and $x_n=[H_n]\in Q_\kappa$.
For each fixed zero-jet block, $J_\Xi H_n$ agrees with $v$
from some index onward; hence $J_\Xi H_n\to v$ in the product.
Equivalently, for each fixed $R$, the finite set in
\eqref{eq:cq-quotient-from-finite-jets} lies in every sufficiently
large $F_n$, so that inequality gives
\[
 q_R(x_n-x_k)=0\qquad(n,k\text{ sufficiently large}).
\]
Thus $(x_n)$ is Cauchy in $Q_\kappa$. If it converged to a
class $[f]$, continuity of $\overline J_\Xi$ and uniqueness of
limits in the product would give $J_\Xi f=v$, already disproved.
For example the product is metrized by the complete
translation-invariant metric
\[
 d(v,w)=\sum_{n=1}^\infty2^{-n}
       \min\!\left(1,\max_{0\leq j<m_{\lambda_n}}
                          |v_{\lambda_n,j}-w_{\lambda_n,j}|\right).
\]
A Cauchy sequence for this metric is Cauchy in each finite
coordinate block, so has a coordinatewise limit; controlling a
finite initial sum and the geometric tail proves convergence in
the metric. This proves product completeness directly. The
isomorphism $\overline J_\Xi$, being linear and continuous
with a continuous linear inverse, preserves the additive
uniformities. Its dense embedding in this complete product
therefore gives the claimed Hausdorff completion.

Finally, we prove the completeness assertion for the stronger
quotient explicitly. The growth seminorms are
\[
 \|f\|_{p,a}=\sup_{s\in\cqC}|f(s)|e^{-a|s|^p},
 \qquad 1<p<2,\quad a>0.
\]
The countable family $p=1+1/j$, $a=1/k$ for integers $j\geq2$
and $k\geq1$ defines the same topology: for arbitrary $1<p<2,a>0$
choose $p_j<p$, and the finite maximum of
$r^{p_j}-ar^p$ controls $\|f\|_{p,a}$ by a constant times
$\|f\|_{p_j,1}$. Each countable seminorm belongs to the original
family, proving both comparisons. Enumerate this countable family
as $(u_\ell)_{\ell\geq1}$ and put
$P_n=\max_{1\leq\ell\leq n}u_\ell$.
If $(f_n)$ is Cauchy in all these seminorms, it is Cauchy uniformly
on each compact set, and has an entire limit $f$. In the weighted
Cauchy inequality for any fixed seminorm, passage to the
pointwise limit in the second index yields the same uniform
weighted bound for $f_n-f$. Thus $f$ has all the finite growth
seminorms, lies in the original $\cqE$, and $f_n\to f$ in
$\gamma$. This proves completeness of $(\cqE,\gamma)$.

The ideal $I$ is closed: the proved kernel description is an
intersection of kernels of derivative evaluations, and Cauchy's
estimate shows that each of these evaluations is continuous for
$\gamma$. The quotient is therefore Hausdorff. Its topology has
the countable increasing seminorm family
\[
 \widetilde P_n([f])=\inf_{h\in I}P_n(f+h).
\]
Let $(y_n)$ be a Cauchy sequence in this quotient. Choose a
strictly increasing subsequence of indices $n_k$ so far out that
\[
 \widetilde P_k(y_{n_{k+1}}-y_{n_k})<2^{-k-1}
 \quad(k\geq1).
\]
This is possible by the Cauchy property and by taking $n_k$ past
the threshold for the $k$th seminorm before selecting the next
index. Choose a representative $d_k\in\cqE$ of each displayed
difference with $P_k(d_k)<2^{-k}$. For any fixed $j$, the tail
$k\geq j$ satisfies $P_j(d_k)\leq P_k(d_k)<2^{-k}$.
Consequently the partial sums of $\sum_{k\geq1}d_k$ are
Cauchy in $\gamma$ and converge in $\cqE$. Adding a
representative of $y_{n_1}$ gives a function whose quotient class
is the limit of $(y_{n_k})$. The full Cauchy sequence $(y_n)$
then has the same limit: for each quotient neighbourhood choose
a smaller additive neighbourhood whose sum with itself is
contained in it, and combine the Cauchy tail with a convergent
subsequence term. The countable seminorm topology is metrizable,
so this proves quotient completeness and the Fr\'echet assertion.

The identity $(\cqE,\gamma)\to(\cqE,\kappa)$ is continuous:
for every $R,p,a$, one has $M_f(R)\leq e^{aR^p}\|f\|_{p,a}$.
It preserves the identical algebraic subspace $I$, so the
quotient universal property proves continuity and bijectivity
of \eqref{eq:cq-topology-comparison-strict}. If its inverse
were continuous, it would be a continuous linear map, hence
uniformly continuous for the additive uniformities. The Cauchy
sequence $(x_n)$ constructed above would then be Cauchy in the
complete $Q_\gamma$, would converge there, and by continuity of
$\iota$ would converge in $Q_\kappa$, a contradiction. This
proves that its inverse is not continuous. No invariance of
metric completeness under arbitrary nonlinear homeomorphisms
has been used.
\end{proof}

The original coordinate action remains visible after completion.
For $v=J_\Xi f$, the exact product rule gives
\begin{equation}\label{eq:cq-original-s-on-completed-jets}
 \bigl(J_\Xi(sf)\bigr)_{\lambda,j}
 =\lambda v_{\lambda,j}+jv_{\lambda,j-1},
 \qquad 0\leq j<m_\lambda,
\end{equation}
where the second term is $0$ when $j=0$. Each output coordinate
depends on finitely many input coordinates, so this extends to
a continuous linear operator on the entire product completion.
The factors $j$ are retained because these coordinates are actual
derivatives. The completion adds precisely the nonrealizable
jet arrays, including \eqref{eq:cq-nonrealizable-jet-array}; it
does not add an entire function of order at most one representing
such an array. Neither analytic quotient topology has been
identified with the source topology $\tau$.


\subsection{Exact derivative domains inside the analytic relation}
Let $I=\Xi\cqE=N_\gamma=N_\kappa$. For $k\geq1$ the exact
representative domain in the relation space is
\begin{equation}\label{eq:cq-Ik}
 \begin{aligned}
 I\cap D^{-k}I=\{\Xi h:\ &h\in\cqE,\quad
 [z^n](A_\lambda(\lambda+z)h(\lambda+z))=0\\
 &\text{for }\max(0,k-m_\lambda)\leq n<k,
                   \text{ at every zero }\lambda\}.
 \end{aligned}
\end{equation}
where the original local unit is
$A_\lambda(s)=\Xi(s)/(s-\lambda)^{m_\lambda}$.
To prove this, write $\Xi h=z^m\sum_{n\geq0}b_nz^n$ at
$\lambda$, with $m=m_\lambda$ and
$b_n=[z^n](A_\lambda h)$. Its $k$th derivative is
\[
 \sum_{m+n\geq k}\frac{(m+n)!}{(m+n-k)!}
               b_nz^{m+n-k}.
\]
Every factorial coefficient is nonzero. Vanishing to order at least
$m$ is equivalent to $b_n=0$ exactly for
$\max(0,k-m)\leq n<k$, as stated. The derivative has the original
order bound, so the vanishing condition is equivalent to membership
in $I$ by \eqref{eq:cq-equalclosures}. This proves both directions
globally, including all multiplicities and local unit coefficients.

For $k=1$, this is simply $h(\lambda)=0$ at every distinct zero,
without changing or discarding its original multiplicity in $\Xi$.
The derivative ambiguity has the exact presentation
\begin{equation}\label{eq:cq-firstambiguity}
 \cqE/\{h:h(\lambda)=0\text{ for every }\Xi(\lambda)=0\}
 \ \xrightarrow{\ \cong\ }\ (DI+I)/I,
 \qquad [h]\longmapsto[\Xi'h].
\end{equation}
Indeed $D(\Xi h)=\Xi h'+\Xi'h$, and the first term lies in $I$.
At each zero of multiplicity $m$, $\Xi'$ has exactly order $m-1$,
so its product with $h$ is divisible by $\Xi$ exactly when
$h(\lambda)=0$ for every zero. This proves the precise kernel,
image, and inverse on the image in \eqref{eq:cq-firstambiguity}.

For the heat generator ($k=2$), at any zero of multiplicity at least
two the condition is $h(\lambda)=h'(\lambda)=0$.
At a simple zero it is instead
\begin{equation}\label{eq:cq-simpleheatdomain}
 A_\lambda(\lambda)h'(\lambda)
        +A_\lambda'(\lambda)h(\lambda)=0.
\end{equation}
It does not require $h(\lambda)=0$. These follow from the index
range in \eqref{eq:cq-Ik}: the simple zero has only the coefficient
$b_1$ constrained. In particular the second-derivative relation
domain need not be a multiplication module: at a simple zero, a
local $\Xi h$ obeying \eqref{eq:cq-simpleheatdomain} and with
$h(\lambda)\ne0$ has
$D^2(s\Xi h)(\lambda)=2(\Xi h)'(\lambda)\ne0$.
This is a local statement with its explicit witness condition; it
does not assume the existence of a global simple zeta zero.

There is also an exact calculation of every finite projection of
the ambiguity $W_k=(D^kI+I)/I$. In the factorial jet coordinates
of any finite set $F$ of distinct zeros, its image is precisely
\begin{equation}\label{eq:cq-ambiguityjets}
 \prod_{\lambda\in F}
 \epsilon^{\max(m_\lambda-k,0)}
       \cqC[\epsilon]/(\epsilon^{m_\lambda}).
\end{equation}
The local dimension is $\min(k,m_\lambda)$. The differentiated
series above proves containment. For surjectivity, prescribe the
allowed output coefficient $y_j$ of $z^j$, where
$\max(m-k,0)\leq j<m$. Set
$b_{j+k-m}=j!y_j/(j+k)!$ and give the other coefficients
$b_0,\ldots,b_{k-1}$ arbitrary values, say zero. Finite Taylor
division by the nonzero unit $A_\lambda$ then specifies the first
$k$ Taylor coefficients of $h$ at $\lambda$.
These finitely many specifications can be achieved by one polynomial:
put $M_\lambda(s)=\prod_{\nu\in F,\nu\ne\lambda}(s-\nu)^k$;
take the Taylor polynomial of degree $k-1$ of the prescribed local
polynomial divided by $M_\lambda$, and sum the resulting
$M_\lambda$ multiples over $F$. At its own node each summand
agrees with the prescribed polynomial modulo $(s-\lambda)^k$,
and at the other nodes it vanishes to that order. This constructs
the required $h\in\cqC[s]\subset\cqE$ and proves
\eqref{eq:cq-ambiguityjets}. No claim about arbitrary infinite jet
interpolation has replaced this exact finite image.

\begin{cqtheorem}[Direct differential and heat obstructions]
\label{thm:cq-obstruction}
Neither $N_\gamma$ nor $N_\kappa$ is invariant under $D$ or $D^2$.
For either frozen quotient, neither representative formula $[f']$
nor $[-f''/4]$ defines an operator on any nonempty class domain
admitting all representatives.
For every $t\ne0$, neither $U_tN_\gamma\subseteq N_\gamma$ nor
$U_tN_\kappa\subseteq N_\kappa$ holds. Ordinary multiplication
by $s$ preserves neither $U_tN_\gamma$ nor $U_tN_\kappa$.
\end{cqtheorem}
\begin{proof}
Choose any zero $\lambda$ of the actual $\Xi$ and let $m\geq1$
be its multiplicity. The existence of its zeros is part of the
classical arithmetic datum in the primary source; no realness or
simplicity is used. Write $\Xi=(s-\lambda)^mA$ with
$A(\lambda)\ne0$. Then $\Xi'=(s-\lambda)^{m-1}
(mA+(s-\lambda)A')$ does not vanish to order $m$.
Since $\Xi\in N_\gamma\subseteq N_\kappa$, both derivative
invariance assertions fail. If $m\geq2$, $\Xi''$ has leading
coefficient $m(m-1)A(\lambda)$ at order $m-2$, proving the
second-derivative failure. If $m=1$, use the actual polynomial
multiple $n=(s-\lambda)\Xi$ instead: $n''(\lambda)=2\Xi'(\lambda)
\ne0$. This retains the product-rule cross term and covers all
multiplicities. The nonempty-domain assertion follows from the
explicit fibre computation in \eqref{eq:cq-relation}.

For $a\in\cqC$, $e^{as}\Xi\in N_\gamma$ by polynomial
approximation of $e^{as}$ and continuous multiplication, and hence
belongs to $N_\kappa$. The conjugation $D(e^{as}f)=e^{as}(D+a)f$
and the absolutely convergent heat and translation series give
\begin{equation}\label{eq:cq-exponential}
 U_t(e^{as}\Xi)(s)
    =e^{as-ta^2/4}Z_t(s-ta/2).
\end{equation}
For completeness, the translation series for $f$ is its Taylor
series about $s$ and sums to $f(s-ta/2)$; its local uniform
convergence and the derivative estimates justify commutation of
$e^{-(ta/2)D}$ with $U_t$. If $U_tN_\bullet\subseteq N_\bullet$
for either choice, evaluate \eqref{eq:cq-exponential} at $\lambda$.
The exponential never vanishes, so $Z_t(\lambda-ta/2)=0$ for
every $a$. When $t\ne0$ these arguments fill $\cqC$; this contradicts
$Z_t\ne0$. Finally if $M$ preserved $U_tN_\bullet$, conjugation
would imply $(M+(t/2)D)N_\bullet\subseteq N_\bullet$.
Since $M$ already preserves $N_\bullet$, this forces
$DN_\bullet\subseteq N_\bullet$, contrary to the first part.
\end{proof}

The analytic moving-divisor quotient, which does carry the zeros of
the actual $Z_t$ under ordinary $s$, is
\begin{equation}\label{eq:cq-divisor}
 I_t=\overline{\cqC[s]Z_t}^{\,\kappa}
 =\{f\in\cqE:f/Z_t\text{ is entire}\},\qquad
 A_t=(\cqE,\kappa)/I_t.
\end{equation}
The division lemma and polynomial approximation further prove
\[
 I_t=Z_t\cqE=\overline{\cqC[s]Z_t}^{\,\gamma}\quad(t\in\cqC).
\]
For a complex $t$ with $Z_t(0)=0$, use the retained nonzero center
$c$ in the division lemma; $Z_t$ is not identically zero.
The analytically transported relation has the different exact
principal presentation for the proved transported multiplication:
\begin{equation}\label{eq:cq-starideal}
 U_t(\Xi\cqE)=Z_t\star_t\cqE.
\end{equation}
Indeed $U_t(\Xi h)=Z_t\star_t U_th$, and $U_t$ is onto.
The map $h\mapsto Z_t\star_t h$ is a continuous linear
isomorphism of $(\cqE,\gamma)$ onto this closed relation;
its inverse on that relation is
$n\mapsto U_t((U_{-t}n)/\Xi)$, continuous by the division and
heat lemmas. Thus the two multiplications and their exact relation
spaces are connected by a proved isomorphism with a proved inverse.
The proof of Theorem~\ref{thm:cq-compact} applies verbatim with
the nonzero entire order-at-most-one function $Z_t$: it uses only
local Taylor division and compact polynomial approximation.
Consequently $\cqSpec(M\text{ on }A_t)=\{\lambda:Z_t(\lambda)=0\}$,
with the full local multiplicity jets, and with the explicit inverse
$[f]\mapsto[(f-f(\lambda)Z_t/Z_t(\lambda))/(s-\lambda)]$
off that set. Equation \eqref{eq:cq-zero} is the evolution of this
original-coordinate spectrum on its simple-zero locus.
This spectrum can move; \eqref{eq:cq-specmove} is the spectrum
of the different, exactly specified transported operator $B_t$.

We now give the promised comparison morphisms with every kernel
retained. Put $J=\cqN\cap N_\kappa$ and
$C=(\cqE,\tau\vee\kappa)/J$, where $\tau\vee\kappa$ is the
topology generated by both topologies. Since each summand is closed
in its own topology, $J$ is closed in their join. The identity maps
give continuous surjections
\begin{equation}\label{eq:cq-span}
 \cqQ\ \longleftarrow\ C\ \longrightarrow\ Q_\kappa,
 \qquad\ker(C\to\cqQ)=\cqN/J,
 \quad\ker(C\to Q_\kappa)=N_\kappa/J.
\end{equation}
Each assertion follows by taking the same representative $f$ and
testing membership in the respective relation subspace; the quotient
universal property proves continuity. An algebraic common quotient is
$\cqE/(\cqN+N_\kappa)$ with the quotient topology from the join;
its Hausdorff version uses
$\overline{\cqN+N_\kappa}^{\,\tau\vee\kappa}$.
The map from $C$ to this further quotient is continuous. Continuity
of maps to it from each of the two separately topologized quotients
is not asserted. No closedness of the sum is presumed. This describes
the precise information lost in that further quotient.

At time $t$ replace the two subspaces by $N_t$ and $I_t$ and the
join by $\tau_t\vee\kappa$. The same construction gives continuous
surjections from their intersection quotient to the genuine transported
arithmetic quotient $Q_t$ and to $A_t$, with kernels
$N_t/(N_t\cap I_t)$ and $I_t/(N_t\cap I_t)$, respectively.
The two generator systems are linked by the explicit formula
$R(B_t)Z_t$ versus $R(s)Z_t$ and the defects
\eqref{eq:cq-generators}--\eqref{eq:cq-unit}.
These are exact morphisms and closure formulas for the actual source
space. Determining that either kernel vanishes for $\tau$ is not
established by the inspected primary loci. Accordingly the
source-frozen statements $D\cqN\subseteq\cqN$ and
$D^2\cqN\subseteq\cqN$ have not been decided solely from that
unnamed topology. The unqualified noninvariance results above have
been proved in the two explicitly specified analytic realizations.

\section{Arithmetic traces and their exact Mellin obstructions}
\label{sec:cq-arithmetic}
\subsection{The full Schwartz trace and arithmetic principal parts}
\label{sec:cq-principalparts}
Retain the source domain
\[
 \mathcal S_0^{\rm ev}
 =\{\psi\in\mathcal S(\mathbb R):
       \psi(-x)=\psi(x),\ \psi(0)=\widehat\psi(0)=0\},
 \qquad
 \widehat\psi(y)=\int_{\mathbb R}\psi(x)e^{-2\pi ixy}\,dx.
\]
Its full trace map takes values in entire functions:
\begin{equation}\label{eq:cq-Schwartz-trace}
 \begin{aligned}
 F_\psi(s)&=\mathcal F_\mu\mathcal E\psi(s),\\
 F_\psi(i(z-\tfrac12))&=\zeta(z)\mathcal M\psi(z),
 &\mathcal M\psi(z)&=\int_0^\infty\psi(x)x^z\,\frac{dx}{x}.
 \end{aligned}
\end{equation}
The Mellin integral is holomorphic for $\operatorname{Re}z>-2$,
and the factorization holds there with its removable value at
$z=1$. This statement does not impose an unproved order bound
on the entire $F_\psi$ for arbitrary Schwartz inputs.

Here is a direct proof with all analytic domains retained.
For fixed $x>0$, periodize $\psi$ by the smooth function
\[
 y\longmapsto\sum_{k\in\mathbb Z}\psi(x(k+y)).
\]
Its Fourier coefficients are $x^{-1}\widehat\psi(\ell/x)$:
integrate on $y\in[0,1]$, substitute $v=x(k+y)$, and sum the
adjacent intervals, justified by Schwartz decay. These coefficients
decrease faster than every power of $\ell$, so their Fourier series
converges with all derivatives and at $y=0$ gives the Poisson
identity. Evenness and the two exact zero conditions yield
\[
 \sum_{n\geq1}\psi(nx)=x^{-1}\sum_{n\geq1}\widehat\psi(n/x),
 \qquad \mathcal E\psi(x)=\mathcal E\widehat\psi(1/x).
\]
For $x\geq1$, a Schwartz bound of arbitrarily high order gives
$|\sum\psi(nx)|\leq C_Ax^{-A}\sum n^{-A}$ for $A>1$.
The same estimate after termwise differentiation follows by choosing
the Schwartz exponent larger than the derivative order plus one.
The Poisson identity supplies the corresponding bounds at zero.
Thus $\mathcal E\psi(e^u)$ decreases faster than $e^{-A|u|}$
for every $A>0$ at both ends. Its Fourier integral and each
derivative in $s$ are dominated locally uniformly in $s$ by an
integrable function. This proves that $F_\psi$ is entire.

Since $\psi$ is even and $\psi(0)=0$, its Taylor formula gives
$\psi(x)=O(x^2)$ at zero. Together with Schwartz decay at infinity
this proves the stated holomorphy of $\mathcal M\psi$ by
dominated differentiation, including all powers of $\log x$.
For $\sigma=\operatorname{Re}z>1$,
\[
 \sum_{n\geq1}\int_0^\infty|\psi(nx)|x^\sigma\,\frac{dx}{x}
 =\Bigl(\sum_{n\geq1}n^{-\sigma}\Bigr)
          \int_0^\infty|\psi(y)|y^\sigma\,\frac{dy}{y}<\infty.
\]
Absolute interchange and the original exponent $x^{-is}$ now give
\eqref{eq:cq-Schwartz-trace} on this half-plane. Moreover
$\mathcal M\psi(1)=\int_0^\infty\psi(x)\,dx
=\widehat\psi(0)/2=0$. This cancels the simple pole of $\zeta$
at $1$, so the right side is holomorphic for
$\operatorname{Re}z>-2$. The identity theorem extends the equality
to that entire half-plane, proving the assertion.

Let $\mathcal Z$ denote the set of distinct nontrivial zeros $\rho$
of $\zeta$, with their original multiplicities $m_\rho$.
The evenness of $\Xi$, immediate from the retained cosine integral,
gives the exact reflected-coordinate identity
\begin{equation}\label{eq:cq-reflected-Xi}
 s=i(z-\tfrac12),\qquad z=\tfrac12-is,\qquad
 \Xi(i(z-\tfrac12))=\xi(1-z)=\xi(z).
\end{equation}
At each $\rho\in\mathcal Z$ the completion factor
$z(z-1)\pi^{-z/2}\Gamma(z/2)/2$ is a nonvanishing holomorphic
unit. Hence $\lambda_\rho=i(\rho-\tfrac12)$ is a zero of
$\Xi$ of exactly multiplicity $m_\rho$, and these are all its
zeros. No reflection of the retained spectral operator $s$ has
been performed: \eqref{eq:cq-reflected-Xi} is the exact Mellin
parameter dictionary for that same operator. The factorization
\eqref{eq:cq-Schwartz-trace} proves
\begin{equation}\label{eq:cq-fulltrace-jets}
 \begin{gathered}
 D^jF_\psi(\lambda_\rho)=0
 \quad(\rho\in\mathcal Z,\ 0\leq j<m_\rho),\\
 F_\psi\in\cqE\ \Longrightarrow\ F_\psi\in\Xi\cqE.
 \end{gathered}
\end{equation}
The first assertion follows because $\mathcal M\psi$ is holomorphic
at every $\rho$ and $s-\lambda_\rho=i(z-\rho)$; the second
then follows from the proved division theorem. It is an exact
intersection assertion about the full Schwartz trace, without
enlarging the order-at-most-one domain or closing it in $\tau$.

For each $\rho$ retain the actual local unit
$a_\rho(z)=\zeta(z)/(z-\rho)^{m_\rho}$ and the space
\[
 \mathcal P_\rho^-=
  \left\{\sum_{j=1}^{m_\rho}c_{-j}(z-\rho)^{-j}:c_{-j}\in\cqC\right\}.
\]
The principal-part map on the original entire-function set is
\begin{equation}\label{eq:cq-principalparts}
 \mathfrak P:\cqE\longrightarrow\prod_{\rho\in\mathcal Z}\mathcal P_\rho^-,
 \qquad
 F\longmapsto\left(\operatorname{PP}_{z=\rho}
                 \frac{F(i(z-\tfrac12))}{\zeta(z)}\right)_\rho.
\end{equation}
The kernel of this map is exactly $\Xi\cqE$. It induces an
injective linear map from the underlying vector space of each
$Q_\gamma,Q_\kappa$ onto the same specified image. Every finite
projection is surjective. The map from $Q_\kappa$ to its image,
with the product subspace topology, is a topological linear
isomorphism; the product of all principal-part spaces is its
completion. From $Q_\gamma$ it is a continuous bijection onto
that image whose inverse is not continuous.

To prove each claim while retaining the analytic units, write
$a_\rho(\rho+w)^{-1}=\sum_{\ell\geq0}\alpha_{\rho,\ell}w^\ell$.
The exact Laurent coefficient at degree $-m_\rho+q$, $0\leq q<m_\rho$,
of the $\rho$ component is
\begin{equation}\label{eq:cq-principalpart-coefficients}
 c_{-m_\rho+q}
 =\sum_{j=0}^{q}\frac{i^j}{j!}F^{(j)}(\lambda_\rho)
                                  \alpha_{\rho,q-j}.
\end{equation}
This follows by multiplying the two convergent local Taylor
series and the retained factor $w^{-m_\rho}$. Its finite matrix
is triangular with nonzero diagonal entries
$i^j/(j!a_\rho(\rho))$. Explicitly, given a prescribed principal
part $P_\rho(w)$, take the coefficients $b_j$ through degree
$m_\rho-1$ of $a_\rho(\rho+w)w^{m_\rho}P_\rho(w)$;
the required original derivatives are $j!i^{-j}b_j$.
Multiplying back by $a_\rho^{-1}$ recovers the prescribed negative
powers, since the discarded error is divisible by $w^{m_\rho}$.
This proves a finite-dimensional linear isomorphism at every zero,
including its inverse and every factorial and power of $i$.

The vanishing of all principal parts is therefore equivalent to
the vanishing of all the original multiplicity jets of $F$.
Theorem~\ref{thm:cq-equalclosures} identifies that kernel exactly
with $\Xi\cqE$. At any finite set of zeros, the required derivatives
can be realized by the explicit Hermite lift
\eqref{eq:cq-finite-Hermite-lift}; this proves surjectivity of
every finite principal-part projection. Each block map and its
inverse are continuous because their finite matrices have fixed
complex coefficients. Their product is therefore a homeomorphism
of the two product spaces: the inverse image of a cylinder fixing
finitely many coordinates is again controlled by those finitely
many blocks. Applying this product isomorphism to
Theorem~\ref{thm:cq-jet-topology} proves all stated topology,
completion, and comparison assertions for $\mathfrak P$.

This map also carries the original spectral action exactly. For
$P_\rho=\mathfrak P_\rho(F)$ one has
\begin{equation}\label{eq:cq-s-principalparts}
 \mathfrak P_\rho(sF)
 =\operatorname{PP}_{w=0}
                  ((\lambda_\rho+iw)P_\rho(w)).
\end{equation}
Indeed the full Laurent quotient for $sF$ is the quotient for $F$
multiplied by $i(z-\tfrac12)=\lambda_\rho+iw$.
Its holomorphic residual contributes no negative power after this
multiplication. Thus the finite block is the retained
$\lambda_\rho$ plus the explicitly specified nilpotent operator
$P\mapsto\operatorname{PP}(iwP)$, whose nilpotency index is
$m_\rho$: iteration on $w^{-m_\rho}$ gives a nonzero multiple
of $w^{-1}$ at step $m_\rho-1$ and zero at step $m_\rho$.

For $I=\Xi\cqE$, the derivative ambiguity
$W_k=(D^kI+I)/I$ has, under $\mathfrak P$, the exact finite image
\begin{equation}\label{eq:cq-principal-ambiguity}
 \prod_{\rho\in F}
 \left\{\sum_{j=1}^{\min(k,m_\rho)}c_{-j}(z-\rho)^{-j}:c_{-j}\in\cqC\right\}
 \quad(F\subset\mathcal Z\text{ finite}).
\end{equation}
The finite image in \eqref{eq:cq-ambiguityjets} consists exactly
of numerator jets divisible by $w^{\max(m_\rho-k,0)}$ after
the invertible substitution $s-\lambda_\rho=iw$.
Multiplication by the nonvanishing unit $a_\rho^{-1}$ preserves
this divisibility and is invertible on the truncated Taylor ring.
Division by $w^{m_\rho}$ therefore yields exactly all negative
powers of order at most $\min(k,m_\rho)$, proving both
containment and surjectivity in the displayed finite image.
In particular the complete local-unit expression, rather than only
a count of vanishing orders, transports the derivative defect to
the arithmetic Mellin poles.

\subsection{The exact Schwartz--Mellin image}
\label{sec:cq-Schwartz-Mellin-image}
Put $\theta=x\partial_x$. For real $a\leq b$, retain the finite polynomial
\[
 P_{a,b}(z)=\prod_{\substack{k\geq1\\a-1\leq-2k\leq b+1}}(z+2k),
\]
with empty product equal to $1$. Define $\mathfrak M_0^{\rm ev}$
to be the vector space of meromorphic functions $M$ on $\cqC$
whose only possible poles are simple poles at $-2k$, $k\geq1$,
which satisfy $M(1)=0$, and for which all the seminorms
\begin{equation}\label{eq:cq-Mellin-strip-seminorm}
 q_{a,b,N}(M)=
 \sup_{\substack{a\leq\sigma\leq b\\v\in\mathbb R}}
 (1+|v|)^N
 \bigl|P_{a,b}(\sigma+iv)M(\sigma+iv)\bigr|,
 \qquad N=0,1,2,\ldots,
\end{equation}
are finite. At a pole of $M$, the product in this formula has its
removable value. The margin $1$ in $P_{a,b}$ is retained in this
definition. It ensures that all poles in a neighborhood of the
closed strip have been removed. These seminorm conditions are
equivalent to rapid decrease of $M$ on every finite vertical strip
for $|\operatorname{Im}z|\geq1$: on that part of the strip each
factor $|z+2k|$ is bounded above by a constant times $1+|v|$
and is at least $|v|$; on the remaining compact rectangle,
$P_{a,b}M$ is holomorphic in a neighborhood and is bounded.
In particular no bound uniform in the pole index $k$ has been
inserted into the definition.
One can equivalently use the residue coordinates together with
the explicit away-from-poles strip seminorms
\[
 r_{a,b,N,\delta}(M)=
 \sup_{\substack{a\leq\operatorname{Re}z\leq b\\
       \operatorname{dist}(z,\{-2,-4,\ldots\})\geq\delta}}
   (1+|\operatorname{Im}z|)^N|M(z)|,
 \qquad 0<\delta\leq\tfrac12.
\]
Here is the equivalence as a statement about topologies, not only
about membership. On the domain of this last supremum every
factor of $P_{a,b}$ has modulus at least $\delta$, so
$r_{a,b,N,\delta}\leq\delta^{-\deg P_{a,b}}q_{a,b,N}$.
Residues are controlled by $q$ as shown explicitly below.
Conversely, outside the radius-$\delta$ disks about the poles,
$|P_{a,b}(z)|\leq C_{a,b}(1+|\operatorname{Im}z|)^{\deg P_{a,b}}$
on the strip, so the corresponding part of $q$ is bounded by
an $r$ seminorm with this degree added to its exponent.
For a pole whose radius-$\delta$ disk meets the strip, that pole
is among the factors of $P_{a,b}$. The function $P_{a,b}M$
is holomorphic throughout the disk; its boundary has distance
at least $\delta$ from every pole, since distinct pole locations
are separated by $2$. That boundary lies in the enlarged strip
$a-1\leq\operatorname{Re}z\leq b+1$. The maximum-modulus
principle bounds the product inside the disk by its boundary
values, and $(1+|\operatorname{Im}z|)^N\leq(1+\delta)^N$
there. An $r$ seminorm on that enlarged strip therefore bounds
the remaining part of $q$. There are finitely many relevant disks.
These estimates prove equality of the two locally convex
topologies; separate residue coordinates may be included without
changing them.

The map
\begin{equation}\label{eq:cq-Mellin-isomorphism}
 \mathcal M:\mathcal S_0^{\rm ev}\longrightarrow
                 \mathfrak M_0^{\rm ev},\qquad
 (\mathcal M\psi)(z)=\int_0^\infty\psi(x)x^z\,\frac{dx}{x}
 \quad(\operatorname{Re}z>-2),
\end{equation}
is a linear bijection. Its inverse, for every fixed real $c>1$, is
\begin{equation}\label{eq:cq-Mellin-reconstruction}
 \psi(x)=\frac1{2\pi}\int_{\mathbb R}
              M(c+iv)x^{-c-iv}\,dv\quad(x>0),
 \qquad \psi(0)=0,\qquad\psi(-x)=\psi(x).
\end{equation}
For every $k\geq1$ its precise residue is
\begin{equation}\label{eq:cq-Mellin-Taylor-residue}
 \operatorname*{Res}_{z=-2k}M(z)
       =\frac{\psi^{(2k)}(0)}{(2k)!}.
\end{equation}
Both maps are continuous when the source has its ordinary Schwartz
seminorms and the target has \eqref{eq:cq-Mellin-strip-seminorm}.

Here is a proof including the strip and reconstruction estimates.
Fix a smooth function $\chi$ on $[0,\infty)$ which equals $1$
on $[0,1]$ and $0$ on $[2,\infty)$. Such a function is obtained,
for example, by setting $\eta(y)=0$ for $y\leq0$ and
$\eta(y)=e^{-1/y}$ for $y>0$, and taking
$\chi(x)=\eta(2-x)/(\eta(2-x)+\eta(x-1))$. Its denominator is
positive at every $x\geq0$. Let
$a_k=\psi^{(2k)}(0)/(2k)!$. For an integer $K\geq1$ write
\[
 r_K(x)=\psi(x)-\chi(x)\sum_{k=1}^K a_kx^{2k}.
\]
Evenness and $\psi(0)=0$ show that $r_K$ vanishes to order
at least $2K+2$ at zero. Taylor's integral remainder, applied
also after differentiating, proves
\[
 |\theta^Jr_K(x)|\leq C_{J,K}(\psi)x^{2K+2}
 \quad(0<x\leq1),\qquad J\geq0,
\]
where $C_{J,K}(\psi)$ is bounded by a constant times finitely
many Schwartz seminorms of $\psi$. To see explicitly that all
$J$ are covered, expand
$\theta^J=\sum_{\ell=0}^J S(J,\ell)x^\ell\partial_x^\ell$,
where $S(0,0)=1$ and
$S(J+1,\ell)=\ell S(J,\ell)+S(J,\ell-1)$. For
$\ell\leq2K+2$, Taylor's remainder gives
\[
 |\partial_x^\ell r_K(x)|
 \leq\frac{x^{2K+2-\ell}}{(2K+2-\ell)!}
              \sup_{0\leq y\leq1}|\psi^{(2K+2)}(y)|.
\]
For $\ell>2K+2$, the differentiated polynomial vanishes on
this interval and $x^\ell\leq x^{2K+2}$, which gives the
claimed estimate using $\sup_{[0,1]}|\psi^{(\ell)}|$.
At infinity $r_K=\psi$, so each $\theta^Jr_K$ decreases
faster than every inverse power.

The cutoff Mellin function
\[
 C(w)=\int_0^\infty\chi(x)x^w\,\frac{dx}{x}
       \quad(\operatorname{Re}w>0)
\]
has the meromorphic continuation
\begin{equation}\label{eq:cq-Mellin-cutoff}
 C(w)=-\frac{L(w)}{w},\qquad
 L(w)=\int_1^2\chi'(x)x^w\,dx.
\end{equation}
The entire function $L$ satisfies $L(0)=\chi(2)-\chi(1)=-1$,
so the residue of $C$ at $0$ is $1$. Repeated integration by
parts with $\theta$ in the expression
$L(w)=\int_1^2(x\chi'(x))x^w\,dx/x$ gives, on every finite
vertical strip, rapid decrease for $|\operatorname{Im}w|\geq1$.
Indeed $x\chi'(x)$ is smooth with support in $[1,2]$ and all
boundary derivatives vanish, and for every integer $J\geq0$
\[
 w^JL(w)=(-1)^J\int_1^2
       \theta^J(x\chi'(x))x^w\,\frac{dx}{x}.
\]
Thus $C$ has the same rapid decrease away from its sole pole.

For any strip $a\leq\operatorname{Re}z\leq b$, choose $K$
with $a+2K+2>0$. On the larger half-plane
$\operatorname{Re}z>-2K-2$, the formula
\begin{equation}\label{eq:cq-Mellin-meromorphic-extension}
 M(z)=\int_0^\infty r_K(x)x^z\,\frac{dx}{x}
                     +\sum_{k=1}^K a_k C(z+2k)
\end{equation}
is meromorphic and agrees with the original Mellin integral where
$\operatorname{Re}z>-2$. The expressions for successive $K$
therefore agree by the identity theorem and define one meromorphic
function on $\cqC$. Formula \eqref{eq:cq-Mellin-cutoff} proves
that its only possible poles and residues are exactly
\eqref{eq:cq-Mellin-Taylor-residue}.
For every $J\geq0$, integration by parts gives
\[
 z^J\int_0^\infty r_K(x)x^z\,\frac{dx}{x}
       =(-1)^J\int_0^\infty\theta^Jr_K(x)x^z\,\frac{dx}{x}.
\]
All boundary terms vanish: at zero they are bounded by a constant
times $x^{a+2K+2}$, and at infinity the Schwartz exponent can
be chosen larger than $b+J+1$. The absolute value of the final
integral, uniformly in this strip, is at most
\begin{equation}\label{eq:cq-Mellin-forward-bound}
 B_{K,J;a,b}(\psi)=
 \int_0^1|\theta^Jr_K(x)|x^{a-1}\,dx
 +\int_1^\infty|\theta^Jr_K(x)|x^{b-1}\,dx<\infty.
\end{equation}
For $|\operatorname{Im}z|\geq1$, divide by $|z|^J\geq
|\operatorname{Im}z|^J$. Taking $J$ arbitrarily large proves
all the required strip estimates, including the polynomial
$P_{a,b}$. At bounded imaginary part, multiply
\eqref{eq:cq-Mellin-meromorphic-extension} by $P_{a,b}$;
every pole meeting the strip cancels and the resulting expression
is bounded on the compact rectangle. The constants in these
bounds involve only finitely many Schwartz seminorms:
the first integral in \eqref{eq:cq-Mellin-forward-bound} is
at most $C_{J,K}(\psi)/(a+2K+2)$; on $[1,2]$ the cutoff
has fixed bounded derivatives; on $[2,\infty)$ the expansion
of $\theta^J$ above uses finitely many bounds
$(1+x)^A|\psi^{(\ell)}(x)|$ with $A>b+J+1$.
This proves continuity of the forward map.
Finally $M(1)=\int_0^\infty\psi(x)\,dx
=\widehat\psi(0)/2=0$, retaining the factor $2$ from evenness.

Conversely, let $M\in\mathfrak M_0^{\rm ev}$ and define
\eqref{eq:cq-Mellin-reconstruction}. The integral and its
ordinary $x$ derivatives converge uniformly on each compact
subinterval of $(0,\infty)$. More precisely, with the rising
factorial $(z)_0=1$ and
$(z)_j=z(z+1)\cdots(z+j-1)$ for $j\geq1$, they are
\begin{equation}\label{eq:cq-Mellin-differentiated-inverse}
 \psi^{(j)}(x)=\frac{(-1)^j}{2\pi}
       \int_{\mathbb R}(c+iv)_jM(c+iv)x^{-c-iv-j}\,dv.
\end{equation}
The factor $(c+iv)_j$ grows at most polynomially in $v$,
so a strip estimate with exponent greater than $j+1$ supplies
an integrable dominating function. This proves smoothness for $x>0$.

The inverse integral is independent of the positive choice of
$c$. In fact integrate $M(z)x^{-z}$ around the rectangle
between two positive vertical lines and heights $\pm T$.
It has no poles there. On each horizontal side, its absolute
value is at most a constant, depending on the fixed $x$ and
the bounded real interval, times $(1+T)^{-N}$ for arbitrary
$N$. Thus the horizontal contributions tend to zero, and the
two upward vertical integrals agree. In particular for every
real $b>0$ the differentiated formula holds with $c$ replaced
by $b$. If
\begin{equation}\label{eq:cq-Mellin-line-integral}
 I_{d,j}(M)=\frac1{2\pi}\int_{\mathbb R}
                |(d+iv)_jM(d+iv)|\,dv
 \quad\bigl(d\notin\{-2,-4,-6,\ldots\}\bigr),
\end{equation}
then $I_{d,j}(M)$ is finite and
\begin{equation}\label{eq:cq-Mellin-infinity-bound}
 |\psi^{(j)}(x)|\leq I_{b,j}(M)x^{-b-j}
 \quad(x>0,\ b>0).
\end{equation}
This gives every required decay estimate at infinity.

For the behavior at zero fix $K\geq1$ and take the actual
left vertical line $d=-2K-1$. The rectangle between $d$ and
$c$ has precisely the possible poles $-2,-4,\ldots,-2K$.
Traverse its right edge upward and its left edge downward;
this is the positive orientation. The residue theorem and
the same horizontal estimates give
\begin{equation}\label{eq:cq-Mellin-contour-shift}
 \begin{aligned}
 \psi(x)&=\sum_{k=1}^K a_kx^{2k}+R_K(x),
 &a_k&=\operatorname*{Res}_{z=-2k}M(z),\\
 R_K(x)&=\frac1{2\pi}\int_{\mathbb R}
                    M(-2K-1+iv)x^{2K+1-iv}\,dv,\\
 |R_K^{(j)}(x)|&\leq I_{-2K-1,j}(M)x^{2K+1-j}.
 \end{aligned}
\end{equation}
The sign of each residue is positive. The last estimate follows
by differentiating the remainder integral with the exact factor
$(-1)^j(-2K-1+iv)_j$; this factor is included in
\eqref{eq:cq-Mellin-line-integral}.
For any given nonnegative integer $j$, choose $K$ so that
$2K+1>j$. Formula \eqref{eq:cq-Mellin-contour-shift} then
proves that all derivatives through order $j$ have limits at zero
given by the displayed even polynomial. These limits define a
$C^j$ extension: inductively, if $\psi^{(\ell)}$ and
$\psi^{(\ell+1)}$ have these limits, integrate
$\psi^{(\ell+1)}$ from a positive $\varepsilon$ to $x$ and
let $\varepsilon\downarrow0$; the resulting integral identity
proves differentiability at zero with the asserted derivative.
Because $j$ was arbitrary, the extension is smooth, has value
zero, all its odd derivatives vanish at zero, and its $2k$th
derivative is $(2k)!a_k$. Reflecting it evenly across zero
preserves smoothness: on the negative half-line its $j$th
derivative is $(-1)^j\psi^{(j)}(-x)$, and the one-sided
values coincide because the odd jets vanish. Together with
\eqref{eq:cq-Mellin-infinity-bound}, this proves
$\psi\in\mathcal S(\mathbb R)$ and the residue assertion.

For completeness the inversion just used also proves that the
Mellin transform of this reconstructed function is the prescribed
$M$. On a line $c>1$ put
$h(u)=e^{cu}\psi(e^u)$ and $m(v)=M(c+iv)$. The formula says
$h(u)=(2\pi)^{-1}\int m(v)e^{-ivu}\,dv$.
Here $m$ is continuous and integrable by the strip estimates,
and $h$ is integrable by the bounds at both ends just proved.
For $\varepsilon>0$, absolute Fubini and the Gaussian integral give
\[
 \int_{\mathbb R}h(u)e^{ivu}e^{-\varepsilon u^2/2}\,du
 =\int_{\mathbb R}m(w)
       \frac{e^{-(v-w)^2/(2\varepsilon)}}{\sqrt{2\pi\varepsilon}}
       \,dw.
\]
The right side tends to $m(v)$: split the integral at
$|v-w|=\delta$, use continuity on the inner interval, and use
the Gaussian tail and boundedness of $m$ on its complement.
The left side tends by dominated convergence to
$\int h(u)e^{ivu}\,du$. Substitution $x=e^u$ proves
$\mathcal M\psi(c+iv)=M(c+iv)$. Both sides are holomorphic
on $\operatorname{Re}z>-2$, so the identity theorem proves
their equality there and hence equality of their meromorphic
continuations. In particular
$\widehat\psi(0)=2\mathcal M\psi(1)=2M(1)=0$.
Reconstruction is injective because its output has Mellin transform
equal to its input. The forward map is injective as well, with
the following direct uniqueness proof. If the Mellin transform of
$\delta\in\mathcal S_0^{\rm ev}$ vanishes, the continuous
integrable function $g(u)=e^{cu}\delta(e^u)$ has Fourier transform
$\int g(u)e^{ivu}\,du=0$. Put
$G_\varepsilon(u)=(2\pi\varepsilon)^{-1/2}
e^{-u^2/(2\varepsilon)}$. Absolute Fubini applied to the Gaussian
integral gives
\[
 (g*G_\varepsilon)(u)=\frac1{2\pi}
 \int_{\mathbb R}\left(\int_{\mathbb R}g(w)e^{ivw}\,dw\right)
                   e^{-\varepsilon v^2/2}e^{-ivu}\,dv=0.
\]
As $\varepsilon\downarrow0$, the Gaussian convolution tends
to $g(u)$ at every $u$ by continuity, the unit total mass of
$G_\varepsilon$, and its vanishing tail mass. Thus $g=0$,
$\delta=0$ on the positive half-line, and evenness gives
$\delta=0$ everywhere. This proves that the two displayed maps
are mutual inverses, not only that reconstructed inputs have the
prescribed transform.

The inverse continuity can be read directly from the estimates.
For nonnegative integers $A,j$, choose $b=A+1$ and $K\geq1$
with $2K+1>j$. On $|x|\geq1$ use
\eqref{eq:cq-Mellin-infinity-bound}; on $|x|\leq1$ use
\eqref{eq:cq-Mellin-contour-shift}. They give
\begin{equation}\label{eq:cq-Mellin-inverse-seminorm}
 \sup_{x\in\mathbb R}(1+|x|)^A|\psi^{(j)}(x)|
 \leq2^A\left(
 I_{b,j}(M)+I_{-2K-1,j}(M)
 +\sum_{\substack{1\leq k\leq K\\2k\geq j}}
           \frac{(2k)!}{(2k-j)!}|a_k|\right).
\end{equation}
Each line integral on the right is bounded by a constant times
a strip seminorm with exponent greater than $j+1$ on a strip
containing its line and no pole on the line. Every one of the
finitely many residues is controlled by a strip seminorm too:
if $-2k$ belongs to $[a,b]$, then
\[
 |a_k|=
 \frac{|(P_{a,b}M)(-2k)|}
 {\left|\displaystyle\prod_{\substack{\ell\geq1,\ \ell\ne k\\
                         a-1\leq-2\ell\leq b+1}}(2\ell-2k)\right|}.
\]
Its denominator is a finite product of nonzero numbers. This
proves continuity without a uniform assumption on the entire
residue sequence and completes the characterization.

Apply this result to the exact arithmetic trace map already proved
in \eqref{eq:cq-Schwartz-trace}. For $F\in\cqE$ define the
meromorphic function
\[
 M_F(z)=\frac{F(i(z-\tfrac12))}{\zeta(z)}.
\]
Its definition uses meromorphic division, including the removable
values determined by any canceled zeros or pole. Then the precise
intersection of the full source trace image with its original
order-at-most-one entire-function set is
\begin{equation}\label{eq:cq-exact-Schwartz-trace-image}
 \begin{aligned}
 \mathcal A_{\rm Sch}
 &:=\{F_\psi:\psi\in\mathcal S_0^{\rm ev}\}\cap\cqE\\
 &=\{F\in\cqE:M_F\in\mathfrak M_0^{\rm ev}\}\\
 &=\left\{\Xi h:h\in\cqE,\quad
     A(z)h(i(z-\tfrac12))\in\mathfrak M_0^{\rm ev}\right\},\\
 A(z)&=\tfrac12z(z-1)\pi^{-z/2}\Gamma(z/2).
 \end{aligned}
\end{equation}
For the first equality after the definition, a trace has the
factorization and Mellin properties already proved. Conversely,
if $M_F\in\mathfrak M_0^{\rm ev}$, reconstruct the unique
$\psi\in\mathcal S_0^{\rm ev}$. In $\operatorname{Re}z>1$
the trace factorization gives
$F_\psi(i(z-\tfrac12))=\zeta(z)M_F(z)
=F(i(z-\tfrac12))$. Both $F_\psi$ and $F$ are entire;
their equality on this open set proves equality everywhere.
This argument retains the original $s$ coordinate and does not
assert that arbitrary Schwartz traces belong to $\cqE$.

For the last equality in
\eqref{eq:cq-exact-Schwartz-trace-image}, the previous
nontrivial-zero jet calculation and the proved entire division
theorem give $F=\Xi h$ with $h\in\cqE$ for every trace in
this intersection. Conversely the retained identity
$\Xi(i(z-\tfrac12))=\xi(1-z)=\xi(z)$ gives exactly
\[
 M_{\Xi h}(z)=
      \tfrac12z(z-1)\pi^{-z/2}\Gamma(z/2)
                      h(i(z-\tfrac12)),
\]
including its meromorphic continuation and all its local units.
The factor $A$ has only simple poles at $-2k$, $k\geq1$:
the gamma poles at $z=0$ and the factor $z$ cancel, the
remaining gamma poles occur at the displayed negative even
integers, and the factor $z-1$ gives $A(1)=0$.
Thus for an arbitrary $h\in\cqE$ the pole locations,
orders and the equality $M_{\Xi h}(1)=0$ are automatic;
the exact additional requirement is the family of strip bounds
in \eqref{eq:cq-Mellin-strip-seminorm}, with the actual factor
$A$ retained. No topology on the unnamed source closure is
identified by this characterization.

This criterion also retains every Taylor coefficient of its
arithmetic preimage. The trivial zero of $\zeta$ at $-2k$ is
simple, as follows directly from the retained completion.
The gamma recurrence gives a simple pole of $\Gamma(z/2)$
there with nonzero residue $2(-1)^k/k!$, while the other factors
in $A(z)$ are nonzero there. On the other hand
$\xi(-2k)=\xi(1+2k)>0$: all factors in the completion at $1+2k$
are positive, including the convergent positive Dirichlet series
for $\zeta(1+2k)$. Thus $\xi=A\zeta$ forces a zero of exactly
order one at $-2k$. Meromorphic division now gives
\begin{equation}\label{eq:cq-Schwartz-arithmetic-Taylor-data}
 \begin{aligned}
 \frac{\psi^{(2k)}(0)}{(2k)!}
 &=\frac{F(i(-2k-\tfrac12))}{\zeta'(-2k)}\\
 &=\frac{2k(2k+1)(-1)^k\pi^k}{k!}
                     h\bigl(-i(2k+\tfrac12)\bigr)
 \qquad(F=\Xi h\in\mathcal A_{\rm Sch}).
 \end{aligned}
\end{equation}
In the first line there is no extra factor $i$: the residue is
taken in the actual coordinate $z$ and the denominator has leading
term $\zeta'(-2k)(z+2k)$. For the second line the gamma
recurrence gives
$\operatorname*{Res}_{z=-2k}\Gamma(z/2)=2(-1)^k/k!$.
Indeed $\Gamma(w)=\Gamma(w+k+1)/
[w(w+1)\cdots(w+k)]$, whose numerator is $1$ at $w=-k$
and whose denominator has derivative $(-1)^kk!$ there;
substitution $w=z/2$ supplies the factor $2$. Multiplying
this exact residue by
$(-2k)(-2k-1)\pi^k/2$ and the value of the entire $h$
gives the asserted coefficient. Thus the inclusion criterion
also reconstructs the full even jet at the origin.

The inclusion $\mathcal A_{\rm Sch}\hookrightarrow\Xi\cqE$
is the actual identity on entire functions, with exact image
specified by \eqref{eq:cq-exact-Schwartz-trace-image}. It contains
$\cqC[s]\Xi$ because the retained finite Hermite trace families
span precisely that polynomial image. Consequently its closures
in each of the explicitly proved analytic topologies satisfy
\begin{equation}\label{eq:cq-Schwartz-image-analytic-closures}
 \overline{\mathcal A_{\rm Sch}}^{\,\gamma}
 =\overline{\mathcal A_{\rm Sch}}^{\,\kappa}
 =\Xi\cqE.
\end{equation}
Indeed the left image lies between $\cqC[s]\Xi$ and the
closed subspace $\Xi\cqE$, and the closures of the smaller
polynomial subspace were proved to equal that subspace.
This constructs the exact inclusion, its image criterion and both
analytic closures; it inserts no equality involving the unnamed
topology $\tau$.

\subsection{An exact exponential orbit and its Schwartz boundary}
\label{sec:cq-exponential-trace}
For each unchanged complex parameter $a=u+iv$ define
\[
 F_a(s)=\Xi(s)e^{as},\qquad
 C(z)=\tfrac12 z(z-1)\pi^{-z/2}\Gamma(z/2).
\]
Every $F_a$ belongs to $\Xi\cqE=N_\gamma=N_\kappa$.
Its membership in the actual unclosed Schwartz trace image has
the exact boundary
\begin{equation}\label{eq:cq-exponential-Schwartz-domain}
 F_a\in\mathcal F_\mu\mathcal E(\mathcal S_0^{\rm ev})
 \quad\Longleftrightarrow\quad |\operatorname{Re}a|<\pi/4.
\end{equation}
On this open strip its unique original Schwartz input is
\begin{equation}\label{eq:cq-exponential-Schwartz-input}
 \psi_a(x)=e^{-ia/2}
   \bigl(4\pi^2e^{-4ia}x^4-6\pi e^{-2ia}x^2\bigr)
                    e^{-\pi e^{-2ia}x^2}\quad(x\in\mathbb R).
\end{equation}
All coefficients and the factor $e^{-ia/2}$ are retained.
We prove both directions, including the two boundary lines.

Write $\Theta=x\partial_x$ and $G(x)=e^{-\pi x^2}$. Direct
differentiation gives
\[
 \psi_0=\Theta(\Theta+1)G
       =(4\pi^2x^4-6\pi x^2)e^{-\pi x^2}.
\]
For $\operatorname{Re}z>0$, substitution $y=\pi x^2$ in the
Euler gamma integral gives
$\mathcal M G(z)=\frac12\pi^{-z/2}\Gamma(z/2)$.
Integration by parts gives
$\mathcal M(\Theta f)(z)=-z\mathcal M f(z)$ when the boundary
terms vanish. They do for $G$ and the needed derivatives on
this half-plane. Consequently $\mathcal M\psi_0=C(z)$, with
the exact multiplier $(-z)(1-z)=z(z-1)$.

Suppose $|u|<\pi/4$ and put $\beta=e^{-2ia}$. Then
$\operatorname{Re}\beta=e^{2v}\cos(2u)>0$, and its logarithm
in the right half-plane is exactly $-2ia$. The Gaussian
integral with coefficient $\beta$ satisfies
\[
 \int_0^\infty e^{-\pi\beta x^2}x^z\,\frac{dx}{x}
 =\tfrac12\pi^{-z/2}e^{iaz}\Gamma(z/2)
 \quad(\operatorname{Re}z>0).
\]
Indeed both sides are holomorphic in the connected right half-plane
of $\beta$, locally dominated there by a positive Gaussian;
they agree for positive real $\beta$ by a real substitution, and
the identity theorem proves equality everywhere in that half-plane.
Applying $\Theta(\Theta+1)$ and retaining $e^{-ia/2}$ gives
\begin{equation}\label{eq:cq-exponential-Mellin}
 \mathcal M\psi_a(z)=C(z)e^{ia(z-1/2)}.
\end{equation}
The function in \eqref{eq:cq-exponential-Schwartz-input} is even,
Schwartz, and zero at the origin, since its Gaussian has positive
real coefficient and its polynomial starts at degree two.
At $z=1$, \eqref{eq:cq-exponential-Mellin} gives
$\widehat\psi_a(0)=2\mathcal M\psi_a(1)=0$.
Thus it belongs to the exact source domain.
The proved trace factorization gives
\[
 F_{\psi_a}(i(z-\tfrac12))
 =\zeta(z)C(z)e^{ia(z-1/2)}
 =\Xi(i(z-\tfrac12))e^{a i(z-1/2)}
 \quad(\operatorname{Re}z>1).
\]
Both sides are entire in $s=i(z-\tfrac12)$, so they agree
everywhere. This proves sufficiency.

We next derive the exact gamma-modulus formula used for necessity.
For real $y$, the absolutely convergent two Euler integrals,
followed by $r=x_1+x_2$, $q=x_1/(x_1+x_2)$, give
\[
 \begin{aligned}
 |\Gamma(1+iy)|^2
 &=\int_0^\infty\!\!\int_0^\infty
       e^{-x_1-x_2}(x_1/x_2)^{iy}\,dx_1\,dx_2\\
 &=\int_0^1(q/(1-q))^{iy}\,dq
 =\int_{\mathbb R}e^{iyw}\frac{e^w}{(1+e^w)^2}\,dw .
 \end{aligned}
\]
The middle equality uses $\int_0^\infty r e^{-r}\,dr=1$,
and the last substitution is $w=\log(q/(1-q))$.
For $y>0$, integrate
$f_y(w)=e^{iyw}e^w/(1+e^w)^2$ on the positively oriented
rectangle with vertices $-R,R,R+2\pi i,-R+2\pi i$.
The integrals on the vertical sides tend to zero, because the
rational exponential factor is $O(e^{-R})$ there, uniformly
in the height. The top horizontal integral is
$-e^{-2\pi y}$ times the bottom integral. The only pole in
the strip is at $w=\pi i$. Writing $w=\pi i+\eta$ gives
\[
 \frac{e^w}{(1+e^w)^2}
 =-\frac{1}{4\sinh^2(\eta/2)}
 =-\eta^{-2}+O(1),
\]
so the residue of $f_y$ is $-iy e^{-\pi y}$. The residue
theorem therefore gives
$(1-e^{-2\pi y})\int_{\mathbb R}f_y(w)\,dw
=2\pi y e^{-\pi y}$. By conjugation for negative $y$ and
continuity at zero, we have proved for every real $y$
\begin{equation}\label{eq:cq-Gamma-modulus}
 |\Gamma(1+iy)|^2=\frac{\pi y}{\sinh(\pi y)},
\end{equation}
with value $1$ at $y=0$. This proof uses no asymptotic remainder.

If a Schwartz input had trace $F_a$, its Mellin transform on
$\operatorname{Re}z>1$ would have to be the quotient
$C(z)e^{ia(z-1/2)}$, by the factorization and the nonvanishing
of $\zeta$ on that half-plane. On the fixed line $z=2-2iy$,
its exact squared modulus is
\begin{equation}\label{eq:cq-exponential-Mellin-modulus}
 \left|C(2-2iy)e^{ia(3/2-2iy)}\right|^2
 = e^{-3v+4uy}\,
       \frac{(1+y^2)(1+4y^2)y}{\pi\sinh(\pi y)}.
\end{equation}
To check every factor, $|z|^2=4(1+y^2)$,
$|z-1|^2=1+4y^2$, $|\pi^{-z/2}|^2=\pi^{-2}$,
and \eqref{eq:cq-Gamma-modulus} supplies the remaining gamma
factor. The real part of $ia(3/2-2iy)$ is $2uy-3v/2$.
At $y=0$ the right side has its removable value $e^{-3v}/\pi^2$.

For $y>0$, $\sinh(\pi y)<e^{\pi y}/2$.
If $u\geq\pi/4$, the right side of
\eqref{eq:cq-exponential-Mellin-modulus} is therefore greater
than
\[
 \frac{2e^{-3v}}{\pi}(1+y^2)(1+4y^2)y
                   e^{(4u-\pi)y},
\]
which tends to infinity, including $u=\pi/4$. If
$u\leq-\pi/4$, put $y=-Y$ and let $Y\to+\infty$; the
same inequality with $-u$ proves divergence. But for every
Schwartz input the defining integral gives the uniform bound
\[
 |\mathcal M\psi(2-2iy)|
 \leq\int_0^\infty|\psi(x)|x\,dx<\infty
 \qquad(y\in\mathbb R).
\]
This contradiction proves necessity and the exact open boundary.

Uniqueness of the input is also exact. If two source inputs have
the same trace, the Mellin transform of their difference vanishes
on $\operatorname{Re}z=2$, by the same factorization. Under
$x=e^w$ this is the Fourier transform, with the retained sign
change of frequency, of the integrable continuous function
$e^{2w}(\psi_1-\psi_2)(e^w)$. Fourier inversion gives zero
on the positive half-line, and evenness gives zero on the negative
half-line and at zero. Thus \eqref{eq:cq-exponential-Schwartz-input}
is the unique original input.

On the strip $|\operatorname{Re}a|<\pi/4$ this is a holomorphic
family in the Schwartz topology. Indeed on a compact subset of
the strip, $\operatorname{Re}(e^{-2ia})$ has a positive lower
bound. Every fixed derivative in $a$ and in $x$ is a polynomial
in $x$ with locally bounded coefficients times this same Gaussian.
It is bounded in each Schwartz seminorm by an integrable Gaussian
majorant; the Cauchy integral formula in $a$, with the same
seminorm bounds, proves complex differentiability in that topology.
Direct differentiation, with no change of the input coordinate,
gives the exact source-generator equation
\begin{equation}\label{eq:cq-exponential-source-generator}
 \partial_a\psi_a=-i(\Theta+\tfrac12)\psi_a,\qquad
 \mathcal F_\mu\mathcal E\bigl(-i(\Theta+\tfrac12)\psi_a\bigr)
                                      =sF_a(s).
\end{equation}
The second equality also follows from
$\mathcal M(-i(\Theta+\tfrac12)\psi)
=i(z-\tfrac12)\mathcal M\psi$ and the same factorization.
This operator is defined on every $\mathcal S_0^{\rm ev}$ input
and preserves that space. Multiplication by $x$ and differentiation
preserve Schwartz functions, and $x\psi'$ is even when $\psi$ is even.
Its value at zero is zero. Finally integration by parts on the
whole real line gives $\int x\psi'(x)\,dx=-\int\psi(x)\,dx$,
with vanishing boundary terms. Hence
$\int-i(\Theta+\tfrac12)\psi\,dx=(i/2)\int\psi\,dx=0$.
This supplies the exact morphism from the original scaling
generator to multiplication by the retained spectral coordinate.
In particular the inclusion of the actual Schwartz trace
intersection into $\Xi\cqE$ is strict: any real $a=\pi/4$
already gives an element of the latter with no Schwartz preimage.
Neither this explicit separating element nor this exact generator
map identifies the source closure in the unnamed topology $\tau$.

\subsection{An exact arithmetic inverse for every heated Hermite trace}
\label{sec:cq-moebius}
Let $\psi$ be any finite linear combination of the original
$\psi_\ell^\pm$, retaining all its coefficients, and write
$F_0=\mathcal F_\mu\mathcal E\psi=p\Xi$. For $x>0$ and
$t\in\cqC$ define
\begin{equation}\label{eq:cq-moebius-data}
 \begin{aligned}
 q(x)&=\sum_{n\geq1}\psi(nx),&
 g_t(x)&=e^{t(\log x)^2/4}\mathcal E\psi(x),\\
 q_t(x)&=x^{-1/2}g_t(x)=e^{t(\log x)^2/4}q(x),&
 \psi_t(x)&=\sum_{n\geq1}\mu(n)q_t(nx).
 \end{aligned}
\end{equation}
Here $\mu(1)=1$, $\mu(n)=0$ when a prime square divides $n$,
and $\mu(n)=(-1)^j$ for a product of $j$ distinct primes.
For the inverse $\psi_t$, the symbol $\mathcal E\psi_t$ means
the absolutely convergent trace-sum extension
$x^{1/2}\sum_{m\geq1}\psi_t(mx)$ on the positive half-line.
Its domain here consists of smooth functions with rapid decrease
at infinity; no Schwartz extension through zero is required.
The transform and inverse assertions, with their domains, are
\begin{equation}\label{eq:cq-moebius-inverse}
 \mathcal E\psi_t=g_t,\qquad
 \mathcal F_\mu g_t=U_t(p\Xi)=:F_t,\qquad
 \int_0^\infty\psi_t(x)x^z\,\frac{dx}{x}
       =\frac{F_t(i(z-\tfrac12))}{\zeta(z)}
       \quad(\operatorname{Re}z>1).
\end{equation}
The series for $\psi_t$ converges locally with all derivatives on
$(0,\infty)$, and $\psi_t$ and all its derivatives decrease faster
than every inverse power at infinity. No behavior at zero has been
assumed for this inverse.

We prove all analytic assertions in these formulas. A finite Hermite
combination and its Fourier transform are polynomials times the
original Gaussian $e^{-\pi x^2}$. Thus at infinity, for every
derivative order $j$, their derivatives are bounded by
$C_j(1+x)^{d_j}e^{-\pi x^2}$. For $x\geq1$, termwise derivatives
of $\sum\psi(nx)$ are bounded by a sum of
$C_j n^j(1+nx)^{d_j}e^{-\pi n^2x^2}$; splitting the Gaussian
into two equal exponential factors bounds that sum by
$C'_j(1+x)^{d_j}e^{-\pi x^2/2}$. This also justifies the
differentiation. Both zero conditions give, by the Poisson identity,
\[
 \mathcal E\psi(x)=\mathcal E\widehat\psi(1/x).
\]
Consequently, in $u=\log x$, the functions
$\mathcal E\psi(e^u)$ and all their $u$ derivatives have bounds
of the form $C\exp(C'|u|-c e^{2|u|})$ on the two ends.
Multiplication by $e^{tu^2/4}$, and any finite number of its
derivatives, preserves integrability against $e^{A|u|}$ for every
$A>0$, locally uniformly for $t$ in compact subsets of $\cqC$.
It follows that $\mathcal F_\mu g_t$ is entire in both $t,s$.
Since $D_s^{2k}e^{-isu}=(-1)^ku^{2k}e^{-isu}$, absolute
termwise integration of the time exponential gives exactly
\[
 \int_{\mathbb R}e^{tu^2/4}\mathcal E\psi(e^u)e^{-isu}\,du
 =\sum_{k\geq0}\frac{(-t/4)^k}{k!}
       D_s^{2k}(p(s)\Xi(s))=F_t(s).
\]
The heat lemma identifies the same series in $\gamma$; hence
$F_t\in\cqE$ with the original coordinate and coefficient.

The proved estimates imply that $q_t$ and all its $x$ derivatives
decrease faster than every power at infinity, and that $q_t$
decreases faster than every power at zero. On a compact interval
$x\in[\delta,L]\subset(0,\infty)$, choose $A>j+1$ to bound
the $j$th derivative of the $n$th inverse summand by
$C n^j(nx)^{-A}\leq C\delta^{-A}n^{j-A}$.
The finitely many arguments below one are harmless. This proves
local convergence with derivatives. On $x\geq1$ the same bound,
with arbitrarily large $A$, proves rapid decrease of each derivative.
For a fixed $x>0$ the double sum used below is absolutely bounded by
\[
 \sum_{m,n\geq1}|q_t(mnx)|
 =\sum_{k\geq1}d(k)|q_t(kx)|<\infty,
\]
where $d(k)\leq k$ and rapid decrease permits an exponent greater
than two. Unique prime factorization gives
$\sum_{n\mid k}\mu(n)=1$ for $k=1$ and
$\prod_{p\mid k}(1-1)=0$ for $k>1$. Therefore
\[
 \sum_{m\geq1}\psi_t(mx)
 =\sum_{k\geq1}\Bigl(\sum_{n\mid k}\mu(n)\Bigr)q_t(kx)
 =q_t(x).
\]
Multiplication by the retained factor $x^{1/2}$ proves the first
identity of \eqref{eq:cq-moebius-inverse}. It also proves uniqueness
of this inverse among functions decreasing faster than every power
at infinity: for any such $v$ with $\sum_m v(mx)=q_t(x)$, the
same absolutely convergent double sum applied to
$\sum_n\mu(n)\sum_m v(mnx)$ equals $v(x)$.

For $\sigma=\operatorname{Re}z>1$, absolute integration gives
\[
 \sum_{n\geq1}|\mu(n)|\int_0^\infty
       |q_t(nx)|x^\sigma\,\frac{dx}{x}
 \leq\Bigl(\sum_{n\geq1}n^{-\sigma}\Bigr)
       \int_0^\infty|q_t(y)|y^\sigma\,\frac{dy}{y}<\infty.
\]
The divisor identity also gives
$(\sum\mu(n)n^{-z})\zeta(z)=1$ by absolute multiplication of
the two Dirichlet series in this half-plane. Finally
\[
 \int_0^\infty q_t(x)x^z\,\frac{dx}{x}
 =\int_0^\infty g_t(x)x^{z-1/2}\,\frac{dx}{x}
 =F_t(i(z-\tfrac12)),
\]
because $-i\,i(z-\tfrac12)=z-\tfrac12$.
These equalities prove the last identity of
\eqref{eq:cq-moebius-inverse}, including its absolute domain.

At a nontrivial zero $\rho$ of $\zeta$ of multiplicity $m$, put
$\lambda=i(\rho-\tfrac12)$ and retain
$\zeta(z)=(z-\rho)^m a_\rho(z)$ with $a_\rho(\rho)\ne0$.
The retained cosine integral makes $\Xi$ even, so its definition gives
\[
 \Xi(i(z-\tfrac12))=\xi(1-z)=\xi(z).
\]
At a nontrivial zero, the factor
$z(z-1)\pi^{-z/2}\Gamma(z/2)/2$ is a nonvanishing holomorphic
unit. Thus this reflected coordinate sends the zero at $\rho$
to the zero at $\lambda$ with exactly the same multiplicity.
The meromorphic continuation of the last quotient has pole order
\begin{equation}\label{eq:cq-moebius-pole}
 \max\{m-\operatorname{ord}_\lambda F_t,0\}.
\end{equation}
For the zero seed $F_t=0$, the quotient is zero and its pole order
is zero; the displayed formula uses $\operatorname{ord}_\lambda0=+\infty$.
Indeed $s-\lambda=i(z-\rho)$, so its numerator has exactly
the original order $\operatorname{ord}_\lambda F_t$, with its
leading coefficient multiplied by the nonzero power of $i$.
The unit $a_\rho$ is retained in every principal coefficient.

For every fixed $t\ne0$ and every zero $\lambda$ of $\Xi$, there
is a polynomial $p$ for which $F_t(\lambda)\ne0$: otherwise all
$U_t(s^n\Xi)$ would vanish there, contrary to the no-common-zero
proof in Section~\ref{sec:cq-moving}. The exact polynomial-image
calculation supplies a finite original Hermite seed for this $p$.
For that seed the unique inverse $\psi_t$ in
\eqref{eq:cq-moebius-data} is unbounded near zero. To prove this
last assertion, boundedness near zero together with its proved rapid
decrease at infinity would make its Mellin integral holomorphic in
$\operatorname{Re}z>0$: on compact subsets, derivatives in $z$
are dominated by powers of $|\log x|x^{\sigma-1}$ near zero
with $\sigma>0$. Its equality in $\operatorname{Re}z>1$ would
then continue as a holomorphic function through every nontrivial
$\rho$, contradicting the pole of order $m$ just proved at
$\rho=\tfrac12-i\lambda$. Thus for each nonzero heat time
some actual finite Hermite trace has an explicit arithmetic inverse
which fails even boundedness at the origin. This proves a concrete
obstruction to Schwartz descent with a unique inverse and its poles;
it does not replace the unresolved closure in $\tau$ by a different
closure.


\subsection{Global meromorphic arithmetic heat transport}
\label{sec:cq-meromorphic-arithmetic-heat}
The original coordinate dictionary permits an exact heat action on
the entire arithmetic Mellin quotient, including its poles. Define
\begin{equation}\label{eq:cq-meromorphic-J}
 \begin{aligned}
 (\mathscr JF)(z)&=\frac{F(i(z-\tfrac12))}{\zeta(z)},
 &\mathfrak R_\zeta&=\mathscr J(\cqE),\\
 H_M(z)&=\zeta(z)M(z),
 &(\mathscr J^{-1}M)(s)&=
       \zeta(\tfrac12-is)M(\tfrac12-is).
 \end{aligned}
\end{equation}
All divisions and products in this display are meromorphic, with
their removable values retained. Thus $\mathfrak R_\zeta$ consists
exactly of meromorphic $M$ for which $H_M$ extends to an entire
function of order at most one. To verify the order statement in
both directions, the affine coordinate maps are
$z\mapsto i(z-\tfrac12)$ and $s\mapsto\tfrac12-is$.
Their moduli are bounded by $|z|+\tfrac12$ and
$|s|+\tfrac12$, respectively. The inequality
$(r+\tfrac12)^p\leq2^{p-1}(r^p+2^{-p})$ for every $p>1$
then transfers each of the defining order-class growth bounds
in both directions. The displayed maps are inverse because
$\tfrac12-i\,i(z-\tfrac12)=z$ and
$i(\tfrac12-is-\tfrac12)=s$.

Give $\mathfrak R_\zeta$ the transported growth topology
$\gamma_\zeta=\mathscr J_*\gamma$, namely the seminorms
\begin{equation}\label{eq:cq-meromorphic-seminorms}
 \|M\|_{\zeta;p,a}
 =\sup_{s\in\cqC}
   |\zeta(\tfrac12-is)M(\tfrac12-is)|e^{-a|s|^p},
 \qquad 1<p<2,\quad a>0.
\end{equation}
The canceled product is used at every apparent singularity.
Consequently $\mathscr J:(\cqE,\gamma)\to
(\mathfrak R_\zeta,\gamma_\zeta)$ is a linear homeomorphism;
completeness and the locally convex structure follow directly
from those of $(\cqE,\gamma)$. No topology has been assigned
to $\mathfrak R_\zeta$ merely by referring to unspecified
convergence of meromorphic functions.

Every $M\in\mathfrak R_\zeta$ is holomorphic away from zeros
of $\zeta$. At a zero of multiplicity $m$ its pole order is at
most $m$. At $z=1$, where $\zeta$ has its simple pole, $M$
is holomorphic and vanishes. These assertions follow by dividing
the entire $H_M$ by the actual local Laurent expansion of
$\zeta$. They concern all complex zeros, including the simple
trivial zeros, and do not discard the local units.

Put
\[
 \nabla_\zeta M=M'+\frac{\zeta'}{\zeta}M
             =\frac{H_M'}{\zeta},
 \qquad \mathcal L_\zeta M=\frac14\nabla_\zeta^2M.
\]
These are globally defined meromorphic operators on
$\mathfrak R_\zeta$; the right expression supplies every
removable value. With $D=\partial_s$, the exact conjugacies are
\begin{equation}\label{eq:cq-meromorphic-generator}
 \begin{aligned}
 \mathscr JD\mathscr J^{-1}&=-i\nabla_\zeta,\\
 \mathcal L_\zeta M
 &=\frac{1}{4\zeta}(\zeta M)''
  =\frac14\left(M''+2\frac{\zeta'}{\zeta}M'
                         +\frac{\zeta''}{\zeta}M\right),\\
 \mathscr J(-D^2/4)\mathscr J^{-1}&=\mathcal L_\zeta.
 \end{aligned}
\end{equation}
Indeed, $H_M(z)=F(i(z-\tfrac12))$ gives
$H_M'(z)=iF'(i(z-\tfrac12))$ and
$H_M''(z)=-F''(i(z-\tfrac12))$. Division by $\zeta$ proves
the first and last equalities. Expanding
$(\zeta M)''=\zeta M''+2\zeta'M'+\zeta''M$ proves every
coefficient in the middle one. In particular the coefficient
$\zeta''/\zeta$ includes both the derivative and the square
of $\zeta'/\zeta$; neither term has been omitted.

For every original complex time $t$, define
\begin{equation}\label{eq:cq-meromorphic-heat}
 \begin{aligned}
 \mathcal V_tM
 &=\mathscr JU_t\mathscr J^{-1}M
   =\frac{H_{M,t}}{\zeta},\\
 H_{M,t}(z)
 &=\exp\!\left(\frac t4\partial_z^2\right)H_M(z)
   =\sum_{k=0}^{\infty}
              \frac{t^k}{4^k k!}H_M^{(2k)}(z).
 \end{aligned}
\end{equation}
The last series is a definition by the proved entire-function
heat calculus, not an operation on unrestricted meromorphic
functions. More explicitly, its $k$th term equals
$(-t/4)^kF^{(2k)}(i(z-\tfrac12))/k!$.
It therefore has exactly the growth-topology convergence already
proved for $U_tF$, after the continuous affine change of argument.
In particular the numerator series converges locally uniformly in
$(t,z)$; after division by $\zeta$, it converges in every seminorm
\eqref{eq:cq-meromorphic-seminorms}, locally uniformly in complex $t$.
Each $\mathcal V_t$ is a continuous automorphism of
$(\mathfrak R_\zeta,\gamma_\zeta)$, with
\[
 \mathcal V_t\mathcal V_v=\mathcal V_{t+v},\qquad
 \mathcal V_t^{-1}=\mathcal V_{-t},\qquad
 \partial_t\mathcal V_tM=\mathcal L_\zeta\mathcal V_tM.
\]
These follow by conjugating the proved heat group and its
differentiated series. Equivalently
$\mathcal V_t=\exp(t\mathcal L_\zeta)$ in this precisely
specified topology. Near a zero of $\zeta$, local uniform
convergence of the numerators also gives convergence of each
Laurent coefficient and of the holomorphic remainders after the
finite principal parts are removed.

The original multiplication coordinate is the actual function
$i(z-\tfrac12)$, because
$\mathscr J(sF)=i(z-\tfrac12)\mathscr JF$. Thus the transported
operator $B_t=s-(t/2)D$ becomes
\begin{equation}\label{eq:cq-meromorphic-spectral-coordinate}
 \mathcal C_tM
 :=\mathscr JB_t\mathscr J^{-1}M
 =i(z-\tfrac12)M+\frac{it}{2}\nabla_\zeta M
 =\mathcal V_t\bigl(i(z-\tfrac12)\bigr)\mathcal V_{-t}M.
\end{equation}
In the last expression the function in parentheses denotes
multiplication. The equality follows by conjugating
$B_t=U_t sU_{-t}$ and proves the retained coordinate and sign.
It does not replace the original spectral value by $z$.

This also has an exact relation to the actual source quotient.
On the same underlying set $\mathfrak R_\zeta$ let
$\sigma_t=\mathscr J_*\tau_t$ and
$\mathfrak N_t=\mathscr JN_t$. The already proved closedness of
$N_t$ in $\tau_t$ gives closedness of $\mathfrak N_t$ in
$\sigma_t$. The map
\[
 [F]_{N_t}\longmapsto[\mathscr JF]_{\mathfrak N_t}
\]
is a homeomorphism of the corresponding quotient spaces, because
it and its inverse are induced by the defining homeomorphism
$\mathscr J$ of the ambient spaces. Equation
\eqref{eq:cq-meromorphic-spectral-coordinate} intertwines the
quotient spectral operators. For every scalar $\lambda$, it
therefore also intertwines $B_t-\lambda$ and
$\mathcal C_t-\lambda$, their kernels and images, and their
continuous inverses whenever those exist. Their spectra are the
same original arithmetic $s$-spectrum. This statement uses
$\sigma_t$, while the entire-time heat estimates above use
$\gamma_\zeta$; it inserts no equality between those topologies.

\subsection{Local pole cancellation with every arithmetic unit retained}
\label{sec:cq-meromorphic-local-heat}
Fix an actual nontrivial zero $\rho$ of $\zeta$, of its full
multiplicity $m=m_\rho$, and use its unchanged local coordinate
$w=z-\rho$. Write
\begin{equation}\label{eq:cq-meromorphic-local-data}
 \zeta(\rho+w)=w^m a_\rho(\rho+w),\qquad
 \frac1{a_\rho(\rho+w)}=\sum_{j\geq0}c_jw^j,\qquad
 H_M(\rho+w)=\sum_{n\geq0}b_nw^n.
\end{equation}
Here $a_\rho(\rho)\ne0$ and $c_0=1/a_\rho(\rho)$.
All series converge on a sufficiently small disk. The principal
coefficients of $M$ and its heat generator are exactly
\begin{equation}\label{eq:cq-meromorphic-principal-generator}
 \begin{aligned}
 [w^{-m+q}]M
   &=\sum_{n=0}^{q}c_{q-n}b_n,\\
 [w^{-m+q}]\mathcal L_\zeta M
   &=\frac14\sum_{n=0}^{q}
                c_{q-n}(n+2)(n+1)b_{n+2},
       &&0\leq q<m.
 \end{aligned}
\end{equation}
In particular neither $M$ nor its generator has a pole of order
greater than $m$, and its potential order-$m$ generator coefficient
is $H_M''(\rho)/(4a_\rho(\rho))=b_2/(2a_\rho(\rho))$.
If that coefficient vanishes, the remaining coefficients in the
display determine the full lower pole order without suppressing
the unit.

There is a direct cancellation calculation for the apparently
higher poles in \eqref{eq:cq-meromorphic-generator}. Put
$v(w)=H_M(\rho+w)/a_\rho(\rho+w)$, so $M=w^{-m}v$,
and temporarily write $a=a_\rho(\rho+w)$ with primes denoting
the same original $z$ derivative. Then
\[
 \begin{aligned}
 M'&=w^{-m}v'-mw^{-m-1}v,\\
 M''&=w^{-m}v''-2mw^{-m-1}v'
                         +m(m+1)w^{-m-2}v,\\
 \frac{\zeta'}{\zeta}&=\frac m w+\frac{a'}a,\\
 \frac{\zeta''}{\zeta}
   &=\frac{m(m-1)}{w^2}+\frac{2m}{w}\frac{a'}a
                              +\frac{a''}a.
 \end{aligned}
\]
In $M''+2(\zeta'/\zeta)M'+(\zeta''/\zeta)M$, the
coefficient multiplying $w^{-m-2}v$ is
$m(m+1)-2m^2+m(m-1)=0$. The coefficients multiplying
$w^{-m-1}v'$ are $-2m+2m=0$, and those multiplying
$w^{-m-1}(a'/a)v$ are $-2m+2m=0$. The complete remaining
expression is
\[
 w^{-m}\left(v''+2\frac{a'}a v'+\frac{a''}a v\right)
 =\frac{w^{-m}}a(av)''=\frac{H_M''}{\zeta}.
\]
This proves the cancellation of both possible extra pole orders
with every cross term and local unit present. Expanding the final
numerator and $1/a$ proves
\eqref{eq:cq-meromorphic-principal-generator}.

The same calculation holds at every time. Set
\[
 b_n(t)=\frac{H_{M,t}^{(n)}(\rho)}{n!},
 \qquad \lambda_\rho=i(\rho-\tfrac12).
\]
The globally convergent numerator heat series gives the exact
entire functions of $t$
\begin{equation}\label{eq:cq-meromorphic-time-jets}
 \begin{aligned}
 b_n(t)
 &=\sum_{k=0}^{\infty}
       \frac{t^k(n+2k)!}{4^k k!n!}b_{n+2k}(0)
  =\frac{i^n}{n!}(D^nU_tF)(\lambda_\rho),\\
 \partial_tb_n(t)&=\frac{(n+2)(n+1)}4 b_{n+2}(t),
   &&F=\mathscr J^{-1}M.
 \end{aligned}
\end{equation}
Differentiation and evaluation are continuous on the entire
growth space by the proved derivative estimates. They therefore
justify termwise differentiation and evaluation of this series,
including local uniform convergence in $t$. If
$\ell_q(t)=[w^{-m+q}]\mathcal V_tM$, then
\begin{equation}\label{eq:cq-meromorphic-principal-evolution}
 \ell_q(t)=\sum_{n=0}^q c_{q-n}b_n(t),\qquad
 \partial_t\ell_q(t)=\frac14\sum_{n=0}^q
           c_{q-n}(n+2)(n+1)b_{n+2}(t),\quad 0\leq q<m.
\end{equation}
Thus the evolution of the arithmetic principal parts is a
calculation from the full entire numerator. It is not a closed
system on only the finitely many coefficients $\ell_q$; the
displayed higher numerator jets are part of the exact map.

For an initial $M$ holomorphic at $\rho$, write
$M(\rho+w)=M_0+M_1w+O(w^2)$. The formula for the generator
then gives the complete possible principal part
\begin{equation}\label{eq:cq-meromorphic-Schwartz-instantaneous-poles}
 \operatorname{PP}_\rho(\mathcal L_\zeta M)
 =\frac{m(m-1)M_0}{4w^2}
  +\frac{m(m+1)M_1+2m(a_\rho'/a_\rho)(\rho)M_0}{4w}.
\end{equation}
For $m=1$ the order-two term is zero. For $m\geq2$ this
principal part vanishes exactly when $M_0=M_1=0$. For $m=1$
it vanishes exactly when $(a_\rho M)'(\rho)=0$. These
equivalences follow from the nonzero constants $m$, $m-1$
in the first case and $a_\rho(\rho)\ne0$ in both cases.
They describe the exact local derivative obstruction, including
multiple zeros, rather than assuming their multiplicity is one.

\subsection{The calculated domain for heat to retain an original Schwartz input}
\label{sec:cq-meromorphic-Schwartz-heat-domain}
Let $\mathfrak M_0^{\rm ev}$, $P_{a,b}$, and $q_{a,b,N}$
be the explicitly proved meromorphic Mellin image and strip
seminorms in \eqref{eq:cq-Mellin-strip-seminorm}. Its intersection
\[
 \mathfrak A_\zeta=\mathfrak R_\zeta\cap\mathfrak M_0^{\rm ev}
                    =\mathscr J\mathcal A_{\rm Sch}
\]
is the exact original Schwartz trace domain in this presentation,
by \eqref{eq:cq-exact-Schwartz-trace-image}. For each unchanged
complex time $t$, its complete domain of retained Schwartz inputs
is the following explicitly calculated subspace:
\begin{equation}\label{eq:cq-meromorphic-Schwartz-heat-domain}
 \begin{aligned}
 \mathfrak D_t=\biggl\{M\in\mathfrak A_\zeta:\quad
 &\sum_{k=0}^{\infty}
       \frac{t^k(n+2k)!}{4^k k!n!}b_{n+2k}^{\rho}(0)=0
       \quad(\rho\in\mathcal Z,\ 0\leq n<m_\rho),\\
 &\sup_{\substack{a\leq\sigma\leq b\\v\in\mathbb R}}
    (1+|v|)^N
    \left|P_{a,b}(\sigma+iv)
       \frac{H_{M,t}(\sigma+iv)}{\zeta(\sigma+iv)}\right|
       <\infty\\
 &\hspace{35mm}\text{for every real }a\leq b
          \text{ and every integer }N\geq0\biggr\}.
 \end{aligned}
\end{equation}
Here $\mathcal Z$ is the actual set of nontrivial zeros,
$b_n^\rho(0)=H_M^{(n)}(\rho)/n!$, and $H_{M,t}$ is the
globally convergent series in \eqref{eq:cq-meromorphic-heat}.
The removable values in the strip expression are the values
forced by the preceding zero-jet equalities and by the factors
$P_{a,b}$. No unknown initial input, omitted pole cancellation,
or assumed decay estimate occurs in this definition.

The exact characterization and reconstruction are
\begin{equation}\label{eq:cq-meromorphic-Schwartz-heat-equivalence}
 \begin{aligned}
 \mathfrak D_t
   &=\mathfrak A_\zeta\cap\mathcal V_{-t}\mathfrak A_\zeta,\\
 M\in\mathfrak D_t\quad&\Longleftrightarrow\quad
 U_t\mathscr J^{-1}M\in\mathcal A_{\rm Sch},\\
 \psi_t(x)
   &=\frac1{2\pi}\int_{\mathbb R}
     \frac{H_{M,t}(c+iv)}{\zeta(c+iv)}x^{-c-iv}\,dv
          \quad(x>0,\ c>1),\\
 \psi_t(0)&=0,\qquad \psi_t(-x)=\psi_t(x).
 \end{aligned}
\end{equation}
In the equivalence $M$ is taken in its declared initial domain
$\mathfrak A_\zeta$; the reconstruction formula is for
$M\in\mathfrak D_t$. We now prove that the conditions in
\eqref{eq:cq-meromorphic-Schwartz-heat-domain} are both necessary
and sufficient, rather than leaving them as proof obligations.

The first family of equalities says exactly
$H_{M,t}^{(n)}(\rho)=0$ for $0\leq n<m_\rho$, by
\eqref{eq:cq-meromorphic-time-jets}. Division by
$(z-\rho)^{m_\rho}a_\rho(z)$ therefore shows that it is
equivalent to holomorphy of $\mathcal V_tM$ at every nontrivial
zero. At every trivial zero $-2k$, the zero of $\zeta$ is
simple, so $H_{M,t}/\zeta$ automatically has at most a simple
pole there. At $z=1$ it automatically is holomorphic and zero,
because the numerator is entire and $1/\zeta$ has a simple zero.
At all remaining points it is holomorphic. Thus the first family
is equivalent to all pole and value conditions in the definition
of $\mathfrak M_0^{\rm ev}$. The second family is exactly that
definition's complete finite-strip decay requirement, with its
original polynomial and margin. This proves
\[
 M\in\mathfrak D_t
 \quad\Longleftrightarrow\quad
 M\in\mathfrak A_\zeta\text{ and }
            \mathcal V_tM\in\mathfrak M_0^{\rm ev}.
\]
Since $\mathcal V_tM$ already belongs to $\mathfrak R_\zeta$,
the last membership is precisely membership in
$\mathfrak A_\zeta$. The group inverse proves the first
identity in \eqref{eq:cq-meromorphic-Schwartz-heat-equivalence},
and the proved exact trace-image characterization proves the
second. The inverse Mellin theorem then proves that the displayed
integral converges with the required differentiated estimates,
is independent of $c>1$, gives the unique original even Schwartz
input with both zero conditions, and satisfies
\[
 \mathcal F_\mu\mathcal E\psi_t
       =\mathscr J^{-1}\mathcal V_tM
       =U_t\mathscr J^{-1}M.
\]
These conclusions use the already proved inverse theorem with
exactly its hypotheses verified above.

All Taylor data at the original input origin are also explicit:
for $M\in\mathfrak D_t$ and every $k\geq1$,
\begin{equation}\label{eq:cq-meromorphic-heated-Taylor-data}
 \frac{\psi_t^{(2k)}(0)}{(2k)!}
   =\operatorname*{Res}_{z=-2k}\frac{H_{M,t}(z)}{\zeta(z)}
   =\frac{H_{M,t}(-2k)}{\zeta'(-2k)}.
\end{equation}
The first equality is the proved Mellin residue theorem for the
reconstructed input. The second divides by the actual leading
term $\zeta'(-2k)(z+2k)$; in particular it introduces no factor
$i$ into the residue coordinate $z$. The right side is defined
as an entire function of $t$ for every initial
$M\in\mathfrak R_\zeta$, even at times outside
$\mathfrak D_t$ when it cannot be interpreted as a Schwartz
input's Taylor coefficient.

At $t=0$ the definition gives all of $\mathfrak A_\zeta$.
For every $t\ne0$ it is a proper subspace. Indeed fix any
nontrivial zero $\rho$ and its retained spectral coordinate
$\lambda_\rho=i(\rho-\tfrac12)$. The proved absence of a
common evaluation zero among $U_t(\cqC[s]\Xi)$ gives a
polynomial $p$ for which
$U_t(p\Xi)(\lambda_\rho)\ne0$. The actual finite Hermite
trace theorem places $p\Xi$ in $\mathcal A_{\rm Sch}$.
For $M=\mathscr J(p\Xi)$ the computed numerator at time $t$
is nonzero at $\rho$, so its quotient by $\zeta$ has a pole
of the full order $m_\rho$ there. This violates the first
family in \eqref{eq:cq-meromorphic-Schwartz-heat-domain} and
proves $M\notin\mathfrak D_t$. The argument is an obstruction
in the actual unclosed Schwartz trace map and its calculated
meromorphic heat domain; it makes no identification of the
unnamed source closure.

\subsection{The exact persistent heat core and all possible return times}
\label{sec:cq-persistent-heat-core}
Retain the original functions, coordinate, and analytic relation
$I=\Xi\cqE=N_\gamma=N_\kappa$. The obstruction can be computed
for an entire orbit, not only for a chosen generator or one time.
Define
\begin{equation}\label{eq:cq-persistent-cores}
 \begin{aligned}
 I_{\rm even}^{\infty}
   &=\{f\in\cqE:\ D^{2k}f\in I\text{ for every integer }k\geq0\},\\
 I_{\rm heat}^{\infty}
   &=\{f\in\cqE:\ U_tf\in I\text{ for every }t\in\cqC\}.
 \end{aligned}
\end{equation}
Then the complete calculation is
\begin{equation}\label{eq:cq-persistent-core-zero}
 I_{\rm even}^{\infty}=I_{\rm heat}^{\infty}=\{0\}.
\end{equation}
More strongly, for every set $T\subset\cqC$ with a finite
accumulation point,
\begin{equation}\label{eq:cq-accumulating-heat-times}
 \bigcap_{t\in T}U_{-t}I=\{0\}.
\end{equation}
We give the complete proof, including the arithmetic obstruction
to an exceptional periodic function.

First suppose $f\in I_{\rm even}^{\infty}$. At every actual
zero $\lambda$ of $\Xi$ of multiplicity $m_\lambda$, its
Taylor coefficients satisfy
\[
 D^{2k+j}f(\lambda)=0
 \qquad(k\geq0,\ 0\leq j<m_\lambda).
\]
If any $m_\lambda\geq2$, the choices $j=0,1$ cover all
nonnegative derivative orders. The entire Taylor series at that
unchanged $\lambda$ is zero, so $f=0$. It remains to prove
the result when all the actual zeros are simple; this is the
other case of the proof, not a simplicity hypothesis.
For each of these zeros, the vanishing of every even derivative
and the entire Taylor expansion give the exact reflection identity
\begin{equation}\label{eq:cq-heat-core-reflections}
 f(2\lambda-s)=-f(s)\qquad(s\in\cqC).
\end{equation}
In particular $f$ is odd about each actual zero without changing
the spectral coordinate. Fix one zero $\lambda_0$.
Composing the reflections at $\lambda_0$ and at $\lambda$ gives
\begin{equation}\label{eq:cq-heat-core-periods}
 f\bigl(s+2(\lambda-\lambda_0)\bigr)=f(s)
 \qquad(s\in\cqC,\ \Xi(\lambda)=0).
\end{equation}

We prove the needed period assertion explicitly. Let $f$ be a
nonconstant entire function and put
\[
 P(f)=\{p\in\cqC:\ f(s+p)=f(s)\text{ for all }s\in\cqC\}.
\]
It is an additive subgroup, closed because pointwise limits of
periods remain periods by continuity of $f$. It has no sequence
of nonzero periods tending to zero: such a sequence $p_n$ would
give
$f'(s)=\lim_n(f(s+p_n)-f(s))/p_n=0$ at every $s$.
Thus some disk about zero contains no other period.
Two periods $p,q$ linearly independent over $\mathbb R$ are
also impossible. Every complex number has the form $ap+bq$
with real $a,b$. Subtracting the integer parts of $a,b$
puts it in the compact parallelogram
$\{xp+yq:0\leq x,y\leq1\}$ without changing the value of $f$.
The resulting global bound and the Cauchy estimate
$|f'(s)|\leq\sup_{\cqC}|f|/R$ for every $R>0$ imply that
$f$ is constant, a contradiction.

If $P(f)$ contains a nonzero period $p$, all its periods
therefore lie on the real line $\mathbb Rp$. The set
$L=\{a\in\mathbb R:ap\in P(f)\}$ is a closed additive subgroup
of $\mathbb R$, contains $1$, and has a positive gap about zero.
Let $d=\inf(L\cap(0,\infty))$. This number is positive; a
sequence approaching the infimum and closedness show $d\in L$.
For any positive $a\in L$, subtract
$\lfloor a/d\rfloor d$. The remainder lies in $L\cap[0,d)$
and must be zero by minimality. Negative elements are handled
by changing sign. Hence $L=d\mathbb Z$ and
\[
 P(f)=\omega\mathbb Z,\qquad\omega=dp\ne0.
\]
This proves the full assertion about the period group used here.

Return to \eqref{eq:cq-heat-core-periods}. A constant satisfying
\eqref{eq:cq-heat-core-reflections} is zero. If $f\ne0$, it is
therefore nonconstant. The set of distinct zeros of $\Xi$ is
infinite, as proved above from its retained product and completion,
so it contains a zero different from $\lambda_0$.
The period group consequently contains a nonzero period and is
$\omega\mathbb Z$ by the preceding argument. Every actual zero
must then lie in the single arithmetic progression
\begin{equation}\label{eq:cq-heat-core-zero-progression}
 \lambda\in\lambda_0+\tfrac12\omega\mathbb Z.
\end{equation}
The number of points in this progression in the disk $|s|\leq r$
is at most $1+4r/|\omega|$: their integer indices lie in an
interval of length at most $4r/|\omega|$, by projecting onto
the line of the progression. This bound does not assume that
$\lambda_0$ or $\omega$ is real.

Here is a direct contradiction with the actual completed zeta
function, using no zero-counting asymptotic. Retain its exact
even Hadamard product from the source,
\[
 \Xi(s)=\Xi(0)\prod_{\alpha\in\Lambda_+}
                    (1-s^2/\alpha^2),
\]
where one member of each pair $\{\lambda,-\lambda\}$ is
retained with its original multiplicity. In the case currently
under consideration every multiplicity is one.
Set $N_+(u)=\#\{\alpha\in\Lambda_+:|\alpha|\leq u\}$.
There is a number $r_*>0$ such that $N_+(u)=0$ for $u<r_*$,
since $\Xi(0)\ne0$ and its zeros are discrete.
The progression bound implies
\[
 N_+(u)\leq Au\quad(u\geq r_*),\qquad
 A=\frac1{r_*}+\frac4{|\omega|}.
\]
For $r>0$, the elementary identity
\[
 \log(1+r^2/v^2)
   =\int_v^\infty\frac{2r^2}{u(u^2+r^2)}\,du
 \qquad(v>0)
\]
and integration of nonnegative summands give
\[
 \begin{aligned}
 \sum_{\alpha\in\Lambda_+}\log(1+r^2/|\alpha|^2)
 &=\int_0^\infty\frac{2r^2N_+(u)}{u(u^2+r^2)}\,du\\
 &\leq2Ar^2\int_0^\infty\frac{du}{u^2+r^2}
 =\pi Ar .
 \end{aligned}
\]
The source product, first for its finite partial products and
then by its locally uniform convergence, therefore gives
\begin{equation}\label{eq:cq-progression-exponential-bound}
 |\Xi(s)|\leq|\Xi(0)|e^{\pi A|s|}\qquad(s\in\cqC).
\end{equation}
No lower bound on an individual product factor was used.
But at the original points $s=-i(2n-\tfrac12)$, $n\geq2$,
the unchanged completion is
\[
 \Xi\bigl(-i(2n-\tfrac12)\bigr)
 =\xi(2n)
 =n(2n-1)\pi^{-n}(n-1)!\zeta(2n).
\]
Every factor is positive and $\zeta(2n)>1$ from its positive
Dirichlet series. For $n\geq4$, let $b_n=\lfloor n/2\rfloor$.
The last $b_n$ factors of $(n-1)!$ are at least $b_n$, so
\[
 \log\left|\Xi\bigl(-i(2n-\tfrac12)\bigr)\right|
 \geq b_n\log b_n-n\log\pi .
\]
Dividing by $2n-\tfrac12$ makes the right side tend to
$+\infty$, in contradiction with
\eqref{eq:cq-progression-exponential-bound}.
This excludes every nonzero $f$ in the remaining case and
proves $I_{\rm even}^{\infty}=\{0\}$.

For the time statements, for every actual zero and every one
of its required derivative orders, define
\[
 h_{\lambda,j}(t)=(D^jU_tf)(\lambda),
 \qquad 0\leq j<m_\lambda.
\]
These are entire functions of the original complex time by the
proved heat series and derivative estimates. Its differentiated
coefficients at zero are exactly
\begin{equation}\label{eq:cq-heat-return-time-jets}
 h_{\lambda,j}^{(k)}(0)=(-1/4)^kD^{2k+j}f(\lambda).
\end{equation}
If $U_tf\in I$ for all $t$ in a set with a finite accumulation
point, every $h_{\lambda,j}$ vanishes there, and the identity
theorem makes every one of them identically zero.
Equation \eqref{eq:cq-heat-return-time-jets} then puts $f$
in $I_{\rm even}^{\infty}$ and proves
\eqref{eq:cq-accumulating-heat-times}.
Taking all complex times proves
$I_{\rm heat}^{\infty}=\{0\}$.
Conversely, membership of $I_{\rm even}^{\infty}$ would make
all Taylor coefficients of each $h_{\lambda,j}$ zero and hence
put every $U_tf$ in $I$, by the already proved entire division
criterion. This also verifies directly the equality of the two
cores, with both directions and the original coefficient retained.

In particular for each nonzero $f\in\cqE$ the exact return set
\begin{equation}\label{eq:cq-heat-return-set}
 \mathcal T_I(f)=\{t\in\cqC:\ U_tf\in I\}
   =\bigcap_{\lambda,j}\{t:h_{\lambda,j}(t)=0\}
\end{equation}
is closed and has no finite accumulation point. Thus it is a
discrete set, possibly empty. More explicitly, at least one
$h_{\lambda,j}$ is a nonzero entire function, since otherwise
the core calculation would give $f=0$; its isolated zero set
already contains $\mathcal T_I(f)$. Every compact disk meets
this return set finitely, so the set is countable by exhaustion
by the integer-radius disks.
Since the exact unclosed Schwartz trace intersection
$\mathcal A_{\rm Sch}$ is a subspace of $I$,
\begin{equation}\label{eq:cq-Schwartz-discrete-returns}
 \mathcal T_{\rm Sch}(f)
   :=\{t\in\cqC:\ U_tf\in\mathcal A_{\rm Sch}\}
   \subseteq\mathcal T_I(f)
 \qquad(f\ne0)
\end{equation}
also has no finite accumulation point and is closed and
discrete. The closedness follows because a subset of a set
with no finite accumulation point is closed in $\cqC$.
Its exact membership, including all strip conditions and
reconstruction of the original input, was computed in
Section~\ref{sec:cq-meromorphic-Schwartz-heat-domain}.
Consequently no nonzero original Schwartz trace follows a heat
orbit through original Schwartz traces on any nontrivial time
interval, even though its entire heat evolution exists for all
complex times.

There is also an exact invariant-subspace conclusion.
If $W$ is a linear subspace of $I$ with $D^2W\subseteq W$,
then $D^{2k}f\in I$ for every $f\in W$ and every $k\geq0$.
The proved core calculation forces $W=\{0\}$.
The same conclusion holds for a subspace of $I$ invariant
under all $U_t$. These conclusions do not assume that $W$
is closed or that its topology has any particular form.
They apply to the analytic relation $I$ and to its unclosed
Schwartz trace subspace, and do not identify either with
the actual source closure $\cqN$ in the still unnamed $\tau$.

\subsection{The original four-term coefficient and its arithmetic return times}
\label{sec:cq-xi4-heat-returns}
Retain the companion mechanism's actual coefficient operator
on the unchanged entire variable $s$:
\begin{equation}\label{eq:cq-xi4-original-operator}
 \begin{aligned}
 (T_Df)(s)={}&-\frac{59}{275184}f(s+3)
             -\frac1{336}f'(s+3)\\
            &+\frac1{1296}f(s+\tfrac92)
             +\frac3{132496}f(s+\tfrac{13}2),\\
 K_t(s)={}&T_DZ_t(s),\qquad Z_t=U_t\Xi=8H_t(2s).
 \end{aligned}
\end{equation}
Here $T_D$ denotes this given four-term operator; it is not
replaced by a differential truncation. We prove directly from
its complete formula that $K_0\ne0$, that $K_t=U_tK_0$,
and that its arithmetic return times have the precise restrictions
just established.

Translations by each of the three original shifts act
continuously on $(\cqE,\gamma)$. For a general fixed shift $h$,
the unchanged growth estimate is
\[
 \|f(\,\cdot+h)\|_{p,a}
 \leq e^{a|h|^p}\|f\|_{p,a/2^{p-1}},
 \qquad 1<p<2,\quad a>0,
\]
because $|s+h|^p\leq2^{p-1}(|s|^p+|h|^p)$.
Differentiation is continuous by Lemma~\ref{thm:cq-heat}.
Thus every displayed term, and $T_D$ itself, maps $\cqE$
continuously to $\cqE$. Differentiation commutes with each
fixed translation. Applying these continuous maps to the
convergent series for $U_t$ therefore proves term by term
\begin{equation}\label{eq:cq-xi4-heat-intertwiner}
 T_DU_t=U_tT_D,\qquad K_t=U_tK_0.
\end{equation}
In particular this calculation includes the derivative coefficient
$-1/336$ at the shift $3$, as well as the three function-value terms.

Nonvanishing of this particular coefficient is not assumed.
Use the retained positive kernel $\Phi$ from
\eqref{eq:cq-realheat}, with its original formula
\[
 \Phi(u)=\sum_{n\geq1}
 (2\pi^2n^4e^{9u}-3\pi n^2e^{5u})e^{-\pi n^2e^{4u}}
 \qquad(u\geq0).
\]
The substitution $v=2u$ in the original cosine integral gives
the exact entire Fourier representation
\begin{equation}\label{eq:cq-xi4-full-Fourier-kernel}
 \begin{aligned}
 Z_t(s)&=2\int_{\mathbb R}
          \Phi(|v|/2)e^{tv^2/4}e^{isv}\,dv,\\
 K_t(s)&=2\int_{\mathbb R}
          \Phi(|v|/2)m(v)e^{tv^2/4}e^{isv}\,dv,\\
 m(v)&=\left(-\frac{59}{275184}-\frac{iv}{336}\right)e^{3iv}
          +\frac{e^{(9/2)iv}}{1296}
          +\frac{3e^{(13/2)iv}}{132496}.
 \end{aligned}
\end{equation}
The coefficient $2$ is required: first $8\,du=4\,dv$ on
the positive half-line, then the even extension of the cosine
integral halves that coefficient. A shift by $h$ multiplies
the kernel by $e^{ihv}$, and the derivative multiplies it by
$iv$ for this original positive Fourier exponent.
These two facts prove the second identity from the first,
with the precise multiplier in the third.
Every interchange and derivative is justified by the previously
proved bound
$|\Phi(u)|\leq Ce^{9u}e^{-ce^{4u}}$ on $u\geq0$.
On a compact set of complex $t,s$, this double exponential
dominates the factors $e^{(\operatorname{Re}t)v^2/4}$,
$e^{-\operatorname{Im}(s)v}$, and every polynomial factor
from $m(v)$ and the derivatives. Thus the integrals converge
absolutely and locally uniformly with all needed derivatives.

The multiplier has the exact value
\begin{equation}\label{eq:cq-xi4-nonzero-multiplier}
 m(0)=-\frac{59}{275184}+\frac1{1296}
                    +\frac3{132496}=\frac{127}{219024}>0.
\end{equation}
Also
\[
 \Phi(0)=\sum_{n\geq1}
       \pi n^2(2\pi n^2-3)e^{-\pi n^2}>0;
\]
every summand is positive and the series is absolutely
convergent. The continuous integrable function
$q(v)=2\Phi(|v|/2)m(v)$ therefore has
$q(0)=2\Phi(0)\,127/219024>0$.
If $K_0$ were identically zero, its Fourier transform in
\eqref{eq:cq-xi4-full-Fourier-kernel} would vanish at every
real $s$. The Gaussian Fourier uniqueness argument already
proved in Section~\ref{sec:cq-Schwartz-Mellin-image} applies
to this continuous integrable $q$: absolute Fubini makes
every Gaussian convolution $q*G_\varepsilon$ zero, while
these convolutions tend to $q$ pointwise. This would force
$q(0)=0$, a contradiction. Hence $K_0\ne0$.
The invertibility of $U_t$ in
\eqref{eq:cq-xi4-heat-intertwiner} also proves $K_t\ne0$
for every original complex time $t$.

The previous theorem now applies to this exact coefficient,
not to an arbitrary replacement. In particular the sets
\begin{equation}\label{eq:cq-xi4-arithmetic-return-times}
 \begin{aligned}
 \mathcal T_I(K_0)
    &=\{t:\ T_D(8H_t(2\,\cdot))\in\Xi\cqE\},\\
 \mathcal T_{\rm Sch}(K_0)
    &=\{t:\ T_D(8H_t(2\,\cdot))\in\mathcal A_{\rm Sch}\}
       \subseteq\mathcal T_I(K_0)
 \end{aligned}
\end{equation}
are closed discrete subsets of $\cqC$.
Their complete individual membership conditions are still the
original multiplicity jets and, for the second set, the
finite-strip Mellin conditions proved above.
At any time outside $\mathcal T_I(K_0)$ there is an actual
nontrivial zero $\rho$ of multiplicity $m_\rho$ at which
\[
 \frac{K_t(i(z-\tfrac12))}{\zeta(z)}
\]
has a pole of positive order
$m_\rho-\operatorname{ord}_{i(\rho-1/2)}K_t$.
Indeed nonmembership in $I$ is exactly failure of at least
one required zero jet, by the proved division theorem; local
division by $(z-\rho)^{m_\rho}a_\rho(z)$ gives this pole order
with the unchanged coordinate and unit.
The full principal coefficients, including their time evolution,
are those of Section~\ref{sec:cq-meromorphic-local-heat}
with the actual entire numerator
$K_t(i(z-\tfrac12))$ inserted.

Finally retain the original incidence factor $a_{ef}$ in
$a_{ef}K_t$. For $a_{ef}\ne0$, scalar multiplication preserves
both subspaces $I$ and $\mathcal A_{\rm Sch}$ and has inverse
multiplication by $a_{ef}^{-1}$, so it leaves both return sets
unchanged and multiplies every principal coefficient by $a_{ef}$.
For $a_{ef}=0$, the coefficient is identically zero and belongs
to both spaces at every time. This proves both incidence cases
without suppressing that factor.
These are exact arithmetic relations for the mechanism coefficient.
They do not equate $I$ or $\mathcal A_{\rm Sch}$ with $\cqN$
in the unnamed source topology.

\section{The inverse-root cover, the branch locus, and arithmetic fibres}
\label{sec:cq-cover}
The retained polynomial inverse-fibre slice is
$P_{a,0,0}(r)=-2r^2-2a$. Its root equation is $a=-r^2$.
The arithmetic coordinate is $a=2s$, so
\begin{equation}\label{eq:cq-cover}
 s=-r^2/2.
\end{equation}
The finite free extension of the actual quotient is
\[
 \widehat Q=\cqC[r]\otimes_{\cqC[s]}\cqQ
           =\cqQ\oplus r\cqQ,
 \qquad
 \mathcal R=\begin{pmatrix}0&-2S\\1&0\end{pmatrix},
 \qquad -\tfrac12\mathcal R^2=\begin{pmatrix}S&0\\0&S\end{pmatrix}.
\]
The even--odd decomposition of a polynomial in $r$ proves freeness
with basis $(1,r)$ and uniqueness coefficient by coefficient.
Multiplying $q_0+rq_1$ by $r$ proves the displayed matrices;
the finite direct-sum topology makes every induced map continuous.
The involution is $(q_0,q_1)\mapsto(q_0,-q_1)$; the trace is
$\operatorname{Tr}(q_0+rq_1)=2q_0$, with kernel $r\cqQ$.
It is $\cqC[s]$-linear. It does not provide a $\cqC[r]$-module
quotient with this kernel unless $S=0$: multiplying $rq$ by $r$
gives $-2Sq$ in the even summand. This proves the exact limitation,
rather than suppressing the odd component. The same construction
at time $t$ uses $S_t$ and is conjugated by
$\operatorname{diag}(\mathcal U_t,\mathcal U_t)$.

The original $s$-spectrum is unchanged by the extension: an inverse
of $S-\lambda$ duplicates to a diagonal inverse, and conversely
projection of an inverse of the doubled matrix to its first component,
following inclusion of that component, gives an inverse of
$S-\lambda$. Both directions are continuous. Moreover
\begin{equation}\label{eq:cq-rspec}
 \cqSpec(\mathcal R)
 =\{\rho\in\cqC:-\rho^2/2\in\cqSpec(S)\}.
\end{equation}
To prove this without an unsupported spectral-mapping citation,
solve $(\mathcal R-\rho)(u,v)=(a,b)$. The second row gives
$u=b+\rho v$ and the first becomes
$-2(S+\rho^2/2)v=a+\rho b$. A continuous inverse of
$S+\rho^2/2$ therefore gives a continuous inverse matrix.
Conversely the right-hand side $(a,0)$ and the second projection
of an inverse matrix supply an inverse of $-2(S+\rho^2/2)$,
with both identities checked by the same two rows.

Before any quotient by $\Xi$, let $\mathcal O_s$ be the sheaf of
holomorphic functions in $s$ and put
$B=\mathcal O_s[r]/(r^2+2s)$, identified with
$\mathcal O_s\oplus r\mathcal O_s$. On the punctured base $s\ne0$,
implicit differentiation gives
$dr/ds=-1/r=r/(2s)$. Thus the actual lifted derivative is
\begin{equation}\label{eq:cq-connection}
 \nabla_s\binom{f_0}{f_1}
 =\binom{f_0'}{f_1'+f_1/(2s)},\quad
 R=\begin{pmatrix}0&-2s\\1&0\end{pmatrix},\quad
 [\nabla_s,R]=-R^{-1}.
\end{equation}
Indeed with $A=\operatorname{diag}(0,1/(2s))$,
$R'+AR-RA=\left(\begin{smallmatrix}0&-1\\1/(2s)&0\end{smallmatrix}\right)
=-R^{-1}$. A second application gives
\begin{equation}\label{eq:cq-connection2}
 \nabla_s^2=\operatorname{diag}
 \left(D^2,D^2+s^{-1}D-\frac1{4s^2}\right).
\end{equation}
Equivalently, on a function of $r$ the derivative is
$-r^{-1}\partial_r$ and the heat generator is
$-\tfrac14r^{-2}\partial_r^2+\tfrac14r^{-3}\partial_r$.
Squaring $-r^{-1}\partial_r$ by the product rule proves this
with both signs preserved.

The branch extension is exact. The meromorphic connection sends
$B$ into $s^{-1}B$ and the connection one-form has a simple pole
with residue $\operatorname{diag}(0,1/2)$. Its maximal domain for
holomorphic output at $s=0$ is
\begin{equation}\label{eq:cq-branchdomain}
 \operatorname{Dom}_{\rm hol}\nabla_s
       =\mathcal O_s\oplus r s\mathcal O_s.
\end{equation}
In fact the odd coefficient is $f_1'+f_1/(2s)$; its only possible
pole has coefficient $f_1(0)/2$, proving necessity and sufficiency.
For the square, write $f_1=\sum_{n\geq0}a_ns^n$. Its image has
coefficient $(n^2-1/4)a_ns^{n-2}$. The negative powers vanish
exactly when $a_0=a_1=0$, so
\[
 \operatorname{Dom}_{\rm hol}\nabla_s^2
       =\mathcal O_s\oplus r s^2\mathcal O_s.
\]
For $k$ iterates the odd coefficient is
$\prod_{h=0}^{k-1}(n+1/2-h)a_ns^{n-k}$, whose prefactor is
never zero for integral $n$. Thus its domain is
$\mathcal O_s\oplus r s^k\mathcal O_s$.
No derivation $B\to B$ extending $\partial_s$ exists at the branch:
derivation of $r^2+2s=0$ would give $2r\delta(r)+2=0$,
forcing $r$ to be a unit although it vanishes at the branch point.
The logarithmic derivative $s\nabla_s$
does extend to all of $B$, with matrix
$sD+\operatorname{diag}(0,1/2)$.

These are sheaf and representative calculations; they do not assert
that $\nabla_s$ descends through the frozen arithmetic relation.
For a holomorphic local multiplier $R_0(s)$ the relation vector
$(R_0\Xi,0)$ acquires $(R_0'\Xi+R_0\Xi',0)$, and the odd
relation $(0,R_0\Xi)$ acquires
$(0,R_0'\Xi+R_0\Xi'+R_0\Xi/(2s))$.
The second iterate includes the full $2R_0'\Xi'$ term, as well as
the odd $s^{-1}D-1/(4s^2)$ terms of
\eqref{eq:cq-connection2}. Jet extension
\eqref{eq:cq-jetextension}, or the moving closed quotient
\eqref{eq:cq-Qt}, retains these relations exactly.

Finally, the branch point is not an arithmetic spectral point at
$t=0$. The source gives $\Xi(0)\ne0$, so its Proposition 3.6
makes $S$ continuously invertible; then $\mathcal R$ is invertible
by \eqref{eq:cq-rspec}. In the analytic divisor realization the
stronger local statement follows from $\Xi$ being a holomorphic
unit on a disk about zero: the quotient stalk is zero there.
For real $t$, $Z_t(0)>0$ directly from the retained kernel:
\[
 2\pi^2n^4e^{9u}-3\pi n^2e^{5u}
 =\pi n^2e^{5u}(2\pi n^2e^{4u}-3)>0\quad(u\geq0),
\]
and the positive integral converges by the bound already proved.
Thus the actual real heat divisor never meets the cover branch
point $s=0$. On a simple zero curve its retained lift satisfies
$r^2=-2\lambda$, $r'=-\lambda'/r$; together with
\eqref{eq:cq-zero} these are the original-coordinate and cover
equations, including the unit contribution. A nonreal $r$ alone
does not specify whether $\lambda=-r^2/2$ is nonreal.

\section{Exact ingress from compact Weil tests}
\label{sec:cq-tests}
The accompanying retained Weil-test calculation supplies the concrete
space $\mathcal E=C_c^\infty(\mathbb R;\cqC)$ with coordinate
operator $T=-i\partial_u$ and Fourier transform
\[
 (\mathscr Ff)(s)=\int_{\mathbb R}f(u)e^{-ius}\,du.
\]
Here we give the complete heat ingress and its inverse domains;
none of the following claims requires identifying the source
quotient topology with a test-space topology. Integration by parts
gives $\mathscr F(Tf)=s\mathscr Ff$ with no endpoint term.
If $\operatorname{supp}f\subset[-L,L]$, then
\[
 |\mathscr Ff(s)|\leq\|f\|_1e^{L|\operatorname{Im}s|},
 \qquad
 \|\mathscr Ff\|_{p,a}\leq
 \|f\|_1\sup_{r\geq0}e^{Lr-ar^p}<\infty.
\]
Thus the Fourier transform maps the actual test space into the
original entire-function set, continuously in $\gamma$ on every
fixed-support smooth-function space and hence in its usual
inductive-limit topology.

Define, for all complex $t$,
\begin{equation}\label{eq:cq-testheat}
 (G_tf)(u)=e^{tu^2/4}f(u),\qquad
 U_t\mathscr Ff=\mathscr F G_tf,\qquad
 T_t=G_tTG_{-t}=T+\frac{it}{2}u.
\end{equation}
The multiplication $G_t$ is a continuous automorphism on each
fixed-support smooth-function space, with inverse $G_{-t}$.
Indeed each derivative of the smooth multiplier is bounded on the
fixed compact support, and Leibniz's finite sum bounds every output
derivative seminorm by finitely many input ones. The same argument
for the inverse proves the assertion. The exponential series and
every $u$-derivative converge uniformly on a fixed compact set,
locally uniformly in $t$; consequently the family is entire in time
in that test topology. In the Fourier integral,
$D^{2k}\mathscr Ff=\mathscr F((-1)^ku^{2k}f)$.
The compact support permits absolute interchange with the
exponential series, proving the middle identity with the exact
factor $1/4$. Finally differentiation of $e^{-tu^2/4}f$ gives
the displayed $T_t$, and
\[
 \mathscr F(T_tf)=B_t\mathscr Ff,
 \qquad D\mathscr Ff=\mathscr F(-iu f).
\]
Both equations follow directly from the integral, including their
signs. The test-space version of the retained cover is therefore
\[
 \mathcal R_t^{\rm test}
   =\begin{pmatrix}0&-2T_t\\1&0\end{pmatrix},\qquad
 (\mathcal R_t^{\rm test})^2=-2\operatorname{diag}(T_t,T_t).
\]
Componentwise Fourier transformation intertwines this matrix with
the representative matrix using $B_t$, and then with the actual
quotient matrix using $S_t$.

To retain every possible arithmetic relation in this ingress, put
\begin{equation}\label{eq:cq-testkernel}
 K=\{f\in\mathcal E:\mathscr Ff\in\cqN\},\qquad
 K_t=\{f\in\mathcal E:\mathscr Ff\in N_t\}.
\end{equation}
Then $K_t=G_tK$: by \eqref{eq:cq-testheat},
$\mathscr Ff\in U_t\cqN$ exactly when
$\mathscr F(G_{-t}f)\in\cqN$. Thus there are injective linear
maps onto their explicitly specified images
\[
 \mathcal E/K\longrightarrow\cqQ,\quad[f]\longmapsto[\mathscr Ff],
 \qquad
 \mathcal E/K_t\longrightarrow Q_t,\quad[f]\longmapsto[\mathscr Ff],
\]
and the square formed by $G_t$ and $\mathcal U_t$ commutes.
The kernel statements follow exactly from
\eqref{eq:cq-testkernel}; no claim that $K=0$ in the unnamed
source topology is inserted. To make the ingress continuous without
that identification, give $\mathcal E$ the join $\eta$ of its
usual test topology with the inverse-image topology of $\tau$
under $\mathscr F$. Then $K$ is closed and the first quotient map
is continuous. Give the time-$t$ test space the transported
topology $(G_t)_*\eta$; the second map and the commutative square
are continuous by \eqref{eq:cq-testheat}. This specifies the exact
topologies rather than claiming an unproved topological embedding
of a quotient image equipped with its subspace topology.

There is a fully explicit inverse-domain calculation in the same
test coordinates. For $\lambda\in\cqC$, $k\geq1$, the image is
\begin{equation}\label{eq:cq-testmoments}
 (T_t-\lambda)^k\mathcal E
 =\left\{g\in\mathcal E:
 \int_{\mathbb R}v^j e^{-tv^2/4-i\lambda v}g(v)\,dv=0
                  \quad(0\leq j<k)\right\}.
\end{equation}
On that precise image the inverse is
\begin{equation}\label{eq:cq-testinverse}
 [(T_t-\lambda)^{-k}g](u)
 =\frac{i^ke^{tu^2/4+i\lambda u}}{(k-1)!}
 \int_{-\infty}^{u}(u-v)^{k-1}
            e^{-tv^2/4-i\lambda v}g(v)\,dv.
\end{equation}
To prove both directions, write
$f=e^{tu^2/4+i\lambda u}h$.
Direct differentiation gives
$(T_t-\lambda)f=e^{tu^2/4+i\lambda u}(-i h')$.
Iteration reduces the inverse equation to a $k$th primitive, giving
\eqref{eq:cq-testinverse} because $(-i)^ki^k=1$.
For $g$ supported in $[a,b]$, that formula vanishes below $a$.
Above $b$, expansion of $(u-v)^{k-1}$ shows it vanishes identically
exactly when all $k$ moments in \eqref{eq:cq-testmoments} vanish;
the binomial coefficients and the nonzero exponential do not change
that equivalence. Its smoothness follows by repeated differentiation
of the integral; smoothness of $g$ with compact support gives
matching derivatives at both endpoints. Necessity for any compact
solution also follows by integrating $h^{(k)}$ against
$1,v,\ldots,v^{k-1}$ and integrating by parts $k$ times.
Uniqueness follows since a compactly supported solution of
$h^{(k)}=0$ is a polynomial of degree at most $k-1$ which vanishes
on a ray, and hence is zero. This proves the exact image, inverse,
kernel zero, and preservation of the original support interval.

In particular the inverse of $\mathcal R_t^{\rm test}$ is
\[
 (g_0,g_1)\longmapsto(g_1,-\tfrac12T_t^{-1}g_0),
 \qquad
 \operatorname{Dom}(\mathcal R_t^{\rm test})^{-1}
 =\left\{(g_0,g_1):\int e^{-tv^2/4}g_0(v)\,dv=0\right\}.
\]
Substitution into its two rows proves both inverse identities.
Thus the branch-locus obstruction in the concrete test realization
is a retained Gaussian-weighted zero-moment condition. The exact
arithmetic quotient instead has invertible $S_t$ and root operator,
as proved above. Their ingress map and its kernels explain precisely
how these statements coexist; one operator's inverse domain is not
silently assigned to the other realization.

\subsection{The source Sobolev cokernel and an exact common test-space map}
\label{sec:cq-sobolev-comparison}
The Sobolev construction cited immediately after Connes--Marcolli
Proposition 2.24 is specified in Chapter 2, Section 9. It is a different,
explicitly defined quotient, and supplies a useful exact comparison.
We retain its weight, its generator, and its limited set of detected jets.
The primary locators are Connes--Marcolli, printed pp.407--408 and
416--420 \cite{cqCM}, and Connes, Section III, pp.10--15 and
Appendix I, pp.62--67 of the author PDF \cite{cqConnesSelecta}.
Here and below $\delta>1$.

For the rational field, choose the splitting
$C_{\mathbb Q}=C_{\mathbb Q,1}\times\mathbb R_+^*$ used in the source,
give the compact factor Haar mass one, and put $x=e^u$ on the second
factor. Its trivial-character component is the Hilbert space
\[
 B_\delta=L^2\bigl(\mathbb R,(1+u^2)^{\delta/2}\,du\bigr),
 \qquad
 \|f\|_\delta^2=\int_{\mathbb R}|f(u)|^2(1+u^2)^{\delta/2}\,du.
\]
The source adelic map is
\[
 E_{\rm ad}\psi(g)=|g|^{1/2}\sum_{q\in\mathbb Q^*}\psi(qg),
 \quad
 \psi\in\mathcal S(\mathbb A_{\mathbb Q})_0
 =\{\psi:\psi(0)=0,\ \int\psi\,dx=0\}.
\]
Let $P_1$ denote averaging over $C_{\mathbb Q,1}$, and define
\[
 R_\delta=\overline{P_1E_{\rm ad}
             (\mathcal S(\mathbb A_{\mathbb Q})_0)}^{\,B_\delta},
 \qquad H_{\delta,1}=B_\delta/R_\delta.
\]
These are exactly the trivial-character relation space and cokernel
of the source. Indeed the source completion uses the norm of its image,
so the extended isometry has closed range equal to the closure of the
uncompleted image. Compact averaging commutes with that equivariant
closed range. If $b$ is a limit of image vectors and satisfies $P_1b=b$,
averaging an approximating sequence gives a sequence in the averaged
image converging to $b$. This proves that the displayed closure is the
trivial-character part of the full relation space, without replacing
any source relation by an orthogonal complement as a definition.

We import explicitly the following result of the source proof:
under the bilinear pairing $\langle f,\eta\rangle=\int f(u)\eta(u)\,du$,
the annihilator of $R_\delta$ in
$B_{-\delta}=L^2(\mathbb R,(1+u^2)^{-\delta/2}\,du)$ is the closed
linear span of
\begin{equation}\label{eq:cq-sobolev-source-annihilator}
 u^a e^{itu},\qquad
 \zeta(\tfrac12+it)=0,\quad t\in\mathbb R,\quad
 0\leq a<m_{-t},\quad a<\frac{\delta-1}{2}.
\end{equation}
Here $m_{-t}$ is the original multiplicity of the zero $-t$ of
$\Xi$, equal to that of $\tfrac12+it$ by the retained reflected
Mellin coordinate. The source proves the finite combinations after a
compact Fourier cutoff and proves that the cutoff vectors are dense
by uniformly bounded approximate units. Thus the closed span, not
only its finite-dimensional projections, is the imported statement.
No such annihilator statement is imported for the unnamed topology
$\tau$ of the entire-function ring.

The precise cutoff can also be checked directly from the defined norm:
\[
 \|u^a e^{itu}\|_{-\delta}^2
   =\int_{\mathbb R}|u|^{2a}(1+u^2)^{-\delta/2}\,du<\infty
 \quad\Longleftrightarrow\quad 2a-\delta<-1.
\]
Near zero the integrand is locally integrable for every integer $a\geq0$;
for $|u|\geq1$ it is bounded above and below by positive constants
times $|u|^{2a-\delta}$. Integration of this power proves the strict
inequality and the logarithmic divergence at equality. Hence set
\begin{equation}\label{eq:cq-sobolev-jetcount}
 q_\delta=\left\lceil\frac{\delta-1}{2}\right\rceil,
 \qquad n_{\delta,\lambda}=\min(m_\lambda,q_\delta).
\end{equation}
There are exactly $n_{\delta,\lambda}$ allowed degrees at a real zero
$\lambda$. The original Connes Theorem 1, Section III, p.13, states
the equivalent multiplicity bound $n<(1+\delta)/2$, $n\leq m_\lambda$.
The inspected Connes--Marcolli printed p.408 instead displays
$n<1+\delta/2$ in Theorem 2.47; its proof, pp.419--420,
equations (2.444) and (2.448), gives precisely the stricter bound
proved above. This records the primary discrepancy rather than
silently identifying the two printed formulas. No assumption that
an actual zeta zero has multiplicity greater than one is used.

Let $\mathcal E=C_c^\infty(\mathbb R;\mathbb C)$ retain its usual
fixed-support smooth-function inductive-limit topology, and let
$\mathscr Ff(s)=\int f(u)e^{-ius}\,du$ as above. Define
\begin{align}
 K_{\rm full}
   &=\{f\in\mathcal E:\ D^a\mathscr Ff(\lambda)=0
               \text{ for every zero }\lambda\text{ of }\Xi,
                         \ 0\leq a<m_\lambda\},
       \label{eq:cq-sobolev-full-kernel}\\
 K_\delta
   &=\{f\in\mathcal E:\ D^a\mathscr Ff(\lambda)=0
               \text{ for every real zero }\lambda\text{ of }\Xi,
                         \ 0\leq a<n_{\delta,\lambda}\}.
       \label{eq:cq-sobolev-kernel}
\end{align}
The already proved entire division theorem identifies
$K_{\rm full}=\{f:\mathscr Ff\in\Xi\cqE\}$ exactly.

\begin{cqtheorem}[Exact dense test ingress into the source Sobolev cokernel]
\label{thm:cq-sobolev-ingress}
The map $j_\delta:\mathcal E\to H_{\delta,1}$ sending a test function
to its Hilbert quotient class is continuous, has dense image, and
has kernel exactly $K_\delta$. It induces the continuous map
\begin{equation}\label{eq:cq-sobolev-quotient-map}
 J_\delta:\mathcal E/K_{\rm full}\longrightarrow H_{\delta,1},
 \qquad [f]\longmapsto[f],
 \qquad \ker J_\delta=K_\delta/K_{\rm full}.
\end{equation}
Its image is dense. These statements neither assert surjectivity
nor identify the topology of its image with a quotient topology.
For the source translation generator $D_1$, this ingress retains
the original spectral coordinate through
\begin{equation}\label{eq:cq-sobolev-generator}
 j_\delta(Tf)=iD_1j_\delta(f),
 \qquad T=-i\partial_u,\qquad \mathscr F(Tf)=s\mathscr Ff.
\end{equation}
\end{cqtheorem}
\begin{proof}
The inclusion $\mathcal E\to B_\delta$ is continuous on every
fixed-support space: on $[-L,L]$,
\[
 \|f\|_\delta^2\leq
   \left(\int_{-L}^L(1+u^2)^{\delta/2}\,du\right)\sup|f|^2.
\]
The inductive-limit universal property and the continuous Hilbert
quotient map give continuity of $j_\delta$.
Smooth compactly supported functions are dense in $B_\delta$:
truncate an arbitrary $f$ outside $[-L,L]$, whose squared norm tail
tends to zero by integrability; on the enlarged interval
$[-L-1,L+1]$, the weight is bounded above and below by positive
constants, and convolution with a smooth compactly supported
approximate identity approximates the truncation in ordinary $L^2$
and therefore in the weighted norm. These convolutions have compact
support and are smooth. Consequently their quotient images are dense.

The continuous dual of $B_\delta$ is exactly $B_{-\delta}$ under
the displayed bilinear pairing. To check this, multiply by
$(1+u^2)^{\delta/4}$ to identify $B_\delta$ isometrically with
ordinary $L^2$; the Riesz representation there, with complex
conjugation absorbed in the representing function, gives exactly
the inverse-weight dual and its norm. Since $R_\delta$ is closed,
a vector belongs to it exactly when every member of its annihilator
vanishes on that vector. One direction is the definition. For the
other, a vector outside a closed Hilbert subspace has a nonzero
orthogonal component, and its inner product defines a continuous
separating functional. By the explicitly imported closed-span
description \eqref{eq:cq-sobolev-source-annihilator}, it is enough
to test the listed log-monomial characters. But
\[
 D^a\mathscr Ff(-t)=(-i)^a\int u^a e^{itu}f(u)\,du,
 \qquad
 \int u^a e^{itu}f(u)\,du=i^aD^a\mathscr Ff(-t).
\]
Every factor $i^a$ is nonzero. The coordinate $\lambda=-t$ and
the proved count \eqref{eq:cq-sobolev-jetcount} now give
$\ker j_\delta=K_\delta$. Every equation defining $K_\delta$
appears among the equations defining $K_{\rm full}$, proving
$K_{\rm full}\subseteq K_\delta$. Thus
\eqref{eq:cq-sobolev-quotient-map} is well defined and continuous
by the quotient universal property; changing the representative
changes its image by zero exactly for a member of $K_\delta$.
This proves the kernel and the dense-image assertion.

Finally the source left translation is
$(V(e^\varepsilon)f)(u)=f(u-\varepsilon)$, so its infinitesimal
generator on $\mathcal E$ is $-\partial_u$. The difference quotient
converges in $B_\delta$ because the supports remain in a fixed
compact set and the smooth-function difference quotient converges
uniformly there. Passing to the quotient proves that
$j_\delta(f)$ belongs to the domain of $D_1$ and
$D_1j_\delta(f)=j_\delta(-f')$. Multiplication by $i$ gives the
first equation of \eqref{eq:cq-sobolev-generator}. Integration by
parts in the compactly supported Fourier integral gives the second.
In particular the source eigenparameter $it$ becomes $-t$ under
$iD_1$, exactly the same $s$ coordinate used in the jet equations;
no sign change is absorbed into the source operator.
\end{proof}

For completeness the unspecified source topology also has an exact
comparison, requiring no unproved inclusion of its kernel in
$K_\delta$. Retain
\[
 K_\tau=\{f\in\mathcal E:\mathscr Ff\in\cqN\},
 \qquad L_{\tau,\delta}=K_\tau\cap K_\delta.
\]
Give $\mathcal E$ the join of its usual test topology, the
inverse-image topology $\mathscr F^{-1}\tau$, and the seminorm
$\|\cdot\|_\delta$. The last seminorm is already continuous in
the test topology, but is displayed to specify every target map.
Then $K_\tau$, $K_\delta$ and $L_{\tau,\delta}$ are closed, and
there are continuous maps
\begin{equation}\label{eq:cq-sobolev-source-span}
 \cqQ\ \longleftarrow\
 \mathcal E/L_{\tau,\delta}
 \ \longrightarrow\ H_{\delta,1},
 \qquad
 [f]\longmapsto[\mathscr Ff],\quad [f]\longmapsto[f].
\end{equation}
Their respective kernels are exactly
$K_\tau/L_{\tau,\delta}$ and $K_\delta/L_{\tau,\delta}$;
the right image is dense. Indeed each composite on $\mathcal E$
is continuous by definition, vanishes on the intersection, and
has the asserted kernel before quotienting. The quotient universal
property proves the assertions, and density follows from the same
set of compact tests as above. The map to the product of the two
targets is injective because its kernel is precisely the
intersection already divided out. No topological embedding of that
image is inferred from injectivity. This is a proved relation to
the actual source quotient while its named topology remains retained.
The comparison also intertwines the coordinate operators: $T$
preserves $K_\tau$ because $\mathscr F(Tf)=s\mathscr Ff$ and
$M\cqN\subset\cqN$; it preserves $K_\delta$ because
\[
 D^a(sF)(\lambda)=\lambda D^aF(\lambda)
                              +aD^{a-1}F(\lambda),
\]
where the second term is zero for $a=0$ and otherwise uses an
already imposed lower jet. Thus $T$ descends to the common span,
whose left map intertwines it with $S$ and whose right map
intertwines it with $iD_1$ on the proved domain of the ingress.

There is also a literal compatibility of the original finite Hermite
traces with the adelic map. For the source's even $\psi$, put
$\psi_{\rm ad}=\psi\otimes\mathbf1_{\widehat{\mathbb Z}}$.
It lies in $\mathcal S(\mathbb A_{\mathbb Q})_0$ because its value
and additive integral are the two original zero constraints.
On the idele represented by $(x,1,1,\ldots)$, a rational number
belongs to every $\mathbb Z_p$ exactly when it is an integer:
write it in lowest terms and use any prime dividing its denominator
for the reverse implication. Therefore
\[
 E_{\rm ad}\psi_{\rm ad}(x)
   =x^{1/2}\sum_{n\in\mathbb Z\setminus\{0\}}\psi(nx)
   =2\mathcal E\psi(x).
\]
Every norm-one rational idele class has a representative consisting
of a real sign and finite units. Indeed, for an idele $g$ put
$q=\prod_p p^{v_p(g_p)}\in\mathbb Q_{>0}$, a finite product.
Then $g_p/q\in\mathbb Z_p^*$ for every prime, while the
norm-one identity gives $|g_\infty|=q$. Dividing by the diagonal
rational $q$ gives the asserted representative. The function
is unchanged by the finite units; evenness removes the real sign.
It therefore lies in the entire trivial-character component.
The Fourier transform in $u=\log x$ is thus exactly twice the
original source trace, retaining the factor $2$. In particular the
finite Hermite generators yield $2R_\ell^\pm(s)\Xi(s)$ in this
explicit realization. Their logarithmic representatives lie in
every $B_\delta$ by their proved rapid decay. This proves a concrete
arithmetic compatibility as well as the quotient comparison; it
does not identify a Sobolev topology with the topology on all of
$\cqE$.

\section{The retained spatial bilaplacian and the actual source quotient}
\label{sec:cq-bilaplacian-source}

The original polynomial spatial map and its inverse-fibre trajectory are
\begin{gather*}
 Q(x,y,w)=1+xy,\qquad
 \sigma(x,y,w)=\tfrac12\{Q^3w+y^2Q(4+3xy)\},\\
 \gamma(z)=(z^{-1},-3z/2,13z^2/2),\quad z>0.
\end{gather*}
Here the Euclidean coordinates are exactly $(x,y,w)$, and
$\Delta=\partial_x^2+\partial_y^2+\partial_w^2$.
The symbol $\gamma(z)$ in this section denotes this spatial curve;
the topology $\gamma$ on $\cqE$ retains its earlier definition.
Substitution gives $Q(\gamma(z))=-1/2$ and
$\sigma(\gamma(z))=-z^2/8$.
The pullback
$\sigma^*:\cqO(\cqC)\longrightarrow\cqO(\cqC^3)$,
$h\longmapsto h\circ\sigma$, is a continuous injective algebra
homomorphism for compact convergence. Indeed,
$\eta(s)=(0,0,2s)$ satisfies $\sigma\circ\eta(s)=s$,
so $\eta^*\sigma^*=1$. A compact set has compact image under each
polynomial map, proving continuity of both pullbacks by their
defining supremum seminorms. This also proves that the inverse on
the image is continuous. Restriction to $\cqE$ with its growth
topology is continuous because that topology dominates compact
convergence.

\subsection{Every retained coefficient and the residue functional}
Write $g=\nabla\sigma$, $H=D^2\sigma$, $N=g\cdot g$,
and $L=\Delta\sigma$. For every entire $h$,
\begin{align}
 \Delta(h\circ\sigma)&=Nh''(\sigma)+Lh'(\sigma),\notag\\
 \Delta^2(h\circ\sigma)&=
 N^2h^{(4)}(\sigma)+(4g^{\mathsf T}Hg+2NL)h^{(3)}(\sigma)\notag\\
 &\quad+(2\|H\|_{\rm F}^2+4g\cdot\nabla L+L^2)h''(\sigma)
       +(\Delta L)h'(\sigma).
 \label{eq:cq-bilaplacian-chain}
\end{align}
These formulas use the bilinear polynomial contractions; on the
real locus $\|H\|_{\rm F}^2=\sum_{i,j}H_{ij}^2$ is the usual
squared Frobenius norm. To prove the second identity, the two
product rules give
\begin{align*}
 \Delta(Nh'')&=(\Delta N)h''+
   (2g\cdot\nabla N+NL)h^{(3)}+N^2h^{(4)},\\
 \Delta(Lh')&=(\Delta L)h'+
   (2g\cdot\nabla L+L^2)h''+LN h^{(3)}.
\end{align*}
All derivatives of $h$ are evaluated at $\sigma$. Differentiating
$N=\sum_i(\partial_i\sigma)^2$ gives
$\nabla N=2Hg$ and
$\Delta N=2\sum_{i,j}H_{ij}^2+2g\cdot\nabla L$.
Substitution proves every term in
\eqref{eq:cq-bilaplacian-chain}.

The original derivatives before restriction are
\begin{align*}
 g_1&=\tfrac32yQ^2w+\tfrac12y^3(7+6xy),\\
 g_2&=\tfrac32xQ^2w+\tfrac12(8y+21xy^2+12x^2y^3),\\
 g_3&=\tfrac12Q^3,\\
 L&=3wQ(x^2+y^2)+3y^4+18x^2y^2+21xy+4.
\end{align*}
Differentiating these polynomials and substituting the unchanged
curve yields
\[
 g(\gamma(z))=
 \begin{pmatrix}-9z^3/32\\-3z/16\\-1/16\end{pmatrix},
 \qquad
 H(\gamma(z))=
 \begin{pmatrix}
 -27z^4/4&3z^2/16&-9z/16\\
 3z^2/16&13/4&3/(8z)\\
 -9z/16&3/(8z)&0
 \end{pmatrix}.
\]
The complete contractions are
\begin{align*}
 N(\gamma(z))&=(81z^6+36z^2+4)/1024,\\
 L(\gamma(z))&=(13-27z^4)/4,\\
 \|H(\gamma(z))\|_{\rm F}^2
  &=729z^8/16+9z^4/128+81z^2/128+169/16+9/(32z^2),\\
 (g\cdot\nabla L)(\gamma(z))
  &=9477z^8/512-405z^4/64+27z^2/128+81/32+3/(32z^2),\\
 (g^{\mathsf T}Hg)(\gamma(z))
  &=-2187z^{10}/4096+81z^6/4096
       -81z^4/4096+117z^2/1024+9/1024,\\
 (\Delta L)(\gamma(z))&=-111z^2+36/z^2.
\end{align*}
Consequently the exact restricted operator is
\begin{equation}
 (\Delta^2\sigma^*h)(\gamma(z))
   =\sum_{j=1}^4 C_j(z)h^{(j)}(-z^2/8),
 \label{eq:cq-bilaplacian-all-coefficients}
\end{equation}
where
\begin{align*}
 C_4(z)&=\bigl((81z^6+36z^2+4)/1024\bigr)^2,\\
 C_3(z)&=-6561z^{10}/2048+243z^6/2048
             -135z^4/1024+351z^2/512+31/512,\\
 C_2(z)&=26973z^8/128-4419z^4/64
             +135z^2/64+669/16+15/(16z^2),\\
 C_1(z)&=-111z^2+36/z^2.
\end{align*}
In particular both off-diagonal Hessian occurrences and all
derivatives of $L$ have been retained.

For every entire $h$, continuity of its first four derivatives at
zero and these four Laurent polynomials prove the limit
\begin{equation}
 \mathfrak r(h):=
 \lim_{z\downarrow0}z^2(\Delta^2\sigma^*h)(\gamma(z))
 =36h'(0)+\frac{15}{16}h''(0).
 \label{eq:cq-retained-spatial-residue}
\end{equation}
Indeed $z^2C_4,z^2C_3$ tend to zero, whereas $z^2C_2$ and
$z^2C_1$ tend to $15/16$ and $36$. No evenness assumption on
$h$ is needed. For any $R>0$, Cauchy's integral formula gives
\[
 |\mathfrak r(h)|\le
 \left(\frac{36}{R}+\frac{15}{8R^2}\right)
       \sup_{|s|\le R}|h(s)| .
\]
Thus $\mathfrak r$ is continuous for compact convergence and for
the growth topology. This estimate asserts no continuity in the
unnamed source topology $\tau$.

For the actual coefficient $K_t=T_{\rm D}Z_t$ of
Section~\ref{sec:cq-xi4-heat-returns}, differentiation of its
four original terms gives the exact value
\begin{align}
 \mathfrak r(K_t)
 &=-\frac{59}{275184}
       \left(36Z_t'(3)+\frac{15}{16}Z_t''(3)\right)\notag\\
 &\quad-\frac1{336}
       \left(36Z_t''(3)+\frac{15}{16}Z_t'''(3)\right)\notag\\
 &\quad+\frac1{1296}
       \left(36Z_t'(9/2)+\frac{15}{16}Z_t''(9/2)\right)\notag\\
 &\quad+\frac3{132496}
       \left(36Z_t'(13/2)+\frac{15}{16}Z_t''(13/2)\right).
 \label{eq:cq-xi4-spatial-residue}
\end{align}
For the unchanged incidence factor $a_{ef}$ and physical
diffusion $\nu\Delta$, its scaled second diffusion is
$\nu^2a_{ef}\mathfrak r(K_t)$ by linearity and operator composition.
The limit probes spatial infinity: $|\gamma(z)|\ge z^{-1}$.
Every $\sigma^*h$ is smooth at every finite real point by its
entire composition. These assertions specify the domain and
meaning of the observable exactly.

\subsection{The full ambiguity on the actual source quotient}
Set $E=\cqE$, $N=\cqN$, and $Q_{\rm ar}=(E,\tau)/N$.
The original even function $\Xi$ has $\Xi(0)>0$, as proved
from its retained positive kernel. Define
\begin{equation}
 e(s)=\frac{8}{15\Xi(0)}s^2\Xi(s),\qquad
 K=\ker\mathfrak r,\qquad P f=f-e\,\mathfrak r(f).
 \label{eq:cq-residue-section}
\end{equation}
Then $e\in\cqC[s]\Xi\subset N$, $e'(0)=0$, and
$e''(0)=16/15$, so $\mathfrak r(e)=1$.
Consequently
\begin{equation}
 \mathfrak r(N)=\cqC,\qquad
 \{([f]_N,\mathfrak r(f)):f\in E\}
       =Q_{\rm ar}\times\cqC .
 \label{eq:cq-source-residue-full-ambiguity}
\end{equation}
Here is the exact representative for every pair on the right.
Given any representative $f$ of $q$ and any $a\in\cqC$, put
$f_a=f+(a-\mathfrak r(f))e$. Then $[f_a]_N=q$ and
$\mathfrak r(f_a)=a$. In particular the zero class itself has
representatives $ae$ with every residue $a$. No nonempty set of
source quotient classes admits this scalar as a value independent
of all its representatives.

The strongest corresponding quotient map is explicit. Algebraically
$P^2=P$, $\operatorname{im}P=K$, $\ker P=\cqC e$, and
\[
 N_0:=N\cap K=P(N),\qquad P^{-1}(N_0)=N .
\]
The first equality for $N_0$ follows because $e\in N$ and
$P$ fixes $K$. If $Pf\in N$, the identity
$f=Pf+e\mathfrak r(f)$ proves $f\in N$, establishing the
inverse-image equality. Give $K$ the quotient topology
$\tau_P=P_*\tau$: a subset $V\subset K$ is open exactly when
$P^{-1}(V)$ is open in $(E,\tau)$. This is the topology of
$(E,\tau)/\cqC e$ under its displayed algebraic isomorphism
with $K$. In particular it is a vector topology. Since $N$ is
closed, the same inverse-image equality proves that $N_0$ is
closed in $(K,\tau_P)$. There is a homeomorphism
\begin{equation}
 (K,\tau_P)/N_0
    \ \xrightarrow{\ \ [k]\mapsto[k]_N\ \ }\
    (E,\tau)/N .
 \label{eq:cq-source-residue-reduced-quotient}
\end{equation}
For completeness, let $q_N:E\to E/N$ and
$q_0:K\to K/N_0$ denote the quotient maps. The proposed map
$A$ satisfies $Aq_0P=q_N$. Since $q_0P$ is a composition
of quotient maps and hence is quotient, this identity proves
continuity of $A$. The inverse is induced by $q_0P$, because
its kernel is $N$; the quotient property of $q_N$ proves its
continuity. These maps are inverse since
$P f-f\in N$ and $P k=k$. This proves the assertion without
using $\tau$-continuity of $\mathfrak r$ or of $P$ as an
endomorphism of $(E,\tau)$.

There is also a proved comparison using the inherited topology:
the identity $(K,\tau|_K)\to(K,\tau_P)$ is continuous.
Indeed for $V$ open in $\tau_P$,
$V=K\cap P^{-1}(V)$, which is inherited-open.
It induces a continuous bijection
$(K,\tau|_K)/N_0\to Q_{\rm ar}$.
No equality of these two topologies on $K$ is asserted.
For the growth topology the stronger result is available:
$\mathfrak r$ and $P$ are continuous by the Cauchy bound,
so $f\mapsto(Pf,\mathfrak r(f))$ and
$(k,a)\mapsto k+ae$ are inverse continuous linear maps
\begin{equation}
 (E,\gamma)\simeq (K,\gamma|_K)\oplus\cqC .
 \label{eq:cq-residue-growth-splitting}
\end{equation}
The continuity of $a\mapsto ae$ follows from each defining
seminorm. This proves the full split topology at precisely
the topology where its hypotheses have been established.

\subsection{Every finite jet at the original coordinate origin}
The residue calculation is a particular finite-jet phenomenon.
For $d\ge0$ retain the derivative coordinates
$J_d f=(f(0),f'(0),\ldots,f^{(d)}(0))$.
Write the convergent germ
$1/\Xi(s)=\sum_{n\ge0}b_ns^n$; its coefficients are given
without any change of $\Xi$ by
\[
 b_0=\Xi(0)^{-1},\qquad
 b_n=-\Xi(0)^{-1}
       \sum_{k=1}^n\frac{\Xi^{(k)}(0)}{k!}\,b_{n-k}.
\]
For $a=(a_0,\ldots,a_d)$ define the entire polynomial trace
\begin{equation}
 E_da(s)=\Xi(s)\sum_{n=0}^d
           \left(\sum_{j=0}^n\frac{a_j}{j!}b_{n-j}\right)s^n.
 \label{eq:cq-origin-jet-section}
\end{equation}
Multiplying its truncated Taylor series by that of $\Xi$
gives $E_da(s)=\sum_{j=0}^d a_js^j/j!+O(s^{d+1})$,
because the defining reciprocal convolution has coefficient
one at zero and zero at each positive degree.
Therefore $J_dE_d=1$ and
$\operatorname{im}E_d\subset\cqC[s]\Xi\subset N$.
This proves $J_d(N)=\cqC^{d+1}$, with its complete
finite-dimensional section, rather than only the scalar image.
For any $q=[f]_N$ and any jet $a$, the representative
$f+E_d(a-J_df)$ realizes $(q,a)$.

Put $P_d=1-E_dJ_d$, $K_d=\ker J_d$, and
$N_d=N\cap K_d$. The same identities, now with
$\ker P_d=\operatorname{im}E_d$, prove
$P_d(N)=N_d$ and $P_d^{-1}(N_d)=N$.
Thus $(K_d,(P_d)_*\tau)/N_d\simeq Q_{\rm ar}$ by the
two quotient-map arguments just given. Cauchy's bounds
$|f^{(j)}(0)|\le j!R^{-j}\sup_{|s|\le R}|f(s)|$
prove the continuous growth-topology decomposition
$(E,\gamma)\simeq (K_d,\gamma|_{K_d})\oplus\cqC^{d+1}$.
All factorials are retained in the derivative coordinates
and in the section.

\subsection{Heat transport of the retained residue and its quotient}
The scalar observable also has an exact heat transport. Define
\[
 \mathfrak r_t=\mathfrak r U_{-t},\qquad
 e_t=U_te,\qquad
 P_t=1-e_t\mathfrak r_t=U_tPU_{-t},\qquad
 K_t^{\rm res}=\ker\mathfrak r_t .
\]
These definitions retain $U_t=\exp(-tD^2/4)$ and hence
give the convergent formula
\begin{equation}
 \mathfrak r_t(f)
 =36\sum_{k=0}^{\infty}\frac{(t/4)^k}{k!}
       f^{(2k+1)}(0)
  +\frac{15}{16}\sum_{k=0}^{\infty}\frac{(t/4)^k}{k!}
       f^{(2k+2)}(0).
 \label{eq:cq-transported-spatial-residue}
\end{equation}
The earlier heat-series proof gives convergence in the growth
topology, locally uniformly in $t$; applying the two continuous
derivative evaluations gives the displayed convergence and
its value. In particular
$\mathfrak r_tU_t=\mathfrak r$,
$\mathfrak r_t(e_t)=1$, and $e_t\in N_t$.
Also $U_tK=K_t^{\rm res}$ and
$U_tN_0=N_t\cap K_t^{\rm res}$ by these identities and
invertibility of $U_t$.

Give $K_t^{\rm res}$ the topology $(P_t)_*\tau_t$.
It equals $(U_t)_*\tau_P$: for $V\subset K_t^{\rm res}$,
the identity $P_tU_t=U_tP$ gives
$U_{-t}(P_t^{-1}(V))=P^{-1}(U_{-t}V)$,
which is exactly the equivalence of the two defining open-set
conditions. Thus $U_t$ gives a homeomorphism between the
residue-zero quotient in
\eqref{eq:cq-source-residue-reduced-quotient} and
\[
 (K_t^{\rm res},(P_t)_*\tau_t)/
       (N_t\cap K_t^{\rm res}) .
\]
The latter is homeomorphic to $(E,\tau_t)/N_t$ by
$[k]\mapsto[k]_{N_t}$, with inverse $[f]\mapsto[P_tf]$;
the two quotient universal properties apply exactly as above.
The two homeomorphisms commute with the already proved heat
homeomorphism of the actual source quotient because both send
$[f]$ to $[U_tf]$. Finally, for each $q=[f]_{N_t}$ and
$a\in\cqC$, the unchanged class has the representative
$f+e_t(a-\mathfrak r_t(f))$ with transported residue $a$.
This computes the full transported scalar relation without
identifying $\mathfrak r_t$ with the fixed physical functional
$\mathfrak r$.

\subsection{Every one-point finite jet of the moving actual relation}
\label{sec:cq-moving-finite-jets}

The complete finite-jet statement persists on the actual transported
source relation, at every point once the time is nonzero.
Throughout this subsection the derivative coordinate is the original
$s$, the time may be complex, and
$N_t=U_t\cqN$ retains its proved closure topology $\tau_t$.

\begin{cqproposition}[Full finite jets on the moving source relation]
\label{prop:cq-moving-finite-jets}
For every $t\in\cqC\setminus\{0\}$, $\lambda\in\cqC$, and
integer $d\ge0$, put
\[
 J_{\lambda,d}f=(D^jf(\lambda))_{j=0}^d.
\]
Then
\begin{equation}
 J_{\lambda,d}\bigl(U_t(\cqC[s]\Xi)\bigr)
   =J_{\lambda,d}(N_t)=\cqC^{d+1}.
 \label{eq:cq-moving-finite-jet-image}
\end{equation}
There is a continuous linear right inverse of this jet map from
the usual finite-dimensional space into $(\cqE,\tau_t)$, whose
image consists of finite linear combinations of the actual
transported polynomial traces.
\end{cqproposition}

\begin{proof}
Suppose a complex linear functional
\[
 Lf=\sum_{j=0}^d c_jD^jf(\lambda)
\]
annihilates $U_t(\cqC[s]\Xi)$. We first justify the full
exponential generating calculation in the established growth
topology, without invoking convergence in $\tau$ or $\tau_t$.
For every $1<p<2$, $a>0$, and $n\ge1$, direct maximization gives
\begin{equation}
 \|s^n\Xi\|_{p,a}
 \le\|\Xi\|_{p,a/2}
       \left(\frac{2n}{ap\mathrm e}\right)^{n/p}.
 \label{eq:cq-exponential-jet-convergence}
\end{equation}
Indeed the maximum of $r^n\exp(-ar^p/2)$ occurs at
$r^p=2n/(ap)$. The $n=0$ term is controlled directly by
$\|\Xi\|_{p,a}$. For every $A>0$, the sum of the right side
times $A^n/n!$ converges: the elementary bound
$n!\ge(n/\mathrm e)^n$ makes its $n$-th root tend to zero.
Each fixed derivative of the parameter power series adds only
a polynomial factor in $n$, together with a fixed shift of its
power. The same root estimate proves locally uniform convergence
of all those differentiated series in every growth seminorm.
Thus
\[
 \alpha\longmapsto e^{\alpha s}\Xi(s)
   =\sum_{n=0}^\infty\frac{\alpha^n}{n!}s^n\Xi(s)
\]
is entire with values in $(\cqE,\gamma)$, with the displayed
parameter derivatives. The proved continuity of $U_t$ in
$\gamma$, and Cauchy's estimate for every derivative evaluation
in $L$, justify applying $LU_t$ to every series. All its Taylor
coefficients vanish by the assumed annihilation. Hence
$LU_t(e^{\alpha s}\Xi)=0$ for every $\alpha\in\cqC$.

The exact heat and exponential identity, retaining its original
sign and coefficient, is
\[
 U_t(e^{\alpha s}\Xi)(s)
  =e^{\alpha s-t\alpha^2/4}Z_t(s-t\alpha/2).
\]
It also follows directly by applying the convergent heat series
to $D(e^{\alpha s}f)=e^{\alpha s}(D+\alpha)f$ and expanding
the commuting operators $D$ and the scalar $\alpha$.
Leibniz's rule therefore gives
\begin{equation}
 0=e^{\alpha\lambda-t\alpha^2/4}
   \sum_{j=0}^d c_j\sum_{k=0}^j\binom jk
      \alpha^{j-k}Z_t^{(k)}(\lambda-t\alpha/2).
 \label{eq:cq-moving-jet-annihilator}
\end{equation}
The exponential is nonzero. Set $w=\lambda-t\alpha/2$,
$b=2/t$, and $x=w-\lambda$. Since $t\ne0$, $w$ ranges
over all of $\cqC$, and $\alpha=-bx$. Substitution in
\eqref{eq:cq-moving-jet-annihilator} gives
\begin{equation}
 \sum_{j=0}^d c_j\mathcal N_jZ_t=0,
 \qquad
 \mathcal N_j=\sum_{k=0}^j\binom jk(-bx)^{j-k}D_w^k.
 \label{eq:cq-moving-jet-normal-order}
\end{equation}
The polynomial factors in this formula multiply after the
indicated derivatives are taken. They do not commute with
$D_w$; in particular this formula is not an uncorrected
power of $D_w-bx$.

Here is the exact correction at every order. Define
\begin{equation}
 A=D_w-bx,\qquad
 H_j(X;b)=j!\sum_{\ell=0}^{\lfloor j/2\rfloor}
        \frac{(b/2)^\ell X^{j-2\ell}}
             {\ell!(j-2\ell)!}.
 \label{eq:cq-moving-jet-Hermite-polynomial}
\end{equation}
Then $\mathcal N_j=H_j(A;b)$. To prove this coefficient
identity on any entire $f$, use
\[
 G(u,w)=e^{-bu(w-\lambda)}f(w+u)
       =\sum_{j\ge0}\frac{u^j}{j!}\mathcal N_jf(w).
\]
Differentiating the left side in both variables gives
$\partial_uG=(A+bu)G$. Comparing its Taylor coefficients yields
\[
 \mathcal N_0=1,\quad \mathcal N_1=A,\quad
 \mathcal N_{j+1}=A\mathcal N_j+jb\mathcal N_{j-1}
 \quad(j\ge1).
\]
The explicitly displayed polynomials have exponential generating
series $\exp(uX+bu^2/2)$ and hence satisfy the same recurrence
and initial values. Induction proves the claimed identity.
For example, $\mathcal N_2=A^2+b$ and
$\mathcal N_3=A^3+3bA$, including the terms created by
noncommutation. This argument compares finite-order differential
operators coefficient by coefficient and makes no claim that
multiplication by an exponential quadratic preserves $\cqE$.

In the coordinate $w$, retain the proved conjugated source
operator $B_t=w-tD_w/2=U_tMU_{-t}$. Then
\[
 A=-\frac2t(B_t-\lambda),\qquad
 U_{-t}AU_t=-\frac2t(M-\lambda).
\]
All the finite-order operators in
\eqref{eq:cq-moving-jet-normal-order} act continuously in
$\gamma$. Applying $U_{-t}$ to that identity and using
$Z_t=U_t\Xi$ therefore gives the exact entire identity
\begin{equation}
 \mathcal P(w)\Xi(w)=0,\qquad
 \mathcal P(w)=\sum_{j=0}^d
   c_jH_j\left(-\frac2t(w-\lambda);\frac2t\right).
 \label{eq:cq-moving-jet-polynomial-contradiction}
\end{equation}
If $c_J$ is the last nonzero coefficient, $H_J(X;b)$ is
monic of degree $J$, while every preceding polynomial has
smaller degree. Thus $\mathcal P$ has leading coefficient
$c_J(-2/t)^J\ne0$. But $\Xi(0)\ne0$, so
\eqref{eq:cq-moving-jet-polynomial-contradiction} forces
$\mathcal P$ to vanish on a disk about zero and hence to
be the zero polynomial, a contradiction. The annihilating
functional is therefore zero.

A proper linear subspace of $\cqC^{d+1}$ has a nonzero
annihilating functional: extend a basis of the subspace to a
basis of the whole space and take a coordinate functional
on one of the added basis vectors. The result just proved
therefore gives the first equality in
\eqref{eq:cq-moving-finite-jet-image}. Since every transported
polynomial trace belongs to $N_t$, it gives the second
equality too.

Here is a finite construction of the promised section. Enumerate
$h_n=U_t(s^n\Xi)$, $n=0,1,\ldots$. Choose $n_0$ to be
the first index whose jet column is nonzero, and successively
choose $n_\ell$ to be the first index whose jet column is
outside the span of the previously selected columns.
Full rank, already proved, makes this selection terminate
at $d+1$ columns. Define
\begin{equation}
 G_{j\ell}=D^jh_{n_\ell}(\lambda)
 \quad(0\le j,\ell\le d),\qquad
 \mathcal S_{t,\lambda,d}(v)
   =\sum_{\ell=0}^d h_{n_\ell}(G^{-1}v)_\ell.
 \label{eq:cq-moving-jet-finite-section}
\end{equation}
The determinant of $G$ is nonzero by its construction;
equivalently $G^{-1}=\operatorname{adj}(G)/\det G$.
Direct multiplication gives
$J_{\lambda,d}\mathcal S_{t,\lambda,d}=1$.
Its image is a finite-dimensional subspace of
$U_t(\cqC[s]\Xi)$. Each scalar coefficient in the finite
sum is a continuous linear function of $v$, and scalar
multiplication and addition are continuous in the source
vector topology $\tau_t$. This proves continuity of the
section into precisely $(\cqE,\tau_t)$, as claimed.
\end{proof}

There is an exact topological quotient for these representatives,
without requiring the jet map itself to be $\tau_t$-continuous.
For the fixed $t\ne0$, $\lambda$, and $d$ above, put
\begin{gather*}
 K_{\lambda,d}=\ker J_{\lambda,d},\qquad
 P_{t,\lambda,d}=1-\mathcal S_{t,\lambda,d}J_{\lambda,d},\qquad
 N_{t,\lambda,d}=N_t\cap K_{\lambda,d},\\
 \theta_{t,\lambda,d}=(P_{t,\lambda,d})_*\tau_t.
\end{gather*}
Here the last topology is the quotient topology on
$K_{\lambda,d}$; a set is open exactly when its inverse
image under $P_{t,\lambda,d}$ is $\tau_t$-open.
Since the section has image in $N_t$ and is a right inverse,
direct substitution gives
\[
 P_{t,\lambda,d}^2=P_{t,\lambda,d},\qquad
 P_{t,\lambda,d}(N_t)=N_{t,\lambda,d},\qquad
 P_{t,\lambda,d}^{-1}(N_{t,\lambda,d})=N_t.
\]
The last equality proves that $N_{t,\lambda,d}$ is closed
in the stated quotient topology. There is a homeomorphism
\begin{equation}
 (K_{\lambda,d},\theta_{t,\lambda,d})/N_{t,\lambda,d}
  \xrightarrow{\ [k]\mapsto[k]_{N_t}\ }Q_t,
 \qquad
 [f]_{N_t}\longmapsto[P_{t,\lambda,d}f]
 \label{eq:cq-moving-jet-source-quotient}
\end{equation}
as its inverse. Well-definedness and both inverse identities
follow from the three displayed relations and
$f-P_{t,\lambda,d}f\in N_t$. The forward map is continuous
because its composition with the quotient map
$P_{t,\lambda,d}$ and the further quotient by
$N_{t,\lambda,d}$ is the original continuous quotient
$\cqE\to Q_t$. The reverse map is continuous because
it is induced by that continuous composite quotient map
out of $(\cqE,\tau_t)$ with kernel $N_t$.
As before, restriction of $P_{t,\lambda,d}$ to its range
is the identity, proving the continuous identity map from
the inherited-topology range to the quotient-topology
range. This induces a continuous bijection of their
quotients; inverse continuity for the inherited topology
has not been inserted.

Finally apply the proposition with $\lambda=0$ and $d=2$.
The explicit choice
\begin{equation}
 \varepsilon_t=\mathcal S_{t,0,2}(0,0,16/15)\quad(t\ne0),
 \qquad \varepsilon_0=\frac{8s^2\Xi(s)}{15\Xi(0)}
 \label{eq:cq-moving-residue-section}
\end{equation}
belongs to $N_t$ and satisfies
$\mathfrak r(\varepsilon_t)
 =36\varepsilon_t'(0)+(15/16)\varepsilon_t''(0)=1$.
At $t=0$ this is the previously computed polynomial trace;
at every nonzero time the section fixes the two required
derivatives to $0$ and $16/15$. Therefore, at every complex
time including zero,
\begin{equation}
 \mathfrak r(N_t)=\cqC,\qquad
 \{([f]_{N_t},\mathfrak r(f)):f\in\cqE\}
      =Q_t\times\cqC.
 \label{eq:cq-all-times-residue-ambiguity}
\end{equation}
For each pair $([f]_{N_t},a)$ the exact representative
$f+(a-\mathfrak r(f))\varepsilon_t$ proves surjectivity onto that
pair. Put $P_t^{\mathfrak r}=1-\varepsilon_t\mathfrak r$ and
retain $K=\ker\mathfrak r$. The identities
\[
 P_t^{\mathfrak r}(N_t)=N_t\cap K,
 \qquad (P_t^{\mathfrak r})^{-1}(N_t\cap K)=N_t
\]
give, by the same two quotient-map arguments, the homeomorphism
\begin{equation}
 \bigl(K,(P_t^{\mathfrak r})_*\tau_t\bigr)/(N_t\cap K)
      \cong Q_t,
 \qquad [k]\mapsto[k]_{N_t},\quad
 [f]_{N_t}\mapsto[P_t^{\mathfrak r}f].
 \label{eq:cq-all-times-residue-source-quotient}
\end{equation}
This retains the actual transported source topology and
the exact physical residue at each fixed time. No
regular dependence of the selected finite pivot columns
on $t$ is needed or asserted. In particular, the full
scalar ambiguity holds without a statement that the
original frozen relation $\cqN$ is invariant under heat,
and without identifying that relation with $\Xi\cqE$.

\section{The Cartan construction and an exact antiunitary reversal}
\label{sec:cq-cartan-heat}

The supplied modular--Cartan manuscript \cite{cqCartan} contains both
a sector conjugation preserving a logarithmic generator and a
standard-form involution reversing the left-minus-right generator.
The following construction retains these two equations explicitly.
It constructs the negative radial direction, its domains, and its
maps to the arithmetic heat space and moving source quotients.

\subsection{The generator equation determines the reflected parameter}

Let $\mathscr H$ be a complex Hilbert space, $K$ a self-adjoint
operator, and $E_K$ its projection-valued spectral measure. Define
\begin{equation}
\begin{split}
 T(\zeta)&=e^{-i\zeta K},\qquad \zeta=a+ib\in\mathbb C,\\
 \operatorname{Dom}T(\zeta)
 &=\left\{f\in\mathscr H:
       \int_{\mathbb R}e^{2b\lambda}\,
                   d\langle E_K(\lambda)f,f\rangle<\infty\right\}.
\end{split}
\label{eq:cq-cartan-complex-domain}
\end{equation}
This is a closed densely defined normal operator for every $b$,
including every negative $b$. Density follows because the spectral
cutoffs $E_K([-n,n])f$ belong to this domain and converge to $f$.
Closedness follows from the spectral multiplication representation.

\begin{cqproposition}
Let $J$ be an antiunitary involution and suppose the equality of
self-adjoint operators $JKJ=\epsilon K$ holds, where
$\epsilon\in\{1,-1\}$; thus it includes equality of domains.
Then
\begin{equation}
 JT(\zeta)J=T(-\epsilon\overline\zeta),\qquad
 J e^{zK}J=e^{\epsilon\overline zK}.
 \label{eq:cq-cartan-parity}
\end{equation}
Both equalities include the transported operator domains.
In particular, with $\Delta=e^{-K}$, the Tomita relation
$J\Delta J=\Delta^{-1}$ is equivalent to $JKJ=-K$ and gives
\begin{equation}
 J\Delta^{i(a+ib)}J=\Delta^{i(a-ib)},\qquad
 J e^{bK}J=e^{-bK}\quad(b\in\mathbb R).
 \label{eq:cq-Tomita-radial-reversal}
\end{equation}
\end{cqproposition}

\begin{proof}
The uniqueness of the spectral resolution gives
$JE_K(B)J=E_{\epsilon K}(B)$ for every Borel set $B$.
For a simple scalar function $\phi$, antilinearity gives
$J\phi(K)J=\overline\phi(\epsilon K)$.
Uniform approximation proves this for bounded Borel functions.
For an arbitrary finite-valued Borel function, apply the bounded
identity to $\phi\,\mathbf1_{\{|\phi|\leq n\}}$ and pass to the
spectral integral. The squared-integrability condition for
$|\phi|$ is transported by $J$, proving equality of the full
domains as well as of the values. Taking
$\phi(\lambda)=e^{-i\zeta\lambda}$ and $e^{z\lambda}$ proves
\eqref{eq:cq-cartan-parity}. Applying the same real-valued
functional calculus to $-\log$ proves the equivalence involving
$\Delta$, whose kernel is zero. Substitution of $\epsilon=-1$
proves \eqref{eq:cq-Tomita-radial-reversal}.
\end{proof}

The Cartan logarithm in the supplied construction is
\[
 \mathbf L_\ell(T)=\frac{1}{2\ell}\log(T^*T),\qquad \ell>0.
\]
For the operator in \eqref{eq:cq-cartan-complex-domain}, multiplication
of the spectral functions, including the product domain, gives
$T(a+ib)^*T(a+ib)=e^{2bK}$. Indeed the product domain is precisely
the set on which the squared multiplier $e^{2b\lambda}$ is
square-integrable; its integrability also implies that of
$e^{b\lambda}$ by Cauchy--Schwarz. Consequently,
\begin{equation}
 \mathbf L_\ell(T(a+ib))=
 \begin{cases}
 (b/\ell)K & b\ne0,\quad\text{with domain }\operatorname{Dom}K,\\
 0 & b=0,\quad\text{with domain }\mathscr H.
 \end{cases}
 \label{eq:cq-cartan-log-domain}
\end{equation}
This follows by composing the real spectral functions
$\log(e^{2b\lambda})=2b\lambda$. The zero slice has been stated
separately: the everywhere-defined $\log I$ must not be replaced
by a zero operator restricted to $\operatorname{Dom}K$.

\subsection{The original heat multiplier and a genuine modular double}

Retain the original Fourier variable $v$ and set
\begin{equation}
\begin{split}
 \mathscr H_0&=L^2(\mathbb R,dv),\qquad
 (C_0f)(v)=\overline{f(-v)},\\
 (Af)(v)&=v^2f(v),\qquad
 \operatorname{Dom}A=\{f:v^2f\in L^2\},\\
 G_t&=e^{tA/4},\\
 \operatorname{Dom}G_t
 &=\left\{f:\int_{\mathbb R}
             e^{(\operatorname{Re}t)v^2/2}|f(v)|^2\,dv<\infty\right\}.
\end{split}
\label{eq:cq-Hilbert-heat-domain}
\end{equation}
The multiplication operator $A$ is nonnegative self-adjoint,
$C_0$ is an antiunitary involution, and $C_0AC_0=A$ with the
displayed domain. Thus
$C_0G_tC_0=G_{\overline t}$, which fixes each real heat parameter.
For $\operatorname{Re}t\leq0$, $G_t$ is bounded with norm one.
For $\operatorname{Re}t>0$, it is densely defined and unbounded:
unit vectors supported in $[n,n+1]$ have image norm at least
$e^{(\operatorname{Re}t)n^2/4}$. Its domain is proper, since
$f(v)=e^{-(\operatorname{Re}t)v^2/8}$ belongs to $L^2$ and not
to that domain. Multiplication gives the exact inverse statement
\begin{equation}
 \operatorname{Ran}G_t=\operatorname{Dom}G_{-t},\qquad
 G_{-t}G_t f=f\quad(f\in\operatorname{Dom}G_t).
 \label{eq:cq-heat-unbounded-inverse}
\end{equation}
For the converse inclusion in the range statement, take
$f=G_{-t}g$ for $g\in\operatorname{Dom}G_{-t}$; then
$G_tf=g\in L^2$, proving $f\in\operatorname{Dom}G_t$.
The common invariant space
$\mathscr D=C_c^\infty(\mathbb R)$ is a graph core for each
$G_t$: truncate a vector in its weighted $L^2$ graph norm,
and approximate the truncations by smooth compact functions,
using equivalence of the weighted and unweighted norms on
each fixed compact interval.

There is no antiunitary $J_0$ on this one sector satisfying
$J_0AJ_0=-A$. Such an equality would make both $A$ and $-A$
nonnegative by preservation of real quadratic forms, forcing
$\langle Af,f\rangle=0$ for every $f\in\operatorname{Dom}A$.
A nonzero smooth function supported in $(1,2)$ contradicts this.
The exact construction retaining this original $A$ is as follows.

\begin{cqtheorem}
On $\widehat{\mathscr H}=\mathscr H_0\oplus\mathscr H_0$, define
\begin{equation}
\begin{aligned}
 \widehat K&=\begin{pmatrix}A&0\\0&-A\end{pmatrix},&
 \widehat\Delta&=\begin{pmatrix}e^{-A}&0\\0&e^A\end{pmatrix},\\
 J(f,g)&=(C_0g,C_0f),&
 \widehat G_t&=e^{t\widehat K/4}
              =\begin{pmatrix}G_t&0\\0&G_{-t}\end{pmatrix}.
\end{aligned}
\label{eq:cq-modular-heat-double}
\end{equation}
Here $\operatorname{Dom}\widehat K=\operatorname{Dom}A\oplus
\operatorname{Dom}A$,
$\operatorname{Dom}\widehat\Delta=\mathscr H_0\oplus
\operatorname{Dom}e^A$, and
$\operatorname{Dom}\widehat G_t=\operatorname{Dom}G_t\oplus
\operatorname{Dom}G_{-t}$. Then
\begin{equation}
 J^2=1,\quad J\widehat KJ=-\widehat K,\quad
 J\widehat\Delta J=\widehat\Delta^{-1},\quad
 J\widehat G_tJ=\widehat G_{-\overline t}.
 \label{eq:cq-modular-double-reversal}
\end{equation}
These are the modular data of an explicitly constructed standard
closed real subspace of $\widehat{\mathscr H}$.
\end{cqtheorem}

\begin{proof}
The block formulas and $C_0AC_0=A$ prove all four identities,
including their domains. For the last assertion put $Q=e^{-A/2}$.
This is a bounded positive injective operator, with dense range,
commuting with $C_0$. Define the closed antilinear operator
\[
 S=J\widehat\Delta^{1/2},\qquad
 \operatorname{Dom}S=\mathscr H_0\oplus\operatorname{Dom}Q^{-1},
 \qquad S(f,g)=(C_0Q^{-1}g,C_0Qf).
\]
Its domain is invariant under $S$, and direct substitution gives
$S^2(f,g)=(f,g)$. Its fixed real subspace is exactly
\begin{equation}
 V=\{(f,C_0Qf):f\in\mathscr H_0\}.
 \label{eq:cq-standard-real-heat-space}
\end{equation}
Both implications follow by substituting the second component
in $S(f,g)=(f,g)$. The graph defining $V$ is closed because
$C_0Q$ is bounded. If a vector in $V$ equals $i$ times a vector
in $V$, its first components give $f=ig$, and its second
components give $-iC_0Qg=iC_0Qg$; injectivity of $Q$ gives
$f=g=0$. Hence $V\cap iV=\{0\}$.

Moreover $V+iV=\mathscr H_0\oplus\operatorname{Ran}Q
=\operatorname{Dom}S$. To prove the nontrivial inclusion,
for $x\in\mathscr H_0$ and $y\in\operatorname{Ran}Q$ put
$b=C_0Q^{-1}y$, $f=(x+b)/2$ and $g=(x-b)/(2i)$.
Then
\[
 (f,C_0Qf)+i(g,C_0Qg)
   =(f+ig,C_0Q(f-ig))=(x,y).
\]
This domain is dense since $\operatorname{Ran}Q$ is dense.
Thus $V$ is standard: it is closed real, has zero intersection
with $iV$, and has dense complex span. On that span
$S(v+iw)=v-iw$ for $v,w\in V$, as follows from antilinearity
and $S|_V=1$. Its polar decomposition is exactly
$S=J\widehat\Delta^{1/2}$. This proves the claimed modular
real-subspace construction, including its canonical involution.
\end{proof}

For $\widehat T(\zeta)=e^{-i\zeta\widehat K}$,
the Cartan logarithm is \eqref{eq:cq-cartan-log-domain} with
$K=\widehat K$, and $\widehat G_t=\widehat T(it/4)$.
In particular the involution sends every real $t$ to $-t$.
For nonzero real $t$, the double is unbounded in one of its
two sectors; its domain remains dense. The inclusion
$I_+f=(f,0)$ retains the original multiplier exactly:
\[
 \widehat G_t I_+f=I_+G_tf
       \quad(f\in\operatorname{Dom}G_t).
\]
No Hilbert completion of the unnamed arithmetic topology is
required for this construction.

\subsection{Exact ingress and involution on the actual moving quotients}

On the original entire-function space $\cqE$ put
$(C_Eh)(s)=\overline{h(\overline s)}$. This is an antilinear
involution and preserves each previously defined growth seminorm,
since $s\mapsto\overline s$ preserves $|s|$.
For the Fourier transform already used in the test ingress,
\[
 (F\phi)(s)=\int_{\mathbb R}\phi(v)e^{-isv}\,dv,
\]
the change of variable $v\mapsto-v$ proves
$C_EF=FC_0$ on $\mathscr D$. Differentiating the integral and
summing the absolutely uniformly convergent exponential series
on the compact support proves
\begin{equation}
 U_tF=FG_t,\qquad C_EU_t=U_{\overline t}C_E.
 \label{eq:cq-cartan-Fourier-ingress}
\end{equation}
The second equality also follows directly from the convergent
entire-space series for $U_t$ and the real coefficient $-1/4$.
Thus on $\cqE\oplus\cqE$ the everywhere-defined maps
\[
 \widehat U_t=(U_t,U_{-t}),\qquad
 \mathcal J(h,k)=(C_Ek,C_Eh)
\]
satisfy $\mathcal J\widehat U_t\mathcal J
=\widehat U_{-\overline t}$ for every complex $t$. The exact
diagram is expressed by
\[
 (F\oplus F)\widehat G_t=\widehat U_t(F\oplus F),\qquad
 (F\oplus F)J=\mathcal J(F\oplus F)
       \quad\text{on }\mathscr D\oplus\mathscr D.
\]
Here $\widehat G_t$ on the test space is the restriction of
the closed Hilbert operator constructed above.

Keep the actual source topology $\tau$ and actual closed
trace image $\cqN$. Define its conjugate presentation by
$\tau^c=(C_E)_*\tau$ and $N^c=C_E\cqN$.
Since $\Xi$ is real entire, $C_E(\mathbb C[s]\Xi)
=\mathbb C[s]\Xi$, and therefore
$N^c=\overline{\mathbb C[s]\Xi}^{\,\tau^c}$.
This defines the conjugate topology explicitly without assuming
$C_E$ is continuous for the unnamed original $\tau$.

\begin{cqproposition}
For each complex $t$ form the topological quotient
\begin{equation}
\begin{split}
 \widehat N_t&=U_t\cqN\oplus U_{-t}N^c,\\
 \widehat Q_t&=
 \big((\cqE,\tau_t)\oplus(\cqE,(\tau^c)_{-t})\big)/\widehat N_t,
 \qquad \tau_t=(U_t)_*\tau,\quad
 (\tau^c)_{-t}=(U_{-t})_*\tau^c .
\end{split}
\label{eq:cq-conjugate-moving-quotient}
\end{equation}
The map $\widehat U_t$ induces a complex-linear homeomorphism
$\overline{\widehat U}_t:\widehat Q_0\to\widehat Q_t$.
The involution $\mathcal J$ induces an antilinear homeomorphism
\begin{equation}
 \overline{\mathcal J}_t:\widehat Q_t\longrightarrow
       \widehat Q_{-\overline t},\qquad
 \overline{\mathcal J}_{-\overline t}\overline{\mathcal J}_t=1,
 \qquad
 \overline{\mathcal J}_t\overline{\widehat U}_t
 =\overline{\widehat U}_{-\overline t}\overline{\mathcal J}_0 .
 \label{eq:cq-actual-quotient-antiunitary-bridge}
\end{equation}
For $\iota_t(\phi,\psi)=[F\phi,F\psi]\in\widehat Q_t$,
\begin{equation}
\begin{split}
 \ker\iota_t&=F^{-1}(U_t\cqN)\oplus F^{-1}(U_{-t}N^c),\\
 \iota_t\widehat G_t&=\overline{\widehat U}_t\iota_0,\qquad
 \overline{\mathcal J}_t\iota_t=\iota_{-\overline t}J .
\end{split}
\label{eq:cq-actual-quotient-test-diagram}
\end{equation}
The inverse images here are taken inside $\mathscr D$.
\end{cqproposition}

\begin{proof}
The direct-sum heat map transports both the specified topology
and the specified closed subspace exactly, proving its quotient
homeomorphism. Next
\[
 C_EU_{-t}N^c=U_{-\overline t}\cqN,\qquad
 C_EU_t\cqN=U_{\overline t}N^c .
\]
The same identities applied to the definitions of the transported
topologies prove that $\mathcal J$ is an antilinear homeomorphism
of the two ambient spaces in \eqref{eq:cq-conjugate-moving-quotient}.
It transports their closed relation spaces exactly, so it induces
the asserted quotient map. Applying it twice is the identity
on representatives. The conjugacy follows from
$\mathcal J\widehat U_t=\widehat U_{-\overline t}\mathcal J$.
The kernel formula is exactly the condition that both Fourier
components lie in the two relation spaces. Finally
\eqref{eq:cq-cartan-Fourier-ingress} proves both equations
in \eqref{eq:cq-actual-quotient-test-diagram} on representatives.
\end{proof}

Thus the constructed Hilbert involution is antiunitary and its
actual source-quotient image is a proved antilinear homeomorphism.
No inner product is asserted on $\widehat Q_t$.
For real $t$, this construction provides the requested negative
parameter and its exact quotient target. It does not assert
the distinct frozen inclusion $U_t\cqN\subseteq\cqN$.

\subsection{The exact viscosity sign and tensor Fourier map}

Let $\nu>0$ be the original viscosity and let $\theta\geq0$ be
physical elapsed time. On $L^2(\mathbb R^3,dq)$ define
$A_\nu=-\nu\Delta_q$ by the unitary Fourier transform
\[
 (\mathcal F_3f)(k)=(2\pi)^{-3/2}
                   \int_{\mathbb R^3}f(q)e^{-ik\cdot q}\,dq .
\]
Its domain is exactly
$\{f:|k|^2\mathcal F_3f\in L^2\}$, and its multiplier is
$\nu|k|^2$. Integration by parts verifies this differential
expression on Schwartz functions, and the spectral multiplication
domain supplies the closed realization. On the frequency tensor
product $L^2(\mathbb R)^{\otimes3}=L^2(\mathbb R^3)$ one has
\begin{equation}
 \mathcal F_3 e^{-\theta A_\nu}\mathcal F_3^{-1}
 =e^{-\nu\theta|k|^2}
 =G_{-4\nu\theta}\otimes G_{-4\nu\theta}\otimes G_{-4\nu\theta}.
 \label{eq:cq-viscosity-Fourier-intertwiner}
\end{equation}
The equality holds first on simple tensors, where multiplication
of the three factors gives the middle multiplier, and then on
all of $L^2$ by their common boundedness and tensor density.
On $\mathscr D^{\otimes3}$ the original unnormalized Fourier
transform $F^{\otimes3}$ is $(2\pi)^{3/2}\mathcal F_3$ with
the variables renamed; this scalar is retained. Each transformed
test is entire in the three coordinates, and
\eqref{eq:cq-cartan-Fourier-ingress} gives
$U_{-4\nu\theta}^{\otimes3}$ on that entire tensor test space.
This proves the sign and coefficient in the viscous bridge.
The multiplier also acts componentwise on vector fields and
preserves the divergence-free subspace: in Fourier coordinates
$k\cdot\widehat u(k)=0$ is unchanged by scalar multiplication.
Its inverse at $\theta>0$ has precisely the domain
$\{f:e^{\nu\theta|k|^2}\mathcal F_3f\in L^2\}$.
These identities identify the linear viscous operator, including
its inverse domain; they do not replace the nonlinear velocity,
pressure, or forcing equations.

\subsection{The public viscous source and its exact proof status}
\label{sec:cq-public-NS-source}

The inspected public manuscript \cite{cqNS2026} has two preserved
revisions. Its earlier 165-page PDF and its later 166-page PDF
have different recorded byte hashes. The later revision is the
one served after 19:09:35 UTC on 8 September 2026; its first
three pages were inspected directly. No equivalence of the
two complete proofs is asserted. The pinned source repository
commit is recorded separately in the bibliography and provenance.

Theorem 1.1 in that manuscript is an imported source claim:
for every $\nu>0$ it asserts existence of
$f\in C_c^\infty(\mathbb R^3\times(0,\infty);\mathbb R^3)$,
a compact spatial set $K$, and smooth $u,p$ on
$\mathbb R^3\times[0,1)$ satisfying
\[
 \partial_\theta u+(u\cdot\nabla)u-\nu\Delta u+\nabla p=f,
 \qquad \nabla\cdot u=0,\qquad u(\cdot,0)=0,
\]
with the support and norm conditions
\[
 \operatorname{supp}u(\cdot,\theta)\cup
 \operatorname{supp}p(\cdot,\theta)\subseteq K,\qquad
 \sup_{0\leq\theta<1}\|u(\cdot,\theta)\|_2<\infty,\qquad
 \limsup_{\theta\uparrow1}\|u(\cdot,\theta)\|_\infty=\infty .
\]
It also asserts nonexistence of a global smooth competitor
with the same force and initial datum and uniformly bounded
kinetic energy. These assertions have not been independently
proved in this fragment or checked here by a Lean kernel.
The identities in \eqref{eq:cq-viscosity-Fourier-intertwiner}
are independently proved above and match the exact ordinary
Laplacian and positive viscosity in this source equation.

The updated manuscript credits the preceding constructions
of Diego C\'ordoba and Luis Mart\'inez-Zoroa, and the
hypodissipative work with Zheng. Separately, the primary
statement \cite{cqABStatement} attributes announced
smooth-forcing results for incompressible porous media,
Boussinesq, and three-dimensional incompressible Euler to
Levent Alp\"oge and Tristan Buckmaster, and credits the
C\'ordoba--Mart\'inez-Zoroa programme. Those are distinct
source attributions. The operator calculations here do not
decide any private-data or appropriation allegation.

\subsection{The retained spatial pullback, exact diffusion leakage, and feedback}
\label{sec:cq-spatial-diffusion-splitting}

Keep the original Cartesian coordinates $q=(x,y,w)$ and polynomials
\[
 Q=1+xy,\qquad
 \sigma(q)=\tfrac12\bigl[Q^3w+y^2Q(4+3xy)\bigr],\qquad
 \eta(s)=(0,0,2s).
\]
Put $\mathscr H=\mathcal O(\mathbb C_s)$ and
$\mathscr G=\mathcal O(\mathbb C_q^3)$, with uniform convergence on compact
sets, and define the actual pullback and section maps
\[
 \mathcal I =\sigma^*:\mathscr H\longrightarrow\mathscr G,\qquad
 E_\eta=\eta^*:\mathscr G\longrightarrow\mathscr H.
\]
Substitution gives $\sigma\eta(s)=s$, so $E_\eta \mathcal I =I$.
If $K\subset\mathbb C^3$ is compact, then $\sigma(K)$ is compact and
$\sup_K|\mathcal I h|=\sup_{\sigma(K)}|h|$. The corresponding identity for $\eta$
proves continuity of both maps. Thus $P=\mathcal I E_\eta$ is a continuous
projection and
\begin{equation}
 \mathscr G=\mathcal I \mathscr H\oplus\ker E_\eta .
 \label{eq:cq-spatial-splitting}
\end{equation}
Indeed $F=\mathcal I E_\eta F+(F-\mathcal I E_\eta F)$; the second summand is in
$\ker E_\eta$, and $\mathcal I h\in\ker E_\eta$ implies $h=0$.
These decomposition maps and their inverse, addition, are continuous.
In particular $\mathcal I \mathscr H=\ker(I-P)$ is closed. This analytic splitting
does not identify the source topology $\tau$ with a compact-open topology.

Retain
\[
 D=\partial_s,\quad
 \Delta=\partial_x^2+\partial_y^2+\partial_w^2,\quad
 g=\nabla\sigma,\quad N=\sum_{i=1}^3g_i^2,\quad L=\Delta\sigma.
\]
The sum defining $N$ is polynomial over $\mathbb C$ and is the Euclidean
squared norm on real $q$. Coordinate differentiation gives
\[
 \partial_i(\mathcal I h)=g_i\mathcal I h',\qquad
 \partial_i^2(\mathcal I h)=g_i^2\mathcal I h''+(\partial_i^2\sigma)\mathcal I h',
 \qquad \Delta \mathcal I h=N\mathcal I h''+L\mathcal I h'.
\]
The full coefficients before any evaluation are
\begin{align*}
 g_1&=\tfrac32yQ^2w+\tfrac12y^3(7+6xy),\\
 g_2&=\tfrac32xQ^2w+\tfrac12(8y+21xy^2+12x^2y^3),\\
 g_3&=\tfrac12Q^3,\\
 L&=3wQ(x^2+y^2)+3y^4+18x^2y^2+21xy+4 .
\end{align*}
At $\eta(s)$ these give $g=(0,0,1/2)$, $N=1/4$, $L=4$.
Consequently the compressed operator and the full complementary
component are
\begin{align}
 A=E_\eta\Delta \mathcal I &=\tfrac14D^2+4D:
                  \mathscr H\longrightarrow\mathscr H,
                  \label{eq:cq-spatial-compression}\\
 W=(I-P)\Delta \mathcal I ,\qquad
 Wh&=(N-\tfrac14)\mathcal I h''+(L-4)\mathcal I h':
                  \mathscr H\longrightarrow\ker E_\eta ,
                  \label{eq:cq-spatial-leakage}\\
 \Delta \mathcal I &=\mathcal I A+W. \notag
\end{align}
Every map is continuous because composition, differentiation, and
multiplication by a fixed polynomial are continuous on the specified
spaces.

\begin{cqproposition}
The exact kernel of $W$ is the constant functions. Equivalently,
\[
 \{h\in\mathscr H:\Delta \mathcal I h\in \mathcal I \mathscr H\}=\mathbb C.
\]
\end{cqproposition}
\begin{proof}
On the entire two-parameter family $(x,0,2s)$, direct substitution gives
\[
 \sigma(x,0,2s)=s,\quad
 N(x,0,2s)=\tfrac14+9x^2s^2,\quad
 L(x,0,2s)=4+6x^2s .
\]
Thus $Wh=0$, evaluated with $x=1$, implies
$9s^2h''+6sh'=0$ for every $s\in\mathbb C$.
Write the convergent entire Taylor series $h(s)=\sum_{n\ge0}a_ns^n$.
The coefficient of $s^n$ in this identity is $3n(3n-1)a_n$.
For each integer $n\ge1$ its scalar factor is nonzero, so every such
$a_n$ vanishes. Conversely all derivatives of a constant vanish and
therefore $Wh=0$. Finally, by the splitting,
$\Delta \mathcal I h\in \mathcal I \mathscr H$ is equivalent to $(I-P)\Delta \mathcal I h=0$.
\end{proof}

The splitting also supplies a derivative extension with an explicit
common domain and codomain:
\[
 B=\mathcal I DE_\eta:\mathscr G\longrightarrow\mathscr G,\qquad B\mathcal I =\mathcal I D .
\]
The commutator with the Euclidean Laplacian is consequently well typed.
For every $h\in\mathscr H$ it is
\begin{align}
 [\Delta,B]\mathcal I h
 &=\Delta \mathcal I h'-\mathcal I D(E_\eta\Delta \mathcal I h)\notag\\
 &=N\mathcal I h'''+L\mathcal I h''-\mathcal I (\tfrac14h'''+4h'')\notag\\
 &=(N-\tfrac14)\mathcal I h'''+(L-4)\mathcal I h''=W(h').
 \label{eq:cq-spatial-commutator}
\end{align}
The just-proved kernel formula gives the exact consequence
$[\Delta,B]\mathcal I h=0$ if and only if $h(s)=a+bs$ for some $a,b\in\mathbb C$.
This conclusion requires no invariance of $\mathcal I \mathscr H$ under $\Delta$.

Compression does not reproduce repeated diffusion. To compute the
feedback exactly, write $H=D_q^2\sigma$. The full product rule gives
\begin{align*}
 \Delta(N\mathcal I h'')&=(\Delta N)\mathcal I h''+
 2(g\cdot\nabla N)\mathcal I h^{(3)}+NL\mathcal I h^{(3)}+N^2\mathcal I h^{(4)},\\
 \Delta(L\mathcal I h')&=(\Delta L)\mathcal I h'+
 2(g\cdot\nabla L)\mathcal I h''+L^2\mathcal I h''+LN\mathcal I h^{(3)},\\
 \partial_jN&=2\sum_i g_iH_{ij},\\
 \Delta N&=2\sum_{i,j}H_{ij}^2+2g\cdot\nabla L,\qquad
 g\cdot\nabla N=2g^{\mathsf T}Hg.
\end{align*}
Adding gives the identity on every entire $h$,
\begin{align*}
 \Delta^2\mathcal I h
 &=N^2\mathcal I h^{(4)}
 +(4g^{\mathsf T}Hg+2NL)\mathcal I h^{(3)}\\
 &\quad+(2\sum_{i,j}H_{ij}^2+4g\cdot\nabla L+L^2)\mathcal I h''
 +(\Delta L)\mathcal I h'.
\end{align*}
Differentiating the original $g_i,L$ and evaluating at the same section
$\eta(s)$ gives
\[
 H(\eta(s))=
 \begin{pmatrix}0&3s&0\\3s&4&0\\0&0&0\end{pmatrix},\qquad
 \sum_{i,j}H_{ij}(\eta(s))^2=18s^2+16,
\]
\[
 (g^{\mathsf T}Hg)(\eta(s))=0,\quad
 (g\cdot\nabla L)(\eta(s))=0,\quad
 (\Delta L)(\eta(s))=24s .
\]
For the last equality, the term $3w(x^2+y^2)$ in $L$ contributes
$12w=24s$ after the two derivatives and substitution; all other terms
give zero there. Therefore
\begin{align}
 E_\eta\Delta^2\mathcal I 
 &=\tfrac1{16}D^4+2D^3+(36s^2+48)D^2+24sD,\notag\\
 E_\eta\Delta^2\mathcal I -A^2
 &=(36s^2+32)D^2+24sD=E_\eta\Delta W .
 \label{eq:cq-spatial-second-feedback}
\end{align}
The last equality also follows directly from $\Delta \mathcal I =\mathcal I A+W$.
It retains the component that leaves the pullback subspace and returns
at the second diffusion. It is an identity of differential operators;
no Euclidean heat semigroup on all of $\mathscr G$ is presumed.

Finally let $h\in\mathcal H_{\le1}$ and use the established entire
evolution $h_t=U_th$, $\partial_th_t=-D^2h_t/4$. For $\nu>0$ put
\[
 f(\tau,q)=h_{-4\nu\tau}(\sigma(q)).
\]
The time chain rule gives
$\partial_\tau f=\nu h_{-4\nu\tau}''(\sigma(q))$.
Thus the exact scalar viscous residual is
\begin{equation}
 (\partial_\tau-\nu\Delta)f
 =\nu\bigl[(1-N(q))h_{-4\nu\tau}''(\sigma(q))
          -L(q)h_{-4\nu\tau}'(\sigma(q))\bigr].
 \label{eq:cq-spatial-viscous-residual}
\end{equation}
\begingroup\emergencystretch=2em
For a specified smooth spatial advection field $V$, its complete
additional term is
\[
 (V\cdot\nabla\sigma)h_{-4\nu\tau}'(\sigma(q)).
\]
In particular the retained material identity $V\cdot\nabla\sigma=1$
gives coefficient $1-\nu L$ for the first derivative.
Setting $h=\Xi$ recovers $Z_{-4\nu\tau}$ with the factor $-4\nu$
and the original Cartesian metric intact. These are scalar equations
with explicitly computed forcing. They do not assert the vector
momentum equation, its pressure, divergence constraint, or the support,
initial-value, and energy conditions of a Navier--Stokes solution.
\par\endgroup

\section{The published fluid input and its complete arithmetic profile map}
\label{sec:cq-NS-profile-transfer}

We now apply the public fluid theorem recorded in
Section~\ref{sec:cq-public-NS-source} through a proved map.
The complete companion calculation \cite{cqNSFlow} was read.
The construction and every estimate used here are proved below,
including the polynomial coordinates, real flow, pressure,
forcing, physical metric, arithmetic Fourier constants, and inverse.
The final existence assertion cites the public Theorem~1.1 as
an external mathematical input; it is not an independent
verification of that theorem's complete proof.

\subsection{The original three coordinates and their full inverse differential}

In the original real coordinates $\mathbf x=(x,y,w)$ put
\begin{equation}
\begin{aligned}
 F_1&=(1+xy)^3w+y^2(1+xy)(4+3xy),\\
 F_2&=y+3x(1+xy)^2w+3xy^2(4+3xy),\\
 F_3&=2x-3x^2y-x^3w,\\
 \mathbf S&=(F_1/2,F_2,F_3)^{\mathsf T},\qquad J_S=D\mathbf S .
\end{aligned}
\label{eq:cq-NS-original-S}
\end{equation}
Thus $S_1=\sigma$ is exactly the earlier spatial coordinate.
The bold notation distinguishes this spatial three-vector
from the original quotient multiplication operator $S$.

Here is a complete determinant calculation. On $x\ne0$ set
$v=1+xy$, $h=2-3xy-x^2w$, and
$A=v^2+v-v^3h$, $B=4v+2-3v^2h$. Direct substitution gives
$F=(A/x^2,B/x,xh)$ and
\[
 \det\frac{\partial(x,v,h)}{\partial(x,y,w)}=-x^3,\qquad
 \det\frac{\partial F}{\partial(x,v,h)}
 =x^{-3}\det\begin{pmatrix}
 -2A&2v+1-3v^2h&-v^3\\
 -B&4-6vh&-3v^2\\h&0&1
 \end{pmatrix}=2x^{-3}.
\]
For the last equality the determinant's numerator is
\[
 (4-6vh)(-2v^2-2v+3v^3h)
 +(2v+1-3v^2h)(4v+2-6v^2h)=2,
\]
as expansion verifies. Thus $\det DF=-2$ on $x\ne0$ and,
by polynomial identity, everywhere. Consequently
$\det J_S=-1$ everywhere. Define the polynomial columns
\begin{equation}
\begin{aligned}
 W_1&=-\tfrac12\nabla F_2\times\nabla F_3,&
 W_2&=-\tfrac12\nabla F_3\times\nabla F_1,&
 W_3&=-\tfrac12\nabla F_1\times\nabla F_2,\\
 (DF)^{-1}&=[\,W_1\ W_2\ W_3\,],&
 J_S^{-1}&=[\,2W_1\ W_2\ W_3\,].
\end{aligned}
\label{eq:cq-NS-original-inverse}
\end{equation}
The scalar triple products with the three rows $\nabla F_i$
give $DF\,W_i=e_i$, proving both inverse identities.
In particular $\ker(d\sigma)_{\mathbf x}
=\operatorname{span}_{\mathbb R}\{W_2(\mathbf x),W_3(\mathbf x)\}$:
write any vector in the basis $(W_1,W_2,W_3)$ and use
$d\sigma(W_i)=\delta_{1i}/2$. No global inverse of $F$ or
$\mathbf S$ is presumed.

\subsection{The complete real flow and its volume-preserving inverse}

Fix $K\subset\mathbb R^3$ compact, $0<T<\infty$ and $\nu>0$.
Let $\mathcal U_{K,T,\nu}$ consist of smooth real triples
$(u,p,f)$ on $[0,T)\times\mathbb R^3$ satisfying
\begin{equation}
 \partial_tu_i+\sum_{\ell=1}^3u_\ell\partial_\ell u_i
 =\nu\Delta_{\mathbf x}u_i-\partial_i p+f_i,\qquad
 \sum_i\partial_i u_i=0,\qquad
 \operatorname{supp}u(t,\cdot)\subseteq K .
 \label{eq:cq-NS-input-class}
\end{equation}
There is no restriction on $p,f$ beyond this definition for
the correspondence; the public theorem supplies the stronger
support and initial conditions in its stated example.
For any such triple define
\[
 \partial_tX_t(q)=u(t,X_t(q)),\qquad X_0(q)=q,\qquad
 A_X(t,q)=D_qX_t(q),\qquad q\in\mathbb R^3 .
\]
For each $T_0<T$, every spatial derivative of $u$ is bounded
on $[0,T_0]\times K$ and vanishes outside $K$.
In particular its global spatial Lipschitz constant is
uniform on this time interval. On an interval of length
less than its reciprocal, the integral map
$Y\mapsto q+\int u(s,Y(s))\,ds$ is a contraction on
continuous curves in the uniform norm. Successive such
intervals construct the solution to $T_0$; boundedness of
$u$ prevents finite escape. The same construction backward
from each endpoint constructs its inverse. Differentiating
the integral equation gives a linear integral equation for
each first label derivative, and then successively for every
higher label and time derivative. The same local contraction
and induction prove smoothness and uniqueness of these
derivatives. Thus $X_t$ is a smooth global diffeomorphism
for every $t<T$.

Differentiation and the determinant product rule give
\[
 \partial_t A_X=Du(t,X_t)A_X,\quad A_X(0,q)=I_3,\qquad
 \partial_t\det A_X=(\operatorname{div}u)(t,X_t)\det A_X=0 .
\]
The determinant identity first holds where the matrix is
invertible, including a neighbourhood of $t=0$; its value
one prevents any first loss of invertibility and continues
the identity. Hence $\det A_X=1$ with positive orientation.
Every $q\notin K$ has its constant trajectory. An endpoint
outside $K$ also has a constant backward trajectory, so
inverse uniqueness proves
\begin{equation}
 X_t(K)=K,\qquad X_t(q)=q\ (q\notin K),\qquad
 X_t(L)\subseteq K\cup L\quad(L\subset\mathbb R^3\text{ compact}).
 \label{eq:cq-NS-flow-support}
\end{equation}
All these conclusions hold on each preterminal interval
without a uniform bound on $|u|$ as $t$ approaches $T$.

\subsection{The full complex fibre and both time derivatives}

Set $a_j(t,q)=S_j(X_t(q))$ and
$v_j(t,q)=(\nabla S_j\cdot u)(t,X_t(q))$.
Then $a_j$ is real and $\partial_ta_j=v_j$.
At fixed $t$ the map on the actual real and complex variables is
\[
 (q,\zeta)\longmapsto(q,a_j(t,q)+\zeta),\qquad
 (q,s)\longmapsto(q,s-a_j(t,q)).
\]
These are exact inverse maps, smooth in $q$ and holomorphic
on each complex fibre. Write
$\mathcal B=C^\infty(\mathbb R^3;\mathcal O(\mathbb C_\zeta))$
with seminorms consisting of suprema of finitely many label
and spectral derivatives on compact products.
The map
\[
 \mathcal P_t:\mathcal O(\mathbb C)^3\longrightarrow\mathcal B^3,
 \qquad \mathcal P_t(h_j)_j=(h_j(a_j(t,q)+\zeta))_j
\]
has the left inverse, for any fixed $q_*$,
$(\mathcal R_tG)_j(s)=G_j(q_*,s-a_j(t,q_*))$.
Substitution proves $\mathcal R_t\mathcal P_t=I$.
Images of compact products have compact spectral argument;
the chain rule involves only finitely many bounded derivatives
of $a_j$ there. This proves continuity of $\mathcal P_t$,
while evaluation and translation prove continuity of
$\mathcal R_t$. Therefore $\mathcal P_t$ is a topological
isomorphism onto its image with the inherited topology.
At $\zeta=0$ its first component is exactly
$h_1(\sigma(X_t(q)))=X_t^*\sigma^*h_1(q)$.

Keep a smooth, explicitly chosen real Newman-time function
$\vartheta:[0,T)\to\mathbb R$, independently of physical $t$.
For either original profile $h_\theta=Z_\theta$ or $K_\theta$,
define
\[
 \Psi_j(t,q,\zeta)=h_{\vartheta(t)}(\zeta+a_j(t,q)).
\]
The previously proved heat identity gives
\begin{equation}
 \partial_t\Psi_j
 =-\frac{\vartheta'}4\partial_\zeta^2\Psi_j
      +v_j\partial_\zeta\Psi_j .
 \label{eq:cq-NS-profile-generator}
\end{equation}
Since $v_j$ has no spectral dependence, another complete
differentiation gives
\begin{equation}
 \partial_t^2\Psi_j
 =\frac{(\vartheta')^2}{16}\partial_\zeta^4\Psi_j
 -\frac{\vartheta'v_j}{2}\partial_\zeta^3\Psi_j
 +\left(v_j^2-\frac{\vartheta''}{4}\right)\partial_\zeta^2\Psi_j
 +b_j\partial_\zeta\Psi_j,
 \label{eq:cq-NS-profile-second-generator}
\end{equation}
where, with every original momentum term,
\begin{equation}
 b_j=\partial_tv_j
 =\left[
 \sum_i(\nu\Delta u_i-\partial_i p+f_i)\partial_iS_j
 +\sum_{i,\ell}u_i u_\ell\partial_i\partial_\ell S_j
 \right](t,X_t(q)).
 \label{eq:cq-NS-profile-acceleration}
\end{equation}
Indeed the material derivative of $\nabla S_j\cdot u$
gives the Hessian term and the full material acceleration
of $u$; substituting \eqref{eq:cq-NS-input-class} gives
\eqref{eq:cq-NS-profile-acceleration}. In the second
derivative the mixed cubic term occurs twice, each with
coefficient $-\vartheta'v_j/4$, proving the factor $-1/2$.

The full Euclidean viscous operator in label coordinates is
\[
 G_X=A_X^{-1}A_X^{-\mathsf T},\qquad
 \mathcal L_t=\sum_{\alpha,\beta}
       \partial_{q_\alpha}(G_X^{\alpha\beta}\partial_{q_\beta}).
\]
For every smooth scalar $b$ one has
$\mathcal L_tX_t^*b=X_t^*\Delta b$. To prove it with all
first-order coefficient derivatives, test against
$\phi\in C_c^\infty(\mathbb R^3)$ and use the gradient chain
rule and the change of variables with determinant one:
\[
\begin{aligned}
 \int\phi\,\mathcal L_tX_t^*b\,dq
 &=-\int\nabla_q\phi\cdot A_X^{-1}(\nabla b)\circ X_t\,dq\\
 &=-\int\nabla_{\mathbf x}(\phi\circ X_t^{-1})\cdot\nabla b\,d\mathbf x
 =\int\phi\,X_t^*\Delta b\,dq .
\end{aligned}
\]
Smoothness makes this distributional identity pointwise.
The ordinary spatial chain rule now proves
\begin{equation}
\begin{aligned}
 \mathcal L_t\Psi_j
 &=(|\nabla S_j|^2\circ X_t)\partial_\zeta^2\Psi_j
        +(\Delta S_j\circ X_t)\partial_\zeta\Psi_j,\\
 (\partial_t-\nu\mathcal L_t)\Psi_j
 &=-\left(\frac{\vartheta'}4+\nu|\nabla S_j|^2\circ X_t\right)
          \partial_\zeta^2\Psi_j
       +(v_j-\nu\Delta S_j\circ X_t)\partial_\zeta\Psi_j .
\end{aligned}
\label{eq:cq-NS-material-viscous-residual}
\end{equation}

\subsection{Original Fourier constants and exact velocity recovery}

Define, with the original arithmetic kernel,
\[
 h_*(r)=\frac{\pi}{2}r^2(2\pi r^2-3)e^{-\pi r^2},\qquad
 k(\omega)=e^{\omega/2}\sum_{n\geq1}h_*(ne^\omega).
\]
For the Fourier convention with exponential $e^{-2\pi ir\xi}$,
the second and fourth Gaussian derivative formulas are
\[
 \widehat{r^2e^{-\pi r^2}}
   =(1/(2\pi)-\xi^2)e^{-\pi\xi^2},\quad
 \widehat{r^4e^{-\pi r^2}}
   =(\xi^4-3\xi^2/\pi+3/(4\pi^2))e^{-\pi\xi^2}.
\]
They follow by differentiating the Gaussian transform twice
and four times and dividing by $(-2\pi i)^2$ and
$(-2\pi i)^4$, respectively. The coefficients
$\pi^2$ and $-3\pi/2$ then give $\widehat h_*=h_*$ and
$h_*(0)=0$. Poisson summation gives
$\sum_{n\geq1}h_*(n\lambda)
=\lambda^{-1}\sum_{n\geq1}h_*(n/\lambda)$, and hence
$k(\omega)=k(-\omega)$. Poisson summation for this Schwartz
function follows directly by periodizing it: its Fourier
coefficients are the displayed Fourier transform at integer
frequencies; the function and coefficient series converge
absolutely with derivatives, so evaluation at zero is valid.

Each fixed number of derivatives of the defining series is
locally uniformly convergent. At $\omega\geq0$ its Gaussian
terms give
$|k^{(r)}(\omega)|\leq C_r e^{C_r\omega}
 e^{-\pi e^{2\omega}}$; factor the $n=1$ exponential and
bound the remaining polynomially weighted series by
$\sum n^{M_r}e^{-\pi(n^2-1)}<\infty$.
Evenness gives the same bound with $|\omega|$ at the other end.
Expanding the definition proves
$\Phi(u)=2k(2u)$ with every factor retained. Consequently
\begin{equation}
\begin{aligned}
 Z_\theta(s)&=4\int_{\mathbb R}k(\omega)
                       e^{\theta\omega^2/4}e^{is\omega}\,d\omega,\\
 K_\theta(s)&=4\int_{\mathbb R}m(\omega)k(\omega)
                       e^{\theta\omega^2/4}e^{is\omega}\,d\omega ,
\end{aligned}
\label{eq:cq-NS-original-Fourier}
\end{equation}
where $m$ is exactly the three-shift, four-term multiplier
in \eqref{eq:cq-xi4-full-Fourier-kernel}, including
$-i\omega/336$ at shift $3$.
All derivatives of these amplitudes are Schwartz for real
$\theta$, locally uniformly in $\theta$; the same is true
of their Fourier transforms by integration by parts.

Set $d_h(\theta)^2=\int_{\mathbb R}|h_\theta'(s)|^2\,ds$.
The exact constants are
\begin{equation}
\begin{aligned}
 d_Z(\theta)^2&=32\pi\int_{\mathbb R}\omega^2 k(\omega)^2
                                   e^{\theta\omega^2/2}\,d\omega,\\
 d_K(\theta)^2&=32\pi\int_{\mathbb R}\omega^2|m(\omega)|^2
                            k(\omega)^2e^{\theta\omega^2/2}\,d\omega .
\end{aligned}
\label{eq:cq-NS-profile-norm-constants}
\end{equation}
For completeness, insert $e^{-\epsilon s^2}$ in the
squared norm of $4\int b(\omega)e^{is\omega}\,d\omega$.
Fubini and the Gaussian integral give
\[
 16\sqrt{\pi/\epsilon}\iint
 b(\omega)\overline{b(\eta)}
             e^{-(\omega-\eta)^2/(4\epsilon)}\,d\omega\,d\eta .
\]
The Gaussian convolution has mass $2\pi$ and tends to
$2\pi b$ in this integral: split at fixed distance from
zero, use uniform continuity on the near part and its
Gaussian tail on the far part, with Schwartz domination.
On the original side dominated convergence applies because
the Fourier transform is Schwartz. The result is
$32\pi\int|b|^2$. Taking
$b=i\omega k e^{\theta\omega^2/4}$ or
$i\omega m k e^{\theta\omega^2/4}$ proves
\eqref{eq:cq-NS-profile-norm-constants}.
Both integrals are strictly positive:
$k(0)=\Phi(0)/2>0$ and $m(0)=127/219024>0$
give an interval containing nonzero frequencies where the
integrands' non-Gaussian factors are positive.
The two-ended bounds prove continuity in $\theta$ by
domination on compact parameter intervals. Thus on each
compact $I\subset\mathbb R$ there are
$0<d_-\leq d_h(\theta)\leq d_+<\infty$.

Define the full arithmetic temporal defect on the real
spectral line by
\[
 \mathscr D_j=\partial_t\Psi_j+
                \frac{\vartheta'}4\partial_s^2\Psi_j
        =v_j h_{\vartheta}'(s+a_j).
\]
Translation by the real $a_j$ preserves Lebesgue measure.
Therefore
\begin{equation}
\begin{aligned}
 \sum_j\|\mathscr D_j(t,q,\cdot)\|_2^2
       &=d_h(\vartheta(t))^2|v(t,q)|^2,\\
 v_j(t,q)&=\frac{\int_{\mathbb R}
       \overline{h_{\vartheta}'(s+a_j)}\,\mathscr D_j(t,q,s)\,ds}
                {d_h(\vartheta)^2},\\
 u(t,X_t(q))&=2W_1(X_t(q))v_1+
                    W_2(X_t(q))v_2+W_3(X_t(q))v_3 .
\end{aligned}
\label{eq:cq-NS-exact-velocity-recovery}
\end{equation}
These formulas prove an inverse correspondence
$(X_t,u\circ X_t)\leftrightarrow(X_t,\mathscr D)$ on the
three real profile lines at each $(t,q)$. The flow and its
inverse sheet are retained in both sides.

For nonempty $K$, put
$M_J=\max_K\|J_S\|_{\rm op}$ and
$N_J=\max_K\|J_S^{-1}\|_{\rm op}$.
They are finite and positive. Define
\[
 \mathcal A_\infty(t)=
 \sup_q\left(\sum_j\|\mathscr D_j(t,q,\cdot)\|_2^2\right)^{1/2},
 \quad
 \mathcal A_2(t)=
 \left(\sum_j\int_{\mathbb R^3}\int_{\mathbb R}
               |\mathscr D_j|^2\,ds\,dq\right)^{1/2}.
\]
When $\vartheta([0,T))\subseteq I$, the complete estimates are
\begin{equation}
 \frac{d_-}{N_J}\|u(t)\|_\infty\leq\mathcal A_\infty(t)
       \leq d_+M_J\|u(t)\|_\infty,\qquad
 \frac{d_-}{N_J}\|u(t)\|_2\leq\mathcal A_2(t)
       \leq d_+M_J\|u(t)\|_2 .
 \label{eq:cq-NS-energy-and-supremum}
\end{equation}
Indeed $v=J_S(X_t)(u\circ X_t)$, both vanish off $K$,
and $|J_Su|\leq M_J|u|$, $|u|\leq N_J|J_Su|$ on $K$.
The supremum is preserved by the diffeomorphism $X_t$.
For the integral, its determinant one gives the exact identity
\[
 \mathcal A_2(t)^2
 =d_h(\vartheta(t))^2
           \int_{\mathbb R^3}|J_S(\mathbf x)u(t,\mathbf x)|^2\,d\mathbf x .
\]
This proves both bounds without discarding any spatial or
spectral component. For a fixed original incidence coefficient
$a_{ef}$, the observation $a_{ef}K_\theta$ multiplies
$\mathscr D^K$ by $a_{ef}$ and its squared norm by
$|a_{ef}|^2$; inversion uses its actual nonzero value, and
the zero incidence has exactly the zero observation.

\subsection{Application to the stated public singularity theorem}

\begin{cqtheorem}[Arithmetic consequence of the cited public input]
\label{thm:cq-NS-public-propagation}
Take the velocity, pressure, force and compact set supplied
by Theorem~1.1 of \cite{cqNS2026}, with its original $\nu>0$.
Choose the explicit Newman time $\vartheta(t)=0$ and the
actual arithmetic profile $h_0=\Xi$, on $0\leq t<1$.
Then the constructed entire spectral fields satisfy
\begin{equation}
 \mathscr D_j=\partial_t\Psi_j,\qquad
 \mathscr D_j(0,q,s)=0,\qquad
 \sup_{t<1}\mathcal A_2(t)<\infty,\qquad
 \limsup_{t\uparrow1}\mathcal A_\infty(t)=\infty .
 \label{eq:cq-NS-public-arithmetic-blowup}
\end{equation}
Every $\mathscr D_j$ is supported in $q\in K$.
For every compact $L\subset\mathbb R^3$, compact
$E_0\subset\mathbb C$ and integer $r\geq0$,
\begin{equation}
 \sup_{t<1,\ q\in L,\ \zeta\in E_0}
             |\partial_\zeta^r\Psi_j(t,q,\zeta)|<\infty .
 \label{eq:cq-NS-bounded-spectral-profiles}
\end{equation}
The analogous temporal-defect conclusions hold for the
original full profile $K_0$ with its exact $d_K(0)$.
\end{cqtheorem}

\begin{proof}
The cited public theorem provides a triple in
\eqref{eq:cq-NS-input-class} with $T=1$, $u(0)=0$,
uniformly bounded $\|u(t)\|_2$, and unbounded terminal
$\|u(t)\|_\infty$. The preceding explicit construction
therefore supplies $X_t,\Psi,\mathscr D$. With the chosen
constant Newman time the heat part of the temporal defect
vanishes exactly, and $u(0)=0$ gives $\mathscr D(0)=0$.
Set $d_-=d_+=d_Z(0)>0$ in
\eqref{eq:cq-NS-energy-and-supremum}. Its two inequalities
give the energy bound and the unbounded terminal supremum
in \eqref{eq:cq-NS-public-arithmetic-blowup}.
Support follows from $v=0$ off $K$.
For the last bound,
\eqref{eq:cq-NS-flow-support} places the argument of
$\Xi^{(r)}$ in the fixed compact set $S_j(K\cup L)+E_0$.
Entire continuity gives its finite maximum.
All arguments apply to $K_0$ using the already proved
positive constant $d_K(0)$ and its entire derivatives.
This proof uses the external existence theorem exactly at
its first step; all transfer formulas and estimates are
proved within this fragment.
\end{proof}

\subsection{The actual cutoff-summed source witness and its arithmetic image}
\label{sec:cq-NS-actual-witness}

\subsubsection{The precise imported witness and its unchanged viscosity}

We now use the particular field constructed in the
166-page revision of \cite{cqNS2026}. The original
cutoff-summed fields, the smooth extension of the force,
and their all-order estimates are the imported source input.
Its actual closing equations (9.21), (10.4)--(10.14) and
(10.20)--(10.23) were read directly. The complete companion
specialization \cite{cqNSWitness} was also read; all arithmetic,
flow, energy and terminal maps used here have their proofs below.
No independent verification of every lemma in the public
166-page existence proof is claimed.

Use a star to retain the source's viscosity-one fields, rather than
overwriting the physical viscosity $\nu>0$. In source equation (9.21),
page 115, the actual summed Cartesian vector potential, azimuthal field,
and pressure are denoted by $A_*,B_*e_\theta,p_{*,\mathrm{loc}}$.
These are the cutoff-summed corrected fields of that equation, including
all stages and the finite initialization.
With the source's own radial-time coordinate denoted $q_{\rm src}$
to distinguish it from our material label $q$, its exact sums are
\[
\begin{aligned}
 A_*&=A_{*,0}+\sum_{j\geq1}\chi(a_jq_{\rm src})A_{*,j},\\
 B_*e_\theta&=B_{*,0}e_\theta+
                        \sum_{j\geq1}\chi(a_jq_{\rm src})B_{*,j},\\
 p_{*,\mathrm{loc}}&=p_{*,0}+
                        \sum_{j\geq1}\chi(a_jq_{\rm src})p_{*,j}.
\end{aligned}
\]
The displayed $a_j,\chi$ and all finite-state representatives
are the actual choices in source (9.21); the subscript star
marks viscosity one and introduces no replacement of them.
 Source equation (10.4), page
118, constructs
\[
 u_* = \operatorname{curl}(c_*A_*)+c_*B_*e_\theta,
 \qquad p_*=c_*p_{*,\mathrm{loc}},\qquad c_*=\chi_x\chi_t.
\]
Here the source spatial cutoff is smooth in $(r^2,z)$, supported in
$r<r_0,\ |z|<z_0$, and equals one on
$r\leq r_0/2,\ |z|\leq z_0/2$. Its time cutoff is zero for
$1-t\geq\tau_0$ and one for $0\leq1-t\leq\tau_0/2$, where
$0<\tau_0<1/2$. Retain the actual source choices of these parameters.
The source force before time one is exactly
\[
 f_* = \partial_tu_*+(u_*\cdot\nabla)u_*
                     -\Delta u_*+\nabla p_*.
\]
Thus every cutoff derivative belongs to the force. Its smooth extension
through time one is source (10.11), page 120, with all its prescribed
Taylor coefficients and cutoff sequence. The source proves that its
support is contained in $K_*\times[0,2]$, where
$K_*=\operatorname{supp}\chi_x$, and it vanishes near time zero.
The pressure is the displayed compactly supported source pressure;
no divergence-free condition is imposed on $f_*$.

For the physical parameter $\nu>0$, take precisely source (10.22):
\begin{align*}
 y&=x/\sqrt\nu, & x&=\sqrt\nu\,y,\\
 u_\nu(x,t)&=\sqrt\nu\,u_*(x/\sqrt\nu,t),&
 p_\nu(x,t)&=\nu p_*(x/\sqrt\nu,t),\\
 f_\nu(x,t)&=\sqrt\nu\,f_*(x/\sqrt\nu,t),&
 K_\nu&=\sqrt\nu\,K_*.
\end{align*}
The inverse multiplies the velocity and force by $\nu^{-1/2}$,
the pressure by $\nu^{-1}$, and substitutes $x=\sqrt\nu\,y$;
physical time is unchanged. The chain rule gives
\[
 \partial_tu_\nu+(u_\nu\cdot\nabla_x)u_\nu
       -\nu\Delta_xu_\nu+\nabla_xp_\nu
 =\sqrt\nu\,[\partial_tu_*+(u_*\cdot\nabla_y)u_*
                    -\Delta_yu_*+\nabla_yp_*](y,t)=f_\nu.
\]
Also $\operatorname{div}_xu_\nu=\operatorname{div}_yu_*=0$,
$u_\nu(0)=0$, and $u_\nu,p_\nu$ have support in $K_\nu$ for every
$t<1$. These are exactly the inputs of Section~\ref{sec:cq-NS-profile-transfer} with
$T=1$, with all coordinate and parameter maps specified.

Let
\[
 F_\nu(t)=\int_0^t\|f_\nu(s)\|_{L^2(\mathbb R^3)}\,ds,
 \qquad E_\nu=F_\nu(1)=\nu^{5/4}F_*(1)<\infty.
\]
The last exponent follows from the factor $\sqrt\nu$ in $f_\nu$ and
$dx=\nu^{3/2}dy$. For completeness the energy estimate used below is
proved directly for this input. Compact support and incompressibility
give, by integrating the equation against $u_\nu$,
\[
 \frac12\frac{d}{dt}\|u_\nu(t)\|_2^2
       +\nu\|\nabla u_\nu(t)\|_2^2
       =\langle f_\nu(t),u_\nu(t)\rangle.
\]
The convection integral is the integral of
$\operatorname{div}(u_\nu|u_\nu|^2/2)$, and the pressure integral is
minus the integral of $p_\nu\operatorname{div}u_\nu$; both vanish.
Divide by $(\|u_\nu\|_2^2+\delta^2)^{1/2}$, use Cauchy--Schwarz,
integrate from zero, and let $\delta\downarrow0$. This proves
$\|u_\nu(t)\|_2\leq F_\nu(t)$. Substitution back into the integrated
identity proves, with its full dissipation coefficient,
\[
 \|u_\nu(t)\|_2^2+2\nu\int_0^t\|\nabla u_\nu(s)\|_2^2ds
 \leq 2\int_0^t F_\nu'(s)F_\nu(s)ds=F_\nu(t)^2.
\]

\subsubsection{The source's actual growing component in two arithmetic channels}

The source fixes $0<h<1/100$, $X_{\rm in}>0$, and $e_0>0$ through
its construction. Its closing proof, (10.20)--(10.21), page 124, gives
the actual corrected swirl along
\[
 x_{*,\tau}=(\sqrt{2X_{\rm in}\tau},0,0),\qquad t=1-\tau,
 \qquad
 (u_*)_2(x_{*,\tau},1-\tau)
   =\tau^{-1/2-h}(e_0+R_*(\tau)).
\]
The original remainder is retained as the function
$R_*(\tau)=\tau^{1/2+h}(u_*)_2(x_{*,\tau},1-\tau)-e_0$.
The source estimate supplies constants $C_*>0,\tau_*>0$ with
$|R_*(\tau)|\leq C_*\tau^{2h}$ for $0<\tau<\tau_*$.
At these points the cylindrical angle is zero, so
$e_\theta=(0,1,0)$ and the displayed Cartesian component is precisely
the source swirl. The annular and higher-order corrections have already
been accounted for in this remainder. The parameter $h$ is the chosen
source parameter; no claim that an arbitrary prescribed $h$ works is
needed or made.

Define
\[
 r_{\nu,\tau}=\sqrt{2\nu X_{\rm in}\tau},\qquad
 x_{\nu,\tau}=(r_{\nu,\tau},0,0).
\]
The exact viscosity map proves
\begin{equation}
 (u_\nu)_2(x_{\nu,\tau},1-\tau)
 =\sqrt\nu\,\tau^{-1/2-h}(e_0+R_*(\tau)).
 \label{eq:cq-nw:actualswirl}
\end{equation}
Use the actual material flow $X_{\nu,t}$ of $u_\nu$ proved to exist
on each interval before one in Section~\ref{sec:cq-NS-profile-transfer}. Put
\[
 X_{\nu,t}(q)=\sqrt\nu\,X_{*,t}(q/\sqrt\nu).
\]
This is an equality of the actual flows: its right side equals $q$ at
zero, and its time derivative is
$\sqrt\nu\,u_*(t,X_{*,t}(q/\sqrt\nu))
 =u_\nu(t,\sqrt\nu\,X_{*,t}(q/\sqrt\nu))$.
The previously proved uniqueness gives the equality, with inverse
$X_{\nu,t}^{-1}(x)=\sqrt\nu\,X_{*,t}^{-1}(x/\sqrt\nu)$.
Define the labels
\[
 q_{\nu,\tau}=X_{\nu,1-\tau}^{-1}(x_{\nu,\tau}).
\]
This is the exact inverse between the Eulerian sampling points and
their material labels. It is a family of labels depending on $\tau$;
the source path has not been asserted to be a single fluid trajectory.

Write $S=\mathbf S=(F_1/2,F_2,F_3)$ and $J=J_S=DS$ in the physical Cartesian
coordinates $x$, not the viscosity-one coordinate $y$. Direct
substitution into the original polynomials gives
\[
 S(r,0,0)=(0,0,2r),\qquad
 J(r,0,0)=\begin{pmatrix}
 0&0&1/2\\0&1&3r\\2&-3r^2&-r^3
 \end{pmatrix}.
\]
Indeed the first row comes from differentiating
$F_1/2=((1+xy)^3w+y^2(1+xy)(4+3xy))/2$;
the second comes from
$F_2=y+3x(1+xy)^2w+3xy^2(4+3xy)$; the third comes from
$F_3=2x-3x^2y-x^3w$, all evaluated at $(x,y,w)=(r,0,0)$.
Every entry, including the off-diagonal powers of $r$, is retained.
Writing $v=J(r,0,0)u$ therefore gives the exact component inverse
\begin{align*}
 u_3&=2v_1, &u_2&=v_2-6rv_1,\\
 u_1&=\tfrac12v_3+\tfrac32r^2v_2-8r^3v_1.
\end{align*}
Substituting these expressions into all three rows proves both
directions of this matrix inverse.

Fix actual arithmetic time zero. Take $H=Z_0=\Xi$ or $H=K_0$ with
the complete four-term coefficient from the preceding proof, and write
\[
 d_H^2=\int_\mathbb R|H'(s)|^2ds>0,\qquad
 \Psi_{\nu,j}(t,q,\zeta)=H(S_j(X_{\nu,t}(q))+\zeta),
 \qquad D_{\nu,j}=\partial_t\Psi_{\nu,j}.
\]
The exact positive constants are those of \eqref{eq:cq-NS-profile-norm-constants} at $\theta=0$,
including $32\pi$ and, for $K_0$, the entire factor $|m|^2$.
At the particular pair $(t,q)=(1-\tau,q_{\nu,\tau})$, the first two
shifts are both zero. Consequently define the actual integral
\[
 I_{\nu,j}(\tau)=\frac{1}{d_H^2}\int_\mathbb R
       \overline{H'(s)}D_{\nu,j}(1-\tau,q_{\nu,\tau},s)ds,
 \qquad j=1,2.
\]
The chain rule gives $D_{\nu,j}=v_jH'(s)$ there. It follows, without
an asymptotic replacement of the profiles, that
\begin{equation}
 I_{\nu,2}(\tau)-6r_{\nu,\tau}I_{\nu,1}(\tau)
 =\sqrt\nu\,\tau^{-1/2-h}(e_0+R_*(\tau)).
 \label{eq:cq-nw:integralblowup}
\end{equation}
There is also a whole-profile identity, stronger than its projection:
\begin{equation}
 \begin{split}
 &D_{\nu,2}(1-\tau,q_{\nu,\tau},s)
      -6r_{\nu,\tau}D_{\nu,1}(1-\tau,q_{\nu,\tau},s)\\
 &\hspace{12mm}=\sqrt\nu\,\tau^{-1/2-h}
                      (e_0+R_*(\tau))H'(s).
 \end{split}
 \label{eq:cq-nw:profileblowup}
\end{equation}
Its squared $L^2(ds)$ norm is exactly
$\nu\tau^{-1-2h}(e_0+R_*(\tau))^2d_H^2$.
Thus the source's actual component and remainder, not just an abstract
possibility of growth, have been carried into the original arithmetic
profile. The complex conjugation and integration above retain the
complex values of $K_0$ when that channel is used.

Choose any $0<\tau_{**}\leq\tau_*$ such that
$C_*\tau_{**}^{2h}\leq e_0/2$ and the source cutoffs equal one on
the path for $\tau<\tau_{**}$. Cauchy--Schwarz on the two component
coefficients $(-6r_{\nu,\tau},1)$ yields
\begin{equation}
 \left(\sum_{j=1}^3\|D_{\nu,j}
              (1-\tau,q_{\nu,\tau},\cdot)\|_2^2\right)^{1/2}
 \geq
 \frac{d_H\sqrt\nu\,e_0}
      {2\sqrt{1+72\nu X_{\rm in}\tau}}\,
              \tau^{-1/2-h}.
 \label{eq:cq-nw:lower}
\end{equation}
This lower bound proves divergence in the spatial-supremum,
spectral-$L^2$ norm for the imported constructed field. It is not a
matching upper asymptotic for its global supremum. All forcing,
pressure, and material diffusion terms remain those of
\eqref{eq:cq-NS-profile-acceleration} and \eqref{eq:cq-NS-material-viscous-residual}, with the explicit
$(u_\nu,p_\nu,f_\nu)$ just given.

\subsubsection{A terminal arithmetic profile with a bounded energy norm}

The same constructed input gives a terminal value in a precise
function space. This assertion can be proved from its energy identity
without controlling its unbounded spatial supremum. Let
$M_\nu=\max_{K_\nu}\|J\|_{\mathrm{op}}$ and put
\[
 \mathcal H_\nu=L^2(K_\nu\times\mathbb R,dq\,ds;\mathbb C^3),\qquad
 B_{\nu,j}(t,q,s)=\Psi_{\nu,j}(t,q,s)-H(S_j(q)+s).
\]
Subtracting the displayed initial value defines an affine coordinate
with the explicit inverse $B\mapsto B+(H(S_j(q)+s))_j$.
It does not discard the initial profile. Its domain and reference
element are specified; outside $K_\nu$ the difference is identically
zero. The full profile on $K_\nu$ is also square integrable because
real translation gives norm $|K_\nu|^{1/2}\|H\|_2$ in each channel.

The already proved exact identity gives
\[
 \|D_\nu(t)\|_{\mathcal H_\nu}^2
  =d_H^2\int_{K_\nu}|J(x)u_\nu(t,x)|^2dx
  \leq d_H^2M_\nu^2E_\nu^2.
\]
The identity follows by integrating the real spectral translation and
then using $\det D_qX_{\nu,t}=1$; no spatial coordinate is removed.
On a closed time interval before one the functions are smooth and the
Schwartz bounds justify their Hilbert-space derivatives. Integrating
$\partial_tB_\nu=D_\nu$ there and using the triangle inequality for
the integral yields, for $0\leq a<b<1$,
\begin{equation}
 \|B_\nu(b)-B_\nu(a)\|_{\mathcal H_\nu}
       \leq d_H M_\nu E_\nu(b-a).
 \label{eq:cq-nw:terminalrate}
\end{equation}
Completeness constructs a unique terminal element $B_\nu(1)$ and
the exact bound
$\|B_\nu(1)-B_\nu(t)\|\leq d_HM_\nu E_\nu(1-t)$.
This constructs a Lipschitz extension on $[0,1]$ in the specified
Hilbert space, while \eqref{eq:cq-nw:lower} proves divergence in the stronger
spatial-supremum norm of its time derivative.

We also identify the terminal element through an actual map on labels.
Volume preservation gives
\[
 \int_{K_\nu}\int_0^1|u_\nu(t,X_{\nu,t}(q))|dt\,dq
 =\int_0^1\|u_\nu(t)\|_{L^1(K_\nu)}dt
 \leq |K_\nu|^{1/2}E_\nu.
\]
The nonnegative integral is finite, so for almost every $q$ the full
trajectory has finite total length before one. Its integral equation
therefore constructs the limit $X_{\nu,1}(q)\in K_\nu$ as $t\uparrow1$.
Set its value arbitrarily in $K_\nu$ on the null exceptional set.
Also the integral triangle inequality and volume preservation give
\[
 \|X_{\nu,b}-X_{\nu,a}\|_{L^2(K_\nu)}
       \leq\int_a^b\|u_\nu(t)\|_2dt
       \leq E_\nu(b-a),
\]
and the same estimate with $b=1$ follows by the constructed limit and
Fatou's lemma. For every continuous $g$ on compact $K_\nu$, dominated
convergence gives
\[
 \int_{K_\nu}g(X_{\nu,1}(q))dq
   =\lim_{t\uparrow1}\int_{K_\nu}g(X_{\nu,t}(q))dq
   =\int_{K_\nu}g(x)dx.
\]
This identifies the pushforward of Lebesgue measure exactly. To see
that continuous tests suffice, for each relatively closed set $C$ use
the continuous decreasing approximations
$g_n(x)=\max(0,1-n\operatorname{dist}(x,C))$ to its indicator, then
monotone or dominated convergence; complements give open sets and
the uniqueness of finite measures on their generated Borel sigma
algebra gives the stated pushforward.

For almost every $q$ the pointwise terminal entire fibre is
\[
 \Psi_{\nu,j}(1,q,\zeta)
      =H(S_j(X_{\nu,1}(q))+\zeta).
\]
It agrees with the Hilbert-space limit: real translations are continuous
in $L^2(ds)$, as follows from
$\|H(\cdot+a)-H(\cdot+b)\|_2\leq d_H|a-b|$ by integrating $H'$.
The shifts converge almost everywhere, remain bounded on $K_\nu$,
and their translated-profile squared differences are bounded by
$4\|H\|_2^2$. Dominated convergence proves the identification in
$\mathcal H_\nu$. Replacing $H$ by $H^{(n)}$ in this proof proves the
same assertions for every pure spectral derivative, with the retained
constant $\|H^{(n+1)}\|_2M_\nu E_\nu$.

The terminal label map has a further exact operator correspondence.
On $L^2(K_\nu)$ define $V_tg=g\circ X_{\nu,t}$, including $t=1$.
For $t<1$ its inverse is composition with $X_{\nu,t}^{-1}$, so it is
unitary. The just-proved measure identity makes $V_1$ an isometry.
For continuous $g$, dominated convergence gives $V_tg\to V_1g$ in
$L^2$. Continuous functions are dense for Lebesgue measure restricted
to a compact set: approximate measurable-set indicators by continuous
distance cutoffs between compact subsets and open supersets using
finite-measure regularity, then approximate simple functions. The
isometric bounds extend the convergence to every $g\in L^2(K_\nu)$.
Thus $V_t\to V_1$ strongly. Its Hilbert adjoint satisfies
$V_1^*V_1=I$ by preservation of inner products, and $V_1V_1^*$ is the
orthogonal projection onto its closed range: it is self-adjoint,
idempotent, acts as identity on that range, and annihilates its
orthogonal complement. These are the proved terminal inverse and
range statements. No pointwise inverse of $X_{\nu,1}$, or surjectivity
of $V_1$, has been supplied by this argument.

\subsubsection{The actual zero image, including the terminal fibres}

For every time before one and almost every terminal label, the shift
$a=S_j(X_{\nu,t}(q))$ is real. The full coordinate map is
$\zeta\mapsto a+\zeta$ with inverse $s\mapsto s-a$. If
$H(s)=(s-z_0)^mA_0(s)$ with $A_0(z_0)\ne0$, the actual image is
\[
 H(a+\zeta)=(\zeta-(z_0-a))^mA_0(a+\zeta).
\]
This retains the whole analytic unit and all its derivatives. In
particular its zero has imaginary part $\operatorname{Im}z_0$ at all
these times, even for the source input whose arithmetic time-derivative
combination obeys \eqref{eq:cq-nw:profileblowup}. With $H=\Xi$ the precise
completed-zeta argument is $\rho=1/2+i(a+\zeta)$, and
$\operatorname{Re}\rho=1/2-\operatorname{Im}\zeta$. Thus the actual
source witness produces the proved derivative singularity and terminal
profile above; this explicit map preserves zero heights. A zero-height
change is not obtained by substituting a faster-growing velocity into
this same real-translation map.

\section{The actual source quotient of a flow-driven arithmetic profile}
\label{sec:cq-flow-source-quotient}

\subsection{Translation, heat, and the original spectral generator}

For $a,\theta\in\mathbb C$ define, with the input coordinate $s$ and
output coordinate $\zeta$ retained,
\[
 (T_ah)(\zeta)=h(\zeta+a),\qquad
 V_{\theta,a}=T_aU_\theta:\cqE\longrightarrow\cqE.
\]
The derivative is the complex derivative in the displayed coordinate.
Translation preserves the original order class. More precisely, for
$1<p<2$ and $c>0$ the already defined growth seminorms satisfy
\[
 \|T_ah\|_{p,c}\leq
 e^{c|a|^p}\|h\|_{p,c\,2^{1-p}},
\]
because $|\zeta+a|^p\leq2^{p-1}(|\zeta|^p+|a|^p)$.
Thus $T_a$ and $T_{-a}$ are continuous for $\gamma$. Ordinary
differentiation gives $DT_a=T_aD$; continuity and the convergent
heat series give $T_aU_\theta=U_\theta T_a$. Consequently
\[
 V_{\theta,a}^{-1}=U_{-\theta}T_{-a},\qquad
 V_{\theta,a}D=DV_{\theta,a}.
\]
These assertions hold on all of $\cqE$ and do not impose continuity
of translation or differentiation for the unnamed source topology.

Retain the actual $\tau$, $\cqN$ and $\cqQ$, and define their exact
transported presentations
\begin{equation}
 \tau_{\theta,a}=(V_{\theta,a})_*\tau,\qquad
 N_{\theta,a}=V_{\theta,a}\cqN,\qquad
 Q_{\theta,a}=(\cqE,\tau_{\theta,a})/N_{\theta,a}.
 \label{eq:cq-flow-moving-source}
\end{equation}
Then $N_{\theta,a}$ is closed, and
\[
 \overline V_{\theta,a}:\cqQ\longrightarrow Q_{\theta,a},
 \qquad[h]\longmapsto[V_{\theta,a}h]
\]
is a complex-linear homeomorphism. Its inverse is induced by the
displayed inverse on representatives. The original arithmetic
multiplication has the exact conjugate
\begin{equation}
 B_{\theta,a}=V_{\theta,a}sV_{\theta,a}^{-1}
              =\zeta+a-\frac{\theta}{2}D_\zeta .
 \label{eq:cq-flow-spectral-generator}
\end{equation}
Indeed the previous heat commutator gives
$U_\theta sU_{-\theta}=s-\theta D/2$; translation sends multiplication
by $s$ to multiplication by $\zeta+a$ and commutes with $D$.
This full combination preserves $N_{\theta,a}$ and is continuous for
$\tau_{\theta,a}$ by its conjugacy. The induced quotient operator is
similar through $\overline V_{\theta,a}$ to the source $S$, so its
spectrum is exactly the original $s$-spectrum. To verify both
directions, conjugate a continuous two-sided inverse of $S-\lambda$
by $\overline V_{\theta,a}$, or conjugate an inverse of the transported
operator back. At $\theta=0$, multiplication by $\zeta$ itself
descends, and its spectrum is $\operatorname{Spec}(S)-a$ because
$B_{0,a}=\zeta+a$.

Along the retained real flow let
\[
 a_j(t,q)=S_j(X_t(q)),\qquad
 v_j(t,q)=\partial_ta_j(t,q),\qquad
 \Psi_j(t,q,\zeta)=V_{\vartheta(t),a_j(t,q)}\Xi(\zeta).
\]
Here $S_j$ denotes the original polynomial target component, whereas
the unindexed $S$ above denotes the arithmetic quotient operator.
The source quotient is taken fibre by fibre at each fixed $(t,q,j)$.
If a topology on the whole family of quotient fibres is needed, give
their disjoint union the topology transported from
$\mathbb C^2\times\cqQ$ by
$((\theta,a),[h])\mapsto((\theta,a),[V_{\theta,a}h])$.
This specifies joint continuity by an exact trivialization and does
not assert joint parameter continuity for the untransported $\tau$.

Because $\Xi\in\cqN$, every profile represents zero in its actual
transported quotient:
\begin{equation}
 [\Psi_j]_{Q_{\vartheta,a_j}}=0.
 \label{eq:cq-flow-profile-zero}
\end{equation}
Its temporal defect is nevertheless a specific full representative,
\[
 \mathcal D_j=\partial_t\Psi_j+
              \frac{\vartheta'}4D_\zeta^2\Psi_j
             =v_jV_{\vartheta,a_j}\Xi'.
\]
This is the complete chain rule with the original heat coefficient.
Write
\[
 \beta=[\Xi']_{\cqQ},\qquad
 \mathfrak a_\beta=\{c\in\mathbb C:c\Xi'\in\cqN\}.
\]
Then the exact quotient statement and its scalar kernel are
\begin{equation}
 [\mathcal D_j]_{Q_{\vartheta,a_j}}
    =\overline V_{\vartheta,a_j}(v_j\beta),\qquad
 \ker(c\mapsto c\beta)=\mathfrak a_\beta .
 \label{eq:cq-flow-source-beta}
\end{equation}
The set $\mathfrak a_\beta$ is a closed complex linear subspace:
it is the inverse image of the closed $\cqN$ under the continuous
scalar map $c\mapsto c\Xi'$ into $(\cqE,\tau)$.
Its equality with $\{0\}$ or $\mathbb C$ is precisely the unresolved
membership question $\Xi'\notin\cqN$ or $\Xi'\in\cqN$.
No evaluation functional on the source quotient is used to decide it.

At a fixed physical point $x=X_t(q)$, retain
$v=J(x)u(t,x)$ and the full inverse
$J^{-1}(x)=[2W_1(x)\ W_2(x)\ W_3(x)]$. Thus the exact kernel of
the complexified velocity-to-source-class map is
\[
 J(x)^{-1}(\mathfrak a_\beta^3)\subseteq\mathbb C^3;
\]
for real velocities take its intersection with $\mathbb R^3$.
This follows componentwise from \eqref{eq:cq-flow-source-beta} and
the invertibility of the unchanged $J$.

The derivative relation explains why a zero profile class does not
determine its representative derivative. At a fixed $(\theta,a)$ put
\[
 \mathfrak D_{\theta,a}
 =\{([h],[Dh]):h\in\cqE\}\subset Q_{\theta,a}^2,\qquad
 \mathcal W_{\theta,a}
 =(DN_{\theta,a}+N_{\theta,a})/N_{\theta,a}.
\]
It has full domain and output fibre $[Dh]+\mathcal W_{\theta,a}$
at $[h]$: replace $h$ by every $h+n$ with $n\in N_{\theta,a}$.
Its presenting map has kernel
$\{h\in N_{\theta,a}:Dh\in N_{\theta,a}\}$.
Since $D$ commutes with $V_{\theta,a}$, conjugating both components
by $\overline V_{\theta,a}$ gives exactly the previous source
derivative relation and its ambiguity space. In particular
$\overline V_{\theta,a}\beta$ is an actual element of its
multivalued part at zero. The displayed temporal connection is an
identity on the full smooth profiles; a derivative on all
untransported source quotient representatives is not presumed.

\subsection{An exact information-preserving refinement and its residue projection}

Retain the explicit analytic ideal $I=\Xi\cqE$, already proved closed
for $\gamma$ by the full division estimate, and put
\begin{equation}
 H_\cap=\cqN\cap I,\qquad
 C_\cap=(\cqE,\tau\vee\gamma)/H_\cap,\qquad
 Q_I=(\cqE,\gamma)/I .
 \label{eq:cq-flow-meet-quotient}
\end{equation}
The join topology is precisely the topology induced by
$h\mapsto(h,h)$ into $(\cqE,\tau)\times(\cqE,\gamma)$.
Both subspaces in the intersection are closed in this join, so
$C_\cap$ is Hausdorff. The identity on representatives induces
continuous surjections
\begin{equation}
 p_N:C_\cap\longrightarrow\cqQ,\qquad
 p_I:C_\cap\longrightarrow Q_I,\qquad
 \ker p_N=\cqN/H_\cap,\quad
 \ker p_I=I/H_\cap .
 \label{eq:cq-flow-meet-maps}
\end{equation}
The kernel assertions are direct membership tests; continuity follows
from the quotient universal property. These kernel identifications
are also homeomorphisms when $\cqN$ and $I$ have their subspace
topologies from $\tau\vee\gamma$. To prove this, let $q_\cap$ be
the open quotient map and let $A$ denote either subspace.
Because $H_\cap\subset A$,
$q_\cap(A\cap O)=q_\cap(A)\cap q_\cap(O)$ for every open ambient
$O$. Thus its restriction to $A$ is an open quotient map onto
the corresponding kernel with its subspace topology.
The joint map $(p_N,p_I)$ is continuous and injective since its
kernel is zero. No inference of a topological embedding from those
two properties is needed.

Let $\lambda$ be any actual zero of $\Xi$, with its actual
multiplicity $m_\lambda\geq1$. Choose a disk containing no other zero
and retain the full unit
\[
 \Xi(s)=(s-\lambda)^{m_\lambda}A_\lambda(s),
 \qquad A_\lambda(\lambda)\ne0.
\]
For $h\in\cqE$ define
\begin{equation}
 R_\lambda(h)
 =\operatorname{Res}_{s=\lambda}\frac{h(s)}{\Xi(s)}
 =\frac1{2\pi i}\int_{|s-\lambda|=r}
                    \frac{h(s)}{\Xi(s)}\,ds ,
 \label{eq:cq-flow-residue-functional}
\end{equation}
where the circle has positive orientation and lies in this disk.
The integral is independent of the permitted $r$ by Laurent
expansion. It is continuous for $\gamma$: its absolute value is
at most $r\sup_{|s-\lambda|=r}|h(s)|/
\min_{|s-\lambda|=r}|\Xi(s)|$, and a compact supremum is bounded
by a fixed growth seminorm times the maximum of its exponential
weight on the circle. It annihilates $I$, since $h/\Xi$ is entire
there, and therefore annihilates $H_\cap$.

Logarithmic differentiation of the full local factor gives
\begin{equation}
 \frac{c\Xi'(s)}{\Xi(s)}
 =\frac{m_\lambda c}{s-\lambda}
        +c\frac{A_\lambda'(s)}{A_\lambda(s)},\qquad
 R_\lambda(c\Xi')=m_\lambda c .
 \label{eq:cq-flow-logarithmic-unit}
\end{equation}
In particular $\Xi'\notin I$: the displayed pole is nonzero.
Thus $\beta_I=[\Xi']_{Q_I}\ne0$, and
$\widetilde\beta=[\Xi']_{C_\cap}\ne0$ as well, independent of the
status of $\beta$ in the actual source quotient.
The residue functional factors continuously through $C_\cap$; write
its induced functional as $\widetilde R_\lambda$.
The formula
\begin{equation}
 \Pi_\lambda(c)
  =\frac{\widetilde R_\lambda(c)}{m_\lambda}\,
                  \widetilde\beta,\qquad c\in C_\cap
 \label{eq:cq-flow-meet-projection}
\end{equation}
defines a continuous projection onto the line
$\mathbb C\widetilde\beta$. Indeed its square is itself by
$\widetilde R_\lambda(\widetilde\beta)=m_\lambda$, and scalar
multiplication on the quotient is continuous. Its inverse coordinate
on that line is $\widetilde R_\lambda/m_\lambda$, while
$c\mapsto c\widetilde\beta$ is its continuous inverse.
Consequently this line has exactly the ordinary complex topology,
and
$C_\cap=\mathbb C\widetilde\beta\oplus
\ker\widetilde R_\lambda$ is a continuous direct sum.
This proves an information-preserving refinement of the actual
source map, including a continuous recovery coordinate.

The full principal-part map
$h\mapsto(\operatorname{PP}_\lambda(h/\Xi))_{\lambda\in Z(\Xi)}$
has kernel exactly $I$. Vanishing of every negative Laurent
coefficient makes $h/\Xi$ entire at every zero; it is already
holomorphic elsewhere, and the proved entire division bound
puts it in $\cqE$. The reverse inclusion is immediate.
Each finite principal-part coordinate is continuous by the same
circle integral with the appropriate power of $s-\lambda$.
On the velocity line its entire image at every actual zero is
the single term $m_\lambda c/(s-\lambda)$ from
\eqref{eq:cq-flow-logarithmic-unit}, with the unchanged analytic
unit in the regular part.

For general $(\theta,a)$ let
$I_{\theta,a}=V_{\theta,a}I$ and transport
\eqref{eq:cq-flow-meet-quotient} by $V_{\theta,a}$.
Its ambient topology is
$\tau_{\theta,a}\vee\gamma$, since $V_{\theta,a}$ is a
$\gamma$-homeomorphism, and its relation is
$N_{\theta,a}\cap I_{\theta,a}=V_{\theta,a}H_\cap$.
Every map, kernel and projection above is transported exactly.
In particular the transported residue coordinate is
\[
 R_{\lambda;\theta,a}(h)
 =\operatorname{Res}_{s=\lambda}
           \frac{(U_{-\theta}T_{-a}h)(s)}{\Xi(s)}.
\]
For the actual temporal defect it gives
$R_{\lambda;\vartheta,a_j}(\mathcal D_j)=m_\lambda v_j$.
The pair of source and analytic classes therefore has the
explicit injective common lift
$v_j[V_{\vartheta,a_j}\Xi']$ in this transported refinement.
For nonzero heat time this formula retains the inverse heat
operator; it does not replace $V_{\theta,a}I$ by an unproved
ordinary ideal generated by $T_aZ_\theta$.

\subsection{The logarithmic connection and the public blowup residue field}

Take the explicit Newman-time choice $\vartheta(t)=0$.
For the original real flow set
\[
 \Psi_j(t,q,\zeta)=\Xi(\zeta+a_j(t,q)),\qquad
 \mathcal D_j=\partial_t\Psi_j
             =v_j(t,q)\Xi'(\zeta+a_j(t,q)).
\]
At every actual zero $\lambda$ the translated divisor is
$\zeta_{\lambda,j}(t,q)=\lambda-a_j(t,q)$, and the full factorization
is
\[
 \Psi_j=
 (\zeta-\zeta_{\lambda,j})^{m_\lambda}
 A_\lambda(\zeta+a_j).
\]
Thus all multiplicities and the actual analytic unit are retained.
Since $a_j$ is real, its imaginary coordinate is exactly
$\operatorname{Im}\lambda$. Differentiating its location gives
$\partial_t\zeta_{\lambda,j}=-v_j$.

On the complement of the divisor in
$[0,T)\times\mathbb R^3_q\times\mathbb C_\zeta$, define the
complex-valued differential form, smooth in the real variables and
holomorphic in $\zeta$,
\begin{equation}
 \omega_j=\frac{d\Psi_j}{\Psi_j}
 =\frac{\Xi'(\zeta+a_j)}{\Xi(\zeta+a_j)}
                         (d\zeta+da_j),\qquad
 da_j=v_j\,dt+\sum_b\partial_{q_b}a_j\,dq_b .
 \label{eq:cq-flow-logarithmic-connection}
\end{equation}
All label and time coefficients are displayed. Put $s=\zeta+a_j$.
Direct exterior differentiation gives
$d\omega_j=(\Xi'/\Xi)'(s)\,ds\wedge ds+
(\Xi'/\Xi)(s)d^2s=0$.
Equivalently the connection $d-\omega_j$ on the trivial complex
line has zero curvature and annihilates the section $\Psi_j$.
This calculation is on the complement of its zeros and retains
the singular divisor explicitly.

At fixed $(t,q)$ a sufficiently small positively oriented circle
around $\zeta_{\lambda,j}$ has
\[
 \frac1{2\pi i}\int\omega_j=m_\lambda,
\]
where only its $d\zeta$ restriction is integrated. The pole term in
\eqref{eq:cq-flow-logarithmic-unit} gives this integral; the
holomorphic unit term has zero integral by Cauchy's theorem.
The time component and its residue are exactly
\begin{align}
 \omega_j(\partial_t)
 &=\frac{\mathcal D_j}{\Psi_j}
 =\frac{m_\lambda v_j}{\zeta-\zeta_{\lambda,j}}
       +v_j\frac{A_\lambda'(\zeta+a_j)}
                       {A_\lambda(\zeta+a_j)},\notag\\
 r_j(t,q)
 &:=\operatorname{Res}_{\zeta=\zeta_{\lambda,j}}
                    \frac{\mathcal D_j}{\Psi_j}
 =m_\lambda v_j
 =-m_\lambda\partial_t\zeta_{\lambda,j}.
 \label{eq:cq-flow-temporal-residue}
\end{align}
The coordinate change $s=\zeta+a_j$ has differential $ds=d\zeta$
on each spectral fibre, so this is precisely
$R_{\lambda;0,a_j}(\mathcal D_j)$, including its orientation
and multiplicity factor.

Write $r=(r_1,r_2,r_3)^{\mathsf T}$, and retain the full original
polynomial target and its matrix inverse:
\[
 S=(F_1/2,F_2,F_3)^{\mathsf T},\quad
 J=DS=\operatorname{diag}(1/2,1,1)DF,\quad
 J^{-1}=[2W_1\ W_2\ W_3].
\]
The literal recovery and pointwise norm identities are
\begin{equation}
 \begin{aligned}
 u(t,X_t(q))
 &=\frac{2W_1(X_t(q))r_1+
              W_2(X_t(q))r_2+W_3(X_t(q))r_3}{m_\lambda},\\
 |r(t,q)|^2&=m_\lambda^2|J(X_t(q))u(t,X_t(q))|^2 .
 \end{aligned}
 \label{eq:cq-flow-residue-velocity-recovery}
\end{equation}
They follow by $r=m_\lambda v$, $v=Ju$ and the stated matrix
inverse, with no selection of a global inverse sheet of $S$.
The flow $X_t$ is retained as part of the data.

For the nonempty compact velocity support $K$ put
\[
 M_K=\max_K\|J\|_{\mathrm{op}},\qquad
 N_K=\max_K\|J^{-1}\|_{\mathrm{op}}.
\]
They are finite and positive by compactness and $\det J=-1$.
Since $X_t(K)=K$, $X_t$ is onto, and $u,v,r$ vanish outside
the corresponding compact labels, \eqref{eq:cq-flow-residue-velocity-recovery}
gives
\begin{equation}
 \frac{m_\lambda}{N_K}\|u(t)\|_\infty
 \leq\sup_q|r(t,q)|
 \leq m_\lambda M_K\|u(t)\|_\infty .
 \label{eq:cq-flow-residue-supremum}
\end{equation}
The volume-preserving change of variables $x=X_t(q)$ gives
\begin{equation}
 \int_{\mathbb R^3}|r(t,q)|^2\,dq
 =m_\lambda^2\int_{\mathbb R^3}|J(x)u(t,x)|^2\,dx,
 \label{eq:cq-flow-residue-energy}
\end{equation}
whose square root is between
$(m_\lambda/N_K)\|u(t)\|_2$ and
$m_\lambda M_K\|u(t)\|_2$.
These are exact spatial estimates for the arithmetic residues.
With the previously proved spectral normalization
$d_Z(0)^2=\int_{\mathbb R}|\Xi'(s)|^2ds>0$, the unchanged
real translations also give
\[
 |r(t,q)|
 =\frac{m_\lambda}{d_Z(0)}
       \left(\sum_j\|\mathcal D_j(t,q,\cdot)\|_{L^2(\mathbb R)}^2
       \right)^{1/2}.
\]

\begin{cqproposition}[Application of the cited public Navier--Stokes theorem]
For every $\nu>0$, take the velocity, pressure and force supplied by
the public Theorem~1.1 of \cite{cqNS2026}. The construction above,
with $\vartheta=0$, gives arithmetic logarithmic residue fields
satisfying
\[
 \sup_{0\leq t<1}\int_{\mathbb R^3}|r(t,q)|^2\,dq<\infty,
 \qquad
 \limsup_{t\uparrow1}\sup_q|r(t,q)|=\infty .
\]
Every residue recovers the full physical velocity by
\eqref{eq:cq-flow-residue-velocity-recovery}, and each translated
arithmetic divisor has the same imaginary coordinates as the
original zero divisor of $\Xi$.
\end{cqproposition}
\begin{proof}
The cited theorem supplies a smooth incompressible velocity on
$[0,1)\times\mathbb R^3$, a fixed compact support, zero initial
velocity, a uniform kinetic-energy bound, and unbounded terminal
velocity supremum, with its stated smooth compactly supported force
and its pressure. These are the inputs of the proved preterminal
flow construction, so its volume and inverse statements apply
on every $[0,T_0]$ with $T_0<1$. Formula
\eqref{eq:cq-flow-residue-energy} and the upper bound using $M_K$
prove the uniform residue energy bound. The lower bound in
\eqref{eq:cq-flow-residue-supremum} proves the unbounded terminal
supremum. Recovery and imaginary-coordinate preservation were
proved above for the actual same fields. The existence step in
this corollary is exactly the cited Theorem~1.1; the correspondence
and every displayed estimate are proved here.
\end{proof}

\subsection{The published witness's exact growing residue}

Take the particular corrected field of the public construction, with
the source's fixed $h\in(0,1/100)$, $X_{\rm in}>0$, $e_0>0$,
and remainder $R_*(\tau)$. Its closing formulas
\cite[equations (10.20)--(10.22), p.~124]{cqNS2026}, as imported
in the preceding witness construction, are
\[
 |R_*(\tau)|\leq C_*\tau^{2h}\quad(0<\tau<\tau_*),\qquad
 (u_\nu)_2(x_{\nu,\tau},1-\tau)
   =\sqrt\nu\,\tau^{-1/2-h}(e_0+R_*(\tau)),
\]
\[
 r_{\nu,\tau}=\sqrt{2\nu X_{\rm in}\tau},\quad
 x_{\nu,\tau}=(r_{\nu,\tau},0,0),\quad
 q_{\nu,\tau}=X_{\nu,1-\tau}^{-1}(x_{\nu,\tau}).
\]
The remainder is the actual source remainder, including its corrected
fields. The parameter $h$ is the one chosen in that construction.
The quantity $r_{\nu,\tau}$ here is the displayed physical radius;
the map to the original polynomial target is retained below.
The labels $q_{\nu,\tau}$ vary with $\tau$ according to the exact
material inverse just displayed.

Substitution into the full original polynomials gives
\[
 S(r,0,0)=(0,0,2r),\qquad
 J(r,0,0)=
 \begin{pmatrix}
 0&0&1/2\\
 0&1&3r\\
 2&-3r^2&-r^3
 \end{pmatrix}.
\]
Here the same original formulas are
\[
 \begin{gathered}
 F_1=Q^3w+y^2Q(4+3xy),\quad
 F_2=y+3xQ^2w+3xy^2(4+3xy),\\
 F_3=2x-3x^2y-x^3w,\quad Q=1+xy.
 \end{gathered}
\]
Differentiating $S=(F_1/2,F_2,F_3)$ and setting $(x,y,w)=(r,0,0)$
gives each matrix entry. Its componentwise inverse gives
$u_3=2v_1$ and $u_2=v_2-6rv_1$, retaining the off-diagonal $3r$.
The first two arithmetic shifts at $(1-\tau,q_{\nu,\tau})$ are zero.
Consequently, as an identity of full entire functions of $\zeta$,
\begin{equation}
\begin{split}
 \mathcal E_{\nu,\tau}(\zeta)
 &:=\mathcal D_{\nu,2}(1-\tau,q_{\nu,\tau},\zeta)
       -6r_{\nu,\tau}
          \mathcal D_{\nu,1}(1-\tau,q_{\nu,\tau},\zeta)\\
 &=\sqrt\nu\,\tau^{-1/2-h}
                    (e_0+R_*(\tau))\Xi'(\zeta).
\end{split}
 \label{eq:cq-actual-witness-entire-defect}
\end{equation}
Indeed each defect is $v_j\Xi'$ because its shift is zero, so their
combination is $(u_\nu)_2\Xi'$, and the stated source identity applies.

For every actual zero $\lambda$ of multiplicity $m_\lambda$, its
complete local logarithmic image is
\begin{equation}
 \frac{\mathcal E_{\nu,\tau}(\zeta)}{\Xi(\zeta)}
 =\sqrt\nu\,\tau^{-1/2-h}(e_0+R_*(\tau))
     \left(\frac{m_\lambda}{\zeta-\lambda}
                    +\frac{A_\lambda'(\zeta)}{A_\lambda(\zeta)}\right).
 \label{eq:cq-actual-witness-meromorphic-defect}
\end{equation}
Therefore the actual unshifted arithmetic residue at the source
sample is exactly
\begin{equation}
 \operatorname{Res}_{\zeta=\lambda}
       \frac{\mathcal E_{\nu,\tau}(\zeta)}{\Xi(\zeta)}
 =m_\lambda\sqrt\nu\,\tau^{-1/2-h}
                           (e_0+R_*(\tau)).
 \label{eq:cq-actual-witness-residue-growth}
\end{equation}
All multiplicities, the whole analytic unit, the positive viscosity,
the original remainder and the exact material labels remain present.
Choose $\tau_{**}>0$ within the source interval and its cutoff region
with $C_*\tau_{**}^{2h}\leq e_0/2$. For $0<\tau<\tau_{**}$ the
absolute value of \eqref{eq:cq-actual-witness-residue-growth} is at
least $m_\lambda\sqrt\nu\,e_0\tau^{-1/2-h}/2$, which diverges.
This is an explicit pole-coefficient singularity of the published
constructed field's arithmetic image.

For the three residue channels $r_{\nu,j}$ of
\eqref{eq:cq-flow-temporal-residue}, the same evaluation and the
two-component Cauchy--Schwarz inequality give the quantitative bound
\begin{equation}
 |r_\nu(1-\tau,q_{\nu,\tau})|
 \geq
 \frac{m_\lambda\sqrt\nu\,e_0}
             {2\sqrt{1+72\nu X_{\rm in}\tau}}\,
                   \tau^{-1/2-h}.
 \label{eq:cq-actual-witness-residue-vector-growth}
\end{equation}
To verify the denominator, the coefficient vector
$(-6r_{\nu,\tau},1)$ has squared norm
$1+36r_{\nu,\tau}^2=1+72\nu X_{\rm in}\tau$.
Its product with the first two residue coordinates is exactly
the residue in \eqref{eq:cq-actual-witness-residue-growth}; the
third coordinate can only increase the Euclidean norm on the left.
If $E_\nu=\nu^{5/4}F_*(1)$ is the full source force-integral energy
bound and $M_\nu=\max_{K_\nu}\|J\|_{\mathrm{op}}$, the earlier
volume identity simultaneously gives
\[
 \int_{\mathbb R^3}|r_\nu(t,q)|^2dq
 \leq m_\lambda^2M_\nu^2E_\nu^2
 =m_\lambda^2M_\nu^2\nu^{5/2}F_*(1)^2\quad(t<1).
\]

At these particular first two channels, their zero shifts mean the
source quotient is the original $\cqQ$ itself. The exact classes are
\[
 [\mathcal E_{\nu,\tau}]_{\cqQ}
 =\sqrt\nu\,\tau^{-1/2-h}(e_0+R_*(\tau))\beta,\qquad
 [\mathcal E_{\nu,\tau}]_{C_\cap}
 =\sqrt\nu\,\tau^{-1/2-h}(e_0+R_*(\tau))\widetilde\beta.
\]
The latter line has the explicitly proved continuous inverse
$\widetilde R_\lambda/m_\lambda$, so its scalar coordinate has
exactly the displayed singularity. In particular these common
classes leave every bounded subset of $C_\cap$: a continuous linear
functional sends bounded sets to bounded scalar sets, whereas their
$\widetilde R_\lambda$ values diverge. Their projection to the
actual source quotient and its exact kernel remain
\eqref{eq:cq-flow-meet-maps}; the undecided source membership is
not replaced by an identification with the analytic ideal.

There is also an exact relation between the changing sample labels
and the time component of the flat logarithmic connection.
Let $A_\tau=D_qX_{\nu,1-\tau}(q_{\nu,\tau})$. Differentiating
$X_{\nu,1-\tau}(q_{\nu,\tau})=x_{\nu,\tau}$ gives
\[
 A_\tau\,\frac{dq_{\nu,\tau}}{d\tau}
 =\frac{dx_{\nu,\tau}}{d\tau}
       +u_\nu(1-\tau,x_{\nu,\tau}),\qquad
 \frac{dx_{\nu,\tau}}{d\tau}
       =\frac{r_{\nu,\tau}}{2\tau}(1,0,0).
\]
The sign follows from $d(1-\tau)/d\tau=-1$.
Since $\nabla_q a_j$ is the $j$th row of $J(x_{\nu,\tau})A_\tau$,
and the first two rows of $J(r,0,0)$ have zero first component,
\[
 \nabla_q a_j\cdot\frac{dq_{\nu,\tau}}{d\tau}=v_j,\qquad j=1,2.
\]
Thus $a_j(1-\tau,q_{\nu,\tau})=0$ has total derivative
$-v_j+\nabla_q a_j\cdot dq_{\nu,\tau}/d\tau=0$, exactly as required.
On the spectral complement of the zeros, the pullback of
\eqref{eq:cq-flow-logarithmic-connection} to this curve and fixed
$\zeta$ is zero in the $\tau$ direction, while its original
$\partial_t$ component has the residue
\eqref{eq:cq-flow-temporal-residue}. These statements are related
by the displayed full tangent-vector identity; the sample is not
silently treated as one fixed material label.

For the full four-term profile $K_\theta=U_\theta K_0$, the corresponding
actual source classes are instead
\[
 \kappa=[K_0]_{\cqQ},\qquad \beta_K=[K_0']_{\cqQ},\qquad
 [T_aK_\theta]=\overline V_{\theta,a}\kappa,\qquad
 [vT_aK_\theta']=\overline V_{\theta,a}(v\beta_K).
\]
These follow by the same commuting translation, heat and derivative
maps and retain $\kappa$ without assigning it the value zero.
The residue corollary above uses the actual $\Xi$ profile.
The displayed four-term quotient formulas concern the whole $K$
profile and its derivative; they do not identify individual zeros
of $K$ with zeros of $\Xi$.

\subsection{The translated modular double}

Keep $C_Eh(s)=\overline{h(\overline s)}$,
$\tau^c=(C_E)_*\tau$ and $N^c=C_E\cqN$ from the explicit modular
construction. Direct conjugation gives
$C_ET_a=T_{\overline a}C_E$, and the established heat identity gives
$C_EU_\theta=U_{\overline\theta}C_E$.
On the two entire sectors define
\[
 \widehat V_{\theta,a}
   =\operatorname{diag}(T_aU_\theta,T_aU_{-\theta}),\qquad
 \mathcal J(h,k)=(C_Ek,C_Eh).
\]
Then the exact all-parameter identity is
\begin{equation}
 \mathcal J\widehat V_{\theta,a}
 =\widehat V_{-\overline\theta,\overline a}\mathcal J .
 \label{eq:cq-flow-translated-cartan}
\end{equation}
Indeed its first component is
$C_ET_aU_{-\theta}k=T_{\overline a}U_{-\overline\theta}C_Ek$,
and its second component is the same calculation with $\theta$.
For the actual real flow, $a\in\mathbb R$; for real Newman time it
therefore sends $(\theta,a)$ to $(-\theta,a)$.

Give the two ambient sectors the topologies
$(T_aU_\theta)_*\tau$ and $(T_aU_{-\theta})_*\tau^c$, and quotient
by the closed relation
\[
 \widehat N_{\theta,a}
   =T_aU_\theta\cqN\oplus T_aU_{-\theta}N^c .
\]
Applying \eqref{eq:cq-flow-translated-cartan} both to these relation
spaces and to the definitions of their transported topologies proves
an antilinear homeomorphism from this actual doubled quotient to
the one indexed by $(-\overline\theta,\overline a)$.
Its square is the identity, and its conjugacy with the quotient heat
and translation maps follows on representatives from the same
displayed equality. Thus the Cartan reversal retains the full real
flow translation while reversing the radial heat parameter, with
both original source and conjugate-source sectors specified.

\section{The exact strength of the frozen source spectral information}
\label{sec:cq-source-spectral-strength}

The primary spectrum statement in \cite[Proposition~3.6(ii)]{cqCCM}
specifies the actual zero set as the spectrum of the Hadamard
quotient's multiplication operator. It gives no multiplicity
assertion. The cited \cite[Proposition~2.24]{cqCM} computes the
finite Hermite polynomial image; the discussion after its proof
explicitly treats function-space choices. The finite multiplicity
calculation in \cite[Theorem~2.47 and equations (2.447)--(2.449)]{cqCM}
belongs to the separately specified weighted Hilbert quotient.
The following exact algebra records both the obstruction available
from such data and the limits of transferring a spectrum set alone.

\subsection{One-sided heat inclusion and the full generalized chain}

The already proved commutator identity
\eqref{eq:cq-source-D-equivalence} and its following heat calculation
show that $D\cqN\subseteq\cqN$ makes every $S-\lambda$ algebraically
onto. The single inclusion $U_t\cqN\subseteq\cqN$ at any $t\ne0$
also makes every $S-\lambda$ onto, using every derivative
of $U_t\Xi$ and the exact source division map.
Here is a stronger consequence that needs no descended derivative.

\begin{cqlemma}
Let $T:V\to V$ be a surjective endomorphism of a complex vector
space and let $0\ne v_0\in\ker T$. Then for every integer $k\geq1$,
$\dim\ker T^k\geq k$. More precisely one can choose
$Tv_j=v_{j-1}$ for every $j\geq1$, and the vectors
$v_0,\ldots,v_n$ are linearly independent for every $n$.
\end{cqlemma}
\begin{proof}
Surjectivity gives the successive preimages. It follows inductively
that $T^jv_j=v_0$ and $T^{j+1}v_j=0$. Applying $T^n$ to a relation
$\sum_{j=0}^n c_jv_j=0$ gives $c_nv_0=0$, hence $c_n=0$.
Repeating with the remaining highest index gives every coefficient
zero. The first $k$ vectors lie in $\ker T^k$, proving the bound.
\end{proof}

Applied to $T=S-\lambda$, this proves that either frozen
inclusion just stated is incompatible with any nonzero
finite-dimensional generalized eigenspace of the actual source
operator. The source-independent implication is
\begin{equation}
\begin{gathered}
 U_t\cqN\subseteq\cqN,\quad t\ne0,\quad
                  0\ne v_0\in\ker(S-\lambda)
 \quad\Longrightarrow\quad
 \dim\ker(S-\lambda)^k\geq k\quad(k\geq1).
\end{gathered}
\label{eq:cq-source-surjective-chain}
\end{equation}
With a descended $d[f]=[Df]$, its more specific ladder is
$v_j=(-1)^jd^jv_0/j!$; indeed
$[d,S]=I$ gives $(S-\lambda)d^jv_0=-j d^{j-1}v_0$.
Thus the complete factorial and sign agree with the general
preimage construction. This is a proved algebraic incompatibility;
it does not assert that the unnamed source quotient has the
finite generalized block needed to apply the incompatibility.

\subsection{The actual finite block in the specified Sobolev source}

Retain the source's real frequency $t$, compact character $\chi_0$,
weight $\delta>1$, and actual zero multiplicity $m$.
The vectors in the dual block of the specified source are
\[
 \eta_{t,a}(x)=\chi_0(x)|x|^{it}(\log|x|)^a,\qquad
 a\in\mathbb Z_{\geq0},\quad a<m,\quad a<(\delta-1)/2 .
\]
The last strict bound is exactly integrability of
$u^{2a}(1+u^2)^{-\delta/2}$ on $\mathbb R$: at infinity
its power is $2a-\delta$, integrable precisely when it is
less than $-1$; near zero it is locally integrable for all
these $a$. Let $r$ be the number of permitted integers.
The source representation at $\lambda>0$ acts by
\[
 \vartheta_m^{\rm tr}(\lambda)\eta_{t,a}
 =\lambda^{it}\sum_{b=0}^a\binom ab(\log\lambda)^b\eta_{t,a-b}.
\]
Differentiating at $\lambda=e^\epsilon$, $\epsilon=0$, proves
\begin{equation}
 (D_{\chi_0}^{\rm tr}-it)\eta_{t,a}=a\eta_{t,a-1},\qquad
 (D_{\chi_0}^{\rm tr}-it)^k\eta_{t,a}
 =\frac{a!}{(a-k)!}\eta_{t,a-k}\quad(0\leq k\leq a).
 \label{eq:cq-source-Sobolev-Jordan-block}
\end{equation}
For $k>a$ the expression is zero. These vectors are independent:
divide a proposed relation by the nonzero character and use
that $\log|x|$ ranges over $\mathbb R$, making the remaining
polynomial identically zero. Thus the block has dimension
$r$, one-dimensional ordinary eigenspace, and nilpotence
index $r$. Its multiplicity is generalized multiplicity.
The common-test morphism to the actual Hadamard source was
proved earlier with its exact two kernels; this finite
block is not silently identified with a block in that
different quotient.

\subsection{An explicit spectral model and the inverse at zero}

Let $\Lambda=Z(\Xi)$ be the actual zero set. On the complex
vector space
\[
 V_\Lambda=\bigoplus_{\lambda\in\Lambda}
                   \bigoplus_{n\geq0}\mathbb C e_{\lambda,n}
\]
every vector has finite support. Give it the finest locally
convex topology, defined by all seminorms on this vector space.
Coordinate functionals separate points, so it is Hausdorff.
Every linear map from it to any locally convex space is
continuous: each pullback of a continuous seminorm is among
these defining seminorms.
With $e_{\lambda,-1}=0$ set
\[
 S_\Lambda e_{\lambda,n}=\lambda e_{\lambda,n}+e_{\lambda,n-1},
 \qquad d_\Lambda e_{\lambda,n}=-(n+1)e_{\lambda,n+1}.
\]
Both are continuous, everywhere-defined, and direct evaluation
on the basis proves $[d_\Lambda,S_\Lambda]=I$.
For $\mu\notin\Lambda$ the full continuous inverse is
\begin{equation}
 (S_\Lambda-\mu)^{-1}e_{\lambda,n}
 =\sum_{j=0}^n
       \frac{(-1)^j}{(\lambda-\mu)^{j+1}}e_{\lambda,n-j}.
 \label{eq:cq-source-model-resolvent}
\end{equation}
Multiplication in either order cancels adjacent terms of
this finite sum and leaves the basis vector. The formula
preserves finite support, so it defines an inverse on
the whole space, continuous by its topology.
For $\mu\in\Lambda$, the $\mu$-summand is a surjective
backward shift and every other summand uses the same finite
inverse formula. Thus $S_\Lambda-\mu$ is onto, its kernel
is exactly $\mathbb C e_{\mu,0}$, and
\begin{equation}
 \ker(S_\Lambda-\lambda)^k
   =\operatorname{span}\{e_{\lambda,0},\ldots,e_{\lambda,k-1}\},
 \quad
 \operatorname{Spec}_{\rm alg}S_\Lambda
  =\operatorname{Spec}_{\rm end}S_\Lambda
  =\operatorname{Spec}_{\rm point}S_\Lambda=\Lambda .
 \label{eq:cq-source-model-spectrum}
\end{equation}
Here $\operatorname{Spec}_{\rm end}$ means failure of a
continuous two-sided inverse in the algebra of continuous
endomorphisms, not the Hilbert spectral subdivision called
continuous spectrum. Since $\Xi(0)\ne0$, this $S_\Lambda$
has a continuous inverse at zero.

There is an exact principal-part realization of the model:
send $e_{\lambda,n}$ to the class of $w^{-n-1}$ in
$\mathbb C[w,w^{-1}]/\mathbb C[w]$ in its $\lambda$-summand.
The finite negative Laurent coefficients give its inverse.
Multiplication by $\lambda+w$ is $S_\Lambda$ and
differentiation in $w$ is $d_\Lambda$, as evaluation on
each monomial proves.
This complete model shows that the correct zero set as
spectrum, geometric multiplicity one, and a continuous
inverse at zero can coexist with a continuous Weyl pair.
It does not identify this model with $\cqQ$.

For the actual quotient, injectivity of $S$ gives exactly
$\cqN\cap s\cqE=s\cqN$: if $sf\in\cqN$, then
$S[f]=0$, so $f\in\cqN$; the converse is polynomial invariance.
The original source inverse at zero is explicitly
\begin{equation}
 S^{-1}[f]=\left[
 \frac{f(s)-f(0)\Xi(s)/\Xi(0)}{s}\right].
 \label{eq:cq-source-explicit-zero-resolvent}
\end{equation}
The singularity is removable, and division by a degree-one
polynomial preserves the original order class. Multiplication
by $s$ gives $[f]$, and injectivity gives its unique inverse.
If $f$ is changed by $n\in\cqN$, the numerator difference
belongs to $\cqN\cap s\cqE=s\cqN$, proving representative
independence directly. The usual topological reading of
the source resolvent also supplies continuity of this
quotient inverse, not continuity of its chosen lift.
More generally,
\[
 S-\lambda\text{ injective}\quad\Longleftrightarrow\quad
 \cqN\cap(s-\lambda)\cqE=(s-\lambda)\cqN
\]
by the same calculation. Saturation at zero does not prove
saturation at other points.
The missing frozen-descent decision can therefore be
approached through an actual finite generalized source block
or a range obstruction, as well as through identification
of the topology. Neither has been supplied by the primary
Hadamard spectrum clauses inspected here.

\section{Actual angular momentum, nonlinear fluxes and arithmetic dilation}
\label{sec:cq-angular-physical-map}

The complete companion proof \cite{cqNSAngular} is retained here with its original coordinates and full corrected source fields. The source existence and correction estimates remain the explicit external input \cite{cqNS2026}.

\subsection{Cartesian input, angular projection and its exact section}

Use the complete source witness $(u_\nu,p_\nu,f_\nu)$ constructed in
Section~\ref{sec:cq-NS-actual-witness}, for every physical time $0\leq t<1$. Retain
the viscosity $\nu>0$ and original Cartesian coordinates $(x_1,x_2,z)$.
For $r>0$ put $x_1=r\cos\theta,x_2=r\sin\theta$,
$e_r=(\cos\theta,\sin\theta,0)$ and
$e_\theta=(-\sin\theta,\cos\theta,0)$. Write
$u_r=e_r\cdot u_\nu$, $u_\theta=e_\theta\cdot u_\nu$ and
$u_z=(u_\nu)_3$. These components may depend on $\theta$; in particular
the complete corrected velocity is not assumed axisymmetric.
Define $\langle a\rangle=(2\pi)^{-1}\int_0^{2\pi}a\,d\theta$ and
retain the radial coordinate and inverse
\[
 \sigma=\frac{r^2}{2},\qquad r=\sqrt{2\sigma},\qquad
 r\,dr=d\sigma.
\]
This radial $\sigma$ is not the earlier polynomial coordinate $F_1/2$.
They coexist as explicit functions on the original Cartesian space;
no equality of their Laplacians is required below.
Their exact common coordinate map, on
$(0,\infty)\times(\mathbb R/2\pi\mathbb Z)\times\mathbb R$, is
$(\sigma,\theta,z)\mapsto
(\sqrt{2\sigma}\cos\theta,\sqrt{2\sigma}\sin\theta,z)$.
Its inverse is $\sigma=(x_1^2+x_2^2)/2$, the direction
$(\cos\theta,\sin\theta)=(x_1,x_2)/\sqrt{x_1^2+x_2^2}$,
and the unchanged $z$. In these coordinates the full polynomial is
\[
 \frac{F_1}{2}=\frac12\bigl[
 (1+\sigma\sin(2\theta))^3z
 +2\sigma\sin^2\theta(1+\sigma\sin(2\theta))
                         (4+3\sigma\sin(2\theta))\bigr].
\]
Substitution of $x_1x_2=\sigma\sin(2\theta)$ and
$x_2^2=2\sigma\sin^2\theta$ proves this morphism; adjoining the
displayed value as an extra coordinate gives its graph map with
inverse projection onto $(\sigma,\theta,z)$. The Cartesian definition
continues over the axis where the angular chart is not defined.

The actual angular-momentum density, fluxes and force torque are
\begin{align*}
 B(\sigma,z,t)&=\langle r u_\theta\rangle,
 &P(\sigma,z,t)&=\langle r^2u_ru_\theta\rangle,\\
 Q(\sigma,z,t)&=\langle r u_zu_\theta\rangle,
 &F(\sigma,z,t)&=\langle r f_\theta\rangle.
\end{align*}
These definitions include every corrected, oscillatory and cutoff term
of the actual source fields before averaging. All four functions have
fixed compact support in $(\sigma,z)$ on the closed half-plane
$\sigma\geq0$. They are smooth there on every closed preterminal
time interval, with
\begin{equation}
 B=O(\sigma),\quad Q=O(\sigma),\quad F=O(\sigma),
 \qquad P=O(\sigma^2)\quad(\sigma\downarrow0).
 \label{eq:cq-am:axis}
\end{equation}
Here and below these estimates hold with all $z,t$ derivatives on
compact preterminal sets. To prove them without polar smoothness
assumptions, observe that $r u_\theta=x_1u_2-x_2u_1$ and
$r u_r=x_1u_1+x_2u_2$. Angular averaging their Cartesian expressions
gives smooth rotation-invariant functions of $(x_1,x_2)$.
A smooth rotation-invariant function is smooth in $\sigma$ on the
closed half-line: its restriction to $(r,0)$ is even and smooth;
Taylor's formula to any even order, with its differentiated remainder,
expresses it as a polynomial in $r^2$ plus a remainder of arbitrarily
high order. The chain rule on $r>0$ then supplies every continuous
right derivative in $\sigma$ at zero and the fundamental theorem of
calculus identifies those derivatives successively. The same Taylor
proof works with the retained $z,t$ parameters.
The constants in the averages defining $B,Q,F$ vanish at $r=0$;
odd radial powers vanish under rotation. They therefore start at
$r^2$. For $P$, the quadratic Cartesian term at zero is
\[
 \big\langle(x_1a+x_2b)(x_1b-x_2a)\big\rangle=0,
 \qquad (a,b)=(u_1,u_2)(0,0,z,t),
\]
because $\langle x_1^2-x_2^2\rangle=
\langle x_1x_2\rangle=0$. Its constant term also vanishes, and
rotation removes the odd powers, so $P$ starts at $r^4$.
This proves \eqref{eq:cq-am:axis}.

The projection $\Pi u=B$ has the explicit linear section
\begin{equation}
 (\mathcal RB)(x_1,x_2,z)
   =\frac{B(\sigma,z)}{2\sigma}(-x_2,x_1,0),
 \qquad \Pi\mathcal R=I.
 \label{eq:cq-am:section}
\end{equation}
The quotient at zero is the actual smooth extension:
$B(\sigma,z)/\sigma=\int_0^1\partial_\sigma B(a\sigma,z)da$.
Thus the section is Cartesian smooth, compactly supported, and
divergence-free. For $r>0$ it has sole component $B/r$ in the
$e_\theta$ direction, proving the section identity. For the exact
remainder $w=u_\nu-\mathcal RB$, one has $\Pi w=0$ and
\begin{align}
 P&=rB\langle w_r\rangle+r^2\langle w_rw_\theta\rangle,
 &Q&=B\langle w_z\rangle+r\langle w_zw_\theta\rangle.
 \label{eq:cq-am:covariance}
\end{align}
These follow by substituting $u_r=w_r$, $u_z=w_z$ and
$u_\theta=B/r+w_\theta$ into the original definitions. They exhibit
both mean transport and the complete two quadratic correlations.
The full inverse of $u\mapsto(B,w)$ is
$(B,w)\mapsto\mathcal RB+w$ on the stated subspace $\Pi w=0$.
In particular the projection has the precise kernel
$\{w:\langle r w_\theta\rangle=0\}$; none of that kernel's
contribution to the fluxes has been discarded.

\subsection{The exact equation before and after angular averaging}

For the unaveraged Cartesian angular momentum
$L=x_1u_2-x_2u_1=r u_\theta$, the product rule in the full
Navier--Stokes equation gives
\begin{equation}
 \partial_tL+\operatorname{div}(u_\nu L)
 =\nu\Delta L-2\nu\omega_z-\partial_\theta p_\nu
                  +(x_1f_2-x_2f_1),
 \qquad \omega_z=\partial_1u_2-\partial_2u_1.
 \label{eq:cq-am:cartesian}
\end{equation}
Indeed the material derivatives of $x_1,x_2$ contribute
$u_1u_2-u_2u_1=0$; incompressibility converts material transport
to the displayed divergence. The Laplacian product rule gives
$\Delta L=x_1\Delta u_2-x_2\Delta u_1+2\omega_z$.
Finally $x_1\partial_2-x_2\partial_1=\partial_\theta$ gives the
negative pressure-torque term with its sign.

In polar coordinates
$\omega_z=r^{-1}\partial_r(ru_\theta)
                  -r^{-1}\partial_\theta u_r$.
Average \eqref{eq:cq-am:cartesian} over a full circle. Every angular
derivative has zero integral, by the fundamental theorem of calculus
and periodicity, including the actual non-axisymmetric pressure.
The convection term is
$r^{-1}\partial_r(r\langle u_rL\rangle)
 +\partial_z\langle u_zL\rangle$; the averaged viscous term is
$\nu(\partial_{rr}-r^{-1}\partial_r+\partial_{zz})B$.
Using $\partial_r=r\partial_\sigma$ now proves
\begin{equation}
 \partial_tB+\partial_\sigma P+\partial_zQ
   =\nu(2\sigma\partial_{\sigma\sigma}B+\partial_{zz}B)+F.
 \label{eq:cq-am:pde}
\end{equation}
The formula extends to $\sigma=0$ by \eqref{eq:cq-am:axis} and smoothness.
The pressure has disappeared here by a proved projection, with its
entire original torque term exhibited in \eqref{eq:cq-am:cartesian}.
No scalar heat equation has replaced the actual nonlinear fluxes.

More generally a Cartesian momentum residual
$f_i+\sum_j\partial_jT_{ij}$ contributes the torque
$\langle x_1\partial_jT_{2j}-x_2\partial_jT_{1j}\rangle$.
For a symmetric stress $T$, the product rule gives exactly
\[
 x_1\partial_jT_{2j}-x_2\partial_jT_{1j}
  =\partial_j(x_1T_{2j}-x_2T_{1j})-(T_{21}-T_{12})
  =\partial_j(x_1T_{2j}-x_2T_{1j}).
\]
Writing its vector flux in polar components therefore turns this
contribution into
\begin{equation}
 \partial_\sigma\langle r^2 T_{\theta r}\rangle
       +\partial_z\langle r T_{\theta z}\rangle,
 \quad T_{\theta r}=e_\theta^{\mathsf T}Te_r,
 \quad T_{\theta z}=e_\theta^{\mathsf T}Te_z.
 \label{eq:cq-am:stressmap}
\end{equation}
For a nonsymmetric tensor the retained extra torque is
$-\langle T_{21}-T_{12}\rangle$. Thus every symmetric stress
cancellation in the source momentum residual has this exact projected
image, with its sign determined by the side of the momentum equation
on which the source places $\operatorname{div}T$.
The map is linear in the full Cartesian tensor; its evaluation is not
an identification of the source's auxiliary-torus average with the
physical angular average used here.

The source cutoff sum has this exact image before the terminal limit.
In its original viscosity-one coordinates let
$\mathbf b_j$ denote the actual Cartesian vector appearing as $B_j$
in source (9.21). With $q_*$ here denoting that source's concentration
coordinate, retain every increment as
\[
 \delta u_{*,j}=\operatorname{curl}
          \bigl(c_*\chi(a_jq_*)A_j\bigr)
                  +c_*\chi(a_jq_*)\mathbf b_j,
 \qquad
 \delta p_{*,j}=c_*\chi(a_jq_*)p_j.
\]
The notation $q_*$ in this display is the source coordinate determined
by $1-t=q_*(1-\eta^2)$ and $z=q_*^{1/2-h}\eta$, not a material
label; those two equations specify its actual dependence on $(z,t)$.
Define each viscosity-$\nu$ increment by the full scaling of Section~\ref{sec:cq-NS-actual-witness}, and denote the correspondingly localized base by
$(u_{\nu,0},p_{\nu,0})$. The actual fields are
$u_\nu=\sum_{j\geq0}u_{\nu,j}$ and
$p_\nu=\sum_{j\geq0}p_{\nu,j}$, where the positive-index terms are
the displayed increments with that scaling. Source local finiteness
on every compact set before time one makes these sums, and every
derivative of them, finite on a neighborhood of each sampled circle.
Indeed the source coordinate is independent of $\theta$, so its
positive lower bound and its finite set of active stages can be
chosen over the entire compact circle. Thus the exact nonlinear fluxes
are
\[
 P=\sum_{a,b\geq0}\langle r^2(u_{\nu,a})_r(u_{\nu,b})_\theta\rangle,
 \qquad
 Q=\sum_{a,b\geq0}\langle r(u_{\nu,a})_z(u_{\nu,b})_\theta\rangle.
\]
Each sum contains every diagonal term and both ordered cross terms
for each pair $a\ne b$. The curl in $u_{\nu,j}$ is taken after
multiplying both cutoffs, so their derivative corrections belong to
these same fluxes. Projection, differentiation and the sums commute
by the stated local finiteness. Hence the source's exact residual
identity, including its flat remainder and final smooth force,
passes to \eqref{eq:cq-am:pde} with all these contributions. This argument
uses the actual summed fields for $t<1$; it does not exchange an
uncontrolled infinite series with the terminal limit.

\subsection{The radial Mellin equation with all shifted fluxes}

For each actual field $a=B,P,Q,F$ put
\[
 \widehat a(s,z,t)=\int_0^\infty a(\sigma,z,t)\sigma^{s-1}d\sigma,
 \qquad \sigma^{s-1}=e^{(s-1)\log\sigma}.
\]
The original measure here is $d\sigma=r\,dr$, with real logarithm;
angular averaging already contains its factor $1/(2\pi)$.
The integral for $B,Q,F$ is absolutely convergent and holomorphic on
$\mathbb Re s>-1$, and that for $P$ on $\mathbb Re s>-2$, by compact support
and \eqref{eq:cq-am:axis}. On compact substrips every complex derivative
adds an integrable power of $|\log\sigma|$, which proves
holomorphy by differentiated domination. All $z,t$ derivatives on
closed preterminal sets obey the same justification.

For $\mathbb Re s>0$, integration by parts gives
\begin{align*}
 \int_0^\infty\partial_\sigma P\,\sigma^{s-1}d\sigma
    &=-(s-1)\widehat P(s-1),\\
 \int_0^\infty 2\sigma\partial_{\sigma\sigma}B\,
                       \sigma^{s-1}d\sigma
    &=2s(s-1)\widehat B(s-1).
\end{align*}
All upper boundary terms vanish by compact support. The lower terms
$P\sigma^{s-1}$, $B'\sigma^s$ and $B\sigma^{s-1}$ vanish by
\eqref{eq:cq-am:axis} and $\mathbb Re s>0$. Consequently the full equation is
\begin{equation}
 \partial_t\widehat B-(s-1)\widehat P(s-1)+\partial_z\widehat Q
 =\nu\{2s(s-1)\widehat B(s-1)+\partial_{zz}\widehat B\}
       +\widehat F,\qquad\mathbb Re s>0.
 \label{eq:cq-am:mellin}
\end{equation}
In particular no shift $s\mapsto s-1$, polynomial factor, or force
moment has been suppressed.

This is an actual injective transform with its inverse. For any real
$c>0$ and fixed $(z,t)$ put $g(u)=e^{cu}a(e^u,z,t)$.
It vanishes for sufficiently large positive $u$ and is
$O(e^{(c+1)u})$ as $u\to-\infty$, with the same bound for every
derivative; for $P$ the exponent improves to $c+2$.
These facts follow from the chain rule and \eqref{eq:cq-am:axis}.
Thus $g$ is Schwartz and
\[
 \widehat a(c-i\eta,z,t)=\int_\mathbb R g(u)e^{-i\eta u}du,
 \qquad
 a(\sigma,z,t)=\frac{\sigma^{-c}}{2\pi}\int_\mathbb R
          \widehat a(c-i\eta,z,t)e^{i\eta\log\sigma}d\eta.
\]
To justify the inverse without a distributional shortcut, twice
integrate by parts in the forward integral to obtain an integrable
$O((1+|\eta|)^{-2})$ bound. Insert $e^{-\epsilon\eta^2}$ in the
inverse. Fubini and the Gaussian integral give convolution by
$(4\pi\epsilon)^{-1/2}e^{-u^2/(4\epsilon)}$, which converges to
$g$ by continuity and its Gaussian tail. Dominated convergence in
the integrable Fourier inverse then proves the displayed identity.

\subsection{The same arithmetic dilation map on the actual corrected fields}

Retain the source's chosen $0<h<1/100$. For any smooth compactly
supported half-line function $a$ with $a(0)=0$, define
\begin{align*}
 \phi_a(x)&=a(x)-2^h a(2x),\\
 \psi_a(x)&=\phi_a(x)-\tfrac12\phi_a(x/2)\\
  &=(1+2^{h-1})a(x)-2^ha(2x)-\tfrac12a(x/2),\\
 (\mathcal A_h a)(x)&=\sum_{n\geq1}\psi_a(nx),\qquad x>0.
\end{align*}
These are the exact three dilation coefficients used by the parent
heat-exterior calculation; here their input is the actual $B,P,Q,F$.
Every sum is locally finite on $x>0$, and the trace vanishes for all
sufficiently large $x$. Substitution gives
$\int_0^\infty\psi_a(x)dx=0$ and $\psi_a(0)=0$.
Its small-$x$ bound has a direct proof. On each interval $[nx,(n+1)x]$,
two integrations by parts give the trapezoidal identity
\[
 \int_{nx}^{(n+1)x}\psi_a(v)dv
 =\frac{x}{2}\{\psi_a(nx)+\psi_a((n+1)x)\}
 -\frac12\int_{nx}^{(n+1)x}
             (v-nx)((n+1)x-v)\psi_a''(v)dv.
\]
Summing, using compact support and the two zero quantities just
proved, yields
\begin{equation}
 |(\mathcal A_h a)(x)|\leq\frac{x}{8}\|\psi_a''\|_1.
 \label{eq:cq-am:tracebound}
\end{equation}
Thus its Mellin integral is holomorphic on $\mathbb Re s>-1$.
Parameter derivatives in $z,t$ have the same bound because the trace
is linear, preserves the zero endpoint and zero integral, and commutes
with those derivatives on preterminal compact sets.

On $\mathbb Re s>1$ absolute Fubini gives
\begin{equation}
 \mathcal M(\mathcal A_h a)(s)
    =W_h(s)\widehat a(s),\qquad
 W_h(s)=\zeta(s)(1-2^{s-1})(1-2^{h-s}).
 \label{eq:cq-am:arithmetic}
\end{equation}
Indeed the sum of absolute Mellin integrals is
$\zeta(\mathbb Re s)\int|\psi_a(x)|x^{\mathbb Re s-1}dx$, and a dilation by
$b>0$ multiplies its Mellin integral by $b^{-s}$.
The only pole of $\zeta$ is its simple pole at one, and the factor
$1-2^{s-1}$ cancels it. Consequently $W_h$ is entire and the identity
theorem extends \eqref{eq:cq-am:arithmetic} to $\mathbb Re s>-1$ for these
actual inputs. Its multiplier is the same arithmetic factor as in the
parent exterior calculation. The gamma product there is specifically
the Mellin transform of that exterior's $H_h$; it is not silently
substituted for the present corrected-field moment $\widehat a$.

The complete inverse preserves the actual field and its endpoint:
\begin{align*}
 \psi_a(x)&=\sum_{n\geq1}\mu(n)(\mathcal A_ha)(nx),\\
 \phi_a(x)&=\sum_{j\geq0}2^{-j}\psi_a(x/2^j),\\
 a(x)&=-\sum_{k\geq1}2^{-kh}\phi_a(x/2^k).
\end{align*}
The first double sum is finite at each $x>0$, so rearranging gives
$\sum_{m\geq1}\psi_a(mx)\sum_{n\mid m}\mu(n)=\psi_a(x)$.
The divisor sum is one for $m=1$ and zero otherwise, since the
squarefree expansion gives $\prod_{p\mid m}(1-1)$ for $m>1$.
The finite second sum through $j=J$ equals
$\phi_a(x)-2^{-J-1}\phi_a(x/2^{J+1})$.
The negative finite third sum through $k=K$ equals
$a(x)-2^{-Kh}a(x/2^K)$. Their tails vanish by boundedness at zero
and $h>0$. After any $x$ derivative, each summand gains the additional
factor $2^{-j\ell}$ or $2^{-k\ell}$; bounded derivatives on compact
half-line intervals prove uniform convergence there. The arithmetic
inverse is locally finite away from zero; the subsequent recovered
functions extend to zero by the displayed smooth inverse. These are
actual inverse maps on this stated class, including its parameter
dependence, with no division by zeta values at their zeros.

Set $\mathfrak B=\mathcal M\mathcal A_hB$ and define
$\mathfrak P,\mathfrak Q,\mathfrak F$ in the identical way.
Multiplication of \eqref{eq:cq-am:mellin} by $W_h(s)$ gives the exact
equation
\begin{align*}
 \partial_t\mathfrak B+\partial_z\mathfrak Q
       -\nu\partial_{zz}\mathfrak B-\mathfrak F
  =(s-1)W_h(s)\{\widehat P(s-1)+2\nu s\widehat B(s-1)\}.
\end{align*}
Both hatted terms on the right are explicitly recovered by the
proved inverses, so this equation retains the full dynamics. There is
also the division-free shifted relation
\begin{equation}
 \begin{split}
 W_h(s-1)\{\partial_t\mathfrak B+\partial_z\mathfrak Q
       -\nu\partial_{zz}\mathfrak B-\mathfrak F\}
  &=(s-1)W_h(s)\{\mathfrak P(s-1)
                        +2\nu s\mathfrak B(s-1)\},\\[-2pt]
 &\hspace{-60mm}\mathbb Re s>0.
 \end{split}
 \label{eq:cq-am:arithmeticpde}
\end{equation}
Initially this follows on $\mathbb Re s>1$; every factor is holomorphic on
$\mathbb Re s>0$ by the previous convergence proofs, so it holds throughout
that half-plane. The preceding inverse form retains the information
even at a zero of $W_h(s-1)$; the cross-multiplied relation alone is
not being asserted to reconstruct information lost by multiplying
by a vanishing scalar.

\subsection{The actual concentrating circle and its recoverable arithmetic image}

In the source's fixed inner similarity region all annular corrections
vanish. The remaining cutoff-summed slow field is axisymmetric, so
the swirl estimate on the positive $x_1$ axis in (10.21) holds on
the full circle with the same $r,z,t$. The corresponding source locators are (9.21), (10.20)--(10.21), and the unchanged inner-region construction of \cite{cqNS2026}. Hence at
$\sigma=\nu X_{\rm in}\tau,z=0,t=1-\tau$ the actual average is
\begin{equation}
 B(\nu X_{\rm in}\tau,0,1-\tau)
   =\nu\sqrt{2X_{\rm in}}\,\tau^{-h}(e_0+R_*(\tau)).
 \label{eq:cq-am:actualsample}
\end{equation}
The factor follows from multiplying the full retained swirl in
\eqref{eq:cq-nw:actualswirl} by $r=\sqrt{2\nu X_{\rm in}\tau}$;
the physical angular average of an axisymmetric value leaves it
unchanged. In particular this actual scalar angular momentum also
diverges, now at the source exponent $h$, with its original remainder.

Applying the proved inverse series to
$\mathcal A_hB(\cdot,0,1-\tau)$ recovers exactly the value in
\eqref{eq:cq-am:actualsample}; it is therefore a source-specific arithmetic
image of the concentrating corrected field. The double dyadic inverse
and the finite-at-each-argument arithmetic inverse retain their stated
order of evaluation. No uniform bound for those inverses as the
sampling point tends to zero has been inferred, so the recovery
identity alone is not a claim that the un-inverted trace has the same
pointwise blowup rate.

For completeness the resulting zeta factor also has a precise local
zero map. On $h<\mathbb Re s<1$, both dilation factors are nonzero.
Writing $\zeta(s)=(s-s_0)^m a_0(s)$ at a zeta zero gives
\[
 \mathfrak B(s,z,t)=(s-s_0)^m a_0(s)
       (1-2^{s-1})(1-2^{h-s})\widehat B(s,z,t).
\]
This is the complete product, including the corrected-field Mellin
factor. If that factor has order $n$ at $s_0$, the product has order
$m+n$, by multiplication of the two Taylor series; if it is nonzero,
the multiplicity is exactly $m$. Identical statements hold for each
flux and force moment. The arithmetic map and inverse are thus exact
for the nonlinear witness, while no off-critical zeta zero or
nonvanishing of the actual corrected-field Mellin factor is assumed.

\section{An exact fixed-axis singularity of the arithmetic trace}
\label{sec:cq-ns-axis-arithmetic-jet}

The complete companion fixed-axis proof \cite{cqNSAxis} is included with a local proof of its endpoint summation formula. The angular-momentum transform can see a fixed-axis derivative of the
actual source field directly, without first applying its inverse. The
reason is a specific feature of the construction: every positive-order
azimuthal coefficient has zero axis value. We retain the source amplitude
\(C\), the parameter \(h\), the physical viscosity \(\nu\), all cutoffs and
the original time \(t=1-\tau\). The calculation uses the constructed
field of Section~\ref{sec:cq-NS-actual-witness}; it does not certify the
entire existence argument of the released manuscript independently.

\subsection{The axis trace of the fully corrected velocity}

Here and only here let \(q_{\rm src}(z,t)\) denote the scalar similarity
coordinate called \(q\) in the source, not the material label in
Section~\ref{sec:cq-NS-profile-transfer}. Its source coordinates are
\[
 \tau=1-t=q_{\rm src}(1-\eta^2),\qquad
 z=q_{\rm src}^{D}\eta,\qquad
 X=\frac{r^2}{2q_{\rm src}},\qquad
 D=\frac12-h,\quad A=\frac12+h.
\]
In particular at \(z=0\) one has exactly \(\eta=0\) and
\(q_{\rm src}=\tau\).
Source (5.1), pp.~45--46, writes each azimuthal coefficient as
\[
 u_{\theta,n}
   =\frac rC q_{\rm src}^{-1-h+2nh}\varphi_n(X,\eta).
\]
Its Lemma~5.1, p.~47, prescribes
\(\varphi_n(0,\eta)=0\) for every \(n\ge1\).
The radial extensions and moment corrections are supported away from
the unchanged inner interval. Source Proposition~5.5 and its proof,
pp.~60--62, realize this series with cutoffs \(\chi(c_nq_{\rm src})\);
these depend on \(z,t\), not \(r\), and the realized angular component is
\begin{equation}\label{eq:cq-naj:base-series}
 \frac{u_{\theta,B}}r
 =\frac{q_{\rm src}^{-1-h}}C
 \left(\varphi_0(X,\eta)+
   \sum_{n\ge1}\chi(c_nq_{\rm src})
                       q_{\rm src}^{2nh}\varphi_n(X,\eta)\right).
\end{equation}
For every fixed preterminal \(z,t\), this sum is finite in a neighborhood
of that point: \(c_n\to\infty\), while \(q_{\rm src}>0\) and \(\chi\)
has compact support. Thus evaluation at \(X=0\), and the smooth extension
of the quotient \(u_{\theta,B}/r\), require no exchange of a terminal
limit with an infinite sum.

The leading axis datum is not an unknown positive remainder. Source
(B.3), p.~145, and Proposition~B.2, pp.~146--147, give
\[
 \varphi_0(0,\eta)=\varphi_*(\eta),\qquad
 \varphi_*(\eta)=
  \exp\left(\Lambda\int_0^\eta\zeta_*(w)\,dw\right),\qquad
 \varphi_*(0)=1.
\]
Here \(\zeta_*\) is the source's auxiliary axis function, not the
Riemann zeta function. Neither its value nor \(C\) is rescaled.

All annular wave and mean potentials, including the finite
initialization and the later cutoff-summed stages, vanish near the axis
for each \(t<1\): source Proposition~9.9, (9.21), pp.~114--117.
Only finitely many stages are present locally at such a time. Their
Cartesian derivatives therefore vanish on a common neighborhood of the
axis. The base potential is azimuthal and independent of angle, so its
curl has only radial and axial components. Consequently it contributes
no additional azimuthal term to \eqref{eq:cq-naj:base-series}.
Finally the spatial and temporal cutoffs in source (10.4), p.~118,
are identically one near the origin for all sufficiently small
positive \(\tau\). Their derivatives vanish there as well.
These are facts about the actual summed source field, not a
replacement of that field by its leading approximation
\cite[Lemma 5.1; Proposition 5.5; (9.21); (10.4); (B.3)]{cqNS2026}.

It follows that there is a source-determined \(\tau_1>0\) such that
\begin{equation}\label{eq:cq-naj:axis-star}
 \left.\frac{u_{\theta,*}(r,0,1-\tau)}r\right|_{r=0}
       =\frac1C\tau^{-1-h}\qquad(0<\tau<\tau_1).
\end{equation}
The quotient denotes its proved smooth extension, not division by
zero. For the physical viscosity map
\(u_\nu(x,t)=\sqrt\nu\,u_*(x/\sqrt\nu,t)\), put
\(r_*=r/\sqrt\nu\). Then, for \(r>0\),
\[
 \frac{u_{\theta,\nu}(r,0,t)}r
 =\frac{\sqrt\nu\,u_{\theta,*}(r_*,0,t)}{\sqrt\nu\,r_*}
 =\frac{u_{\theta,*}(r_*,0,t)}{r_*}.
\]
Passing to the smooth axis extension proves the same identity
\eqref{eq:cq-naj:axis-star} for \(u_\nu\). This cancellation follows from the
actual coordinate map; viscosity has not been set to one.

Use the original physical radial coordinate
\(\sigma=r^2/2\) and angular average \( (2\pi)^{-1}\int_0^{2\pi}\).
Let
\[
 B_\nu(\sigma,z,t)=
    \frac1{2\pi}\int_0^{2\pi} r u_{\theta,\nu}(r,\theta,z,t)\,d\theta .
\]
In the axis neighborhood just established the full field is
axisymmetric and \(B_\nu=2\sigma(u_{\theta,\nu}/r)\).
The exact consequence is
\begin{equation}\label{eq:cq-naj:B-axis}
 B_\nu(0,0,1-\tau)=0,\qquad
 \partial_\sigma B_\nu(0,0,1-\tau)=\frac2C\tau^{-1-h}.
\end{equation}
No \(O(\tau^{2h})\) error is present in this first axis derivative:
each positive-order axis value is exactly zero. This does not remove
the actual remainder at the distinct nonzero \(X_{\rm in}\) sample
of Section~\ref{sec:cq-NS-actual-witness}.

\subsection{The endpoint summation identity with its complete remainder}
\label{sec:cq-angular-EM-proof}

The endpoint calculation below uses an exact summation formula.
We give its proof here so that the fixed-axis result has no missing
local dependency. Define the Bernoulli polynomials by the convergent
power series, for the auxiliary variable $v$ near zero,
\[
 \frac{ve^{xv}}{e^v-1}=\sum_{n=0}^\infty B_n(x)\frac{v^n}{n!},
 \qquad B_n=B_n(0).
\]
The singularity at $v=0$ is removable. Differentiation in $x$ gives
$B_n'=nB_{n-1}$; subtracting the series at $x=0$ from that at
$x=1$ gives $B_n(1)=B_n(0)$ for $n\ne1$ and
$B_1(1)-B_1(0)=1$. Expansion to second degree gives
$B_1(x)=x-1/2$ and $B_2=1/6$.
The function $v/(e^v-1)+v/2$ is even, by direct substitution
of $-v$. Therefore $B_{2k+1}=0$ for every $k\geq1$.
Write $\widetilde B_n(y)=B_n(y-\lfloor y\rfloor)$ away
from integers, with the common endpoint value for $n\geq2$.
For these $n$ the function is continuous, bounded and periodic.

Let $f$ be smooth and compactly supported on $[0,\infty)$,
let $x>0$, and let $M\geq1$ be an integer. The exact identity is
\begin{equation}
\begin{split}
 \sum_{n\geq1}f(nx)
 &=\frac1x\int_0^\infty f(v)\,dv-\frac{f(0)}2
 -\sum_{k=1}^M\frac{B_{2k}}{(2k)!}
                         x^{2k-1}f^{(2k-1)}(0)+\mathcal R_M(f,x),\\
 \mathcal R_M(f,x)
 &=-\frac{x^{2M-1}}{(2M)!}
    \int_0^\infty\widetilde B_{2M}(v/x)f^{(2M)}(v)\,dv,\\
 |\mathcal R_M(f,x)|
 &\leq\frac{\|B_{2M}\|_{L^\infty([0,1])}}{(2M)!}
                          x^{2M-1}\|f^{(2M)}\|_{L^1(0,\infty)}.
\end{split}
\label{eq:cq-angular-EM}
\end{equation}
To prove it first take mesh one on $[0,N]$, with integer $N$
beyond the support. On every unit interval integration by parts gives
\[
 \int_n^{n+1} B_1(v-n)f'(v)\,dv
 =\frac{f(n)+f(n+1)}2-\int_n^{n+1}f(v)\,dv.
\]
Summing produces the trapezoidal sum minus the integral.
Repeated integration by parts using $B_j'=jB_{j-1}$ yields
the endpoint terms $B_{2k}(f^{(2k-1)}(N)-f^{(2k-1)}(0))/(2k)!$
and the remaining integral
$-\int_0^N\widetilde B_{2M}f^{(2M)}/(2M)!$.
The internal boundary terms cancel because the endpoint values
of $B_j$ agree for $j\geq2$; the odd endpoint coefficients
above degree one vanish by the already proved parity.
All upper endpoint derivatives vanish by the choice of $N$.
Apply this identity to $v\mapsto f(xv)$ and substitute $u=xv$
in each integral. This gives every power and sign in
\eqref{eq:cq-angular-EM}, including the stated norm bound.

For completeness the value of zeta needed at the first axis pole
also follows from this formula. Apply its finite-interval proof
on $[1,N]$ to $v^{-w}$ and initially let $N\to\infty$
for $\operatorname{Re}w>1$. For $M=2$ the result is
\[
 \zeta(w)=\frac1{w-1}+\frac12
    +\frac{B_2}{2!}w+\frac{B_4}{4!}w(w+1)(w+2)
    -\frac{(w)_4}{4!}\int_1^\infty
                     \widetilde B_4(v)v^{-w-4}\,dv.
\]
Here $(w)_4=w(w+1)(w+2)(w+3)$.
Boundedness of the periodic polynomial proves that the last
integral is holomorphic for $\operatorname{Re}w>-3$, with
all complex derivatives justified by the integrable logarithmic
weights on compact subsets. Thus the equality continues there
apart from the pole at one. Substitution $w=-1$ cancels the
first two terms, and the two higher products vanish. Hence
$\zeta(-1)=-B_2/2=-1/12$, with its sign fixed by the actual
summation formula. These classical formulas agree with
\cite[\S2.10(i)]{cqDLMFEM} and \cite[\S25.6(i)]{cqDLMFZeta}.


\subsection{The arithmetic operator on every endpoint derivative}

We prove the endpoint map on its full smooth compactly supported
domain. Let \(a\in C^\infty([0,\infty))\) have compact support and
\(a(0)=0\). For the same fixed \(h\) define
\[
 \phi_a(x)=a(x)-2^h a(2x),\qquad
 \psi_a(x)=\phi_a(x)-\tfrac12\phi_a(x/2),\qquad
 (\mathcal A_h a)(x)=\sum_{n\ge1}\psi_a(nx)\quad(x>0).
\]
This is exactly the arithmetic operator used on the nonlinear fluxes.
There is no radius change hidden in its argument \(x\).
Every sum and its derivatives are locally finite for \(x>0\).
Substitution in the two integrals proves
\(\int_0^\infty\psi_a=0\), and \(\psi_a(0)=0\).
The full endpoint derivatives of the input are
\begin{equation}\label{eq:cq-naj:psi-jet}
 \psi_a^{(j)}(0)
  =(1-2^{-j-1})(1-2^{h+j})a^{(j)}(0)\qquad(j\ge0).
\end{equation}

For clarity the passage from the sum to actual endpoint derivatives
is justified here, rather than differentiating a formal asymptotic
series. We use the Euler--Maclaurin identity proved with its remainder
in \eqref{eq:cq-angular-EM}; its classical provenance is
\cite[\S2.10(i), (2.10.1)]{cqDLMFEM}. With
\(f_j(v)=v^j\psi_a^{(j)}(v)\), direct differentiation on \(x>0\)
gives
\[
 x^j(\mathcal A_ha)^{(j)}(x)=\sum_{n\ge1}f_j(nx).
\]
All these \(f_j\) are smooth and compactly supported. Integration by
parts \(j\) times gives
\(\int f_j=(-1)^j j!\int\psi_a=0\); the factors \(v^j\) make
each zero-end boundary term vanish. Their axis derivatives are
\[
 f_j^{(r)}(0)=0\quad(r<j),\qquad
 f_j^{(r)}(0)=\frac{r!}{(r-j)!}\psi_a^{(r)}(0)\quad(r\ge j).
\]
Choose the Euler--Maclaurin order \(M\) so that \(2M-1>j\).
After division by \(x^j\), its remainder is bounded by
\[
 \frac{\|B_{2M}\|_\infty}{(2M)!}\,
 x^{2M-1-j}\|f_j^{(2M)}\|_1,
\]
and hence tends to zero. Terms below degree \(j\) vanish, so every
differentiated series has a finite limit at zero. The only constant
term for \(j\ge1\) is the Bernoulli term with \(2k-1=j\); for even
\(j\ge2\) it is absent. The fundamental theorem of calculus on
\([\epsilon,x]\), followed by \(\epsilon\downarrow0\), now identifies
these derivative limits with successive derivatives of the extended
function. Boundedness of the next derivative justifies the passage
under each integral. Thus \(\mathcal A_ha\) is smooth at zero, and
\begin{equation}\label{eq:cq-naj:all-jets}
 \begin{split}
 (\mathcal A_ha)(0)&=0,\\
 (\mathcal A_ha)^{(j)}(0)&=
 -\frac{B_{j+1}}{j+1}
 (1-2^{-j-1})(1-2^{h+j})a^{(j)}(0),\qquad j\ge1.
 \end{split}
\end{equation}
The Bernoulli convention is \(B_2=1/6\) and odd \(B_{j+1}=0\)
for even \(j\ge2\).
In particular
\begin{equation}\label{eq:cq-naj:first-map}
 (\mathcal A_ha)'(0)=
       \frac{2^{h+1}-1}{16}\,a'(0).
\end{equation}
The coefficient is strictly positive for the source \(h>0\).

The preceding proofs apply also after any fixed \(z,t\) derivative of
a smooth compactly supported parameter family with common support on
compact preterminal parameter sets. Every bound then holds uniformly
on those sets, with constants depending on the derivative order.
This proves compatibility of the endpoint map with those parameter
derivatives. It supplies no uniform estimate as \(t\uparrow1\).

\subsection{A directly visible singularity at a fixed arithmetic coordinate}

Apply this same operator to the actual compactly supported
\(a(\sigma)=B_\nu(\sigma,z,t)\). Write its arithmetic image as
\[
 \mathcal B_\nu(x,z,t)=\mathcal A_hB_\nu(x,z,t).
\]
The domain and all nonlinear fluxes are those already given in the
preceding angular-momentum section. Equations \eqref{eq:cq-naj:B-axis} and
\eqref{eq:cq-naj:first-map} now prove the exact formula
\begin{equation}\label{eq:cq-naj:exact-blowup}
 \boxed{\displaystyle
 \partial_x\mathcal B_\nu(0,0,1-\tau)
    =\frac{2^{h+1}-1}{8C}\,\tau^{-1-h},
       \qquad 0<\tau<\tau_1,\quad \nu>0 .}
\end{equation}
Here \(x=0,z=0\) are fixed coordinates, not a shrinking sampling path,
and \(\mathcal B_\nu(0,z,t)=0\) for every preterminal \(t\).
Thus the value is zero while its first \(x\)-derivative diverges
positively. For any fixed \(\epsilon>0\), the \(C^1([0,\epsilon])\)
seminorm is at least the displayed endpoint derivative. In particular
no extension continuous at \(t=1\) with values in
\(C^1([0,\epsilon])\) can agree with this family.
This assertion follows from an exact derivative of the original
arithmetic sum, not merely a lower bound obtained after its inverse.

For every integer \(m\ge0\), with \((1+h)_m=\prod_{j=0}^{m-1}(1+h+j)\) and empty product one, differentiation of the exact equality gives
\[
 \partial_t^m\partial_x\mathcal B_\nu(0,0,t)
  =\frac{2^{h+1}-1}{8C}(1+h)_m(1-t)^{-1-h-m},
 \qquad 1-\tau_1<t<1 .
\]
Indeed \(d_t(1-t)^{-\alpha}=\alpha(1-t)^{-\alpha-1}\).
These are equalities on an open preterminal interval. They make no
assignment of finite terminal derivatives.

The projection onto \(B_\nu\) has not discarded the other velocity
data in the full correspondence. Its exact smooth section and kernel
are
\[
 R B_\nu=\frac{B_\nu(r^2/2,z,t)}{r^2}(-x_2,x_1,0),\qquad
 w_\nu=u_\nu-RB_\nu,\qquad \langle r(w_\nu)_\theta\rangle=0,
\]
where \(B_\nu/\sigma\) has its smooth integral extension at zero.
The preceding section proves the inverse
\((B_\nu,w_\nu)\mapsto RB_\nu+w_\nu\) and retains every covariance
in both momentum fluxes. The signal \eqref{eq:cq-naj:exact-blowup} concerns
the \(B_\nu\) component of that complete map. It does not replace the
full Navier--Stokes equation by a closed scalar equation for \(B_\nu\).

\subsection{The exact Mellin pole carrying this axis signal}

For fixed preterminal parameters, put
\[
 \widehat B_\nu(s)=\int_0^\infty
                  B_\nu(\sigma,z,t)\sigma^{s-1}\,d\sigma,\qquad
 \mathfrak B_\nu(s)=\int_0^\infty
                  \mathcal B_\nu(x,z,t)x^{s-1}\,dx .
\]
Both integrals are initially holomorphic on \(\Re s>-1\).
The full transform calculation in the preceding section gives
\[
 \mathfrak B_\nu(s)=W_h(s)\widehat B_\nu(s),\qquad
 W_h(s)=\zeta(s)(1-2^{s-1})(1-2^{h-s}).
\]
The pole of \(\zeta\) at \(1\) is canceled by the displayed first
dilation factor, making \(W_h\) entire.

Here is the continuation needed at \(s=-1\). For any fixed \(b>0\),
split the integral at \(b\) and subtract the linear Taylor term on
\([0,b]\). Since \(B_\nu(0)=0\), the exact formula is
\[
 \widehat B_\nu(s)=
 \int_0^b [B_\nu(\sigma)-B_{\nu,\sigma}(0)\sigma]\sigma^{s-1}d\sigma
 +\int_b^\infty B_\nu(\sigma)\sigma^{s-1}d\sigma
 +\frac{B_{\nu,\sigma}(0)b^{s+1}}{s+1}.
\]
Taylor's integral remainder is \(O(\sigma^2)\); the first integral
is holomorphic for \(\Re s>-2\), including after any complex
\(s\) derivative on a compact subset, because the additional powers
of \(\log\sigma\) remain integrable. The second is entire by compact
support away from zero. Hence
\(\operatorname{Res}_{s=-1}\widehat B_\nu=B_{\nu,\sigma}(0)\).
The identical argument applies to the proved smooth
\(\mathcal B_\nu\), so its residue is \(\partial_x\mathcal B_\nu(0)\).

The classical value \(\zeta(-1)=-1/12\)
\cite[\S25.6(i)]{cqDLMFZeta} gives
\[
 W_h(-1)=-\frac1{12}\left(1-\frac14\right)(1-2^{h+1})
         =\frac{2^{h+1}-1}{16}>0.
\]
Equality on \(\Re s>-1\) therefore continues to the meromorphic
functions just constructed and agrees with the direct endpoint proof.
At \(z=0,t=1-\tau\) the pole is genuine and its residue is exactly
\begin{equation}\label{eq:cq-naj:residue}
 \operatorname{Res}_{s=-1}\mathfrak B_\nu(s,0,1-\tau)
     =\frac{2^{h+1}-1}{8C}\tau^{-1-h}.
\end{equation}
There is no zero-factor compensation at this pole, since \(W_h(-1)>0\).
The sign of \(\zeta(-1)\) alone is not the sign of the arithmetic
residue: both dilation factors have been evaluated explicitly.

This identifies what the construction propagates in this channel:
an exact axis derivative, equivalently a growing residue at the fixed
Mellin point \(-1\). The identity retains the full
\(\widehat B_\nu\), its flux equations, and its inverse maps.
The nontrivial zeros of \(\zeta\) remain the zero factors specified
in the full product formula; no off-critical zero is inferred from
this axis pole or from its positive residue.

\section{Exact entrance of the angular Mellin profiles into the source quotient}
\label{sec:cq-angular-ingress}

\subsection{The original direct dilation and every endpoint pole}

In this section $\sigma=r^2/2\geq0$ is the physical radial coordinate,
$z$ is the unchanged axial coordinate, $t$ is physical time, $w$ is a
Mellin variable, and $s$ is the original Connes spectral variable.
The radial Mellin coordinate called $s$ in Sections~\ref{sec:cq-angular-physical-map} and~\ref{sec:cq-ns-axis-arithmetic-jet} is denoted by $w$ here. The original Connes spectral coordinate remains $s$. Their exact affine relation is $w=1/2-is$, with inverse $s=i(w-1/2)$. The functional equation gives $\Xi(s)=\xi(1/2+is)=\xi(w)$.
Fix $(z,t)$ before terminal time and let
\[
 a\in\mathscr A=\{a\in C_c^\infty([0,\infty)):a(0)=0\}.
\]
The actual inputs $B,P,Q,F$ of the angular construction belong to
this space. Dependence on $(z,t)$ is retained as a smooth parameter
on every compact preterminal set. Keep the source's chosen
$0<h<1/100$ and its exact dilation operators
\[
 \phi_a(\sigma)=a(\sigma)-2^ha(2\sigma),\quad
 \psi_a(\sigma)=\phi_a(\sigma)-\tfrac12\phi_a(\sigma/2),\quad
 A_ha(x)=\sum_{n\geq1}\psi_a(nx).
\]
The letter $x>0$ in this last expression is the argument of the
direct arithmetic dilation; it has the same numerical value as
the input radial variable and is not the Cartesian first coordinate.
The source identity, with its unchanged measure, is
\begin{equation}
 \begin{gathered}
 \mathcal M(A_ha)(w)=W_h(w)\widehat a(w),\\
 W_h(w)=\zeta(w)(1-2^{w-1})(1-2^{h-w}),\qquad
 \widehat a(w)=\int_0^\infty a(\sigma)\sigma^{w-1}d\sigma .
 \end{gathered}
 \label{eq:cq-angular-direct-Mellin}
\end{equation}
It initially holds on $\operatorname{Re}w>1$ and extends to
$\operatorname{Re}w>-1$ by the proved endpoint bound.

Write $a_n=a^{(n)}(0)/n!$, $n\geq1$. The full meromorphic
continuation of $\widehat a$ has only possible simple poles at
$w=-n$, with
\begin{equation}
 \operatorname{Res}_{w=-n}\widehat a(w)=a_n.
 \label{eq:cq-angular-Mellin-residues}
\end{equation}
Here is a direct proof for arbitrary smooth input. Choose a smooth
compact cutoff $\chi$ equal to one near zero. For any $N\geq1$,
$r_N=a-\chi\sum_{n=1}^Na_n\sigma^n$ vanishes to order $N+1$.
Thus $\int r_N(\sigma)\sigma^{w-1}d\sigma$ is holomorphic on
$\operatorname{Re}w>-N-1$. The cutoff Mellin transform
$C_\chi(v)=\int\chi(\sigma)\sigma^{v-1}d\sigma$ satisfies
$C_\chi(v)=-v^{-1}\int\chi'(\sigma)\sigma^v d\sigma$.
The numerator is entire and has value $-1$ at $v=0$, giving
residue one. Consequently
\[
 \widehat a(w)=
 \int_0^\infty r_N(\sigma)\sigma^{w-1}d\sigma
                  +\sum_{n=1}^Na_nC_\chi(w+n).
\]
These expressions agree on their original common half-plane and
therefore on their overlaps. They prove the continuation and every
residue in \eqref{eq:cq-angular-Mellin-residues}. Euler-operator
integration by parts in the remainder, with its exact endpoint
vanishing order, also proves rapid decrease on every finite
vertical strip away from the finitely many poles, exactly as in
the previous full even-Schwartz Mellin proof.

The dilation factor
$d_h(w)=(1-2^{w-1})(1-2^{h-w})$ is nonzero at every $w=-n$,
$n\geq1$. At negative even integers the simple zero of $\zeta$
cancels the possible Mellin pole. At negative odd integers
$\zeta(-n)\ne0$: in the retained completion
$\xi(w)=w(w-1)\pi^{-w/2}\Gamma(w/2)\zeta(w)/2$, the other
factors are finite and nonzero there, while
$\xi(-n)=\xi(1+n)>0$. The gamma recurrence makes
$\Gamma(-n/2)$ finite and nonzero for odd $n$, starting from
$\Gamma(1/2)>0$; the positivity at $1+n$ follows from its
positive Dirichlet series. At $w=1$ the zero of $1-2^{w-1}$
cancels the sole zeta pole. Thus
\begin{equation}
 W_h\widehat a\text{ is entire}
 \quad\Longleftrightarrow\quad
 a^{(2k+1)}(0)=0\text{ for every }k\geq0.
 \label{eq:cq-angular-direct-entire-domain}
\end{equation}
Its uncanceled poles and residues are precisely
\[
 \operatorname{Res}_{w=-n}W_h(w)\widehat a(w)
 =\zeta(-n)(1-2^{-n-1})(1-2^{h+n})a_n,
 \qquad n\geq1\text{ odd}.
\]
For example its first pole coefficient is
$\zeta(-1)(3/4)(1-2^{h+1})a'(0)$; all factors multiplying
$a'(0)$ are nonzero.

This condition is also the exact parity condition for the direct
original trace. Taylor differentiation gives
\[
 \frac{\psi_a^{(n)}(0)}{n!}
  =(1-2^{-n-1})(1-2^{h+n})a_n .
\]
The multiplier is nonzero for every integer $n\geq1$.
Therefore $\psi_a(|x|)$ is smooth and even on $\mathbb R$
exactly when all odd jets of $a$ vanish. The even extension is then
compactly supported, zero at zero, and has integral zero, because
$\int_0^\infty\psi_a=0$. It belongs to the precise
$\mathcal S_0^{\rm ev}$ used in the source. Its original arithmetic
trace is the entire continuation
\[
 \mathcal F_a(s)=W_h(1/2-is)\widehat a(1/2-is).
\]
The identity with the source trace follows first for
$\operatorname{Re}w>1$ by
$F_\psi(i(w-1/2))=\zeta(w)\mathcal M\psi(w)$, and then everywhere
by analytic continuation. This proves the exact trace identity
without calling a general meromorphic Mellin profile entire.

For the actual angular momentum the first obstruction is a familiar
but fully specified Cartesian observable:
\begin{equation}
 \partial_\sigma B(0,z,t)
 =(\partial_{x_1}u_2-\partial_{x_2}u_1)(0,0,z,t)
 =\omega_z(0,0,z,t).
 \label{eq:cq-angular-axis-vorticity}
\end{equation}
Indeed expand the smooth Cartesian $u_1,u_2$ to first order at
the axis. In
$\langle x_1u_2-x_2u_1\rangle$, constant velocity terms have zero
average; the mixed quadratic terms have zero average; and
$\langle x_1^2\rangle=\langle x_2^2\rangle=r^2/2=\sigma$.
The remaining coefficient is exactly
$\partial_{x_1}u_2-\partial_{x_2}u_1$.
The higher-order remainder is $o(r^2)$ after angular averaging,
and smoothness in $\sigma$ identifies the displayed derivative.
Thus the angular source's $B=O(\sigma)$ alone does not enforce the
first, or any of the higher odd-jet entrance conditions.

\subsection{Order-one growth is an additional exact entrance condition}

For either entire profile below, order at most one is tested in the
original coordinate by the previously defined growth seminorms.
The full direct domain is therefore the explicit linear space
\begin{equation}
\begin{split}
 \mathscr A_{\rm dir}
 =\biggl\{a\in\mathscr A:\;&a^{(2k+1)}(0)=0\ (k\geq0),\\
 &\sup_{w\in\mathbb C}|W_h(w)\widehat a(w)|
       e^{-c|i(w-1/2)|^p}<\infty
       \quad(1<p<2,\ c>0)\biggr\}.
\end{split}
 \label{eq:cq-angular-direct-order-domain}
\end{equation}
The product in this expression uses its proved removable values.
This is a necessary and sufficient condition expressed entirely
through the original input's Mellin integrals and endpoint
continuation above: the parity statement is exactly entireness,
and the displayed bounds are exactly the proved characterization
of the original set $\cqE$. No derivative-size bound has been
silently imposed on arbitrary smooth compact input.
For $a\in\mathscr A_{\rm dir}$, the original Mellin homeomorphism
and trace intersection theorem give
\begin{equation}
 \mathcal F_a\in\mathcal A_{\rm Sch}\subset I=\Xi\cqE,\qquad
 \frac{\mathcal F_a(i(w-1/2))}{\zeta(w)}
       =d_h(w)\widehat a(w)
       \in\mathfrak M_0^{\rm ev}.
 \label{eq:cq-angular-direct-trace-intersection}
\end{equation}
All its Mellin strip bounds and the unique even Schwartz preimage
are therefore the ones already proved, rather than an unproved
new source topology.

Here is a useful sufficient entrance class with a complete growth
argument. Suppose $a$ is analytic at the endpoint, meaning that
$a(\sigma)=\sum_{n\geq1}a_n\sigma^n$ on $[0,\delta_0]$ for a
convergent power series of positive radius. Choose
$0<\delta<\delta_0$ inside its convergence disk, and let
$R_0>\delta$ contain the support. Exact integration gives
\begin{equation}
 \widehat a(w)=
  \int_\delta^{R_0}a(\sigma)\sigma^{w-1}d\sigma
       +\sum_{n\geq1}\frac{a_n\delta^{w+n}}{w+n}.
 \label{eq:cq-angular-analytic-endpoint-Mellin}
\end{equation}
The series converges normally away from its indicated poles because
$\sum|a_n|\delta^n<\infty$. Outside the radius-$1/4$ disks about
the negative integers it gives
$|\widehat a(w)|\leq C e^{C|w|}$: bound the integral over its
fixed positive interval and use $|w+n|^{-1}\leq4$ in the series.

For clarity the zeta factor used in the remaining growth argument
also has order at most one after its pole is removed.
Let $R_\Gamma(z)=1/\Gamma(z)$. The entire product
\[
 R_\Gamma(z)=
 z e^{\gamma_{\rm E}z}
       \prod_{n=1}^\infty(1+z/n)e^{-z/n}
\]
has logarithmic modulus at most $C R\log(R+2)$ on $|z|\leq R$
for $R\geq1$. To verify both its identity and this bound, repeated
integration by parts in the beta integral gives
\[
 \int_0^1t^{z-1}(1-t)^n dt
     =\frac{n!}{z(z+1)\cdots(z+n)}
              \quad(\operatorname{Re}z>0).
\]
After $u=nt$, multiplication by $n^z$ makes the integral converge
to $\Gamma(z)$ by domination with $u^{\operatorname{Re}z-1}e^{-u}$,
locally uniformly in that half-plane. Rewriting its reciprocal
using $H_n-\log n\to\gamma_{\rm E}$ gives the displayed product.
The limit of $H_n-\log n$ exists by comparison of $1/n$ with
the integral of $1/x$ over each unit interval. Its product tails
converge locally uniformly because
$|\log(1+v)-v|\leq2|v|^2$ for $|v|\leq1/2$.
For $n>2R$ their absolute logarithmic sum is at most
$2R^2\sum_{n>2R}n^{-2}\leq C R$.
For $n\leq2R$, use
$(1+R/n)e^{R/n}$ to bound each factor; the logarithmic sum is
at most $2R\log(1+R)+R H_{\lfloor2R\rfloor}$.
The initial factors give the claimed bound. The finite beta-integral
expression times its rewritten reciprocal is identically one;
passing to the two locally uniform limits proves
$\Gamma(z)R_\Gamma(z)=1$ on the positive half-plane, without
presuming nonvanishing of a limit. Analytic continuation proves
the reciprocal-gamma identity everywhere, including its zeros.

The retained completion now gives
\[
 (w-1)\zeta(w)
     =2\pi^{w/2}\xi(w)\frac{R_\Gamma(w/2)}w .
\]
The last quotient has its removable value $1/2$ at zero.
The original $\xi(w)=\Xi(i(w-1/2))$ has order at most one;
the product bound just proved and multiplication by an exponential
show the same for $(w-1)\zeta(w)$. Division by $w$ is bounded
outside a fixed disk, and removable holomorphy bounds its inside.
This proves the assertion with the actual completion constants.
The two dyadic factors have exponential type. Combining these
bounds with \eqref{eq:cq-angular-analytic-endpoint-Mellin} proves
the order-one bounds for $W_h\widehat a$ outside the fixed
negative-integer disks and a fixed disk about one.
When the parity conditions hold, that product is entire. The
maximum-modulus principle on every omitted disk supplies the same
order-one bounds throughout its interior, with radius enlarged
by at most one. Hence every endpoint-analytic $a$ with the stated
odd jets zero belongs to $\mathscr A_{\rm dir}$.
In particular every smooth compact input vanishing on a whole
neighborhood of zero belongs to this domain.

\subsection{A square-root radial ingress retaining every input}

There is a second exact map which enforces the source's even
test-space condition and retains the original radial coordinate.
Define
\[
 b_a(\rho)=a(\rho^2/2),\quad
 \chi_a(\sigma)=\phi_a(\sigma)-2^{-1/2}\phi_a(\sigma/2),\quad
 \psi_a^{\rm rad}(\rho)=\chi_a(\rho^2/2).
\]
Here $\rho$ is a signed real test coordinate; $\sigma=\rho^2/2$
and $|\rho|=\sqrt{2\sigma}$. This is exactly the earlier physical
radius with both signs supplied for the even test. It does not
assert equality with the polynomial coordinate $F_1/2$.
Each $b_a$ and $\psi_a^{\rm rad}$ is smooth, compactly supported
and even, with value zero at zero. Its integral satisfies
\[
 \int_0^\infty\psi_a^{\rm rad}(\rho)d\rho
 =2^{-1/2}\int_0^\infty\chi_a(\sigma)\sigma^{-1/2}d\sigma=0.
\]
Indeed the second dilation in $\chi_a$ multiplies the last integral
by $2^{1/2}$, exactly canceled by its coefficient $2^{-1/2}$.
Thus $\psi_a^{\rm rad}\in\mathcal S_0^{\rm ev}$ for every
$a\in\mathscr A$.

With no measure factor removed, substitution $\rho=\sqrt{2\sigma}$
gives the exact Mellin and source trace formulas
\begin{align}
 M_a^{\rm rad}(w)
 &=2^{w/2-1}(1-2^{w/2-1/2})(1-2^{h-w/2})
                    \widehat a(w/2),\notag\\
 \mathcal F_a^{\rm rad}(s)
 &=\zeta(w)M_a^{\rm rad}(w),\qquad w=1/2-is.
 \label{eq:cq-angular-radial-source-map}
\end{align}
Initially derive them on $\operatorname{Re}w>2$ by absolute
integration and the source trace formula. The proved even
Schwartz trace is entire, so it supplies their continuation
with all removable values. At $w=-2n$ the only possible pole of
$M_a^{\rm rad}$ has the full residue
\begin{equation}
 \operatorname{Res}_{w=-2n}M_a^{\rm rad}(w)
 =2^{-n}(1-2^{-n-1/2})(1-2^{h+n})a_n,\qquad n\geq1 .
 \label{eq:cq-angular-radial-Taylor-data}
\end{equation}
The factor $2$ from the residue of $\widehat a(w/2)$ combines
with $2^{w/2-1}$ to give $2^{-n}$, as displayed.
These are precisely the Taylor coefficients of
$\chi_a(\rho^2/2)$ at powers $\rho^{2n}$; all its original
radial jets have been retained.

This map is injective and has a full inverse on its image.
From its entire source trace recover the meromorphic function
$M_a^{\rm rad}$ by division by $\zeta$, with the removable values
specified by the identity, and apply the already proved even
Schwartz Mellin inverse to recover $\psi_a^{\rm rad}$.
Then
\begin{align}
 \chi_a(\sigma)&=\psi_a^{\rm rad}(\sqrt{2\sigma}),\notag\\
 \phi_a(\sigma)&=\sum_{j\geq0}2^{-j/2}\chi_a(\sigma/2^j),\notag\\
 a(\sigma)&=-\sum_{k\geq1}2^{-kh}\phi_a(\sigma/2^k).
 \label{eq:cq-angular-radial-exact-inverse}
\end{align}
The finite second sum leaves remainder
$2^{-(J+1)/2}\phi_a(\sigma/2^{J+1})$ and the finite third
sum leaves remainder $2^{-Kh}a(\sigma/2^K)$.
Both tend to zero by boundedness at zero. After $\ell$
derivatives, their terms gain $2^{-j\ell}$ and $2^{-k\ell}$;
bounded derivatives on compact half-line intervals give uniform
convergence for every derivative. The first recovered function
has its smooth extension at zero because its data lie in this
image. Thus these formulas recover the full original $a$, not
only its zero set or finitely many moments.

In the original radial Mellin variable the exact relation between
the two arithmetic transforms is the common input $\widehat a$:
\begin{align*}
 \mathcal F_a(i(w-\tfrac12))
 &=\zeta(w)(1-2^{w-1})(1-2^{h-w})\widehat a(w),\\
 \mathcal F_a^{\rm rad}(i(2w-\tfrac12))
 &=2^{w-1}\zeta(2w)(1-2^{w-1/2})(1-2^{h-w})\widehat a(w).
\end{align*}
The first line uses its meromorphic continuation for arbitrary
input. The second is entire in its original spectral variable.
Both retain the distinct zeta arguments, dyadic coefficients
and factor $2^{w-1}$. Their proved inverse constructions recover
the same input even at zeros of these factors. Neither formula
divides one entire trace by the other at an unexamined zero.

The full order-one entrance domain of this second map is
\begin{equation}
 \mathscr A_{\rm rad}=
 \left\{a\in\mathscr A:
 \sup_{w\in\mathbb C}
 |\zeta(w)M_a^{\rm rad}(w)|
 e^{-c|i(w-1/2)|^p}<\infty
       \quad(1<p<2,\ c>0)\right\}.
 \label{eq:cq-angular-radial-order-domain}
\end{equation}
Entireness is already proved for every input; the displayed
condition is exactly the remaining order condition.
Every endpoint-analytic $a$ belongs to this domain, without
any parity restriction. To prove it, apply
\eqref{eq:cq-angular-analytic-endpoint-Mellin} with argument
$w/2$, the just-proved zeta growth bound and the exponential
dyadic factors. The only possible Mellin poles are at negative
even $w$, and all are canceled by the simple trivial zeros.
Maximum modulus fills their fixed disks as before.
For every $a\in\mathscr A_{\rm rad}$ the trace is in
$\mathcal A_{\rm Sch}\subset\Xi\cqE$, and its precise source
Schwartz input is \eqref{eq:cq-angular-radial-source-map} with
the inverse \eqref{eq:cq-angular-radial-exact-inverse}.

\subsection{An exact full-input family in the original order class}

Arbitrary smooth endpoint data still have an information-preserving
entrance into a family of actual order-one trace functions.
Fix one smooth function $\chi:[0,\infty)\to[0,1]$ equal to zero
on $[0,1]$ and one on $[2,\infty)$, and set
\[
 \epsilon_n=2^{-n},\qquad a_n(\sigma)=\chi(\sigma/\epsilon_n)a(\sigma),
                  \qquad n\geq1.
\]
These are specified functions of the full input, not replacements
of it. Every $a_n$ vanishes near zero, so both $\mathcal F_{a_n}$
and $\mathcal F_{a_n}^{\rm rad}$ lie in the original $\cqE$ by
the sufficient entrance theorem. Retain the whole sequence
$\mathfrak I(a)=(\mathcal F_{a_n}^{\rm rad})_{n\geq1}$.
It is injective: invert each entry by
\eqref{eq:cq-angular-radial-exact-inverse}; for every fixed
$\sigma>0$ all sufficiently large $n$ have
$a_n=a$ on a neighborhood of $\sigma$. Hence the recovered
values and all derivatives there are eventually exact.
At zero the full original jets are recovered by
$a^{(k)}(0)=\lim_{\sigma\downarrow0}a^{(k)}(\sigma)$,
using the known smooth half-line domain. Also $a(0)=0$ is retained.
This proves the full inverse on the image, including endpoint
information; no finite truncation of the sequence has been
asserted sufficient.

This gives a direct actual entrance for all the corrected angular
fields without an unproved endpoint analyticity assertion.
For the source sample
$\sigma_\tau=\nu X_{\rm in}\tau$ choose any integer $n(\tau)$
with $2\epsilon_{n(\tau)}<\sigma_\tau$. Then inversion of that
specified sequence entry recovers the exact sample value
\[
 B(\sigma_\tau,0,1-\tau)
 =\nu\sqrt{2X_{\rm in}}\tau^{-h}(e_0+R_*(\tau)).
\]
All sample-dependent index choices recover the same full value,
since the cutoff is exactly one there. This is a proved inverse
evaluation of the source's actual sample, not an assertion that
an uninverted trace or its source quotient class has this same
pointwise rate.

\subsection{The source kernels, arithmetic refinement, and derivative defect}

For $\bullet\in\{{\rm dir},{\rm rad}\}$ give
$\mathscr A_\bullet$ the initial topology of its inclusion in
the usual compact-test space together with
$a\mapsto\mathcal F_a^\bullet\in(\cqE,\tau\vee\gamma)$.
Here the compact-test topology is the inductive limit of the spaces
supported in $[0,R]$, $R=1,2,\ldots$, with seminorms
$\max_{0\leq\sigma\leq R}|a^{(j)}(\sigma)|$, $j\geq0$.
This explicitly defines the topology required for the following
maps; it does not assert continuity into the unnamed $\tau$
from the compact-test topology alone.
The actual source quotient entrance and its exact kernel are
\begin{equation}
 q_\bullet:\mathscr A_\bullet\longrightarrow\cqQ,\quad
 a\longmapsto[\mathcal F_a^\bullet],\qquad
 \ker q_\bullet
  =\{a\in\mathscr A_\bullet:\mathcal F_a^\bullet\in\cqN\}.
 \label{eq:cq-angular-source-kernel}
\end{equation}
Each map is continuous by the topology just specified.
Its underlying entire-trace map is injective: divide by zeta on
its meromorphic image, use the proved Mellin inverse to recover
the exact test, and then the convergent dilation inverses to
recover $a$. For the direct map those inverses are the original
$\phi_a=\sum_{j\geq0}2^{-j}\psi_a(\sigma/2^j)$ and
$a=-\sum_{k\geq1}2^{-kh}\phi_a(\sigma/2^k)$;
their telescoping proof is identical with the retained exponent
$1$ instead of $1/2$ in the second dilation.
Thus \eqref{eq:cq-angular-source-kernel} describes exactly the
information lost on subsequently imposing the actual source
relation, not an ambiguity of the Mellin reconstruction.

Let $H_\cap=\cqN\cap I$ and
$C_\cap=(\cqE,\tau\vee\gamma)/H_\cap$ be the previously proved
common refinement. The maps
$a\mapsto[\mathcal F_a^\bullet]_{C_\cap}$ are continuous and
project to \eqref{eq:cq-angular-source-kernel}. Their kernels
are equal to those in that equation, since all the trace images
already lie in $I$. Their images lie in the exact kernel
$I/H_\cap$ of $C_\cap\to Q_I$. This proves the precise
comparison with the analytic realization; it does not identify
$\cqN$ with $I$ or assert $\mathcal A_{\rm Sch}\subseteq\cqN$.

The raw traces alone therefore have zero analytic principal
parts. Their derivative supplies a different explicitly
calculated coordinate. Write its unique entire quotient as
$\mathcal F_a^\bullet=\Xi g_a^\bullet$, where
$g_a^\bullet\in\cqE$ by the proved division theorem.
At each actual zero $\lambda$ of multiplicity $m_\lambda$,
\begin{equation}
 \frac{D\mathcal F_a^\bullet}{\Xi}
  =g_a^{\bullet\,\prime}
        +g_a^\bullet\frac{\Xi'}{\Xi},\qquad
 \operatorname{PP}_{\lambda}
       \frac{D\mathcal F_a^\bullet}{\Xi}
  =\frac{m_\lambda g_a^\bullet(\lambda)}{s-\lambda}.
 \label{eq:cq-angular-derivative-principal-part}
\end{equation}
To prove the complete local assertion, insert
$\Xi=(s-\lambda)^{m_\lambda}A_\lambda$ into the first equality.
Then $\Xi'/\Xi=m_\lambda/(s-\lambda)+A_\lambda'/A_\lambda$,
and
$(g_a^\bullet(s)-g_a^\bullet(\lambda))/(s-\lambda)$
is holomorphic. All remaining terms, including the actual unit,
are holomorphic; the displayed simple principal part is exact.

The source class of this derivative is still its class modulo
the actual $\cqN$. Its entire presenting relation has the same
already calculated ambiguity $(D\cqN+\cqN)/\cqN$.
For the specific element $\beta=[\Xi']$, equality with any
angular derivative or raw entrance is exactly the respective
membership test
\[
 \Xi'-D\mathcal F_a^\bullet\in\cqN
 \quad\text{or}\quad
 \Xi'-\mathcal F_a^\bullet\in\cqN.
\]
Neither follows from a nonzero angular sample or a zeta
factorization. At the common refinement, raw images belong to
$I/H_\cap$ while $[\Xi']_{C_\cap}$ does not, because its
analytic residue is $m_\lambda\ne0$. Their relation to the
source quotient remains the proved projection with kernel
$\cqN/H_\cap$, not a claimed nonrelation.
Likewise the full heat domain is the inverse image of
$\mathcal A_{\rm Sch}$ under $U_\theta$ as already computed;
its zero-jet and Mellin-strip conditions apply to these exact
traces. No frozen-source inclusion $U_\theta\cqN\subseteq\cqN$
is inferred from their arithmetic dilation or from the
Navier--Stokes angular momentum sample.

\subsection{The radial and polynomial coordinates are an explicit common map}

On the real angular chart retain
\[
 (x_1,x_2,z)
   =(\sqrt{2\sigma}\cos\theta,\sqrt{2\sigma}\sin\theta,z).
\]
Its inverse is $\sigma=(x_1^2+x_2^2)/2$ and the unit direction
$(\cos\theta,\sin\theta)=(x_1,x_2)/\sqrt{x_1^2+x_2^2}$,
with unchanged $z$. The original polynomial coordinate is
\[
 \sigma_{\rm pol}=\frac12\left[
 (1+\sigma\sin2\theta)^3z+
 2\sigma\sin^2\theta(1+\sigma\sin2\theta)
                        (4+3\sigma\sin2\theta)\right].
\]
Indeed $x_1x_2=\sigma\sin2\theta$ and
$x_2^2=2\sigma\sin^2\theta$ substitute directly into the
unchanged $F_1/2$. Adjoining $\sigma_{\rm pol}$ gives an
injective graph map with inverse projection onto
$(\sigma,\theta,z)$; in particular no division or inverse
Jacobian is needed to relate the two observables. The
square-root test ingress above uses $\sigma=\rho^2/2$, whereas
the earlier fluid-profile translation uses this displayed
$\sigma_{\rm pol}$. Their exact common Cartesian chart and
the graph formula retain the full coordinate dictionary.

\subsection{Analytic radial entrance for the actual corrected central field}
\label{sec:cq-actual-angular-analytic}

The preceding general entrance criteria can be applied to the actual
central field without an additional analyticity assumption. We retain
the profiles and all cutoff choices of the cited source. The argument
below proves the required radial analyticity locally at the axis; it
does not claim that a general smooth cutoff is analytic on its
transition region.

First consider the original viscosity-one similarity variables
\[
 \tau=1-t,\qquad \tau=q(1-\eta^2),\qquad
 z=q^{1/2-h}\eta,\qquad X=\frac{r^2}{2q}.
\]
For every preterminal point in the source's local domain, $q>0$.
The source's Proposition B.2 supplies a leading profile jointly
analytic in $X$ and $\eta$ on a neighborhood of its initial radial
rectangle. Every positive-order profile actually selected by the
source is analytic in $X$ near $X=0$, as we now prove from its
displayed recurrence.

For an explicit definition of every coefficient in the ensuing
linear system, retain the complete original recurrence
\cite[equations (4.1)--(4.2), (5.1)--(5.6)]{cqNS2026}.
In the present recurrence $A,D$ are the source's scalar
exponents, and $D_X$ is its logarithmic radial derivative:
\begin{equation}
\begin{gathered}
 A=\tfrac12+h,\qquad D=\tfrac12-h,\qquad
 \lambda_n=2nh,\qquad b=-A-\tfrac12,\qquad c=-A,\\
 d=1-\eta^2,\qquad L=1-2h\eta^2,\qquad D_X=X\partial_X,\\
 \tau=1-t=q(1-\eta^2),\qquad z=q^D\eta,\qquad
 \sigma=\frac{r^2}{2},\qquad X=\frac{\sigma}{q}.
\end{gathered}
 \label{eq:cq-source-recurrence-parameters}
\end{equation}
The source denotes the radial $\sigma$ in this line by $s$;
the equation explicitly identifies that source variable with
our retained radial coordinate, while leaving the Connes
spectral $s$ distinct. Every other source coefficient in the
recurrence is unchanged. For a real exponent $a$ define
\begin{align}
 T_a f&=L^{-1}(-af+D\eta\,\partial_\eta f+D_Xf),\notag\\
 Z_a f&=L^{-1}(2a\eta f+d\,\partial_\eta f-2\eta D_Xf),\notag\\
 T_{a,n}&=T_{a+\lambda_n},\qquad
 Z_{a,n}=Z_{a+\lambda_n},\qquad
 Z^{[2]}_{a,n}=Z_{a+\lambda_n-D}Z_{a+\lambda_n}.
 \label{eq:cq-source-profile-time-axial-operators}
\end{align}
The rightmost operator acts first in the last expression.
These are the exact physical derivative identities
\[
 \partial_t(q^af)=q^{a-1}T_af,\qquad
 \partial_z(q^af)=q^{a-D}Z_af.
\]
Indeed implicit differentiation of the unchanged similarity
coordinates gives
\[
 q_t=-L^{-1},\quad \eta_t=\frac{D\eta}{qL},\quad
 X_t=\frac{X}{qL},\qquad
 q_z=\frac{2\eta q^{1-D}}L,\quad
 \eta_z=\frac{d}{q^DL},\quad
 X_z=-\frac{2\eta X}{q^DL};
\]
substitution in the ordinary chain rule proves both formulas.
A second axial derivative changes the exponent by $-D$,
which proves the stated composition order for $Z^{[2]}$.

At each nonnegative order the physical profiles are
\begin{equation}
\begin{aligned}
 u_{\theta,n}
  &=q^{-A+\lambda_n}E_n
    =r q^{-A-1/2+\lambda_n}\phi_n/C,
 &E_n&=\sqrt{2X}\phi_n/C,\\
 u_{z,n}&=q^{-A+\lambda_n}U_n,
 &r u_{r,n}&=q^{\lambda_n}V_n,\\
 p_n&=q^{-2A+\lambda_n}\Pi_n,
 &A_Xf&=\int_0^1 f(b_0X,\eta)\,db_0 .
\end{aligned}
 \label{eq:cq-source-physical-coefficients}
\end{equation}
The letter $b_0$ in the average is an integration variable,
distinct from the exponent $b=-A-1/2$.
The leading profiles have index zero; every negative-index
field is zero. Primes in the following recurrences mean
$\partial_X$ at fixed $\eta$. Incompressibility gives
\begin{equation}
 \partial_XV_n=-Z_{-A,n}U_n,\qquad
 \frac{V_n}{X}=
 \frac{2\eta U_n-2\eta(D+\lambda_n)A_X(U_n)
                      -d\,\partial_\eta A_X(U_n)}L .
 \label{eq:cq-source-full-incompressibility-recurrence}
\end{equation}
The original positive-order equations are, for $n\geq1$,
\begin{align}
 2(X\phi_n''+2\phi_n')
 &=T_{b,n}\phi_n+
   \sum_{\substack{i,j\geq0\\i+j=n}}
    \left\{V_i(\phi_j'+\phi_j/X)+U_iZ_{b,j}\phi_j\right\}
                         -Z^{[2]}_{b,n-1}\phi_{n-1},
 \label{eq:cq-source-full-angular-recurrence}\\
 2(XU_n''+U_n')
 &=T_{c,n}U_n+
   \sum_{\substack{i,j\geq0\\i+j=n}}
           \left\{V_iU_j'+U_iZ_{c,j}U_j\right\}
   +Z_{-2A,n}\Pi_n-Z^{[2]}_{c,n-1}U_{n-1},
 \label{eq:cq-source-full-axial-recurrence}\\
 \Pi_n'
 &=C^{-2}\sum_{\substack{i,j\geq0\\i+j=n}}\phi_i\phi_j
                          -\frac{\Omega_{n-1}}{2X},
 \label{eq:cq-source-full-pressure-recurrence}\\
 \Omega_k
 &=T_{0,k}V_k+
   \sum_{\substack{i,j\geq0\\i+j=k}}
       \left\{V_i(V_j'-V_j/(2X))+U_iZ_{0,j}V_j\right\}
             -2XV_k''-Z^{[2]}_{0,k-1}V_{k-1}.
 \label{eq:cq-source-full-radial-recurrence}
\end{align}
In particular the pressure term is
$-\Omega_{n-1}/(2X)$, with the factor $2$ in its denominator.
These formulas define each lower-order source without leaving
a pressure, radial-viscous or axial-viscous contribution
implicit. All ordered pairs in the sums are retained.

The explicit six-row reduction can also be stated without an
unnamed source matrix. At fixed $n\geq1$ define the three known
lower-order functions
\begin{align*}
 f_n^\phi
 &=\sum_{i=1}^{n-1}
     \{V_i(\phi_{n-i}'+\phi_{n-i}/X)
                         +U_iZ_{b,n-i}\phi_{n-i}\}
                   -Z^{[2]}_{b,n-1}\phi_{n-1},\\
 f_n^U
 &=\sum_{i=1}^{n-1}
     \{V_iU_{n-i}'+U_iZ_{c,n-i}U_{n-i}\}
                   -Z^{[2]}_{c,n-1}U_{n-1},\\
 f_n^\Pi
 &=C^{-2}\sum_{i=1}^{n-1}\phi_i\phi_{n-i}
                                     -\Omega_{n-1}/(2X).
\end{align*}
Empty sums are zero. Put $\xi=\sqrt X$, $v_0=V_0/X$, and
write the six current unknowns as
\[
 W_n=(\phi,U,K,\Pi,p_\xi,q_\xi)^{\mathsf T},\qquad
 p_\xi=\partial_\xi\phi,\quad q_\xi=\partial_\xi U,\quad
 K=A_X(U)-U .
\]
Here $p_\xi,q_\xi$ are labels for the last two components;
they are not the physical pressure or the similarity
coordinate $q$. The current radial-velocity coefficient is
the following explicit linear expression:
\[
 \mathcal V_n[U,K]=
 \frac{2\eta U-2\eta(D+\lambda_n)(U+K)
                       -d\,\partial_\eta(U+K)}L .
\]
Thus $V_n=X\mathcal V_n[U,K]$. With every known coefficient
evaluated at $(X,\eta)=(\xi^2,\eta)$, the first four rows are
\begin{equation}
\begin{aligned}
 \partial_\xi\phi&=p_\xi,&
 \partial_\xi U&=q_\xi,\\
 \partial_\xi K+\frac2\xi K&=-q_\xi,&
 \partial_\xi\Pi
      &=4\xi C^{-2}\phi_0\phi+2\xi f_n^\Pi .
\end{aligned}
 \label{eq:cq-source-explicit-first-four-rows}
\end{equation}
The complete last two rows are
\begin{equation}
\begin{split}
 \partial_\xi p_\xi+\frac3\xi p_\xi
 =2\biggl\{&
 \frac{-(b+\lambda_n)\phi+D\eta\,\partial_\eta\phi
                                     +\xi p_\xi/2}{L}\\
 &+v_0(\xi p_\xi/2+\phi)
 +U_0\frac{2(b+\lambda_n)\eta\phi+d\,\partial_\eta\phi
                                           -\eta\xi p_\xi}{L}\\
 &+\mathcal V_n[U,K](X\phi_0'+\phi_0)
                    +U Z_{b,0}\phi_0+f_n^\phi\biggr\}.
\end{split}
 \label{eq:cq-source-explicit-angular-row}
\end{equation}
\begin{equation}
\begin{split}
 \partial_\xi q_\xi+\frac1\xi q_\xi
 =2\biggl\{&
 \frac{-(c+\lambda_n)U+D\eta\,\partial_\eta U
                                     +\xi q_\xi/2}{L}
                       +v_0\xi q_\xi/2\\
 &+U_0\frac{2(c+\lambda_n)\eta U+d\,\partial_\eta U
                                           -\eta\xi q_\xi}{L}
 +X\mathcal V_n[U,K]U_0'+U Z_{c,0}U_0\\
 &+\frac{2(-2A+\lambda_n)\eta\Pi+d\,\partial_\eta\Pi
              -2\eta X(2C^{-2}\phi_0\phi+f_n^\Pi)}L
                                      +f_n^U\biggr\}.
\end{split}
 \label{eq:cq-source-explicit-axial-row}
\end{equation}
These rows follow directly from
\eqref{eq:cq-source-full-incompressibility-recurrence}--%
\eqref{eq:cq-source-full-radial-recurrence}.
For the radial conversion,
$\partial_X\phi=p_\xi/(2\xi)$ and
$2(X\phi''+2\phi')=(\partial_\xi p_\xi+3p_\xi/\xi)/2$;
the analogous axial identity has $q_\xi/\xi$ instead of
$3p_\xi/\xi$. The apparent divisions in the transport rows
cancel exactly with $V_0=Xv_0$ and
$V_n=X\mathcal V_n[U,K]$, producing the displayed regular
coefficients. In the pressure contribution to the last row,
the substitution is the full equation
$\Pi'=2C^{-2}\phi_0\phi+f_n^\Pi$; hence its lower-order
$f_n^\Pi$ contribution is included.

Equations \eqref{eq:cq-source-explicit-first-four-rows}--%
\eqref{eq:cq-source-explicit-axial-row} define $A_0W_n+
A_1\partial_\eta W_n+f_n$ uniquely by collecting, in this
specified order of the six components, the coefficients of
$W_n$, $\partial_\eta W_n$, and the remaining known terms.
Thus the Picard argument's coefficient functions are
completely specified. Extracting the $\partial_\eta$ terms
from these rows gives exactly the sparse $A_1$ displayed
in that argument: the coefficient of $\partial_\eta\phi$
is $2(D\eta+dU_0)/L$, the $\partial_\eta U$ and
$\partial_\eta K$ coefficients in the angular row are both
$-2d(\phi_0+X\phi_0')/L$, and the axial row gives
$2(D\eta+dU_0-dXU_0')/L$, $-2dXU_0'/L$, and $2d/L$.
This supplies a direct check of every factor in the sparse
matrix and preserves every lower-order forcing term.


Set $\xi=\sqrt X$. Source equations (5.2)--(5.7) give at each
positive order $n$ the linear regular-singular system
\begin{equation}
\begin{split}
 \partial_\xi W_n+\xi^{-1}C_{\rm diag}W_n
       &=A_0(\xi,\eta)W_n+
             A_1(\xi,\eta)\partial_\eta W_n+f_n(\xi,\eta),\\
 C_{\rm diag}&=\operatorname{diag}(0,0,2,0,3,1),\\
 W_n&=(\phi_n,U_n,K_n,\Pi_n,\partial_\xi\phi_n,
                                      \partial_\xi U_n)^{\mathsf T},\\
 K_n&=A_X(U_n)-U_n,\qquad W_n(0,\eta)=0 .
\end{split}
 \label{eq:cq-actual-axis-linear-system}
\end{equation}
The matrices and source here are the unchanged ones from that
recurrence. In particular the only possibly nonzero entries of
$A_1$ are exactly
\begin{align*}
 (A_1)_{51}&=2H_c/L,&
 (A_1)_{52}=(A_1)_{53}
     &=-2d(\phi_0+X\partial_X\phi_0)/L,\\
 (A_1)_{62}&=2(H_c-dX\partial_XU_0)/L,&
 (A_1)_{63}&=-2dX\partial_XU_0/L,&
 (A_1)_{64}&=2d/L,
\end{align*}
where $d=1-\eta^2$, $L=1-2h\eta^2$ and
$H_c=(1/2-h)\eta+dU_0$. Thus $A_1$ and every parameter derivative
map the first four components to the last two and annihilate
the last two.

Assume inductively that all lower-order profiles are analytic near
the axis. This is true at order zero by the just-cited proposition.
The source's radial extensions agree with their inner profiles on
one fixed initial radial interval; their moment corrections have
support to the right of that interval. Thus those extensions do not
change any lower-order germ used here. The equations forming
$A_0,A_1,f_n$ use finitely many sums, products, radial and parameter
derivatives, and the original radial averages, with divisions only
by $L$, nonzero leading factors on the chosen neighborhood, and
the powers of $X$ whose zero factors are explicitly canceled in
the source recurrence. These operations preserve the analytic
germs. For the average the exact formula is
$A_Xg=\int_0^1g(bX,\eta)\,db$. For the apparently divided sources,
$V_j=Xv_j$ by (5.2), and every summand of $\Omega_j$ in (5.6)
is divisible by $X$, including
\[
 V_i(\partial_XV_j-V_j/(2X))
       =Xv_i(v_j/2+X\partial_Xv_j).
\]
Analytic divisibility follows from the vanishing Taylor coefficient,
so no singular coefficient is discarded. Once the displayed
diagonal singular part is separated, $A_0,A_1,f_n$ consequently
extend holomorphically to a product
$|\xi|<R$, $\eta\in\Omega$ about the chosen real axis point.
Their norms are bounded on a smaller closed product. The radii
may depend on $n$; only positivity of each radius is used below.

For $c_i=(C_{\rm diag})_{ii}$ extend the exact source inverse
to complex $\xi$ by
\[
 (Gg)_i(\xi,\eta)=
       \xi\int_0^1b^{c_i}g_i(b\xi,\eta)\,db,\qquad
 K=G(A_0+A_1\partial_\eta).
\]
It is holomorphic on the product, has value zero at $\xi=0$,
and termwise integration of a convergent Taylor series proves
$(\partial_\xi+c_i/\xi)(Gg)_i=g_i$, with its holomorphic
interpretation at zero. Its kernels are diagonal, independent
of $\eta$, and have absolute value at most one along each
straight integration segment.

For every such diagonal kernel $D_0$, at arbitrary radii,
\[
 A_1(\xi)D_0A_1(\xi')=0,\qquad
 A_1(\xi)D_0\partial_\eta A_1(\xi')=0.
\]
Thus two adjacent derivative factors in an expansion of $K^k$
vanish, including the term in which the left derivative strikes
the right matrix. Each remaining word contains at most
$p_k=\lceil k/2\rceil$ parameter derivatives. Let a smaller
parameter neighborhood have distance $\Delta>0$ from the
boundary of the larger one. Split that loss evenly among the
at most $p_k$ derivatives and use Cauchy's integral estimate at
each differentiation. The ordered radial integration simplex
has volume $|\xi|^{k+1}/(k+1)!$; this follows by parametrizing
each nested segment by its ratio to the preceding complex
radius, so the same positive real simplex appears. Summing
the at most $2^k$ words and absorbing the bounded matrix and
source norms gives
\begin{equation}
 \|K^kGf_n\|_{|\xi|\leq R',\,\Omega'}
 \leq C_n^{k+1}\frac{(R')^{k+1}}{(k+1)!}
           \max\{1,p_k/\Delta\}^{p_k},\qquad R'<R .
 \label{eq:cq-actual-complex-axis-Picard-bound}
\end{equation}
The estimate applies to complete derivative operators, so every
Leibniz derivative of a coefficient is included.
Its $k$th root tends to zero: the factorial has root comparable
to $k$, whereas the derivative loss has root at most a constant
times $k^{1/2+1/(2k)}$. Hence
$\sum_{k\geq0}K^kGf_n$ converges normally to a holomorphic
solution of \eqref{eq:cq-actual-axis-linear-system}. Applying
the same estimate to the iterated homogeneous integral equation
proves uniqueness among bounded parameter-holomorphic solutions
on the real restricted radial interval. The source's actual
solution has that parameter regularity by Lemma 5.1; its
restriction therefore equals this holomorphic solution.

The source equations preserve even parity in the first four
components and odd parity in the last two, as stated and proved
by its uniqueness argument on page 49. This also follows by
reflecting the just-constructed integral equation and using the
same uniqueness. The first four convergent complex Taylor
series therefore contain only powers $\xi^{2j}=X^j$.
Their radius in $X$ is the square of a positive radius in
$\xi$, proving analytic $X$ germs for $\phi_n,U_n,K_n,\Pi_n$.
Formula (5.2) and the analytic radial average prove the same
for $V_n/X$. This closes the induction for the actual profiles,
including their full parameter derivatives.

For completeness the locally summed slow field includes the entire
meridional velocity obtained by differentiating the source's base
potential, including every derivative of the summation cutoffs.
The following formulas make this point explicit. Retain
$A=1/2+h$, $D=1/2-h$, $\lambda_n=2nh$, $X=\sigma/q$,
$\sigma=r^2/2$, and define
\[
 \rho_0(q)=1,\qquad \rho_n(q)=\chi(c_nq)\quad(n\geq1),
 \qquad F_n^{\rm src}(X,\eta)=\int_0^X U_n(v,\eta)\,dv .
\]
The superscript distinguishes the source's radial primitive from
the angular force flux $F$. Source (10.1) gives exactly
\[
 S=\sum_{n\geq0}\rho_n(q)q^{1-A+\lambda_n}
                         F_n^{\rm src}(X,\eta),
 \qquad A_{\rm base}=\frac Sr e_\theta .
\]
On each sufficiently small preterminal parameter neighborhood the
sum is finite. Since $F_n^{\rm src}=X A_XU_n$, put
\begin{equation}
\begin{aligned}
 \mathfrak a(\sigma,z,t)
 &=\frac12\sum_{n\geq0}\rho_n(q)q^{-A+\lambda_n}
                     A_XU_n(\sigma/q,\eta),\\
 \mathfrak b(\sigma,z,t)
 &=\frac1C\sum_{n\geq0}\rho_n(q)q^{-A-1/2+\lambda_n}
                          \phi_n(\sigma/q,\eta),\\
 \mathfrak p(\sigma,z,t)
 &=\sum_{n\geq0}\rho_n(q)q^{-2A+\lambda_n}
                          \Pi_n(\sigma/q,\eta).
\end{aligned}
 \label{eq:cq-actual-full-slow-germ-coefficients}
\end{equation}
Every sum here retains order zero. The radial-average formula proves
that $S=2\sigma\mathfrak a$, including its analytic extension to
$\sigma=0$. Consequently the original potential and its complete
curl have the Cartesian representatives
\begin{equation}
\begin{aligned}
 A_{\rm base}&=\mathfrak a(-x_2,x_1,0),\\
 \operatorname{curl} A_{\rm base}
 &=(-x_1\partial_z\mathfrak a,
    -x_2\partial_z\mathfrak a,
       2\mathfrak a+2\sigma\partial_\sigma\mathfrak a).
\end{aligned}
 \label{eq:cq-actual-full-meridional-curl}
\end{equation}
Indeed the first two curl components differentiate the two displayed
potential components with respect to $z$, with their indicated
signs, while the last is
$\partial_{x_1}(x_1\mathfrak a)+
  \partial_{x_2}(x_2\mathfrak a)
 =2\mathfrak a+2\sigma\partial_\sigma\mathfrak a$.
This proves the formula at positive radius and, because its
coefficients are analytic, also at the axis.

In particular $\partial_z\mathfrak a$ differentiates
$\rho_n(q(z,t))$ as well as the profiles. To check its cutoff
contribution against the original construction, note that
$q_z=2\eta q^{1-D}/L$ and $A+D=1$. The source's
incompressibility formula yields
$-Z_{D+\lambda_n}F_n^{\rm src}=V_n$, so
$-\partial_z(q^{1-A+\lambda_n}F_n^{\rm src})
 =q^{\lambda_n}V_n$.
Differentiating the full cutoff-weighted streamfunction therefore
gives, for every $n\geq1$,
\begin{equation}
\begin{aligned}
 r u_{r,n}^{\rm cut}
 &=q^{\lambda_n}\left[
       \chi(c_nq)V_n-
        \frac{2\eta}{L}(c_nq)\chi'(c_nq)F_n^{\rm src}\right],\\
 u_{z,n}^{\rm cut}
 &=q^{-A+\lambda_n}\chi(c_nq)U_n .
\end{aligned}
 \label{eq:cq-actual-exact-streamfunction-cutoff-term}
\end{equation}
For the second equality use
$r^{-1}\partial_r S_n=\partial_\sigma S_n$ and the
independence of $q,\eta,\rho_n$ from $\sigma$.
These are exactly source (5.45), including the negative cutoff
derivative and its full factor $2\eta(c_nq)/L$.

The analytic profile germs already proved, their analytic radial
averages, and local finiteness show that all three functions in
\eqref{eq:cq-actual-full-slow-germ-coefficients} have analytic
$\sigma$ germs. This remains true after every fixed $z,t$
derivative: the cutoffs and their derivatives are smooth functions
of $q(z,t)$ independent of $\sigma$, and all other factors are
derivatives of the already established analytic profile germs.
Source (9.21) and the support statement accompanying (10.1) say
that every added wave and mean representative, including the
finite initialization block, vanishes near the axis. The active
sum is finite on the chosen preterminal neighborhood. Intersecting
their vanishing neighborhoods proves that the actual local field
there consists of the complete curl
\eqref{eq:cq-actual-full-meridional-curl}, the angular component
$r\mathfrak b\,e_\theta$, and the pressure $\mathfrak p$.
Thus the assertion about the actual slow field includes its full
radial and axial components, rather than only its angular component.

Now apply the final localization exactly as in source (10.4).
On the central slab $|z|<z_0/2$ and sufficiently near the axis,
the spatial cutoff equals one on a neighborhood. Put
$\kappa(t)=\chi_t(t)$. There $c(x,t)=\kappa(t)$ and
$\nabla c=0$, whence
\[
 u_*=\operatorname{curl}(\kappa A)+\kappa B_{\rm src}e_\theta
     =\kappa u_{\rm slow},\qquad
 p_*=\kappa\mathfrak p .
\]
Here $B_{\rm src}$ denotes the source's azimuthal velocity
coefficient in (9.21), not our angular-momentum field $B$.
The complete localized velocity has the exact form
\begin{equation}
\begin{gathered}
 u_*=(\alpha_*x_1-\beta_*x_2,
       \alpha_*x_2+\beta_*x_1,\gamma_*),\qquad
 \pi_*=p_*,\\
 \alpha_*=-\kappa\partial_z\mathfrak a,\qquad
 \beta_*=\kappa\mathfrak b,\qquad
 \gamma_*=\kappa(2\mathfrak a+
                            2\sigma\partial_\sigma\mathfrak a),
 \qquad \pi_*=\kappa\mathfrak p .
\end{gathered}
 \label{eq:cq-actual-localized-all-velocity-coefficients}
\end{equation}
The source's support margin in (10.3) places a neighborhood of
the temporal cutoff support within its local coordinate domain;
outside that support the actual localized fields vanish. Thus
these formulas establish the asserted germs for every
$0\leq t<1$ in the central slab, including the temporal transition.

No time derivative of this cutoff is omitted from the force.
For example its exact viscosity-one residual satisfies
\[
 f_* =\kappa f_{\rm slow}
        +\kappa' u_{\rm slow}
        +(\kappa^2-\kappa)
                   (u_{\rm slow}\cdot\nabla)u_{\rm slow},
\]
where $f_{\rm slow}$ is the full local residual of the preceding
velocity and pressure. This follows by substituting
$u_* =\kappa u_{\rm slow}$ and $p_* =\kappa\mathfrak p$
into the original momentum equation; the spatial derivatives of
$\kappa$ vanish, while its time derivative is retained.

The physical viscosity map gives, with
$\widehat\sigma=\sigma/\nu$ and
$\widehat z=z/\sqrt\nu$,
\[
 \alpha_\nu=\alpha_*(\widehat\sigma,\widehat z,t),\qquad
 \beta_\nu=\beta_*(\widehat\sigma,\widehat z,t),\qquad
 \gamma_\nu=\sqrt\nu\,
                  \gamma_*(\widehat\sigma,\widehat z,t),\qquad
 \pi_\nu=\nu\pi_*(\widehat\sigma,\widehat z,t).
\]
For an explicit check of the force, let subscripts $\sigma,z,t$
denote partial derivatives of these complete localized physical
coefficients and suppress their subscript $\nu$ on the right.
Then
\[
 f_\nu=(\alpha_f x_1-\beta_f x_2,
        \alpha_f x_2+\beta_f x_1,\gamma_f),
\]
where the original residual, with physical viscosity $\nu$, gives
\begin{equation}
\begin{aligned}
 \alpha_f={}&\alpha_t+2\sigma\alpha\alpha_\sigma
       +\gamma\alpha_z+\alpha^2-\beta^2
       -\nu(2\sigma\alpha_{\sigma\sigma}
                    +4\alpha_\sigma+\alpha_{zz})+\pi_\sigma,\\
 \beta_f={}&\beta_t+2\sigma\alpha\beta_\sigma
       +\gamma\beta_z+2\alpha\beta
       -\nu(2\sigma\beta_{\sigma\sigma}
                    +4\beta_\sigma+\beta_{zz}),\\
 \gamma_f={}&\gamma_t+2\sigma\alpha\gamma_\sigma
       +\gamma\gamma_z
       -\nu(2\sigma\gamma_{\sigma\sigma}
                    +2\gamma_\sigma+\gamma_{zz})+\pi_z.
\end{aligned}
 \label{eq:cq-actual-complete-analytic-force-coefficients}
\end{equation}
To verify every coefficient, use
$u\cdot\nabla\sigma=2\sigma\alpha$ and
$u\cdot\nabla z=\gamma$; differentiating the two transverse
Cartesian coordinate factors contributes $\alpha^2-\beta^2$
and $2\alpha\beta$. For a scalar radial coefficient $g$,
\[
 \Delta g=2\sigma g_{\sigma\sigma}+2g_\sigma+g_{zz},
 \qquad
 \Delta(x_jg)=x_j(2\sigma g_{\sigma\sigma}
                                  +4g_\sigma+g_{zz})
       \quad(j=1,2).
\]
The pressure gradient is
$(x_1\pi_\sigma,x_2\pi_\sigma,\pi_z)$.
Substitution proves
\eqref{eq:cq-actual-complete-analytic-force-coefficients}
with every sign and factor. All its coefficients have analytic
radial germs because the formulas use the full differentiated
localized coefficients from
\eqref{eq:cq-actual-localized-all-velocity-coefficients}.

Finally the actual axisymmetric central angular projections,
with the original factor $(2\pi)^{-1}$ in the angular average,
are exactly
\[
 B_\nu=2\sigma\beta_\nu,\qquad
 P_\nu=4\sigma^2\alpha_\nu\beta_\nu,\qquad
 Q_\nu=2\sigma\gamma_\nu\beta_\nu,\qquad
 F_\nu=2\sigma\beta_f .
\]
They prove analytic radial germs for all four actual fields on
$|z|<\sqrt\nu\,z_0/2$, $0\leq t<1$, and retain every
meridional and cutoff contribution. The separate exact
terminal-axis coefficient and its fixed spectral-value consequence
require $0<\tau<\tau_1$, when the temporal cutoff equals one;
the all-preterminal analyticity assertion does not enlarge that
range of the exact terminal formulas.



In the retained physical variables the domain just proved is
\begin{equation}
 |z|<\sqrt\nu\,z_0/2,\qquad 0\leq t<1.
 \label{eq:cq-actual-central-analytic-domain}
\end{equation}

Thus the
radial entrance theorem applies to the actual untruncated
four fields on \eqref{eq:cq-actual-central-analytic-domain}:
\[
 \mathcal F_B^{\rm rad},\quad\mathcal F_P^{\rm rad},\quad
 \mathcal F_Q^{\rm rad},\quad\mathcal F_F^{\rm rad}
           \in\mathcal A_{\rm Sch}\subset\Xi\cqE .
\]
The compact radial support was already proved, so the endpoint
analyticity just established is all that was needed for the
order-one sufficient criterion. At other axial points, the
source only prescribes a smooth spatial cutoff; the previously
constructed full-input packet ingress retains those points
without imposing an unsupported analytic cutoff.

\subsection{The full angular equation in the actual entire radial traces}
\label{sec:cq-radial-entire-dynamics}

Retain the original physical variables, viscosity and source exponent,
and set
\begin{equation}
\begin{split}
 K_h(w)&=\zeta(w)2^{w/2-1}
       (1-2^{w/2-1/2})(1-2^{h-w/2}),\\
 \kappa_h(s)&=K_h(1/2-is),\qquad w=1/2-is,\\
 (T_ha)(s)&=\mathcal F_a^{\rm rad}(s)
                  =K_h(w)\widehat a(w/2).
\end{split}
\label{eq:cq-radial-dynamics-multiplier}
\end{equation}
All products use their previously proved entire continuations.
The multiplier itself is entire and has order at most one. Indeed
\[
 \frac{1-2^{(w-1)/2}}{w-1}
   =-\frac{\log2}{2}\int_0^1
                  e^{v(w-1)\log2/2}\,dv
\]
is entire of exponential type. Multiplying it by the already proved
order-one entire function $(w-1)\zeta(w)$ and the two remaining
exponential factors proves the assertion. Moreover
\[
 K_h(1)=-\frac{\log2}{2\sqrt2}(1-2^{h-1/2})\ne0,
\]
since the retained $0<h<1/100$ makes $h-1/2<0$.

\begin{cqproposition}[Actual entire angular dynamics]
\label{prop:cq-radial-entire-dynamics}
On the central domain \eqref{eq:cq-actual-central-analytic-domain},
each of the four actual untruncated profiles
$\mathcal F_B^{\rm rad},\mathcal F_P^{\rm rad},
\mathcal F_Q^{\rm rad},\mathcal F_F^{\rm rad}$ is a smooth
$(\cqE,\gamma)$-valued function of $(z,t)$ before terminal time.
Every mixed physical parameter derivative is the radial trace of
the corresponding mixed derivative of its full input. They satisfy,
at every $s\in\mathbb C$,
\begin{equation}
\begin{split}
 &\kappa_h(s-2i)
  \bigl(\partial_t\mathcal F_B^{\rm rad}(s)
       +\partial_z\mathcal F_Q^{\rm rad}(s)
       -\nu\partial_{zz}\mathcal F_B^{\rm rad}(s)
       -\mathcal F_F^{\rm rad}(s)\bigr)\\
 &\qquad=(w/2-1)\kappa_h(s)
  \bigl(\mathcal F_P^{\rm rad}(s-2i)
                  +\nu w\mathcal F_B^{\rm rad}(s-2i)\bigr).
\end{split}
\label{eq:cq-radial-entire-angular-pde}
\end{equation}
On the exact radial trace image with its recovered smooth inputs
and endpoint orders \eqref{eq:cq-am:axis}, this entire identity is
equivalent to the full physical equation \eqref{eq:cq-am:pde}.
\end{cqproposition}

\begin{proof}
Fix a compact preterminal parameter set inside the stated central
slab, and slightly enlarge it there. The actual source construction
has a positive lower bound for $q$ on this set. Its cutoff sums
(9.21), (10.1) of \cite{cqNS2026} therefore have a common finite
set of active profiles on a neighborhood of this compact set.
The analytic germs proved in Section~\ref{sec:cq-actual-angular-analytic}
give each of those finitely many profiles a holomorphic radial
extension jointly with its similarity parameter on smaller product
neighborhoods. A finite cover of the real parameter set and a
further radial shrink give one complex $\sigma$-disk for all of
them. The change $X=\sigma/(\nu q)$ preserves such a common
disk because $q$ is bounded away from zero. Its parameter changes
are smooth: in viscosity-one variables the determinant of
$(q,\eta)\mapsto(q(1-\eta^2),q^{1/2-h}\eta)$ is
$q^{1/2-h}(1-2h\eta^2)>0$ on the retained domain.
The finite annular terms vanish on a common smaller radial
interval, and all summation cutoffs depend only on the physical
parameters there. Their arbitrary smooth parameter derivatives
multiply the same holomorphic radial germs. The final spatial
cutoff is identically one near this portion of the axis.
Consequently the Cartesian products and derivatives forming the
full $B,P,Q,F$, including the force residual, have one common
holomorphic radial disk with smooth parameter dependence and
mixed parameter derivatives jointly continuous and bounded on the
compact parameter set times every smaller closed disk. Radial
derivatives are taken on a larger disk before this last shrink.
The disk can be chosen once for all derivative orders: differentiation
changes the bounds, while the finite holomorphic profile domains
and the parameter-only cutoff factors stay the same. The fixed
compact radial support supplies the corresponding bounds away
from the axis. At $t=0$ these statements use the smooth one-sided
derivatives of the source field.

Here is the resulting growth-topology estimate, including the
parameter limit it controls. Choose $0<\delta<\rho$ so that
$|\sigma|\leq\rho$ lies inside this common holomorphic disk,
and choose $R>\delta$ containing the fixed radial support. Write
$\widetilde a$ for the holomorphic germ and put
\[
 N(a)=\sup_{\delta\leq\sigma\leq R}|a(\sigma)|
             +\sup_{|\sigma|\leq\rho}|\widetilde a(\sigma)|.
\]
For every parameter derivative of any of the four fields, Cauchy's
coefficient formula gives $|a_n|\leq N(a)\rho^{-n}$.
The exact expansion \eqref{eq:cq-angular-analytic-endpoint-Mellin}
at argument $w/2$ therefore gives, outside the disks
$|w+2n|<1/2$, $n\geq1$,
\[
 |\widehat a(w/2)|\leq C_{\delta,\rho,R}N(a)e^{C_{\delta,R}|w|}.
\]
In detail, the integral on $[\delta,R]$ is bounded by
$(R-\delta)\delta^{-1}N(a)
 e^{|w|\max\{|\log\delta|,|\log R|\}/2}$; the series is
bounded by
$4N(a)e^{|w||\log\delta|/2}
 \sum_{n\geq1}(\delta/\rho)^n$.
For $1<p<2$ choose $1<p_0<p$. The proved order bound on $K_h$
gives $|K_h(w)|\leq C_{p_0} e^{|w|^{p_0}}$. The entire product
$K_h(w)\widehat a(w/2)$ then has the preceding bound multiplied
by $C_{p_0} e^{|w|^{p_0}}$ off the disks. Maximum modulus on each disk
extends it inside with $|w|$ replaced by $|w|+1$.
Since $|w|\leq|s|+1/2$ and $p_0<p$, for every $c>0$ this proves
\begin{equation}
 p_{p,c}(T_ha):=\sup_{s\in\mathbb C}
         |T_ha(s)|e^{-c|s|^p}
       \leq C_{\delta,\rho,R,p,c}N(a).
\label{eq:cq-radial-parameter-growth-bound}
\end{equation}
Every difference and parameter derivative considered here has
the same support and holomorphic germ domain, so the same estimate
applies to it. For a coordinate direction $e_j$ in $(z,t)$,
the integral Taylor formula in the two sup norms gives
\[
 N\!\left(\frac{a(y+\epsilon e_j)-a(y)}{\epsilon}
                        -\partial_j a(y)\right)
 \leq\frac{|\epsilon|}{2}
       \sup_{0\leq v\leq1}N(\partial_j^2a(y+v\epsilon e_j)).
\]
The bound tends to zero. Applying
\eqref{eq:cq-radial-parameter-growth-bound} proves differentiation
in every defining $\gamma$ seminorm. The first-order integral
formula applied to each $\partial^\alpha a$ also proves continuity
of its trace in those seminorms. Repeating this argument proves
all mixed derivatives and the asserted smoothness.

Put $v=w/2$. On $\operatorname{Re}w>2$ the Mellin integrations
by parts already justified in \eqref{eq:cq-am:mellin} give
\[
 \partial_t\widehat B(v)+\partial_z\widehat Q(v)
       -\nu\partial_{zz}\widehat B(v)-\widehat F(v)
  =(v-1)\{\widehat P(v-1)+2\nu v\widehat B(v-1)\}.
\]
The parameter differentiations commute with $T_h$ by the preceding
proof. The exact spectral change is
$1/2-i(s-2i)=w-2$, so the two shifted traces are
$K_h(w-2)\widehat P(v-1)$ and
$K_h(w-2)\widehat B(v-1)$. Multiplying the last displayed
equation by $K_h(w)K_h(w-2)$, and retaining $2\nu v=\nu w$,
proves \eqref{eq:cq-radial-entire-angular-pde} on that half-plane.
All its terms are entire, hence the identity theorem proves it
at every complex $s$. They also lie in $\cqE$: finite products,
polynomial factors and a fixed translation preserve order at
most one. No radial or nonlinear flux has been removed.

For the claimed exact converse retain the full inverse domain:
each of the four entire functions is the radial trace of its
specified smooth compact input, recovered by
\eqref{eq:cq-angular-radial-exact-inverse}, with the common local
germ and parameter regularity just proved. Define its full
physical residual
\[
 E=\partial_tB+\partial_\sigma P+\partial_zQ
                  -\nu(2\sigma B_{\sigma\sigma}+B_{zz})-F.
\]
The endpoint orders make $E(0)=0$; it has the same compact support
and an analytic endpoint germ. Repeating the preceding integrations
without setting $E$ to zero shows that the left side minus the
right side of \eqref{eq:cq-radial-entire-angular-pde} is exactly
\begin{equation}
 \kappa_h(s-2i)\mathcal F_E^{\rm rad}(s).
\label{eq:cq-radial-entire-residual-defect}
\end{equation}
If the entire identity holds, this product vanishes identically.
The multiplier is nonzero on an open set because $K_h(1)\ne0$,
so the identity theorem gives $\mathcal F_E^{\rm rad}=0$.
The proved injective radial inverse gives $E=0$. Conversely,
$E=0$ makes the product zero. Thus even the isolated zeros of the
multiplier lose no information on this specified full entire
image. This conclusion uses the full functions and their inverse
domain, rather than values at a single multiplier zero.
\end{proof}

The codomain of this proved smooth physical evolution is
$(\cqE,\gamma)$, with every individual trace in
$\mathcal A_{\rm Sch}\subset\Xi\cqE$. Its subsequent map to
the original source quotient is the already specified
$q_{\rm rad}(a)=[\mathcal F_a^{\rm rad}]_{\cqQ}$, with exact
kernel \eqref{eq:cq-angular-source-kernel}. That map is continuous
for the previously defined initial topology involving
$\tau\vee\gamma$. The new seminorm estimate proves no additional
continuity into the unnamed source topology $\tau$. Nor does the
entire equation define its multipliers or spectral translation on
source classes without invariance of the actual relation space
$\cqN$. The exact relation furnished here is therefore the full
smooth entire evolution, its injective radial inverse, and the
specified projection and kernel to the source quotient; the
frozen differential and heat descent questions remain those of
the original $\cqN$.


\subsection{A fixed entire spectral value of the actual concentrating field}

Throughout the exact terminal identities below, $0<\tau<\tau_1$, with the source-determined $\tau_1>0$ from Section~\ref{sec:cq-ns-axis-arithmetic-jet}; the final central cutoffs equal one on this range. At $z=0$, $q=\tau$, $\eta=0$ in the original source coordinates.
The source's equation (B.3) gives $\phi_0(0,0)=1$; Lemma 5.1
gives $\phi_n(0,0)=0$ for every $n\geq1$.
Its locally finite sum, the vanishing of all annular corrections
near the axis, and the final cutoffs equal to one near the
specified terminal central point yield the exact axial identity
\begin{equation}
 \frac{u_{\nu,\theta}}r\bigg|_{\sigma=0,z=0,t=1-\tau}
          =\frac{\tau^{-1-h}}C,\qquad
 \partial_\sigma B(0,0,1-\tau)=\frac2C\tau^{-1-h}.
 \label{eq:cq-actual-axis-angular-coefficient}
\end{equation}
The first coefficient has the original amplitude $C>1$ from
the source and the original power
$q^{-A-1/2}=q^{-1-h}$. The second identity follows from
$B=r u_\theta=2\sigma(u_\theta/r)$.
Under viscosity scaling $u_\nu=\sqrt\nu\,u_*(x/\sqrt\nu,t)$,
both the numerator $u_\theta$ and radius $r$ have the same
factor $\sqrt\nu$, which cancels in this ratio. Thus no
viscosity factor is absent from the displayed coefficient.

For the direct arithmetic Mellin profile this gives the genuine
uncanceled pole
\[
 \operatorname{Res}_{w=-1}W_h(w)\widehat B(w,0,1-\tau)
       =\frac{2^{h+1}-1}{8C}\tau^{-1-h}.
\]
Here $W_h(-1)=(2^{h+1}-1)/16$ follows from
$\zeta(-1)=-1/12$ and the two retained dyadic factors.
The pole is compatible with the entire radial ingress because
the two precise maps have different Mellin arguments and
different moment-canceling dilation coefficients.

For the latter, set the fixed original spectral point
$s_0=-5i/2$, so that $w=1/2-is_0=-2$.
Formula \eqref{eq:cq-angular-radial-Taylor-data} gives
\[
 \operatorname{Res}_{w=-2}M_B^{\rm rad}(w,0,1-\tau)
   =-\frac{(1-2^{-3/2})(2^{h+1}-1)}C\,\tau^{-1-h}.
\]
The completed zeta function determines the exact canceled
zero coefficient:
\[
 \xi(-2)=\xi(3)=\frac{3\zeta(3)}{2\pi},\qquad
 \Gamma(w/2)=-\frac2{w+2}+O(1),\qquad
 \zeta'(-2)=-\frac{\zeta(3)}{4\pi^2}.
\]
For the first equality evaluate the retained completion at
$3$, using $\Gamma(3/2)=\sqrt\pi/2$, and use the functional
equation $\xi(-2)=\xi(3)$. For the second, the gamma recurrence
gives the residue $-1$ at its argument $-1$, and the change
of argument $w/2$ supplies the factor two. Multiplication by
$w(w-1)\pi^{-w/2}/2$, whose value at $-2$ is $3\pi$, shows
$\xi(-2)=-6\pi\zeta'(-2)$, proving the third equality.

Multiplying the simple zeta zero by the exact Mellin pole
therefore gives the actual entire trace value
\begin{equation}
 \mathcal F_B^{\rm rad}(s_0,0,1-\tau)
 =\frac{\zeta(3)}{4\pi^2C}
       (1-2^{-3/2})(2^{h+1}-1)\tau^{-1-h}.
 \label{eq:cq-actual-entire-trace-fixed-value}
\end{equation}
The function here is the untruncated actual radial source trace
already proved to lie in $\Xi\cqE$ for each preterminal time.
Writing its unique entire quotient as
$\mathcal F_B^{\rm rad}=\Xi g_B^{\rm rad}$ and retaining
$\Xi(s_0)=\xi(3)=3\zeta(3)/(2\pi)$ gives
\begin{equation}
 g_B^{\rm rad}(s_0,0,1-\tau)
 =\frac{(1-2^{-3/2})(2^{h+1}-1)}{6\pi C}\tau^{-1-h}.
 \label{eq:cq-actual-entire-quotient-fixed-value}
\end{equation}
All factors are strictly positive, so both displayed values
diverge with exactly the source exponent. This is a
fixed-spectral-point consequence for actual order-one entire
functions, not an uninverted trace claim at a moving sample.

Its precise relation to the source quotient retains the original
relation space. Since $\Xi(s_0)\ne0$, the function
$e(s)=\Xi(s)/\Xi(s_0)$ belongs to
$\mathbb C[s]\Xi\subseteq\cqN$ and satisfies $e(s_0)=1$.
For any representative $f$ and any prescribed scalar $b$,
$f+(b-f(s_0))e$ has the same source quotient class and
value $b$ at $s_0$. Thus the representative relation
$\{([f]_{\cqQ},f(s_0)):f\in\cqE\}$ has full scalar fibres.
This proves the exact information lost by imposing the source
quotient on the now-constructed growing trace value.
The full trace itself and its inverse radial input remain
specified by the preceding injective maps and kernels;
the calculation does not decide the distinct membership of
$\Xi'$ in $\cqN$.

\section{Original-space Fr\'echet rigidity and the ring obstruction}
\label{sec:cq-original-ring-rigidity}

We retain the original ring $\cqE=\mathcal H_{\leq1}$, ordinary
pointwise multiplication, its coordinate $s$, and
$Mf=sf$. The source relation remains
$\cqN=\overline{\mathbb C[s]\Xi}^{\,\tau}$ and
$\cqQ=\cqE/\cqN$.
The word ``topological ring'' in
\cite[\S3.6, p.19, typeset line 552 and Proposition~3.6]{cqCCM}
does not itself assert completeness or metrizability.
The cited \cite[Proposition~2.24 and the following discussion,
pp.379--381]{cqCM} computes the finite Hermite image and discusses
the choice of function spaces. Those inspected clauses do not
state that the Hadamard topology is Fr\'echet or that polynomials
are dense in it. The theorem below classifies a precisely specified
class of topologies on the original space; it does not insert
Fr\'echet structure into the primary source.

\subsection{The exact division map before an open-mapping argument}

\begin{cqlemma}[Original division and closed evaluation hyperplanes]
\label{thm:cq-ring-division-map}
For each $\lambda\in\mathbb C$, the map
\begin{equation}
 \mathsf L_\lambda:\cqE\times\mathbb C\longrightarrow\cqE,
 \qquad \mathsf L_\lambda(g,c)=(s-\lambda)g+c
 \label{eq:cq-ring-division-map}
\end{equation}
is a linear bijection, with inverse
\begin{equation}
 \mathsf L_\lambda^{-1}f=(\mathsf A_\lambda f,f(\lambda)),
 \qquad
 \mathsf A_\lambda f(s)=\frac{f(s)-f(\lambda)}{s-\lambda}.
 \label{eq:cq-ring-divided-difference}
\end{equation}
In any complex topological-vector-space topology $\tau$ with
continuous $M$, $\mathsf L_\lambda$ is continuous. For a Hausdorff
such topology, writing
$\mathfrak m_\lambda=(s-\lambda)\cqE$, one has
\begin{equation}
 \operatorname{ev}_\lambda\text{ is continuous}
 \quad\Longleftrightarrow\quad
 \mathfrak m_\lambda\text{ is closed}.
 \label{eq:cq-ring-evaluation-closed}
\end{equation}
\end{cqlemma}
\begin{proof}
The singularity in \eqref{eq:cq-ring-divided-difference} is
removable by the Taylor expansion at the retained point $\lambda$.
For $|s|\geq2(|\lambda|+1)$ one has $|s-\lambda|\geq|s|/2\geq1$,
so
$|\mathsf A_\lambda f(s)|\leq|f(s)|+|f(\lambda)|$.
The remaining compact disk contributes a finite maximum. The
filled quotient therefore belongs to the same order class $\cqE$.
Evaluation of $\mathsf L_\lambda(g,c)=0$ at $\lambda$ gives $c=0$;
then $(s-\lambda)g=0$ gives $g=0$. Substitution verifies both
inverse identities. Continuity of $\mathsf L_\lambda$ follows
from continuity of $M$, scalar multiplication, addition, and
the constant embedding $c\mapsto c1$.

The kernel of evaluation is exactly $\mathfrak m_\lambda$,
which proves the forward implication of
\eqref{eq:cq-ring-evaluation-closed}. If this hyperplane is
closed, $\cqE/\mathfrak m_\lambda$ is a one-dimensional
Hausdorff complex topological vector space. The map
$\mathbb C\to\cqE/\mathfrak m_\lambda$, $c\mapsto[c1]$,
is a continuous linear bijection. Its inverse is continuous:
the image of the unit circle is compact and excludes zero;
choose a balanced zero neighbourhood avoiding that compact
set. It contains no $c[1]$ with $|c|\geq1$, since balance
would put $(c/|c|)[1]$ in it. This proves inverse continuity
at zero, and scalar rescaling and translation prove it everywhere.
Composing with the quotient map gives continuous evaluation.
\end{proof}

The lemma does not infer continuity of the inverse of
\eqref{eq:cq-ring-division-map} from its bijectivity.
Even when evaluation is continuous, the continuous bijection
$\cqE\to\mathfrak m_\lambda$, $g\mapsto(s-\lambda)g$,
has a further inverse-continuity question. The next argument
settles both inverse components in its stated class.

\subsection{The original growth space and the open-mapping step}

\begin{cqlemma}[Complete growth topology on the unchanged order class]
\label{thm:cq-ring-gamma-Frechet}
Retain every original seminorm
\[
 p_{p,a}(f)=\|f\|_{p,a}
   =\sup_{s\in\mathbb C}|f(s)|e^{-a|s|^p},
 \qquad 1<p<2,\quad a>0.
\]
These seminorms define a Hausdorff Fr\'echet topology $\gamma$
on exactly $\cqE$. The countable subfamily
\begin{equation}
 q_n=p_{1+1/n,1},\qquad n=2,3,\ldots,
 \label{eq:cq-ring-countable-growth}
\end{equation}
is cofinal for this topology. Ordinary multiplication and $M$
are continuous.
\end{cqlemma}
\begin{proof}
For $f\in\cqE$ and $1<r<p$, the order bound gives
$|f(s)|\leq C\exp(|s|^r)$ for a finite $C$.
Since $\sup_{t\geq0}\exp(t^r-at^p)<\infty$, every displayed
seminorm is finite. Conversely, finiteness of all these
seminorms gives
$\log M_f(R)\leq\log p_{p,a}(f)+aR^p$ for every $p>1$
in the displayed range, hence order at most one. The zero
function and constants satisfy the same membership conclusion.

For a given $p,a$, choose $n$ with $r_n=1+1/n<p$. Then
\begin{equation}
 \begin{aligned}
 p_{p,a}(f)&\leq C_{p,a,n}q_n(f),\\
 C_{p,a,n}&=\exp\left(\sup_{t\geq0}(t^{r_n}-at^p)\right)<\infty.
 \end{aligned}
 \label{eq:cq-ring-cofinal-bound}
\end{equation}
Every $q_n$ is already one of the original seminorms, proving
cofinality without replacing their definition. They separate
entire functions, giving a Hausdorff metrizable locally convex
space.

Let $(f_k)$ be Cauchy in all $q_n$. On each disk, a fixed
$q_n$ bounds the compact supremum by a finite factor, so
$f_k$ converges uniformly on compact disks to an entire $f$.
For fixed $n$ and $\epsilon>0$, choose $k_0$ such that
$q_n(f_k-f_l)\leq\epsilon$ for $k,l\geq k_0$.
Sending $l\to\infty$ pointwise and taking the supremum gives
$q_n(f_k-f)\leq\epsilon$. Thus all $q_n(f)$ are finite;
\eqref{eq:cq-ring-cofinal-bound} and the order calculation
put $f$ in the original $\cqE$. This proves completeness
and convergence in the entire original seminorm family.
Finally
\[
 \begin{aligned}
 p_{p,a}(fg)&\leq p_{p,a/2}(f)p_{p,a/2}(g),\\
 p_{p,a}(sf)&\leq
 \left(\sup_{t\geq0}t e^{-at^p/2}\right)p_{p,a/2}(f)
 \end{aligned}
\]
prove the two continuity assertions.
\end{proof}

\begin{cqlemma}[The open-mapping step used here]
\label{thm:cq-ring-open-mapping}
A continuous linear surjection between Fr\'echet spaces is
open. In particular a continuous linear bijection between
such spaces is a homeomorphism.
\end{cqlemma}
\begin{proof}
Let $T:F\to G$ be the surjection and let $U$ be an arbitrary
zero neighbourhood in $F$. Using countably many defining
seminorms, choose balanced convex zero neighbourhoods $U_n$
such that every series $x_n\in U_n$ converges in $F$ and
has sum in $U$. Explicitly, bound the first $n$ seminorms
by $2^{-n}$, and at every stage impose the finitely many
bounds of a smaller neighbourhood inside $U$ with summable
coefficients. Each seminorm then has summable tails, and
completeness gives convergence; the finite bounds keep the
sum inside $U$.

Each $U_n$ is absorbing, so
$G=\bigcup_{k\geq1}kT(U_n)$. The complete-metric Baire
theorem implies that $\overline{T(U_n)}$ has nonempty
interior. Balance and convexity put a zero neighbourhood
inside it: average a neighbourhood about an interior point
with its reflected neighbourhood. Choose zero neighbourhoods
$V_n\subseteq\overline{T(U_n)}$ whose diameters tend to
zero in a compatible metric of $G$.
Given $y\in V_1$, choose $x_1\in U_1$ with
$y-Tx_1\in V_2$. Recursively choose $x_n\in U_n$ so that
the remainder after $n$ terms lies in $V_{n+1}$; membership
in the preceding closure gives each choice. The series
converges to $x\in U$, its image converges to $y$, and
continuity gives $Tx=y$. Thus $V_1\subseteq T(U)$, proving
openness. The Baire property used here follows by the
nested-closed-ball proof in a complete compatible metric.
\end{proof}

\begin{cqtheorem}[Original-space Fr\'echet rigidity]
\label{thm:cq-ring-Frechet-rigidity}
On the original vector space $\cqE$, there is exactly one
Hausdorff Fr\'echet complex-vector-space topology in which
$Mf=sf$ is continuous: it is $\gamma$.
In particular the same uniqueness holds for Hausdorff
Fr\'echet topological-algebra structures with the original
ordinary multiplication and usual complex scalars.
\end{cqtheorem}
\begin{proof}
Let $\tau$ belong to the stated class. The map
$\mathsf L_\lambda$ in \eqref{eq:cq-ring-division-map} is a
continuous linear bijection between the Fr\'echet spaces
$(\cqE,\tau)\times\mathbb C$ and $(\cqE,\tau)$.
Lemma~\ref{thm:cq-ring-open-mapping} makes it a homeomorphism.
Thus evaluation at every $\lambda$, and the divided
difference $\mathsf A_\lambda$, are $\tau$-continuous.

Fix $p,a$ from the unchanged growth family. Its seminorm
$p_{p,a}$ is a supremum of the continuous functions
$f\mapsto|f(s)|e^{-a|s|^p}$, hence is lower semicontinuous
for $\tau$. It is finite everywhere by
Lemma~\ref{thm:cq-ring-gamma-Frechet}. The closed sets
$F_m=\{f:p_{p,a}(f)\leq m\}$, $m\geq1$, cover $\cqE$.
Baire supplies $m,f_0$ and a zero neighbourhood $V$ with
$f_0+V\subseteq F_m$. Taking $0\in V$ gives $f_0\in F_m$,
so for $v\in V$,
\[
 p_{p,a}(v)\leq p_{p,a}(f_0+v)+p_{p,a}(f_0)\leq2m.
\]
Scaling $V$ proves continuity of $p_{p,a}$. All growth
seminorms are therefore $\tau$-continuous, and
$\mathrm{id}:(\cqE,\tau)\to(\cqE,\gamma)$ is a continuous
linear bijection. Both spaces are Fr\'echet, the latter by
Lemma~\ref{thm:cq-ring-gamma-Frechet}; the open-mapping
lemma proves inverse continuity. Hence $\tau=\gamma$.
Conversely $\gamma$ belongs to the stated class by that
same lemma, completing existence and uniqueness.
\end{proof}

No family of seminorms for $\tau$, evaluation continuity,
polynomial density, or principal-ideal description was
assumed in this theorem. They are consequences within its
explicit Fr\'echet domain. The inspected primary topology
has not yet been proved to lie in that domain.

\subsection{The complete closure and blocks in the classified domain}

\begin{cqproposition}[Closure, generalized blocks, and frozen failure]
\label{thm:cq-ring-classified-closure}
For every topology classified in
Theorem~\ref{thm:cq-ring-Frechet-rigidity}, its original
source-form closure is
\begin{equation}
 \begin{aligned}
 \overline{\mathbb C[s]\Xi}^{\,\tau}
 &=\overline{\mathbb C[s]\Xi}^{\,\gamma}=I:=\Xi\cqE,\\
 I&=\left\{f\in\cqE:\begin{gathered}
 f^{(j)}(\lambda)=0\text{ for every }\lambda\in\mathbb C
                      \text{ with }\Xi(\lambda)=0,\\
 \text{and every }j\in\mathbb Z_{\geq0}\text{ with }0\leq j<m_\lambda
 \end{gathered}\right\}.
 \end{aligned}
 \label{eq:cq-ring-classified-closure}
\end{equation}
At every actual zero $\lambda$ of multiplicity $m$, put
\[
 v_{\lambda,j}=\left[\frac{\Xi(s)}{(s-\lambda)^j}\right]_I,
 \qquad 1\leq j\leq m.
\]
For the original multiplication $S$ on $\cqE/I$ and every integer $k\geq1$,
\begin{equation}
 \begin{aligned}
 \ker(S-\lambda)^k
 &=\operatorname{span}\{v_{\lambda,1},\ldots,
                  v_{\lambda,\min(k,m)}\},\\
 \dim\ker(S-\lambda)^k&=\min(k,m).
 \end{aligned}
 \label{eq:cq-ring-full-generalized-block}
\end{equation}
In this classified domain neither $D$ nor $D^2$ descends
to the frozen quotient by all representatives, and
$U_tI\nsubseteq I$ for every $t\ne0$.
\end{cqproposition}
\begin{proof}
The topology equality has just been proved.
Lemma~\ref{thm:cq-density} proves Taylor-polynomial density
in $\gamma$, and Lemma~\ref{thm:cq-division} proves entire
division with the retained estimate
\[
 \|f/g\|_{p,a}\leq
 \left(\frac{\max(1,\|g\|_{p,b})}{|g(0)|}\right)^3
 \|f\|_{p,b},\qquad b=\frac{a}{6\cdot2^p}.
\]
The original center applies since $\Xi(0)>0$.
The complete closure calculation in
Proposition~\ref{thm:cq-equalclosures} therefore gives
\eqref{eq:cq-ring-classified-closure}; it uses closed
vanishing jets for one containment, entire division for
the quotient's order, and Taylor approximation followed
by multiplication for the reverse containment. The
principal-ideal equality is thus a calculated closure.

All $v_{\lambda,j}$ belong to the original order class.
Writing $T=S-\lambda$, one has $Tv_{\lambda,1}=0$ and
$Tv_{\lambda,j}=v_{\lambda,j-1}$ for $j>1$.
A relation among them would make a finite sum of the
principal parts $(s-\lambda)^{-j}$ entire after division
by $\Xi$, forcing every coefficient to vanish.
If $T^k[f]=0$, then $(s-\lambda)^kf=\Xi h$ for some
$h\in\cqE$. For $k\leq m$, subtract the Taylor polynomial
of $h$ through degree $k-1$ at $\lambda$. Division of the
remainder by $(s-\lambda)^k$ stays in $\cqE$ by iterating
Lemma~\ref{thm:cq-ring-division-map}. Modulo $I$, this
expresses $[f]$ in the span of the first $k$ displayed
vectors. For $k>m$, entire regularity of $f$ forces $h$
to vanish to order at least $k-m$ at $\lambda$. Divide by
that exact factor and reduce to $k=m$. This proves the
full equality and dimension in
\eqref{eq:cq-ring-full-generalized-block}.

Finally Theorem~\ref{thm:cq-obstruction} proves the direct
frozen failures for this calculated $I$. Its witnesses
retain every multiplicity: $\Xi'$ loses one order at
an actual zero; for $D^2$ use $\Xi''$ when $m\geq2$,
and $((s-\lambda)\Xi)''(\lambda)=2\Xi'(\lambda)\ne0$
when $m=1$. For heat the complete identity
\eqref{eq:cq-exponential} gives
$U_t(e^{as}\Xi)(s)=e^{as-ta^2/4}Z_t(s-ta/2)$.
Evaluation at an actual zero, under a proposed
$U_tI\subseteq I$, would make nonzero $Z_t$ vanish
at every $\lambda-ta/2$, $a\in\mathbb C$. For $t\ne0$
these fill $\mathbb C$, proving the contradiction.
The full representative-ambiguity calculation in
\eqref{eq:cq-relation} gives the asserted all-representative
domain failure.
\end{proof}

\subsection{The exact multiplier and annihilator maps without completeness}

Return now to a general Hausdorff complex topological vector
algebra $(\cqE,\tau)$ with the original multiplication, and
keep its actual closed $\cqN$. Define
\begin{equation}
 \begin{aligned}
 R_\tau&=\overline{\mathbb C[s]}^{\,\tau},\\
 \operatorname{Mult}(\cqN)&=\{a\in\cqE:a\cqN\subseteq\cqN\},\\
 J_\tau&=\{a\in\cqE:a\cqE\subseteq\cqN\}.
 \end{aligned}
 \label{eq:cq-ring-multiplier-annihilator}
\end{equation}

\begin{cqproposition}[The original-ring representation]
\label{thm:cq-ring-multiplier-representation}
$R_\tau$ and $\operatorname{Mult}(\cqN)$ are closed unital
subalgebras, and
$R_\tau\subseteq\operatorname{Mult}(\cqN)$.
The space $J_\tau$ is the largest $\cqE$-ideal contained
in $\cqN$, and it is closed. There is an exact faithful
algebra representation
\begin{equation}
 \operatorname{Mult}(\cqN)/J_\tau
 \longrightarrow\operatorname{End}_{\rm cont}(\cqQ),
 \qquad [a]\longmapsto\rho(a),\quad
 \rho(a)[f]=[af],\qquad S=\rho(s).
 \label{eq:cq-ring-quotient-representation}
\end{equation}
If $R_\tau=\cqE$, then $\cqN$ is an $\cqE$-ideal and
$\cqN=\overline{\Xi\cqE}^{\,\tau}$.
\end{cqproposition}
\begin{proof}
For every polynomial $p$, continuous multiplication by
$p$ preserves $\mathbb C[s]\Xi$, and hence its closure
$\cqN$. This proves the original polynomial-module
property. Joint multiplication of convergent polynomial
nets proves that $R_\tau$ is a subalgebra; the corresponding
addition and scalar maps prove the remaining algebra
operations, and its definition gives closedness.

The set $\operatorname{Mult}(\cqN)$ is a unital subalgebra
by its definition. It is closed because
\[
 \operatorname{Mult}(\cqN)
 =\bigcap_{n\in\cqN}\{a:an\in\cqN\},
\]
an intersection of continuous inverse images of the closed
space $\cqN$. It contains all polynomials and therefore
their closure $R_\tau$. Similarly $J_\tau$ is closed by
its intersection-of-inverse-images definition.
If $a\in J_\tau$ and $b,f\in\cqE$, then
$(ba)f=a(bf)\in\cqN$, so it is an ideal. Setting $f=1$
shows $J_\tau\subseteq\cqN$. Every ideal contained in
$\cqN$ has this multiplication property and is therefore
contained in $J_\tau$.

For $a\in\operatorname{Mult}(\cqN)$ the map $[f]\mapsto[af]$
is representative-independent. Its composite with the
quotient map is continuous, so the quotient universal
property proves its continuity. Its vanishing on all
classes is exactly $a\cqE\subseteq\cqN$. Thus its kernel
is $J_\tau$, giving the faithful representation
\eqref{eq:cq-ring-quotient-representation}. No topology
on the algebra of endomorphisms is needed for this
algebra representation.

If $R_\tau=\cqE$, take polynomials $p_i\to a$.
For $n\in\cqN$, $p_in\in\cqN$ and continuity gives
$an\in\cqN$, so $\cqN$ is an ideal. In particular
$\Xi\cqE\subseteq\cqN$. The other containment after
closure follows from
$\mathbb C[s]\Xi\subseteq\Xi\cqE$.
\end{proof}

Thus ordinary topological-ring structure already supplies
\eqref{eq:cq-ring-multiplier-annihilator} and
\eqref{eq:cq-ring-quotient-representation}. It does not
justify replacing $R_\tau$ by $\cqE$ without a density
proof. The abstract Weyl model of
Section~\ref{sec:cq-source-spectral-strength} was a model
of vector-space and operator data. It did not construct
the original ring, this closure, or these multiplier maps;
it cannot refute additional deductions from them.

\subsection{The closed-ideal spectral obstruction}

\begin{cqlemma}[Surjectivity gives an explicit continuous ring inverse]
\label{thm:cq-ring-surjective-inverse}
Let $\cqN$ be an actual closed $\cqE$-ideal in the preceding
topological algebra. If $S-\lambda$ on $\cqQ=\cqE/\cqN$
is algebraically onto, it has a continuous inverse.
More precisely, a preimage $b\in\cqQ$ of $[1]$ gives
$(S-\lambda)^{-1}[f]=b[f]$ on all of $\cqQ$.
\end{cqlemma}
\begin{proof}
The quotient map is open: the preimage of the image of an
open set $U$ is $U+\cqN$, a union of translates of $U$.
Its product with itself is therefore an open quotient
map, so ordinary multiplication induces continuous
multiplication on the quotient algebra. Surjectivity gives
$b$ with $([s]-\lambda[1])b=[1]$. Multiplication by $b$
is continuous and, by commutativity, is both a left and
right inverse of $S-\lambda$. No operator-series inverse
or open-unit-group hypothesis enters this proof.
\end{proof}

\begin{cqproposition}[Nonempty source spectrum excludes frozen ideal descent]
\label{thm:cq-ring-ideal-spectral-obstruction}
Among the original source-form closed $\cqE$-ideals whose
multiplication operator has the nonempty spectrum stated
in \cite[Proposition~3.6(ii)]{cqCCM}, one has
\begin{equation}
 D\cqN\nsubseteq\cqN,\qquad
 D^2\cqN\nsubseteq\cqN,\qquad
 U_t\cqN\nsubseteq\cqN\quad(t\ne0).
 \label{eq:cq-ring-ideal-obstruction}
\end{equation}
No finite-multiplicity assertion is required for this conclusion.
\end{cqproposition}
\begin{proof}
The original calculations
\eqref{eq:cq-source-D-equivalence}--\eqref{eq:cq-source-heatconsequence}
prove that either derivative invariance or a single
nonzero one-sided heat inclusion makes every $S-\lambda$
algebraically onto. Their exact mechanism is retained:
derivative invariance puts every $D^j\Xi$ in $\cqN$;
the heat commutator polynomials have recurrence
\[
 P_0=1,\qquad P_{k+1}(X)=\frac t2XP_k(X)-P_k'(X),
\]
with leading coefficient $(t/2)^k$, and put every
$D^kU_t\cqN$ in $\cqN$. At each $\lambda$ at least one
derivative of the relevant nonzero entire function
$\Xi$ or $Z_t$ is nonzero. The actual division formula
\eqref{eq:cq-source-division} then proves surjectivity.

Lemma~\ref{thm:cq-ring-surjective-inverse} would make every
$S-\lambda$ continuously invertible on this quotient
algebra, giving empty spectrum and contradicting the
stated nonempty arithmetic spectrum. This excludes
derivative and heat inclusion. The exact identity
$D^2(sn)-sD^2n=2Dn$, with $s\cqN\subseteq\cqN$,
gives the equivalence of the first and second derivative
invariance questions and excludes the second inclusion.
\end{proof}

There is a direct value form of the same ring obstruction.
If $n\in\cqN$ has $n(\lambda)\ne0$ and $\cqN$ is an
ideal, write
\[
 1=\frac{n(s)}{n(\lambda)}+(s-\lambda)b(s),\qquad
 b(s)=\frac{1-n(s)/n(\lambda)}{s-\lambda}\in\cqE.
\]
Then multiplication by $[b]$ is the continuous inverse of
$S-\lambda$. Thus every spectral point forces every
relation element to vanish there. This conclusion is
algebraic and uses no assumed evaluation continuity.
The actual ideal property of the unnamed primary closure
has not been established merely by proving this obstruction.

\subsection{Why unit openness cannot be added to the original ring}

\begin{cqproposition}[Nonopenness of the original unit set]
\label{thm:cq-ring-unit-nonopenness}
For any complex topological-vector-algebra topology on
the original $\cqE$, its group of units $\cqE^\times$
is not open. In fact
\begin{equation}
 1\in\overline{\cqE\setminus\cqE^\times}^{\,\tau}.
 \label{eq:cq-ring-unit-nonopenness}
\end{equation}
\end{cqproposition}
\begin{proof}
The map $\epsilon\mapsto1-\epsilon s$ is continuous and
tends to $1$ as $\epsilon\to0$ by the usual complex
scalar multiplication. For every $\epsilon\ne0$ it
vanishes at $s=1/\epsilon$, so it has no entire reciprocal
and is not a unit of $\cqE$. This proves
\eqref{eq:cq-ring-unit-nonopenness}, and in particular
no neighbourhood of the unit $1$ is contained in the unit set.
\end{proof}

Thus unit openness is impossible on this original ring,
not a technical condition that can be supplied by a
choice of its complex-vector topology. The explicit
quotient inverse above avoids that premise.
The bounded primary-source reading recorded at the
start of this section has supplied neither Fr\'echet
structure for $\tau$ nor the actual ideal property of
$\cqN$. These remain specific unresolved source assertions.
The original-space classification, complete consequences
within its stated class, exact ring representation,
closed-ideal spectrum obstruction, and unit-set
nonopenness have all been proved without identifying
that source topology by assumption.

\section{Proof scope and reproducibility}
The new source-topology assertions proved here are the complete
transported closure, homeomorphic quotient, conjugated spectral
operator, source-domain differential correspondence, exact coordinate
kernels and cokernels, correlated jet extension, finite free cover,
compact-test heat ingress with its weighted-moment inverse domains,
and exact comparison diagrams. The compact and growth relation
subspaces both equal $\Xi\cqE$ by the complete division proof.
Their quotient topologies are computed: the compact quotient has
its exact jet-image topology and product completion, whereas the
growth quotient is complete and the comparison inverse is not
continuous. All derivative relation domains and finite ambiguity
projections retain the full zero multiplicities and analytic units.
The arithmetic principal-part realization has the exact kernel
$\Xi\cqE$, every finite image, the same product completion,
and explicit original-coordinate spectral blocks. The full Schwartz
trace factorization and its order-class intersection are proved
without assigning every Schwartz trace to that order class.
Its exact Mellin image has a proved inverse and complete
finite-strip estimates, including every Taylor residue at the
input origin. The exponential trace orbit has its exact open
Schwartz boundary. The meromorphic arithmetic heat action
retains the full zeta connection, its pole cancellations and
the convergent numerator jets determining its Schwartz domain.
The full persistent heat and even-derivative cores of $\Xi\cqE$
are zero. Every nonzero input returns to that analytic relation,
or to the original Schwartz trace intersection, only at a closed
discrete set of times; the complete proof retains every actual
zero multiplicity and uses the original Hadamard product.
The actual four-term $\xi^4$ coefficient is proved nonzero
from its complete Fourier multiplier, commutes exactly with
heat, and has those same discrete arithmetic return restrictions,
with its original incidence factor retained.
The retained Cartesian bilaplacian has its complete four
coefficients and exact scaled residue. That original residue
has full complex ambiguity on every actual moving source
quotient, including time zero. Its finite trace sections and
the precise quotient topology give an explicit homeomorphic
residue-zero presentation. Every finite origin jet has a
reciprocal-$\Xi$ polynomial section, and at every nonzero time
every one-point finite jet is realized by transported original
polynomial traces. The complete normal-order correction,
nonzero-polynomial obstruction, finite section construction
and both quotient maps are proved. The separately transported
residue has its full convergent derivative series and the
exact heat-conjugate quotient.
The Cartan construction now has a proved standard-real-subspace
modular double for the original heat multiplier, full spectral
domains at both signs of the parameter, and an exact conjugate
moving-source quotient diagram. Its involution reverses the radial
parameter. The threefold Fourier map identifies the positive
Euclidean viscous generator with the original heat coefficient.
For the retained nonlinear spatial pullback, the complete splitting,
compressed diffusion, constant leakage kernel, affine commutator
kernel, and second-diffusion feedback are computed explicitly.
The separate public Navier--Stokes theorem is recorded as an
imported source claim with distinct revision receipts; its full
proof has not been independently verified in this fragment.
Its actual cutoff-summed witness is now an explicit imported
input to the full arithmetic flow map. The original polynomial
target, inverse differential, pressure, forcing, viscosity scaling,
material metric, both time generators, and the two exact Fourier
norm constants are retained. The source's particular growing
swirl produces a complete arithmetic derivative channel and an
unbounded residue with its whole source remainder. The joint
profile and residue energies stay bounded, and the full profile
has a proved terminal limit in its stated Hilbert space.
The actual moving source class and its exact scalar kernel are
calculated. A common quotient by $\cqN\cap\Xi\cqE$ maps to both
source and analytic quotients with proved kernels and a continuous
residue recovery projection. The full logarithmic connection,
varying source sample labels, and translated Cartan reversal
are proved. A separate exact spectral model and generalized-chain
argument identify what additional actual source spectral data
would decide frozen descent without imposing a replacement topology.
The specified source Sobolev cokernel has an exact dense
compact-test ingress and a common-test comparison with $\cqQ$,
with both kernels and all generator signs proved.
Both explicit realizations have proved frozen differential and
heat noninvariance and the original-coordinate divisor spectrum.
Each heated finite Hermite trace also has its proved unique
M\"obius inverse, precise Mellin domain and full nontrivial-zero
pole orders, with actual failures of boundedness at the origin
for suitable original seeds at every nonzero time.
The missing identification of the paper's unnamed topology is
recorded explicitly and is not replaced by a conditional arithmetic
theorem. Nothing in these results establishes an off-critical zero,
a negative Weil value, or an arithmetic endpoint by analogy with
the inverse-fibre collision. The exact transport through that
collision cover and the derivative information it requires have
been calculated above.


The angular-momentum map retains the exact Cartesian torque and
symmetric-stress divergence. Its full nonlinear fluxes, half-line
endpoint jets and inverse arithmetic dilation are proved.
The source's actual fixed-axis coefficient gives a genuine Mellin
pole at $-1$. The separate radial lift has a proved inverse and
entire-function growth criterion. The full source recurrence proves
radial analyticity of the actual corrected central field, including
its pressure and force. Thus its untruncated radial arithmetic trace
belongs to $\Xi\cqE$ on the stated central slab, and its value at the
fixed original spectral point $-5i/2$ grows with exactly the source
power $\tau^{-1-h}$. Evaluation on the actual source quotient has
full scalar ambiguity, with an explicit trace section. Whole input
sequences provide the exact ingress elsewhere without presuming
analyticity of arbitrary smooth spatial cutoffs.


The central meridional streamfunction and every summation-cutoff
derivative are written explicitly, as are all three physical-force
coefficients. The full radial entire evolution has a proved
growth-topology smoothness estimate and an exact shift by $-2i$
in the original spectral coordinate. Its equation defect is the
nonzero entire multiplier times the radial trace of the complete
physical momentum residual. Injectivity proves equivalence of
the two full equations on the stated inverse domain, including
at the isolated multiplier zeros.

The original entire-function ring has a unique Hausdorff Fr\'echet
vector-space topology with continuous coordinate multiplication,
proved here to be the complete growth topology. The actual primary
source has not been proved to be in that class. Without that
identification, its exact closed multiplier algebra and largest
annihilator ideal still give a faithful quotient representation.
The full closed-ideal spectral obstruction and nonopenness of the
original unit set are proved; the actual source ideal property,
and therefore its remaining frozen-descent decision, remain open
in the inspected source data.

The companion checker uses rational polynomial arithmetic to verify
the heat conjugation on finite polynomial inputs, the retained
product terms, the matrix commutator, the iterated branch-domain
coefficients, the original-coordinate relation, and sample finite
jet multiplication, the multiplicity-sensitive derivative domains,
the source heat commutator recurrence, exact M\"obius divisor sums,
and Mellin principal-part coordinate inverses.
The continuation checks the exponential Gaussian input and its
generator, inverse Mellin differentiation and arithmetic residues,
the Sobolev cutoff and Fourier moments, and the meromorphic
heat conjugacy with all local pole terms.
It also checks the unchanged reflection composition, time-jet
coefficients and positive Hadamard counting integrand used to
compute the persistent core.
The final spatial checks retain all Cartesian derivatives,
all four bilaplacian coefficients, the residue limit and
original four-term coefficient, the reciprocal-jet sections,
and every tested normal-order and heat-conjugation correction.
The Cartesian replay also checks the full spatial compression,
leakage and its returning component, and the original scalar
viscous residual. The modular domains and quotient diagram have
an independent written proof audit.
The new original three-coordinate replay checks its full
determinant and cofactor inverse, the actual source sample
matrix and inverse, the swirl channel, both time derivatives,
all three material accelerations, the translated spectral
conjugacy, and every tested logarithmic unit and residue
coefficient. The varying material-label tangent signs and
the quantitative source lower-bound denominator are retained.
The angular replay checks the full Cartesian torque identities,
physical radial diffusion and graph coordinate, dyadic endpoint
coefficients, exact radial Mellin residues, similarity derivative
maps, sparse recurrence products and the fixed entire trace value.
These are reproducible checks of the displayed
identities, not a substitute for their written proofs. Source
hashes and exact reading locators are in the companion provenance
receipt; no numerical zero search or Lean run is part of this work.
